% arithmetic_shadows.tex -- Arithmetic content of the shadow obstruction tower
%
% Governing constants are imported from chapters/examples/landscape_census.tex.

\chapter{Arithmetic content of the shadow obstruction tower}
\label{chap:arithmetic-shadows}

\providecommand{\PhiFA}{\Phi^{\mathrm{FA}}}
\providecommand{\SpCh}{\mathrm{Sp}^{\mathrm{ch}}}

The shadow tower $\Theta_\cA$ of a chiral algebra $\cA$ is a
sequence of rational numbers $\kappa, S_3, S_4, \ldots$ that
encode the obstruction to formality of the bar complex. Over
$\bC$, these are just numbers. Over $\bQ$, they are arithmetic
invariants: their denominators detect Bernoulli primes, their
$p$-adic valuations encode local Euler factors, and their
congruence classes see Hecke eigenforms. The question is not
whether $\kappa$ is rational (it always is) but what structure
its rationality reveals.

Genus~$1$ forces the issue. The free energy
$F_1 = \kappa/24$ is the leading term of the $q$-expansion
on $\cM_{1,1}$. For Heisenberg at level $k$:
$\kappa = k$, so $F_1 = k/24$.
For Virasoro at central charge $c$: $\kappa = c/2$,
so $F_1 = c/48$; at $c = 26$ this gives $F_1 = 13/24$.
For affine Kac--Moody at level $k$:
$\kappa = \dim(\fg)(k + h^\vee)/(2h^\vee)$,
so $F_1 = \dim(\fg)(k + h^\vee)/(48 h^\vee)$.
The denominator $24$ is the Bernoulli factor
$B_2/2 = 1/12$ inverted; the same $24$ governs
$\eta(\tau) = q^{1/24}\prod(1 - q^n)$. This coincidence
is the first indication that the shadow tower sees modular
arithmetic together with rational arithmetic.

The genus-$1$ propagator is the quasi-modular Eisenstein series
$E_2^*(\tau) = 1 - 24\sum_{n \ge 1} \sigma_1(n)\,q^n$
(Proposition~\ref{prop:quasi-modular-propagator}): weight $2$,
depth $1$, and \emph{not} a holomorphic modular form.
The degree-$2$ shadow amplitude at genus~$1$ is
$\mathrm{Sh}_2^{(1)}(\tau) = \kappa \cdot E_2^*(\tau)$,
and the full genus-$1$ amplitude is a polynomial in
$\{E_2^*, E_4, E_6\}$: the ring of quasi-modular forms,
not the ring $\bC[E_4, E_6]$ of holomorphic modular forms.
The quasi-modular anomaly of $E_2^*$ under
$\tau \mapsto -1/\tau$ produces the non-holomorphic
completion $\widehat{E}_2 = E_2^* - 3/(\pi\,\mathrm{Im}\,\tau)$,
and this completion is the source of the modular
anomaly in the shadow tower at higher degree.
Graph sums at degree $r$ produce degree-$r$ polynomials
in $E_2^*$, pushing the arithmetic into the full
quasi-modular ring.

The Leech lattice vertex algebra makes the arithmetic
structure explicit at a single degree. Its shadow
coefficient is $S_{12} = -65520/691$, where $691$ is
the numerator of $B_{12}/12$ and the index of the
Eisenstein congruence $\tau(n) \equiv \sigma_{11}(n)
\pmod{691}$. The prime $691$ appears on the chiral
side for the same reason it appears on the arithmetic
side: the Hecke decomposition
$\Theta_{\mathrm{Leech}} = E_{12} - (65520/691)\,\Delta_{12}$
in $M_{12}(\mathrm{SL}_2(\bZ))$ projects the lattice theta
function onto the Ramanujan cusp form, and the coefficient
$65520/691$ is determined by the vanishing of the
representation number $r_{\mathrm{Leech}}(2) = 0$ (no vectors
of norm $2$ in the Leech lattice).

This is the data of a general phenomenon. The shadow
obstruction tower, classified by shadow depth into the hierarchy
$\mathbf{G}/\mathbf{L}/\mathbf{C}/\mathbf{F}_d/\mathbf{M}$
(Definition~\ref{def:shadow-depth-classification},
Theorem~\ref{thm:shadow-archetype-classification}), detects Hecke
eigenforms through its graph-sum amplitudes. Each degree resolves
one eigenspace of the constrained Epstein zeta
$\varepsilon^c_s(\cA) = \sum_{\Delta \in \mathcal{S}}(2\Delta)^{-s}$,
which factors into $L$-functions. This factorization defines the
arithmetic component $d_{\mathrm{arith}}$ of shadow depth.
A second component is purely homotopy-theoretic: the $A_\infty$
non-formality depth $d_{\mathrm{alg}}$, measuring how many
transferred operations $m_r$ on bar cohomology are nonzero.
For Virasoro, $d_{\mathrm{arith}}$ is finite but
$d_{\mathrm{alg}} = \infty$: the tower detects the algebra
itself via the self-referential OPE $T \in T_{(1)} T$,
not its representation theory. The depth decomposition
$d = 1 + d_{\mathrm{arith}} + d_{\mathrm{alg}}$
(Theorem~\ref{thm:depth-decomposition}) separates arithmetic
from homotopy. The discriminant, the quadratic extension
cut out by $\kappa$, and the $p$-adic locality of the
shadow ring of integers organize the arithmetic side below.

\smallskip
\noindent\emph{Epistemic scope.}\;
Two regimes of the level parameter $k$ produce two
kinds of arithmetic. At the critical level $k = -h^\vee$,
$\kappa$ vanishes, the bar complex becomes
uncurved, and $H^*(\barB(\fg_{-h^\vee}))$
recovers the Feigin--Frenkel center with its oper
differential-form package
(Theorem~\ref{thm:langlands-bar-bridge}).
This is the Langlands regime: the arithmetic lives in the
opers, and the shadow tower has collapsed.
The complementary regime $\kappa \ne 0$ is where the
tower does \emph{not} collapse, and where its arithmetic
content is visible.

The depth decomposition
$d = 1 + d_{\mathrm{arith}} + d_{\mathrm{alg}}$
(Theorem~\ref{thm:depth-decomposition}) separates the
two sources of tower non-triviality.
The arithmetic component $d_{\mathrm{arith}}$
counts the critical lines of the constrained Epstein
zeta $\varepsilon^c_s$: one for each $L$-function
factor in the Hecke decomposition of the lattice theta
function. The algebraic component $d_{\mathrm{alg}} \ge 0$
is the $A_\infty$ non-formality depth
(Theorem~\ref{thm:ainfty-formality-depth}):
$d_{\mathrm{alg}} = \max\{r : m_r \ne 0\} - 2$ for finite
shadow depth, $d_{\mathrm{alg}} = \infty$ when all
transferred operations are nonzero. For lattice vertex algebras,
$d_{\mathrm{arith}} = 2 + \dim S_{r/2}$ is finite because
the Hecke decomposition terminates. The shadow--spectral
correspondence
(Theorem~\ref{thm:shadow-spectral-correspondence})
identifies each Hecke eigenform $f_j$ with a depth stratum.

The lattice-VOA Ramanujan comparison criterion
(Corollary~\ref{cor:unconditional-lattice}) follows the conditional
route
$\mathrm{MC} \Rightarrow \text{prime-locality}
 \Rightarrow \text{CPS} \Rightarrow \mathrm{Sym}^{r-1}
 \Rightarrow \text{Ramanujan}$.
The chain passes through the converse theorem and Langlands
functoriality; it compares the shadow tower with the same eigenvalue
bounds proved by Deligne from the Weil conjectures.
The arithmetic comparison conjecture
(Conjecture~\ref{conj:arithmetic-comparison}), that
$\Theta_\cA$ canonically determines the arithmetic packet
connection $\nabla^{\mathrm{arith}}_\cA$, remains open.

\smallskip
\noindent\emph{Limitation: scalar blindness for Niemeier lattices.}\;
The scalar tower assigns $\kappa = 24$ and
total depth $d = 4$ to all $24$ Niemeier lattice
vertex algebras $V_\Lambda$, regardless of root system.
The scalar invariants $(\kappa, S_3, S_4, \ldots)$ depend
on $c = \operatorname{rank}(\Lambda) = 24$ alone; they
cannot distinguish $V_{E_8^3}$ from
$V_{D_{16}^+ \oplus D_{16}^+}$ or the Leech lattice.
The discrimination requires the charged-sector
components of $\Theta_\cA$: the decomposition of the
bar complex by the lattice grading
$\gamma \in \Lambda^*/\Lambda$ produces per-sector
invariants that encode the root system structure.
The full genus-$2$ Siegel theta function distinguishes all
$24$ lattices (King, 2003);
the multi-channel bar complex provides an independent
complete discrimination already at genus $1$.
The single B\"ocherer coefficient $c_2$
(the cusp projection onto $\chi_{12}$)
separates only $20$ of $24$; the remaining $4$ pairs
require other genus-$2$ Fourier coefficients.

The passage from genus-$1$ quasi-modular forms to higher-genus
arithmetic requires the moduli space
$\mathrm{SL}(2,\bZ)\backslash\mathfrak{H}$. On a fixed elliptic
curve, the sewing operator $K_q$ acts diagonally on the Fock space
with eigenvalues $q^n$; its Fredholm determinant
$\det(1-K_q) = \prod(1-q^n)$ recovers
$\eta(\tau) = q^{1/24}\prod(1-q^n)$, and the spectral zeta
function $\sum n^{-s} = \zeta(s)$ is the Riemann zeta function.
None of this constrains the zeros of $\zeta$: the constraint
requires variation over moduli. There the primary-counting
function
$\widehat{Z}^c = y^{c/2}\lvert\eta\rvert^{2c} Z$ enters the
Roelcke--Selberg spectral decomposition, and the real-analytic
Eisenstein series $E_s(\tau)$ carries the scattering matrix
\begin{equation}\label{eq:scattering-matrix}
 \varphi(s) = \frac{\Lambda(1-s)}{\Lambda(s)},
 \qquad
 \Lambda(s) := \pi^{-s}\,\zeta(2s)\,\Gamma(s),
\end{equation}
whose poles are the rescaled nontrivial zeros of $\zeta$.
The scattering matrix encodes the zeros; the shadow tower
constrains the spectral coefficients. The question is
whether these constraints are strong enough to reach the
zeros. The answer, for lattice algebras, is that
the shadow depth counts exactly the critical lines of the
constrained Epstein zeta:

\begin{theorem}[Shadow--spectral correspondence]
\label{thm:shadow-spectral-correspondence}
\ClaimStatusProvedElsewhere
Let $V_\Lambda$ be an even lattice vertex algebra of rank~$r$ with
theta function $\Theta_\Lambda\in M_{r/2}(\Gamma)$. Write
$\Theta_\Lambda = c_E\, E_{r/2} + \sum_j c_j f_j$ for the
decomposition into the Eisenstein series~$E_{r/2}$ and Hecke
eigenforms $f_j\in S_{r/2}(\Gamma)$. Then:
\begin{enumerate}[label=\textup{(\roman*)}]
\item The constrained Epstein zeta factors as
\begin{equation}\label{eq:epstein-factorization}
 \varepsilon^r_s(V_\Lambda)
 = 2^{-s}\Bigl(
 c_E\,C_E(s)\,\zeta(s)\,\zeta(s-\tfrac{r}{2}+1)
 + \sum_j c_j\,C_j(s)\,L(s,f_j)
 \Bigr),
\end{equation}
where $C_E$, $C_j$ are explicit products of gamma factors and powers
of~$\pi$.
\item The number of critical lines of~$\varepsilon^r_s$ (the
distinct values of $\mathrm{Re}(s)$ on which $\varepsilon^r_s$ has
infinitely many zeros) equals the shadow depth $d(V_\Lambda)$
minus~$1$.
\end{enumerate}
\end{theorem}

Equality holds for lattice algebras; for non-lattice algebras the
shadow depth exceeds the critical-line count
(\S\ref{sec:non-lattice-theories}).

\begin{remark}[Lattice specificity of the Hecke pipeline]
\label{rem:lattice-specificity}
\index{lattice VOA!Hecke pipeline specificity}
\index{Hecke decomposition!lattice-specific}
For lattice VOAs $V_\Lambda$, the partition function
$Z_1(V_\Lambda, q) = \Theta_\Lambda(q) / \eta(q)^{\operatorname{rank}(\Lambda)}$
decomposes via Hecke theory into Eisenstein and cuspidal
$L$-functions. This decomposition relies on three properties
special to the lattice setting:
\begin{enumerate}[label=\textup{(\roman*)}]
\item the theta function $\Theta_\Lambda \in M_{r/2}(\Gamma)$ is a
 classical modular form of known weight and level;
\item the space $M_{r/2}(\Gamma)$ is finite-dimensional, so the
 Hecke decomposition into eigenforms is a finite sum;
\item the Fourier coefficients of $\Theta_\Lambda$ are lattice
 representation numbers $r_\Lambda(n) = |\{v \in \Lambda :
 |v|^2 = 2n\}|$, giving the Dirichlet series
 $\sum r_\Lambda(n)\,n^{-s}$ a direct arithmetic meaning.
\end{enumerate}
None of these properties holds for non-lattice algebras in
general. For the Virasoro algebra, the character is
$\chi_0(q) = q^{-c/24}\prod(1-q^n)^{-1}$, which is not a
classical modular form of any finite weight. The shadow tower
still exists with $r_{\max} = \infty$
(Theorem~\ref{thm:shadow-archetype-classification}), but the
arithmetic sieve of~\S\ref{sec:five-lattice-verifications} does
not apply: there is no theta function to decompose.
The non-lattice extension is treated in
\S\ref{sec:non-lattice-theories}.
\end{remark}

% ============================================================
\section{The sewing--Rankin--Selberg bridge}
\label{sec:sewing-RS-bridge}

The genus-$1$ amplitude is quasi-modular; the Epstein zeta is
a Dirichlet series; the Hecke decomposition factorizes it into
$L$-functions. The bridge from sewing to spectral theory
passes through eight steps, each proved in the sections that
follow:
\begin{enumerate}[label=\textup{(\arabic*)}]
\item \emph{Sewing}:
 $\det(1{-}K_q) = \prod(1{-}q^n)$
 (Fredholm determinant of the sewing operator on Fock space).
\item \emph{Eta}:
 $\eta(\tau) = e^{\pi i\tau/12}\det(1{-}K_q)$
 (Dedekind $\eta$ from the sewing determinant).
\item \emph{Partition function}:
 $Z_\cA(\tau) = \Theta/\eta^r$
 (lattice theta divided by oscillator contribution).
\item \emph{Bar cohomology}:
 $\widehat{Z}^c = y^{c/2}\lvert\eta\rvert^{2c}Z$
 (primary-counting function strips descendants).
\item \emph{Constrained Epstein}:
 $\varepsilon^c_s = \sum(2\Delta)^{-s}$
 (Dirichlet series over scalar primaries $= $Eisenstein spectral
 coefficient of $\widehat{Z}^c$).
\item \emph{Functional equation}:
 $\varepsilon^c_{c/2-s} = F(s,c)\cdot\varepsilon^c_{c/2+s-1}$
 (from the modular invariance of~$Z$; the factor~$F$ involves
 $\zeta(2s)/\zeta(2s{-}1)$).
\item \emph{$L$-function factorization}:
 Hecke decomposition of~$\Theta$ gives
 $\varepsilon^c_s = \sum c_j\,L(s,f_j)$
 (each eigenform contributes one critical line).
\item \emph{Shadow--spectral correspondence} (lattice case):
 $d(\cA) = 1 + {}$critical-line count
 (the shadow depth counts the eigenforms; equality proved for
 lattice VOAs, inequality $\ge$ in general).
\end{enumerate}

\subsection{The sewing operator on a single curve}

The arithmetic enters through the Fredholm determinant.
Fix $\tau\in\mathfrak{H}$ and set $q=e^{2\pi i\tau}$. The sewing
operator $K_q$ acts on the Fock space $\cF=\bigoplus_{n\ge0}\cF_n$
diagonally: $K_q\lvert n\rangle = q^n\lvert n\rangle$ for $n\ge1$.
It is self-adjoint, positive, and trace class:
$\operatorname{tr}(K_q)=q/(1-q)$, $\det(1-K_q)=\prod(1-q^n)$.
The logarithm of this determinant is a $q$-series whose
coefficients are divisor sums, and divisor sums are
multiplicative: the Euler product is already present
in the Fredholm determinant of a single free boson.

\begin{proposition}[Divisor-sum decomposition]
\label{prop:divisor-sum-decomposition}
\ClaimStatusProvedHere
\begin{equation}\label{eq:log-det-sigma}
 \log\det(1-K_q)
 = -\sum_{N\ge1}\sigma_{-1}(N)\,q^N,
 \qquad
 \sigma_{-1}(N) = \sum_{d\mid N}\frac{1}{d}.
\end{equation}
The function $\sigma_{-1}$ is multiplicative, and its Dirichlet series
satisfies the Ramanujan identity
\begin{equation}\label{eq:sigma-minus-1-dirichlet}
 \sum_{N\ge1}\sigma_{-1}(N)\,N^{-s}
 = \zeta(s)\,\zeta(s+1).
\end{equation}
\end{proposition}

\begin{proof}
Expand
$\log\det(1-K_q)=\sum_n\log(1-q^n)=-\sum_n\sum_m q^{nm}/m$
and reindex $N=nm$. Multiplicativity of $\sigma_{-1}$ is standard;
the Dirichlet series identity is the $k=-1$ case of
$\sum\sigma_k(N)N^{-s}=\zeta(s)\zeta(s-k)$.
\end{proof}

\begin{corollary}[Euler product of the sewing determinant]
\label{cor:sewing-euler-product}
\ClaimStatusProvedElsewhere
The identity~\eqref{eq:sigma-minus-1-dirichlet} gives a prime
factorization of the sewing Fredholm determinant:
\[
 \zeta(s)\,\zeta(s+1)
 = \prod_p \frac{1}{(1-p^{-s})(1-p^{-s-1})}.
\]
Each prime~$p$ contributes an independent factor to the sewing
amplitude.
\end{corollary}

\begin{proposition}[Sewing trace formula]
\label{prop:sewing-trace-formula}
\ClaimStatusProvedHere
\index{sewing operator!trace formula}
The trace of the $N$-th power of the sewing operator is
$\operatorname{tr}(K_q^N) = q^N/(1-q^N)$, and
\begin{equation}\label{eq:sewing-trace}
 \sum_{N\ge1} \frac{\operatorname{tr}(K_q^N)}{N}
 = -\log\det(1-K_q)
 = \sum_{M\ge1} \sigma_{-1}(M)\,q^M.
\end{equation}
The divisor-counting function $\sigma_0(N) = \sum_{d\mid N}1$ satisfies
\begin{equation}\label{eq:sigma0-zeta-squared}
 \sum_{N\ge1}\sigma_0(N)\,N^{-s} = \zeta(s)^2.
\end{equation}
This is the Dirichlet series of the \emph{tensor square}
$K_q\otimes K_q$: the number of ordered pairs $(a,b)$ with
$ab = N$ equals~$\sigma_0(N)$.
\end{proposition}

\begin{proof}
$\operatorname{tr}(K_q^N) = \sum_n q^{nN} = q^N/(1-q^N)$.
Then $\sum\operatorname{tr}(K_q^N)/N
= -\sum_n\sum_m q^{nm}/m = -\log\prod(1-q^n)$, recovering the
divisor-sum expansion. For~\eqref{eq:sigma0-zeta-squared}:
$\sigma_0 = \mathbf{1}\ast\mathbf{1}$ (Dirichlet convolution of
the constant function~$1$ with itself), so
$\sum\sigma_0 N^{-s} = (\sum N^{-s})^2 = \zeta(s)^2$.
\end{proof}

\subsection{The sewing--Selberg formula}

The divisor-sum expansion of $\log\det(1-K_q)$ encodes
the arithmetic of a single curve. Integration over
$\cM_{1,1}$ lifts this to the Rankin--Selberg integral,
and the Mellin transform converts the $q$-expansion
into a product of zeta values.

\begin{theorem}[Sewing--Selberg formula]
\label{thm:sewing-selberg-formula}
\ClaimStatusProvedHere
The Rankin--Selberg integral of the sewing Fredholm determinant
against the Eisenstein series over $\cM_{1,1}$ is
\begin{equation}\label{eq:sewing-selberg}
 \int_{\cM_{1,1}}
 \log\det(1-K(\tau))\cdot E_s(\tau)\,d\mu(\tau)
 = -2(2\pi)^{-(s-1)}\,\Gamma(s-1)\,\zeta(s-1)\,\zeta(s).
\end{equation}
\end{theorem}

\begin{proof}
Rankin--Selberg unfolding replaces the integral over the fundamental
domain by $\int_0^\infty a_0(y)\,y^{s-2}\,dy$, where
$a_0(y) = \int_0^1 \log\det(1-K(\tau))\,dx$ is the zeroth Fourier
mode. On the imaginary axis,
$a_0(y) = \log\lvert\eta(iy)\rvert^2+\pi y/6
= 2\log\det(1-K_q)$, and the Mellin transform of
$-\sum\sigma_{-1}(N)e^{-2\pi Ny}$ against $y^{s-2}$ gives
$(2\pi N)^{-(s-1)}\Gamma(s-1)$ per term. Summing via
\eqref{eq:sigma-minus-1-dirichlet} yields~\eqref{eq:sewing-selberg}.
\end{proof}

\begin{remark}[Structural obstruction]
\label{rem:structural-obstruction}
The scattering matrix $\varphi(s)=\Lambda(1{-}s)/\Lambda(s)$ has
poles at $s=\rho/2$ and zeros at $s=(1{-}\bar\rho)/2$, where $\rho$
ranges over the nontrivial zeros of~$\zeta$. In the spectral
decomposition, the parameter~$t$ is
\emph{real} ($s=\tfrac12+it$), and $\varphi(\tfrac12+it)$
is regular for all $t\in\bR$: its poles are at
$\mathrm{Re}(s)=\mathrm{Re}(\rho)/2\ne\tfrac12$
(assuming $\rho$ is not on the boundary of the critical strip),
hence at \emph{complex} spectral parameters. The shadow tower and
the MC equation constrain the spectral coefficients~$c(t)$ on the
real $t$-axis; the zeta zeros live at complex~$t$. This separation
is \emph{structural}: algebraic constraints on the spectral line
cannot reach the scattering poles without analytic continuation of
the spectral data off the real axis.

The completed Selberg function $\Lambda(s)=\pi^{-s}\Gamma(s)\zeta(2s)$
satisfies $\Lambda(s_0)\ne 0$ and $\Lambda(1{-}s_0)=0$ at
$s_0=(1{+}\rho)/2$, because $\zeta(2{-}2s_0)=\zeta(1{-}\rho)=0$
(the conjugate zero). The quotient
$\varphi(s_0)=\Lambda(1{-}s_0)/\Lambda(s_0)=0$: the ``pole'' of
the uncompleted functional-equation factor
$\zeta(2s)/\zeta(2s{-}1)$ is absorbed into a
\emph{zero} of the completed scattering matrix. No zero-forcing
mechanism for $\varepsilon^c_s$ operates at these points.
\end{remark}

% ============================================================
\section{Four identities}\label{sec:four-identities}

\subsection{The Narain universality theorem}

\begin{theorem}[Narain universality]
\label{thm:narain-universality}
\ClaimStatusProvedHere
For the rank-$1$ Narain lattice vertex algebra at compactification
radius~$R$:
\begin{equation}\label{eq:narain-universality}
 \varepsilon^1_s(R)
 = 2\bigl(R^{2s}+R^{-2s}\bigr)\,\zeta(2s).
\end{equation}
The character factor $2\cosh(2s\log R)$ is positive for
real~$s$ and $R>0$. Consequently:
\begin{enumerate}[label=\textup{(\roman*)}]
\item The zeros of $\varepsilon^1_s(R)$ are the zeros of
$\zeta(2s)$, independent of~$R$.
\item $T$-duality $R\leftrightarrow 1/R$ preserves
$\varepsilon^1_s$ identically, corresponding to Koszul
self-duality at the self-dual radius.
\item The Narain moduli provide no bootstrap constraint on
the zero locations of~$\zeta$.
\end{enumerate}
\end{theorem}

\begin{proof}
The scalar primaries at radius~$R$ have dimensions
$\Delta_n=n^2/(2R^2)$ (momentum, multiplicity~$2$) and
$\Delta_w=w^2R^2/2$ (winding, multiplicity~$2$). Then
$\varepsilon^1_s(R) = 2R^{2s}\zeta(2s)+2R^{-2s}\zeta(2s)$.
\end{proof}

\subsection{The $E_8$ factorization}

\begin{theorem}[$E_8$ Epstein factorization]
\label{thm:e8-epstein}
\ClaimStatusProvedHere
For the $E_8$ lattice vertex algebra at level~$1$ ($c=8$):
\begin{equation}\label{eq:e8-epstein}
 \varepsilon^8_s(V_{E_8})
 = 240\cdot 2^{-s}\,\zeta(s)\,\zeta(s-3).
\end{equation}
The zeros lie on two critical lines:
$\mathrm{Re}(s) = \tfrac12$ and
$\mathrm{Re}(s)=\tfrac72$.
\end{theorem}

\begin{proof}
The theta function of~$E_8$ is the Eisenstein series:
$\Theta_{E_8}(\tau)=E_4(\tau)=1+240\sum\sigma_3(n)q^n$. The
representation count is $r_8(2n)=240\,\sigma_3(n)$. The Epstein zeta
$E_{E_8}(s)=\sum r_8(2n)(2n)^{-s}=240\cdot
2^{-s}\sum\sigma_3(n)n^{-s}=240\cdot 2^{-s}\zeta(s)\zeta(s-3)$,
using the Ramanujan identity $\sum\sigma_k(n)n^{-s}=\zeta(s)\zeta(s-k)$.
Then $\varepsilon^8_s=E_{E_8}(s)$.
\end{proof}

\subsection{The square-lattice Dedekind zeta}

\begin{proposition}[$\bZ^2$ Epstein zeta]
\label{prop:z2-epstein}
\ClaimStatusProvedHere
For the rank-$2$ square lattice:
\begin{equation}\label{eq:z2-epstein}
 E_{\bZ^2}(s)
 = 4\,\zeta(s)\,L(s,\chi_{-4})
 = 4\,\zeta_{\bQ(i)}(s),
\end{equation}
where $\chi_{-4}$ is the Kronecker character and
$\zeta_{\bQ(i)}$ is the Dedekind zeta function of the Gaussian
integers. Both factors have their nontrivial zeros on
$\mathrm{Re}(s)=\tfrac12$: one critical line.
\end{proposition}

\begin{proof}
The representation count $r_2(n)=4\sum_{d\mid n}\chi_{-4}(d)$ is
standard; the Dirichlet series identity follows from the Euler product
of the Dedekind zeta of~$\bQ(i)$.
\end{proof}

\subsection{The Leech factorization}

\begin{proposition}[Leech Epstein factorization]
\label{prop:leech-epstein}
\ClaimStatusProvedHere
For the Leech lattice ($r=24$), the Hecke decomposition of
$\Theta_{\mathrm{Leech}} = E_4^3 - 720\,\Delta$ in $M_{12}$
is
\begin{equation}\label{eq:leech-hecke}
 \Theta_{\mathrm{Leech}}
 = E_{12} - \frac{65520}{691}\,\Delta_{12},
\end{equation}
where $E_{12}$ is the weight-$12$ Eisenstein series and $\Delta_{12}$
is the Ramanujan cusp form. The Epstein zeta factors accordingly:
\begin{equation}\label{eq:leech-epstein}
 E_{\mathrm{Leech}}(s)
 = C_E(s)\,\zeta(s)\,\zeta(s-11)
 - \frac{65520}{691}\,C_\Delta(s)\,L(s,\Delta_{12}),
\end{equation}
where $C_E$, $C_\Delta$ are products of gamma factors and powers
of~$\pi$. Assuming GRH, the zeros lie on three critical lines:
$\mathrm{Re}(s)=\tfrac12$,
$\mathrm{Re}(s)=6$, and
$\mathrm{Re}(s)=\tfrac{23}{2}$.
\end{proposition}

\begin{proof}
The space $M_{12}(\mathrm{SL}(2,\bZ))$ is two-dimensional, spanned by
$E_{12}$ and $\Delta_{12}=\eta^{24}$. Since
$\Theta_{\mathrm{Leech}}$ and $E_{12}$ both have constant term~$1$,
we have $\Theta_{\mathrm{Leech}} - E_{12} \in S_{12} = \bC\,\Delta_{12}$.
Comparing $q$-coefficients:
$E_{12} = 1 + \tfrac{65520}{691}\sum\sigma_{11}(n)q^n$,
$\Theta_{\mathrm{Leech}} = 1 + 0\cdot q + 196560\,q^2 + \cdots$\,.
The $q$-coefficient gives
$0 - \tfrac{65520}{691} = c_\Delta\cdot\tau(1) = c_\Delta$,
so $c_\Delta = -\tfrac{65520}{691}$.
Verification: the $q^2$-coefficient is
$\tfrac{65520}{691}\sigma_{11}(2) - \tfrac{65520}{691}\tau(2)
= \tfrac{65520}{691}(2049+24)
= \tfrac{65520\cdot 2073}{691}
= 196560$. \checkmark

The Hecke correspondence gives
$\sum\sigma_{11}(n)n^{-s}=\zeta(s)\zeta(s{-}11)$ and
$\sum\tau(n)n^{-s}=L(s,\Delta_{12})$, yielding the
factorization~\eqref{eq:leech-epstein}. The critical lines are:
$\mathrm{Re}(s)=\tfrac12$ from $\zeta(s)$,
$\mathrm{Re}(s)=\tfrac{23}{2}$ from $\zeta(s{-}11)$, and
$\mathrm{Re}(s)=6$ from $L(s,\Delta_{12})$ (assuming GRH;
Deligne's theorem gives $|\tau(p)| \leq 2p^{11/2}$
but does not imply that zeros of~$L(s,\Delta_{12})$
lie on $\mathrm{Re}(s)=6$).
\end{proof}

\begin{remark}[Borcherds weight and chiral conductor are distinct
arithmetic faces]
\label{rem:arith-shad-11}%
\index{Phi functor@$\Phi$ functor!universal trace identity}%
\index{Borcherds singular theta correspondence!Koszul conductor bridge}%
\index{universal Phi-trace identity!K3 BKM}%
The Leech factorisation~\eqref{eq:leech-epstein} places the rational
prefactor $65520/691 = -S_{12}$ on the bar--cobar side of the
arithmetic shadow. Borcherds' singular-theta correspondence
\cite{Borcherds1998,Bruinier2002} supplies the automorphic
normalisation with which such rational coefficients must be compared:
for an even Lorentzian lattice
$\Lambda$ of signature $(2, n)$ the Borcherds regularised theta lift
$\Phi_\Lambda \colon M^!_{1 - n/2, \rho_\Lambda}(\mathrm{Mp}_2(\bZ))
\to \operatorname{Aut}(G_\Lambda \backslash \cH_\Lambda)$
produces automorphic products whose leading Fourier coefficient
$c_\Lambda(0)/2$ is the weight of the output automorphic form.
For the K3-fibered class-A threefold
$\cC = \mathrm{K3} \times E$, Vol~III records the Borcherds
weight as
\[
 \kappa_{\mathrm{BKM}}(\Delta_5)
 = \operatorname{wt}(\Delta_5)
 = c_1(0)/2
 = 5.
\]
This is not the Vol~I Koszul conductor. The conductor attached to
the Mukai-enhanced K3 Heisenberg row is
\[
 K^{\kappa_{\mathrm{ch}}}(\mathrm{Mukai}(K3))
 = 2c_+(\mathrm{II}_{4,20})
 = 8,
\]
the $\mathsf B$-row ceiling of Theorem~C. The two integers live on
different lanes: $5$ is the automorphic weight of the Borcherds
product, while $8$ is the Mukai-doubled chiral Koszul conductor.

Vol~III's two-stage CY-to-chiral factorisation
$\Phi = \SpCh_{\Sigma_2, E_\tau} \circ \PhiFA_3$
(Stage~1 canonical $E_3$-holomorphic factorisation algebra on $\cC$,
Stage~2 specialisation to the elliptic reference curve) gives a
comparison map between these lanes, not an equality of their scalar
invariants:
\begin{equation}\label{eq:phi-trace-universal}
 \Phi:
 \bigl(K^{\kappa_{\mathrm{ch}}}(\cA),\Theta_\cA\bigr)
 \longmapsto
 \bigl(\Phi_\Lambda(F_\cA),\operatorname{wt}\Phi_\Lambda(F_\cA)\bigr),
 \qquad
 \operatorname{wt}\Phi_\Lambda(F_\cA)=\tfrac12 c_\Lambda(0),
\end{equation}
where $F_\cA \in M^!_{1 - n/2, \rho_\Lambda}(\mathrm{Mp}_2(\bZ))$ is
the vector-valued nearly-holomorphic input determined by the chiral
OPE data of $\cA$ (its component at the trivial coset is the reduced
graded character of $\cA$ minus the Eisenstein part).
The proof obligation for a stronger statement is precise:
construct a natural transformation from the chiral trace on
$Z^{\mathrm{der}}_{\mathrm{ch}}(\cA)$ to the Borcherds theta lift
whose effect on the K3 base point sends the conductor $8$ to the
weight-$5$ automorphic form by a named normalisation map. No such
normalisation is used here as an equality.

The identity specialises at the Leech lattice to
$\operatorname{wt}(\Phi_{12})=12=c_{II_{25,1}}(0)/2$ through the
$24$-dimensional sublattice embedding
$\Lambda_{24} \hookrightarrow II_{25,1}$ that underlies the
Borcherds construction of the fake Monster Lie algebra: the
$-65520/691$ prefactor of the Hecke decomposition of
$\Theta_{\Lambda_{24}}$ and the $\chi_{12}$ projection coefficient of
\eqref{eq:leech-chi12-decomposition} are Hecke projection data.
They are not, by themselves, a supertrace identity on
$Z(V_{\Lambda_{24}})$. The true statement is the comparison of
normalisations:
\[
 S_{12}(V_{\Lambda_{24}})=-65520/691,\qquad
 \operatorname{wt}(\Phi_{12})=12,\qquad
 \operatorname{wt}(\Delta_5)=5,\qquad
 K^{\kappa_{\mathrm{ch}}}(\mathrm{Mukai}(K3))=8.
\]
The arithmetic bridge is the requirement that all four numbers be
kept on their native objects.
\end{remark}

% ============================================================
\section{Five lattice discriminants}
\label{sec:five-lattice-verifications}

\begin{center}
\renewcommand{\arraystretch}{1.3}
\begin{tabular}{lccccc}
\hline
\textbf{Algebra} & $r$ & $k{=}r/2$ & $\Theta$ decomposition
 & \textbf{Critical lines} & \textbf{Depth}\\
\hline
$V_{\bZ}$ & 1 & $\tfrac12$ & $\theta_3$
 & $1\ (\mathrm{Re}=\tfrac14)$ & 2\\
$V_{\bZ^2}$ & 2 & 1 & $\theta_3^2$
 & $1\ (\mathrm{Re}=\tfrac12)$ & 2\\
$V_{A_2}$ & 2 & 1 & hex $\theta$
 & $1\ (\mathrm{Re}=\tfrac12)$ & 2\\
$V_{E_8}$ & 8 & 4 & $E_4$
 & $2\ (\tfrac12,\,\tfrac72)$ & 3\\
$V_{\mathrm{Leech}}$ & 24 & 12 & $E_4^3-720\Delta$
 & $3\ (\tfrac12,\,6,\,\tfrac{23}{2})$ & 4\\
\hline
\end{tabular}
\end{center}

In every row, the depth equals $1$ plus the critical-line count.

% ============================================================
\subsection{The charged-sector resolution: beyond the scalar shadow}
\label{sec:charged-sector-resolution}

The five lattice rows expose an intrinsic deficiency of the scalar tower.
All even unimodular lattice vertex algebras $V_\Lambda$ of rank~$r$
share $\kappa(V_\Lambda) = r$, class~G shadow depth $r_{\max} = 2$,
and genus expansion $F_g(V_\Lambda) = r \cdot \lambda_g^{\mathrm{FP}}$
\textup{(}UNIFORM-WEIGHT\textup{)} at every genus.
For the $24$ Niemeier lattices ($r = 24$), the scalar invariants
$(\kappa, S_3, S_4, \ldots) = (24, 0, 0, \ldots)$ are identical:
the scalar tower is blind to the root
system~$R(\Lambda)$, the theta function~$\Theta_\Lambda$, and
the lattice geometry.
The Niemeier discrimination problem
admits resolution at three levels of the MC~hierarchy, each strictly
refining the previous.

\medskip
\noindent\textbf{Level~0: scalar shadow (trivial).}\;
The modular characteristic $\kappa(V_\Lambda) = 24$ for all
$24$ lattices.
The full scalar tower $(S_r)_{r \geq 2} = (24, 0, 0, \ldots)$
is universal.
Among the $\tbinom{24}{2} = 276$ lattice pairs, exactly $0$
are distinguished.

\medskip
\noindent\textbf{Level~1: genus-$1$ theta (partial).}\;
The genus-$1$ theta series $\Theta_\Lambda \in M_{12}(\mathrm{SL}(2,\bZ))$
is determined by a single integer, the root count~$|R(\Lambda)|$:
\[
 \Theta_\Lambda(\tau) \;=\; E_{12}(\tau)
 + \frac{691\cdot|R| - 65520}{691}\,\Delta(\tau),
\]
since $\dim M_{12}(\mathrm{SL}(2,\bZ)) = 2$.
The coefficient $|R(\Lambda)|$ is a genus-$1$ theta-series root count.
It is not a Borcherds weight $c(0)/2$ and not a root-multiplicity
count in a denominator algebra.
This distinguishes $271$ of $276$ pairs.
Five collision pairs remain, each consisting of lattices with
identical root counts:
\begin{center}
\small
\renewcommand{\arraystretch}{1.2}
\begin{tabular}{rcl}
\toprule
$|R|$ & \textbf{Pair} & \textbf{Root systems}\\
\midrule
$144$ & $4A_5 D_4$ vs.\ $6D_4$
 & $(5,5,5,5,4)$ vs.\ $(4^6)$\\
$240$ & $2A_9 D_6$ vs.\ $4D_6$
 & $(9,9,6)$ vs.\ $(6^4)$\\
$288$ & $A_{11} D_7 E_6$ vs.\ $4E_6$
 & $(11,7,6)$ vs.\ $(6^4)$\\
$432$ & $A_{17} E_7$ vs.\ $D_{10}(E_7)^2$
 & $(17,7)$ vs.\ $(10,7,7)$\\
$720$ & $D_{16} E_8$ vs.\ $3E_8$
 & $(16,8)$ vs.\ $(8^3)$\\
\bottomrule
\end{tabular}
\end{center}
The genus-$1$ theta series of each pair is
\emph{exactly} identical to all Fourier orders: the identity
$c_\Delta(\Lambda_1) = c_\Delta(\Lambda_2)$ is
an algebraic consequence of $|R(\Lambda_1)| = |R(\Lambda_2)|$.
The class-$\mathcal S$ $A_{N-1}$-on-$\Sigma_{0,24}$ labelling
singles out $A_5^4 D_4$ at $N = 6$ (Theorem
\ref{thm:walgdeep-N6-reanchor-A5-4-D4}) and the Niemeier root system
$2 A_7 D_5^2$ at $N = 8$ (where $D_5$ is the type-$D$ root lattice,
not the monic Borcherds product $D_5=64^{-1}\Delta_5$ of the OP
normalisation; Theorem~\ref{thm:walgdeep-N7-N8-re-anchor}): the
$(N-1)$-dimensional summand in the fugacity data selects the
Niemeier containing $A_{N-1}$ among the collision candidates
at common Coxeter slot.

\medskip
\noindent\textbf{Level~2: per-factor bar complex
(complete).}\;
The MC element $\Theta_{V_\Lambda}$
(Theorem~\ref{thm:mc2-bar-intrinsic})
decomposes along the root system.
For a Niemeier lattice with root system
$R(\Lambda) = R_1 \oplus \cdots \oplus R_k$,
each factor $V_1(\fg_i)$ (level-$1$ affine Kac--Moody of
type~$R_i$) contributes an independent summand to the
convolution algebra~$\fg^{\mathrm{mod}}_{V_\Lambda}$
(Proposition~\ref{prop:independent-sum-factorization}).
The per-factor data visible in the bar complex is:
\begin{enumerate}[label=\textup{(\alph*)}]
\item $\operatorname{rank}(R_i)$: the per-factor curvature
 $\kappa_i = \operatorname{rank}(R_i)$ (from the genus-$1$
 obstruction of each affine factor);
\item $h(R_i)$: the Coxeter number, which controls the
 PBW filtration on $H^1(\barB(V_1(\fg_i)))$ (the
 highest root has height $h - 1$, placing a bar generator
 at filtration degree~$h$).
\end{enumerate}

\begin{proposition}[Multi-channel Niemeier discrimination]
\label{prop:niemeier-multichannel}
\ClaimStatusProvedHere
\index{Niemeier lattice!multi-channel discrimination}
For the $24$ Niemeier lattices, the sorted per-factor
tuple $\{(\operatorname{rank}(R_i), h(R_i))\}_{i=1}^{k}$,
extracted from the bar complex $\barB(V_\Lambda)$,
is a complete invariant: it distinguishes all
$276$ pairs.
The per-factor rank vector alone distinguishes $267$
of $276$ pairs, leaving $9$ undistinguished pairs
in $5$ collision groups%
\textup{:}
\begin{center}
\small
\renewcommand{\arraystretch}{1.2}
\begin{tabular}{cl}
\toprule
\textbf{Rank partition} & \textbf{Collision}\\
\midrule
$(24)$ & $A_{24}$ vs.\ $D_{24}$\\
$(12,12)$ & $2A_{12}$ vs.\ $2D_{12}$\\
$(8,8,8)$ & $3A_8$,\; $3D_8$,\; $3E_8$\\
$(6,6,6,6)$ & $4A_6$,\; $4D_6$,\; $4E_6$\\
$(4^6)$ & $6A_4$ vs.\ $6D_4$\\
\bottomrule
\end{tabular}
\end{center}
Adjoining the common Coxeter number~$h$ resolves
every collision, since the map
$(n, h) \mapsto \text{ADE type}$ is injective\textup{:}
$A_n \mapsto (n, n{+}1)$, $D_n \mapsto (n, 2(n{-}1))$,
$E_6 \mapsto (6, 12)$, $E_7 \mapsto (7, 18)$,
$E_8 \mapsto (8, 30)$.
By Niemeier's theorem, all factors of a Niemeier root system
share a single Coxeter number~$h$, and the pair
$(h, \text{rank partition of } 24)$ determines the root system
type, hence the lattice.
\end{proposition}

\begin{proof}
The bar complex $\barB(V_\Lambda)$ decomposes as the tensor
product of per-factor bar complexes
(Proposition~\ref{prop:independent-sum-factorization}).
The first bar cohomology
$H^1(\barB(V_1(\fg_i))) \cong \fg_i^*$ carries the Lie
algebra type via its dimension and PBW filtration degree.
The ADE injectivity of $(n, h)$ is verified by
inspection of the Coxeter numbers:
$h(A_n) = n+1$ and $h(D_n) = 2(n-1)$
never coincide for $n \geq 4$
($n+1 \ne 2(n-1)$ for $n \ne 3$, and $D_3 \cong A_3$ is excluded);
the $E$-types have $h \in \{12, 18, 30\}$, each
exceeding $2(\operatorname{rank}-1)$ and $\operatorname{rank}+1$.
By Niemeier's classification
\textup{(}$1973$\textup{)},
a rank-$24$ even unimodular lattice is determined by its
root system, and root systems with uniform Coxeter number~$h$
and rank partition summing to~$24$ are classified:
there are exactly $23$ non-Leech types, each with a
distinct $(h, \text{partition})$ pair.
The Leech lattice (empty root system) is the unique lattice with
$R(\Lambda) = \varnothing$.
\end{proof}

\begin{remark}[Discrimination hierarchy]
\label{rem:discrimination-hierarchy}
The discrimination cascade
$\text{scalar} \subset \text{per-factor rank}
 \subset \text{per-factor (rank, Coxeter)}$
is strict: $0/276 < 267/276 < 276/276$.
At Level~$2$ (per-factor rank alone), the $9$
undistinguished pairs all arise from the coincidence
that $A_n$ and $D_n$ have the same rank~$n$.
Adjoining $h$ costs nothing (the
PBW filtration is already present in the bar complex)
and gives complete resolution.
The multi-channel bar invariant is stronger than
the genus-$1$ theta series ($271/276$) because
it sees the \emph{type} of each factor and the total root count.
\end{remark}

\medskip
\noindent\textbf{The oscillator--lattice factorization.}\;
The partition function of a lattice vertex
algebra decomposes canonically into oscillator and lattice contributions.

\begin{proposition}[Shadow--arithmetic factorization]
\label{prop:shadow-arithmetic-factorization}
\ClaimStatusProvedHere
\index{lattice VOA!oscillator--lattice factorization}
For an even lattice vertex algebra $V_\Lambda$ of rank~$r$,
the genus-$g$ partition function on a principally polarised
abelian variety $(\Sigma, \Omega) \in \cA_g$ factors as
\begin{equation}\label{eq:shadow-arithmetic-factorization}
 Z_g(V_\Lambda;\, \Omega)
 \;=\;
 Z_g^{\mathrm{osc}}(\Omega)
 \;\cdot\;
 \Theta_\Lambda^{(g)}(\Omega),
\end{equation}
where\textup{:}
\begin{enumerate}[label=\textup{(\roman*)}]
\item $Z_g^{\mathrm{osc}}(\Omega)
 = \bigl(\det \operatorname{Im}(\Omega)\bigr)^{r/2}
 \cdot |\eta_g(\Omega)|^{-2r}$
 depends only on the rank~$r$ and is controlled by
 the scalar shadow tower
 \textup{(}$F_g^{\mathrm{scal}}
 = r \cdot \lambda_g^{\mathrm{FP}}$,
 UNIFORM-WEIGHT\textup{)};
\item $\Theta_\Lambda^{(g)}(\Omega)
 = \sum_{v_1, \ldots, v_g \in \Lambda}
 e^{\pi i\, \sum_{j,k} \Omega_{jk}\, (v_j, v_k)}$
 is the genus-$g$ Siegel theta function of~$\Lambda$,
 a Siegel modular form of weight~$r/2$ for
 $\mathrm{Sp}(2g, \bZ)$, encoding all lattice-specific
 arithmetic.
\end{enumerate}
The scalar tower controls
$Z_g^{\mathrm{osc}}$ and is $\Lambda$-independent;
all lattice-specific information resides in
$\Theta_\Lambda^{(g)}$.
\end{proposition}

\begin{proof}
This is the standard boson--lattice factorization
for lattice vertex algebras
\textup{(}cf.~Polchinski, \emph{String Theory},
Vol.~I, Ch.~$8$\textup{)}.
The Fock space of the rank-$r$ Heisenberg
subalgebra contributes $Z_g^{\mathrm{osc}}$
(an $r$-fold product of free-boson determinants),
and the lattice sum over $\Lambda^g \subset \bR^{rg}$
contributes $\Theta_\Lambda^{(g)}$.
The independence of $Z_g^{\mathrm{osc}}$ from~$\Lambda$
is a consequence of the universality of the Heisenberg
subsector: $\kappa(H_r) = r$ for every rank-$r$
Heisenberg algebra, regardless of the lattice embedding.
\end{proof}

\begin{remark}[Genus-$2$ theta and the five collision pairs]
\label{rem:genus2-theta-niemeier}
\index{Niemeier lattice!genus-2 discrimination}
The genus-$1$ theta $\Theta_\Lambda^{(1)} \in M_{12}(\mathrm{SL}(2,\bZ))$
depends only on $|R(\Lambda)|$, leaving $5$ collision pairs
unresolved.
At genus~$2$, the Siegel modular form
$\Theta_\Lambda^{(2)} \in M_{12}(\mathrm{Sp}(4,\bZ))$
acquires a new degree of freedom: the space
$M_{12}(\mathrm{Sp}(4,\bZ))$ is three-dimensional,
spanned by the Siegel Eisenstein series~$E_{12}^{(2)}$,
the Klingen Eisenstein series~$E_{12,1}^{\mathrm{Kling}}$
\textup{(}induced from $\Delta \in S_{12}(\mathrm{SL}(2,\bZ))$
via the Klingen parabolic\textup{)},
and the unique cuspidal Hecke eigenform~$\chi_{12}$.
The decomposition
\[
 \Theta_\Lambda^{(2)}
 \;=\;
 E_{12}^{(2)}
 + c_1(\Lambda)\, E_{12,1}^{\mathrm{Kling}}
 + c_2(\Lambda)\, \chi_{12}
\]
with $c_1 = c_\Delta$ (from the diagonal restriction)
introduces one new invariant, the B\"ocherer
coefficient~$c_2(\Lambda)$.
King \textup{(}$2003$\textup{)} computed that the $24$ values
of $(c_\Delta, c_2)$ are pairwise
distinct:
the genus-$2$ Siegel theta series
is a \emph{complete} invariant for the
$24$ Niemeier lattices.
In particular, all five genus-$1$ collision pairs are resolved
at genus~$2$.

The multi-channel bar invariant
\textup{(}Proposition~\textup{\ref{prop:niemeier-multichannel}}\textup{)}
achieves complete discrimination at genus~$1$ by accessing
non-scalar sectors of the MC element.
It is strictly stronger than the genus-$1$ theta
\textup{(}$276/276$ vs.\ $271/276$\textup{)}
and orthogonal to the theta comparison:
the bar complex sees the root system type directly
through the per-factor PBW filtration,
while the theta function sees only aggregate lattice
statistics
\textup{(}representation numbers summed over all lattice
vectors\textup{)}.
\end{remark}

\begin{remark}[The MC element as complete invariant]
\label{rem:mc-complete-niemeier}
\index{Maurer--Cartan element!completeness for lattice VOAs}
The per-factor bar data is a lossy projection
of the full MC element $\Theta_{V_\Lambda}
\in \mathrm{MC}(\fg^{\mathrm{mod}}_{V_\Lambda})$.
By bar-cobar inversion
\textup{(}Theorem~\textup{\ref{thm:bar-cobar-inversion-qi}}\textup{)},
$\Theta_{V_\Lambda}$ is a \emph{complete} invariant
of the chiral algebra $V_\Lambda$: it encodes the full
OPE structure, the lattice grading by
$\gamma \in \Lambda^*/\Lambda$, the glue vectors,
and the modular completion.
The discrimination hierarchy
\[
 \text{scalar tower}
 \;\subsetneq\;
 \text{per-factor bar}
 \;\subsetneq\;
 \text{full } \Theta_{V_\Lambda}
 \;\cong\;
 V_\Lambda
\]
shows that the scalar projection loses the root system,
the per-factor projection loses the glue and inter-factor
coupling, and the full MC element loses nothing.
\end{remark}

\begin{remark}[Three comparison targets for $8$: Mukai doubling,
Humbert monodromy, Lusztig specialisation]
\label{rem:arith-shad-8}%
\index{Mukai doubling!chiral Koszul conductor}%
\index{Humbert divisor!Heegner monodromy}%
\index{Lusztig specialisation!root of unity}%
\index{three faces of eight!K3 BKM}%
The integer $8$ occurs in three comparison targets. The Vol~I result
used below is only the first one: the scalar Koszul conductor of the
Mukai-enhanced K3 row. The Humbert and Lusztig appearances are target
normalisations for a source-level comparison theorem; the scalar
conductor does not by itself identify a Humbert monodromy class or a
root-of-unity BKM category.
\begin{enumerate}[label=\textup{(\roman*)}]
\item \emph{Koszul conductor face.} For the $\Delta_5$-family of
chiral algebras parametrised by the Lorentzian lattice $2U \oplus
E_8(-1)$, the bar--cobar Koszul conductor evaluates to
\begin{equation}\label{eq:three-faces-8-conductor}
 K^{\kappa_{\mathrm{ch}}}(\cA_{\Delta_5})
 \;=\;
 2 c_+\bigl(\mathrm{Mukai}(\mathrm{K3})\bigr)
 \;=\;
 8,
\end{equation}
where $c_+(\mathrm{Mukai}(\mathrm{K3})) = 4$ is the positive-signature
rank of the Mukai lattice $H^{2*}(\mathrm{K3}, \bZ) = U^{\oplus 4}
\oplus E_8(-1)^{\oplus 2}$ and the factor~$2$ is the Mukai doubling
(the self-pairing $\langle \alpha, \alpha \rangle_{\mathrm{Muk}} =
2\,\mathrm{ch}_2 \cdot \mathrm{ch}_0 - \mathrm{ch}_1^2$).
\item \emph{Humbert monodromy target.} The automorphic comparison target
is a weight-$5$ line bundle $\cL^{\Delta_5}$ on the Heegner divisor of
Humbert discriminant~$1$, $\cH_1\subset\cA_2$, whose restriction has
monodromy order
\begin{equation}\label{eq:three-faces-8-monodromy}
 \operatorname{ord}_{H_1(\cH_1, \bZ)}
 \bigl(\cL^{\Delta_5}\big|_{\cH_1}\bigr)
 \;=\;
 8,
\end{equation}
after the automorphic normalisation, local monodromy convention, and
Bruinier Chern-class map
\cite[Prop.~5.1]{Bruinier2002} have been fixed on the same line
bundle. This is a comparison criterion, not a consequence of
\eqref{eq:three-faces-8-conductor}.
\item \emph{Lusztig specialisation target.} The quantum-group target is
a root-of-unity model at $\ell=8$ for the finite height truncation of
the $\Delta_5$ BKM datum. Lifting the scalar statement to this face
requires a PBW-finite truncation, a parity/root fixture, a comparison
functor to the chiral Koszul side, and a proof that no extra
root-of-unity blocks appear.
\end{enumerate}
The comparison problem is therefore the construction of a class
$\xi \in \mathrm{CH}^1(\cH_1)$, after the two target normalisations have
been supplied, with prescribed scalar images
\begin{equation}\label{eq:three-faces-8-reciprocity}
 \xi
 \longmapsto
 \bigl(
 K^{\kappa_{\mathrm{ch}}}(\cA_{\Delta_5}),
 \operatorname{ord}(\cL^{\Delta_5}|_{\cH_1}),
 \ell_{\mathrm{Lusztig}}
 \bigr)
 =
 (8,8,8).
\end{equation}
Only the first component is used here. The Humbert and Lusztig
components become theorems only after the displayed comparison class,
with its local monodromy and PBW hypotheses, has been constructed.
\emph{Enlarged Vol~I Theorem~C bucket.} The classical four-archetype
$\mathsf{G}/\mathsf{L}/\mathsf{C}/\mathsf{M}$ landmark ceiling
$\kappa + \kappa^! \in \{0, 13, 250/3, 98/3\}$ of
Theorem~\ref{thm:complementarity-root-datum} enlarges, on the
Lorentzian-lattice-parametric $\mathsf{B}$-family, to the canonical
five-element bucket
\begin{equation}\label{eq:theorem-c-bucket-enlarged-arithmetic}
 \kappa + \kappa^!
 \;\in\;
 \{0,\; 8,\; 13,\; 250/3,\; 98/3\},
\end{equation}
with the new entry $8$ witnessed by the scalar conductor computation
\eqref{eq:three-faces-8-conductor}. The entry is not certified here by
the Humbert or Lusztig targets. It becomes a Heegner-divisor family
only after the comparison class~\eqref{eq:three-faces-8-reciprocity}
has been constructed; until then the Vol~I statement is the Mukai
conductor value and the target normalisations it must match.
See Part~\ref{part:characteristic-datum} for the complementarity
framework and Part~\ref{part:standard-landscape} for the enlarged
landscape census.
\end{remark}

% ============================================================
\section{Proof of the shadow--spectral correspondence
(lattice case)}\label{sec:proof-lattice-case}

\begin{proof}[Proof of
Theorem~\textup{\ref{thm:shadow-spectral-correspondence}}]
\mbox{}

\emph{Step~1 (Epstein--Hecke factorization).}
The Epstein zeta $E_\Lambda(s) = \sum_{0\ne\lambda\in\Lambda}
\lvert\lambda\rvert^{-2s}$ is the Mellin transform of
$\Theta_\Lambda(\tau)-1$:
\[
 E_\Lambda^*(s)
 := (2\pi)^{-s}\,\Gamma(s)\,E_\Lambda(s)
 = \int_0^\infty \bigl(\Theta_\Lambda(iy)-1\bigr)\,y^{s-1}\,dy.
\]
The functional equation
$E_\Lambda^*(s)=E_\Lambda^*(r/2-s)$ follows from Poisson summation.

\emph{Step~2 (Modular-form decomposition).}
With $k=r/2$, decompose
$\Theta_\Lambda = c_E\,E_k + \sum_j c_j\,f_j$ in
$M_k(\Gamma)$. The Eisenstein series $E_k$ has Fourier coefficients
$\sigma_{k-1}(n)$; each $f_j$ is a Hecke eigenform.

\emph{Step~3 ($L$-function factorization).}
By the Hecke correspondence:
$\sum\sigma_{k-1}(n)n^{-s}=\zeta(s)\zeta(s-k+1)$, and
$\sum a_n(f_j)n^{-s}=L(s,f_j)$.
So $E_\Lambda(s)$ factors as in~\eqref{eq:epstein-factorization}.

\emph{Step~4 (Critical-line identification).}
The zeros of $\zeta(s)$ lie on $\mathrm{Re}(s)=\tfrac12$; of
$\zeta(s-k+1)$ on $\mathrm{Re}(s)=k-\tfrac12$; of $L(s,f_j)$
on $\mathrm{Re}(s)=\tfrac{k}{2}$ (assuming GRH for~$f_j$;
Deligne's theorem gives $|a_p(f_j)| \leq 2p^{(k-1)/2}$
but does not imply the Riemann hypothesis for $L(s,f_j)$).
The distinct values are
$\{\tfrac12,\,\tfrac{k}{2},\,k-\tfrac12\}$
(for $k\ge12$; fewer for small~$k$).

\emph{Step~5 (The arithmetic sieve).}
Each degree detects exactly one
arithmetic object in the Hecke decomposition of~$\Theta_\Lambda$:
\begin{center}
\small
\renewcommand{\arraystretch}{1.2}
\begin{tabular}{cll}
\toprule
\textbf{Degree} & \textbf{Detects} & \textbf{$L$-function} \\
\midrule
$2$ & curvature $\kappa$ (total spectral weight)
 & $\zeta(s)$ \\
$3$ & Eisenstein shift (root structure)
 & $\zeta(s - k + 1)$ \\
$3+j$ ($j\ge1$) & $j$-th cusp eigenform $f_j \in S_k$
 & $L(s, f_j)$ \\
\bottomrule
\end{tabular}
\end{center}
The obstruction $o_{3+j+1}$ vanishes iff
$\dim S_k < j{+}1$, i.e., iff there is no $(j{+}1)$-th cusp
form. The tower terminates at degree $d = 3 + \dim S_k$ for
$k \ge 4$. Hence for even unimodular lattices of rank
$r \ge 8$:
\begin{equation}\label{eq:depth-cusp-formula}
 d(V_\Lambda) = 3 + \dim S_{r/2}.
\end{equation}
The number of critical lines is $d - 1 = 2 + \dim S_{r/2}$:
two from the Eisenstein $\zeta(s)\,\zeta(s-k+1)$, and one
for each cusp eigenform.
\end{proof}

\begin{remark}[Chriss--Ginzburg identification]
\label{rem:shadow-spectral-CG}
For $V_\Lambda$, the graph amplitudes factor through the Hecke algebra
of $\mathrm{SL}(2,\bZ)$ acting on $M_{r/2}$:
$\ell_\Gamma^{(g)}(V_\Lambda) \in
\mathrm{Image}(T_n \colon M_{r/2} \to M_{r/2})$.
The shadow depth counts Hecke eigenspaces receiving nonzero
graph-sum contributions.
\end{remark}

\begin{proposition}[Leading-order Hecke identification]
\label{prop:leading-hecke-identification}
\ClaimStatusProvedElsewhere
At leading order in graph complexity, the shadow amplitude at degree~$r$
is the chain-graph trace
\begin{equation}\label{eq:chain-trace}
 \operatorname{Sh}_r
 = P^{r-2}\,\operatorname{tr}(m_2^{r-1})
 + (\text{graphs with loops}),
\end{equation}
where $P = H_\cA^{-1}$ is the complementarity propagator and
$m_2$ is the binary $A_\infty$ product on~$H^*(\barB(\cA))$.
For $V_\Lambda$ with Hecke decomposition
$\Theta_\Lambda = c_0 + \sum_i c_i f_i$, the trace decomposes:
\[
 \operatorname{tr}(m_2^{r-1})
 = \sum_i c_i^r\,\lVert f_i \rVert^{r-1}_{\mathrm{Pet}}
 \cdot (\text{combinatorial factor}).
\]
The obstruction $o_{r+1}$ vanishes iff no Hecke eigenform of
weight~$r/2$ contributes a new summand. For Virasoro, the
single-generator trace $\operatorname{tr}(m_2^{r-1})=2^{r-1}$
and the chain-graph formula reproduces
Theorem~\textup{\ref{thm:shadow-tower-asymptotics}} with effective
coupling $\lambda_{\mathrm{eff}}=-3P$.
\end{proposition}

\begin{proposition}[Exactness of the Hecke identification]
\label{prop:hecke-all-orders}
\label{prop:loop-corrections}
\ClaimStatusProvedHere
\index{Hecke decomposition!loop corrections}
The critical-line count $d_{\mathrm{arith}}$ is an \emph{exact}
invariant:
\begin{enumerate}[label=\textup{(\roman*)}]
\item On the Virasoro primary line
 \textup{(}$\dim H^*(\barB) = 1$\textup{)}: the chain graph is the
 \emph{only} connected graph at each degree \textup{(}a
 one-dimensional space admits no loops\textup{)}. The chain-graph
 formula is exact \textup{(}verified symbolically through $r = 12$\textup{)}.
\item For lattice VOAs with multi-dimensional bar cohomology: loop
 graphs contribute trace corrections
 $\ell^{\mathrm{loop}}_\Gamma =
 \operatorname{tr}_{H^*(\barB)}(m_k \circ P \circ \cdots)$
 that are multilinear in the Fourier coefficients of
 $\Theta_\Lambda = \sum c_j f_j$. The trace decomposes in the
 \emph{same} Hecke basis $\{f_j\}$ as the tree amplitudes:
 loops modify the coefficients~$c_j$ but \emph{not} the eigenform
 set. Therefore $d_{\mathrm{arith}}$ is unchanged by loop
 corrections.
\end{enumerate}
\end{proposition}

\begin{proof}
For~(i): on a $1$-dimensional generator space $\bC\cdot\bar{T}$,
every connected degree-$r$ graph is a chain (no room for loops).
For~(ii): the lattice sum
$\sum_\lambda(\text{vertex labels around loop})$ involves
$r_\Lambda(n)=\sum c_j a_n(f_j)$ at each vertex,
where $a_n(f_j)$ are Fourier coefficients of the eigenform~$f_j$.
Products of such sums decompose in the tensor powers
$f_j^{\otimes p}$, which live in $\mathrm{Sym}^p M_k$. Since
$\mathrm{Sym}^p M_k$ is spanned by products of the
\emph{same} eigenforms $\{f_j\}$, no new eigenspace appears.
\end{proof}

% ============================================================
\section{The period--shadow dictionary}
\label{sec:period-shadow-dictionary}

The arithmetic sieve (Step~5 of
Theorem~\ref{thm:shadow-spectral-correspondence}) assigns to each
shadow degree a specific arithmetic period. The complete correspondence for even unimodular lattice
VOAs of rank~$r$, where $k=r/2$ and
$\Theta_\Lambda = c_E\,E_k + \sum_{j=1}^{g_k} c_j\,f_j$ is
the Hecke decomposition in $M_k(\mathrm{SL}(2,\bZ))$
with $g_k = \dim S_k$.

\begin{proposition}[Period--shadow dictionary]
\label{prop:period-shadow-dictionary}
\ClaimStatusProvedHere
For an even unimodular lattice VOA $V_\Lambda$ of rank~$r\ge8$,
the shadow tower $\Theta_{V_\Lambda}^{\le r}$ decomposes
the Epstein zeta function $E_\Lambda(s)$ into its spectral
constituents, one degree at a time:
\begin{center}
\renewcommand{\arraystretch}{1.35}
\begin{tabular}{cllll}
\toprule
\textbf{Degree}
 & \textbf{Shadow}
 & \textbf{Modular source}
 & \textbf{$L$-function}
 & \textbf{Period type} \\
\midrule
$2$
 & $\kappa$ \textup{(curvature)}
 & constant term of $\Theta_\Lambda$
 & $\zeta(s)$
 & Riemann \\
$3$
 & $C$ \textup{(cubic)}
 & Eisenstein part $E_k$
 & $\zeta(s)\,\zeta(s{-}k{+}1)$
 & Dedekind \\
$3+j$\, ($j\ge1$)
 & $\operatorname{Sh}_{3+j}$
 & $j$-th cusp eigenform $f_j \in S_k$
 & $L(s, f_j)$
 & Hecke \\
\bottomrule
\end{tabular}
\end{center}
The tower terminates at degree $d = 3 + g_k$, and
\begin{equation}\label{eq:depth-cusp-full}
 d(V_\Lambda)
 = 3 + \dim S_{r/2}(\mathrm{SL}(2,\bZ)).
\end{equation}
The number of independent $L$-functions in the factorization of
$E_\Lambda(s)$ equals $d-1 = 2 + g_k$: two from the Eisenstein
product $\zeta(s)\,\zeta(s{-}k{+}1)$, and one for each cusp
eigenform.
\end{proposition}

\begin{proof}
Combine Steps~1--5 of the proof of
Theorem~\ref{thm:shadow-spectral-correspondence}
with Proposition~\ref{prop:loop-corrections}.
At degree~$2$, the curvature $\kappa = \operatorname{rank}(\Lambda)$
carries only the constant term of $\Theta_\Lambda$, whose Mellin
transform is $\zeta(s)$ (the spectral zeta of the sewing operator).
At degree~$3$, the cubic shadow~$C$ detects the root structure
encoded in the divisor sums $\sigma_{k-1}(n)$, whose generating
series is $\zeta(s)\zeta(s{-}k{+}1)$. At degree $3+j$
($j\ge1$), the obstruction $o_{3+j}$ is nonzero iff
$\dim S_k \ge j$, and the shadow amplitude
$\operatorname{Sh}_{3+j}$ projects onto the $j$-th Hecke
eigenform~$f_j$ with generating series $L(s,f_j)$
(Proposition~\ref{prop:leading-hecke-identification}).
By Proposition~\ref{prop:loop-corrections}, loop corrections
do not introduce new eigenforms, so the correspondence is exact.
\end{proof}

\begin{remark}[The depth formula for non-unimodular lattices]
\label{rem:depth-non-unimodular}
For a general even lattice of rank~$r$ and level~$N$, the
theta function $\Theta_\Lambda \in M_{r/2}(\Gamma_0(N), \chi)$
and the depth formula generalizes to
\[
 d(V_\Lambda)
 = 3 + \dim S_{r/2}(\Gamma_0(N), \chi),
\]
where $\chi$ is the character determined by the discriminant
of~$\Lambda$. The dictionary table extends with the
$j$-th newform $f_j \in S_{r/2}^{\mathrm{new}}(\Gamma_0(N),\chi)$
contributing $L(s,f_j)$ at degree $3+j$.
\end{remark}

\subsection{Five worked examples}
\label{subsec:period-shadow-examples}

\begin{computation}[$V_{\bZ}$: the Gaussian archetype]
\label{comp:period-shadow-vz}
\ClaimStatusProvedHere
The rank-$1$ lattice has
$\Theta_{\bZ}(\tau)=\theta_3(2\tau)=\sum_{n\in\bZ}q^{n^2}$, a
modular form of weight~$\frac12$ and level~$4$. The space
$S_{1/2}(\Gamma_0(4))=0$ contains no cusp forms, so
the constrained Epstein zeta is purely Eisenstein:
\[
 \varepsilon^1_s
 = 2\,\zeta(2s).
\]
Depth~$2$ (class~$\mathbf{G}$). One critical line at
$\mathrm{Re}(s)=\frac14$. The tower terminates at
the curvature: $\kappa = 1$, and $o_3 = 0$. No $L$-function
content beyond the Riemann zeta.
\end{computation}

\begin{computation}[$V_{E_8}$: the Lie/tree archetype]
\label{comp:period-shadow-ve8}
\ClaimStatusProvedHere
The $E_8$ root lattice has
$\Theta_{E_8}(\tau)=E_4(\tau)=1+240\sum\sigma_3(n)q^n$,
a pure Eisenstein series in $M_4(\mathrm{SL}(2,\bZ))$. Since
$S_4(\mathrm{SL}(2,\bZ))=0$, there are no cusp forms and the
Epstein zeta factors as
(Theorem~\textup{\ref{thm:e8-epstein}}):
\[
 \varepsilon^8_s
 = 240\cdot 2^{-s}\,\zeta(s)\,\zeta(s-3).
\]
Depth~$3$ (class~$\mathbf{L}$). Two critical lines at
$\mathrm{Re}(s)=\frac12$ and $\mathrm{Re}(s)=\frac72$.
The degree-$3$ cubic shadow detects the root structure via the
divisor sums $\sigma_3(n)$, and $o_4 = 0$: the \emph{absence} of
cusp forms in weight~$4$ means that degree~$\ge4$ shadows carry
no $L$-function content. The $240$ root vectors of $E_8$ are
entirely visible to the Eisenstein projection.
\end{computation}

\begin{computation}[$V_{\mathrm{Leech}}$: the Ramanujan archetype]
\label{comp:period-shadow-leech}
\ClaimStatusProvedHere
The Leech lattice has rank~$24$ and
$\Theta_{\mathrm{Leech}} = E_{12} - \frac{65520}{691}\,\Delta_{12}$
(Proposition~\textup{\ref{prop:leech-epstein}}),
where $\Delta_{12} = \sum \tau(n)\,q^n$ is the Ramanujan cusp form,
the unique normalized eigenform spanning
$S_{12}(\mathrm{SL}(2,\bZ))$. The Epstein zeta factors into
three $L$-functions:
\begin{equation}\label{eq:leech-three-L}
 E_{\mathrm{Leech}}(s)
 = C_E(s)\,\zeta(s)\,\zeta(s{-}11)
 \;-\;
 \frac{65520}{691}\,C_\Delta(s)\,L(s,\Delta_{12}).
\end{equation}
Depth~$4$ (class~$\mathbf{C}$). Three critical lines:
\begin{center}
\renewcommand{\arraystretch}{1.2}
\begin{tabular}{clcl}
\toprule
\textbf{Degree} & \textbf{Source} & \textbf{Critical line} & \textbf{Period} \\
\midrule
$2$ & curvature $\kappa = 24$
 & $\mathrm{Re}(s) = \tfrac12$ & $\zeta(s)$ \\
$3$ & Eisenstein $E_{12}$
 & $\mathrm{Re}(s) = \tfrac{23}{2}$ & $\zeta(s{-}11)$ \\
$4$ & Ramanujan $\Delta_{12}$
 & $\mathrm{Re}(s) = 6$ & $L(s,\Delta_{12})$ \\
\bottomrule
\end{tabular}
\end{center}
The degree-$4$ shadow \emph{is} the cusp form period: the quartic
resonance class $Q(V_{\mathrm{Leech}})$ extracts $L(s,\Delta_{12})$
from the Epstein zeta function. This is the tower's arithmetic content at degree~$4$. The coefficient
$\tau(n)$ is governed by the Ramanujan conjecture
$\lvert\tau(p)\rvert \le 2p^{11/2}$ (Deligne's theorem),
which places the zeros of $L(s,\Delta_{12})$ on
$\mathrm{Re}(s) = 6$. The tower terminates at degree~$4$ because
$\dim S_{12} = 1$: a single cusp form exhausts the cuspidal
spectrum, and $o_5 = 0$.

The Leech lattice is distinguished among unimodular lattices:
the coefficient
$-65520/691$ is the \emph{unique} ratio making the $q^1$-coefficient
of $\Theta_{\mathrm{Leech}}$ vanish (no vectors of squared
length~$2$). This arithmetic accident (the vanishing of the
shortest nonzero shell) is the lattice-theoretic origin of the
contact shadow class.
\end{computation}

\begin{computation}[$D_{16}^+ \oplus D_{16}^+$ vs.\ $E_8^3$: same depth,
different coefficient]
\label{comp:period-shadow-rank24-comparison}
\ClaimStatusProvedHere
Both $V_{D_{16}^+ \oplus D_{16}^+}$ and $V_{E_8^3}$
have rank~$24$ and theta functions in $M_{12}(\mathrm{SL}(2,\bZ))$,
so their depth is $3 + \dim S_{12} = 4$ in both cases.
They share the same Hecke eigenform basis $\{E_{12}, \Delta_{12}\}$
but differ in the Eisenstein--cuspidal decomposition
coefficients~$c_E, c_\Delta$. Since $E_8^3$ has $3\times 240 = 720$
root vectors contributing $720\,q^2 + \cdots$ while
$D_{16}^+\oplus D_{16}^+$ has $2\times 480 = 960$
vectors of squared length~$2$, their quartic shadows $Q$ carry distinct
arithmetic data from \emph{the same} $L$-function
$L(s,\Delta_{12})$ at \emph{different} cuspidal weights.
The period--shadow dictionary assigns identical period types but
distinguishable shadow amplitudes: the dictionary is functorial
in the lattice.
\end{computation}

\begin{computation}[$V_{\bZ^2}$ and $V_{A_2}$: Dedekind zeta at depth~$2$]
\label{comp:period-shadow-rank2}
\ClaimStatusProvedHere
For the rank-$2$ square lattice $\bZ^2$, the theta function is
$\Theta_{\bZ^2} = \theta_3^2$, a modular form in
$M_1(\Gamma_0(4))$. At weight~$1$, the Eisenstein--cuspidal
decomposition gives
$E_{\bZ^2}(s) = 4\,\zeta(s)\,L(s,\chi_{-4})
= 4\,\zeta_{\bQ(i)}(s)$
(Proposition~\textup{\ref{prop:z2-epstein}}): the Epstein zeta is
a Dedekind zeta function. Both $\zeta(s)$ and $L(s,\chi_{-4})$
have zeros on $\mathrm{Re}(s) = \frac12$: a single critical line.
Depth~$2$ (class~$\mathbf{G}$).

For the hexagonal lattice $A_2$, $\Theta_{A_2} \in
M_1(\Gamma_0(3),\chi_{-3})$ and
$E_{A_2}(s) = 6\,\zeta_{\bQ(\omega)}(s)$ with
$\omega = e^{2\pi i/3}$.
Again depth~$2$, one critical line. These rank-$2$ examples
illustrate the boundary between Gaussian (no root structure)
and Lie (root structure visible at degree~$3$): the weight is too
low for the Eisenstein product to produce a second critical line.
\end{computation}

\begin{theorem}[Spectral decomposition principle]
\label{thm:spectral-decomposition-principle}
\ClaimStatusProvedHere
For a lattice vertex algebra $V_\Lambda$ with Hecke decomposition
$\Theta_\Lambda = c_E\,E_k + \sum_{j=1}^{g_k} c_j\,f_j$,
the shadow tower
$\Theta_{V_\Lambda}^{\le 2} \hookrightarrow
\Theta_{V_\Lambda}^{\le 3} \hookrightarrow
\cdots \hookrightarrow
\Theta_{V_\Lambda}^{\le 3+g_k} = \Theta_{V_\Lambda}$
is a spectral filtration: the degree-$r$ extension
$\Theta^{\le r}/\Theta^{\le r-1}$ isolates the $L$-function
$L(s,f_{r-3})$ of a single Hecke eigenform. The depth
$d(V_\Lambda)$ counts the number of independent $L$-functions
in the Epstein zeta factorization:
\begin{equation}\label{eq:depth-L-count}
 d(V_\Lambda)
 = 1
 + \underbrace{2}_{\text{Eisenstein}}
 + \underbrace{g_k}_{\text{cuspidal}}
 = 3 + \dim S_{r/2}.
\end{equation}
\end{theorem}

\begin{proof}
The depth formula~\eqref{eq:depth-cusp-full} gives
$d = 3 + g_k$. It remains to verify the spectral filtration
structure. At degree~$2$, the curvature $\kappa$ extracts only
the constant term, which is the residue of $E_\Lambda(s)$ at
$s = k$ and involves $\zeta(s)$ alone. At degree~$3$, the
cubic shadow~$C$ detects $\sigma_{k-1}(n)$, the Eisenstein
coefficients, contributing $\zeta(s)\zeta(s{-}k{+}1)$. At
degree $3+j$, the obstruction $o_{3+j}(V_\Lambda) \ne 0$ iff
$j \le g_k$, and the leading-order amplitude
(Proposition~\ref{prop:leading-hecke-identification}) projects
onto~$f_j$. Since the Hecke eigenforms $\{E_k, f_1,\ldots,f_{g_k}\}$
form a basis for $M_k$, the filtration exhausts all spectral
content and $o_{3+g_k+1} = 0$.
\end{proof}

\begin{remark}[Homotopy becomes arithmetic]
\label{rem:homotopy-becomes-arithmetic}
\index{period--shadow dictionary}%
The period--shadow dictionary concretely manifests
the principle that homotopy-algebraic data (the $A_\infty$
structure on bar cohomology) and arithmetic data (the $L$-functions
of Hecke eigenforms) are unified by the single Maurer--Cartan
element $\Theta_\cA$. For the Leech lattice, this unification
is maximally visible: the quartic resonance class
$Q(V_{\mathrm{Leech}})$, defined by the homotopy-theoretic
obstruction $o_4$ in the cyclic deformation complex
$\mathrm{Def}_{\mathrm{cyc}}^{\mathrm{mod}}(\cA)$, is
\emph{equal} to the Hecke $L$-function period $L(s,\Delta_{12})$
of the Ramanujan cusp form under the Mellin correspondence.
The vanishing of~$o_5$ (the termination of the shadow
tower) is equivalent to the one-dimensionality of $S_{12}$,
Ramanujan's observation that the cusp form $\Delta$ is the unique
eigenform of weight~$12$.

The dictionary identifies the shadow depth
classification
$\mathbf{G}/\mathbf{L}/\mathbf{C}/\mathbf{F}_d/\mathbf{M}$
with the spectral complexity of the theta function: Gaussian
means no Eisenstein content beyond the constant term, Lie/tree
means pure Eisenstein, contact means one cusp form, and the
$\mathbf{F}_d$ classes correspond to lattices whose theta
functions probe the cuspidal spectrum at higher levels.
The first lattice of depth~$5$ would require a weight~$k$
with $\dim S_k \ge 2$, which occurs at $k=24$
(rank~$48$, $\dim S_{24} = 2$).
\end{remark}

% ============================================================
\section{The cusp-singularity theorem}
\label{sec:cusp-singularity}

\begin{definition}[Cusp growth rate]
\label{def:cusp-growth-rate}
For a vertex algebra~$\cA$ at central charge~$c$, the
\emph{cusp growth rate} $\alpha(\cA)\ge0$ is defined by the asymptotic
\[
 \widehat{Z}^c_\cA(iy)
 \sim y^{c/2}\,e^{\alpha(\cA)\cdot y}
 \qquad\text{as } y\to\infty.
\]
\end{definition}

\begin{proposition}[Growth-rate dictionary]
\label{prop:growth-rate-dictionary}
\ClaimStatusProvedHere
\mbox{}
\begin{enumerate}[label=\textup{(\roman*)}]
\item \textup{(Lattice, class~G):}
 $\alpha(V_\Lambda)=0$; the growth is polynomial.
\item \textup{(Free fields, class~C):}
 $\alpha(\beta\gamma)=\pi/6$, from the
 factor $1/\eta(\tau)^2$ in the partition function.
\item \textup{(Virasoro, class~M):}
 $\alpha(\mathrm{Vir}_c)=\pi c/12$, from the vacuum energy
 $q^{-c/24}$.
\end{enumerate}
\end{proposition}

\begin{proof}
For lattice algebras,
$\widehat{Z}^c=y^{c/2}\lvert\Theta_\Lambda\rvert^2$ and
$\Theta_\Lambda(iy)\to1$ as $y\to\infty$. For~$\beta\gamma$,
$\widehat{Z}^1=y^{1/2}\theta_3^2/\eta^2$ and
$\eta(iy)\sim e^{-\pi y/12}$. For Virasoro, the vacuum character
$\chi_0(q)=q^{-c/24}/\eta(\tau)\cdot(\text{null corrections})$
dominates.
\end{proof}

\begin{theorem}[$\beta\gamma$--Virasoro rate coincidence]
\label{thm:bg-vir-coincidence}
\ClaimStatusProvedElsewhere
The cusp growth rate of the $\beta\gamma$ system at $c=2$
equals the Virasoro rate at $c=2$:
\[
 \alpha(\beta\gamma)
 = \frac{\pi}{6}
 = \frac{\pi\cdot 2}{12}
 = \alpha(\mathrm{Vir}_{c=2}).
\]
The $\beta\gamma$ tower terminates at depth~$4$ because the OPE is
free; the Virasoro tower does not terminate because the OPE is
nonlinear.
\end{theorem}

\subsection{The self-referentiality criterion}

\begin{proposition}[Self-referentiality criterion; \ClaimStatusProvedHere]
\label{prop:self-referentiality-criterion}
\label{prop:free-interacting-dichotomy}
The OPE nonlinearity type controls the shadow depth $d(\cA)$:
\begin{enumerate}[label=\textup{(\roman*)}]
\item \textup{(Abelian, Tier~$1$):}
 If the OPE of all strong generators produces only vacuum and
 derivatives of inputs, then $d(\cA) \leq 2$.
\item \textup{(Controlled, Tier~$2$, proved families):}
 In the lattice-VOA and rank-one $\beta\gamma$ families, if the OPE
 creates new fields but no strong generator~$a$ satisfies
 $a \in a_{(k)}a$ for any~$k$ \textup{(}no direct self-loop\textup{)},
 then $d(\cA) < \infty$. For lattice VOAs this follows from the
 arithmetic sieve \textup{(}$d = 3 + \dim S_{r/2}$,
 Theorem~\textup{\ref{thm:shadow-spectral-correspondence})}; for the
 rank-one $\beta\gamma$ system from abelian rigidity
 \textup{(}Corollary~\textup{\ref{cor:nms-betagamma-mu-vanishing})}.
 Across these proved families the depth is \emph{unbounded}: it grows
 with the lattice rank or with the complexity of the root system.
\item \textup{(Self-referential, Tier~$3$):}
 If some strong generator~$a$ satisfies $a \in a_{(k)}a$ for some~$k$,
 then $d(\cA) = \infty$.
\end{enumerate}
\end{proposition}

\begin{proof}
For~(i), the bar differential is strictly coassociative ($m_k = 0$
for $k \geq 3$ in the transferred $A_\infty$-structure), so all
obstruction sources vanish.

For~(ii), the finiteness of depth is established in the named
families. For lattice VOAs: the depth equals $3 + \dim S_{r/2}$
(Theorem~\ref{thm:shadow-spectral-correspondence}), which is
finite because $M_{r/2}(\Gamma)$ is finite-dimensional.
For rank-one $\beta\gamma$: the quintic obstruction vanishes by
abelian rigidity, giving $d = 4$.

For~(iii), the direct self-loop $a \in a_{(k)}a$ generates an
infinite $A_\infty$ cascade: the shadow obstruction
$o_{r+1} = \{o_r, \Theta^{\leq r}\}_H$ is nonzero for all~$r$
because the self-loop prevents stabilization.
\end{proof}

\begin{remark}[General controlled regime remains open]
The general statement that absence of direct self-loops forces finite
shadow depth is still Conjecture~\ref{conj:general-finiteness}. The
proposition records only the families where this controlled regime has
been proved.
\end{remark}

\begin{conjecture}[General finiteness conjecture]
\label{conj:general-finiteness}
\ClaimStatusConjectured
\index{shadow depth!general finiteness conjecture}
If~$\cA$ is strongly finitely generated and no strong generator~$a$
satisfies $a \in a_{(k)}a$ for any~$k$, then $d(\cA) < \infty$.
\end{conjecture}

The conjecture holds for lattice VOAs (via the Hecke
decomposition, $d = 3 + \dim S_{r/2}$), for $\beta\gamma$ (via
rank-one abelian rigidity), and for all other standard families
(case-by-case). The conjecture is equivalent to the statement that
the $A_\infty$-Massey products of a non-self-referential
transferred algebra are eventually cohomologically trivial, a
condition on the rational homotopy type of~$H^*(\barB(\cA))$.
A spectral-radius approach fails: for affine~$\mathfrak{sl}_2$,
the recursion operator has nonzero spectral radius (from the cycle
$E^\alpha \to J \to E^\alpha$), yet depth~$= 3$ because the
Jacobi identity cancels the amplitudes at degree~$\geq 4$.

\begin{corollary}[Conformal vector implies infinite shadow depth]
\label{cor:conformal-vector-infinite-depth}
\ClaimStatusProvedHere
If the stress tensor~$T$ of~$\cA$ is a strong generator
\textup{(}not a composite like the Sugawara construction\textup{)},
then $d(\cA) = \infty$.
\end{corollary}

\begin{proof}
$T_{(1)}T = 2T$ holds universally (from the Virasoro OPE).
If $T$ is a strong generator, this is a direct self-loop on a
strong generator, so
Proposition~\ref{prop:self-referentiality-criterion}(iii) applies.
\end{proof}

This covers all $\mathcal{W}$-algebras, all $\mathcal{N} = 1, 2, 3, 4$
superconformal algebras, and every VOA whose conformal vector
is a primitive generator. It excludes Heisenberg
(where, at nonzero level, $T=(2\kappa)^{-1}
\mathord{:}\alpha\alpha\mathord{:}$ is composite; at $\kappa=0$
there is no inverse-metric Sugawara conformal vector) and affine
algebras (where $T$ is the Sugawara
composite $T = \tfrac{1}{2(k+h^\vee)}\sum\mathord{:}J^a J^a
\mathord{:}$).

\begin{theorem}[Shadow obstruction tower leading asymptotics]
\label{thm:shadow-tower-asymptotics}
\ClaimStatusProvedHere
For the Virasoro algebra at central charge~$c\neq0$, the Laurent
expansion at $c^{-1}=0$ of the shadow coefficient
$S_r$ in $\operatorname{Sh}_r = S_r\,x^r$ satisfies, for $r \geq 4$:
\begin{equation}\label{eq:shadow-leading}
 S_r = \frac{2}{r}\,(-3)^{r-4}\,P^{r-2}
 + O(c^{-(r-1)}),
 \qquad P = \frac{2}{c}.
\end{equation}
The displayed monomial is the exact leading Laurent coefficient:
$\lim_{c\to\infty} S_r \cdot r / [(-3)^{r-4} P^{r-2}] = 2$
holds for all $r \geq 4$.
\end{theorem}

\begin{proof}
The master equation $\nabla_H(\operatorname{Sh}_r) + \mathfrak{o}^{(r)} = 0$
gives, at leading order in $1/c$, the recursion
$S_r = -3(r{-}1)\,P\,S_{r-1}/r$
(from the dominant bracket
$\{\mathfrak{C}, \operatorname{Sh}_{r-1}\}_H$ with
$\mathfrak{C} = 2x^3$, scalar propagator $P=P_0=2/c$, and
$\nabla_H^{-1}(x^r) = x^r/(2r)$). Iterating with initial condition
$S_4 = P^2/2$:
\[
 S_r = S_4 \cdot \prod_{k=5}^{r} \frac{-3(k{-}1)P}{k}
 = \frac{P^2}{2} \cdot (-3P)^{r-4} \cdot \frac{4}{r}
 = \frac{2}{r}\,(-3)^{r-4}\,P^{r-2}.\qedhere
\]
\end{proof}

\begin{corollary}[Rigorous Virasoro infinite shadow depth]
\label{cor:rigorous-infinite-depth}
\ClaimStatusProvedHere
For the universal Virasoro algebra at $c\neq0$,
$d(\mathrm{Vir}_c)=\infty$. In the depth decomposition of
Theorem~\ref{thm:depth-decomposition}, this infinite contribution is
the algebraic summand $d_{\mathrm{alg}}$.
\end{corollary}

\begin{proof}
$T$ is the strong generator of the universal Virasoro algebra and
$T_{(1)}T=2T$. Thus $T$ has a direct self-loop, so
Proposition~\ref{prop:self-referentiality-criterion}\textup{(iii)}
gives $d(\mathrm{Vir}_c)=\infty$. This source is not a
Roelcke--Selberg Hecke eigenspace of a lattice theta series; it is the
self-OPE of the primitive stress tensor, hence belongs to the
homotopy-defect summand $d_{\mathrm{alg}}$ after the decomposition
theorem is invoked. The Laurent calculation~\eqref{eq:shadow-leading}
is the independent large-$c$ shadow of the same mechanism: its leading
term is nonzero whenever $c\neq0$. On the nonsingular
Zamolodchikov surface
$c(5c+22)\neq0$, the canonical low-degree coefficients
$S_2=c/2$, $S_3=2$, $S_4=10/[c(5c+22)]$, and
$S_5=-48/[c^2(5c+22)]$ agree with
Proposition~\ref{prop:virasoro-shadow-canonical}.
\end{proof}

\begin{remark}[The effective coupling constant]
\label{rem:effective-coupling}
The recursion $S_r \propto (-3P)\,S_{r-1}$ identifies
$\lambda_{\mathrm{eff}} = -3P = -6/c$. The shadow generating function
$G(t) = -\log(1 + 3Px\,t)$ has convergence radius
$|t| < |c|/3 = |\lambda_{\mathrm{eff}}|^{-1}$.
\end{remark}

% ============================================================
\section{Non-lattice theories}\label{sec:non-lattice-theories}

The non-lattice regime is where the tower
$\Theta_\cA^{\le r}$ detects arithmetic data beyond
the theta-function decomposition.

\subsection{The $\beta\gamma$ system}

\begin{proposition}[$\beta\gamma$ primary-counting function]
\label{prop:bg-primary-counting}
\ClaimStatusProvedElsewhere
The primary-counting function of the $\beta\gamma$ system,
viewed as a $\mathrm{U}(1)^1$ theory at the diagonal point, is
\begin{equation}\label{eq:bg-Zhat}
 \widehat{Z}^{\beta\gamma}(\tau)
 = y^{1/2}\,\frac{\theta_3(\tau)^2}{\eta(\tau)^2}.
\end{equation}
This differs from the rank-$1$ Narain function
$\widehat{Z}^{V_{\bZ}}=y^{1/2}\theta_3^2$ by the factor~$1/\eta^2$.
\end{proposition}

\begin{theorem}[Refined shadow--spectral correspondence]
\label{thm:refined-shadow-spectral}
\ClaimStatusProvedElsewhere
For a vertex algebra $\cA$ with $\mathrm{U}(1)^c$ symmetry:
\begin{equation}\label{eq:refined-inequality}
 d(\cA) \ge 1 + (\text{number of critical lines of
 $\varepsilon^{c,\mathrm{U}(1)}_s$}).
\end{equation}
Equality holds for lattice vertex algebras. For the $\beta\gamma$
system, $d=4$ and the $\mathrm{U}(1)$ constrained Epstein has
$1$~critical line, giving a gap of~$2$. The gap originates from
the weight-changing deformations at shadow degrees~$3$ and~$4$, which
are invisible to the $\mathrm{U}(1)$ spectral decomposition.
\end{theorem}

\subsection{Virasoro at arbitrary central charge}

The Virasoro shadow tower provides an
infinite sequence of moment constraints on the spectral data:
\begin{align}
 \alpha_2 &= c/2, \label{eq:vir-alpha2}\\
 \alpha_3 &= 0 \quad\text{(cubic vanishing from diagonal metric)},
 \label{eq:vir-alpha3}\\
 \alpha_4 &= \frac{10}{c(5c+22)} \quad\text{(quartic contact
 $Q^{\mathrm{contact}}_{\mathrm{Vir}}$)}, \label{eq:vir-alpha4}\\
 \alpha_5 &\ne 0 \quad\text{(quintic forced)},
 \label{eq:vir-alpha5}\\
 &\;\;\vdots \notag
\end{align}

\begin{proposition}[Ising model: $d_{\mathrm{arith}} = 0$]
\label{prop:ising-d-arith}
\ClaimStatusProvedHere
For the Ising model $\cM(3,4)$ at $c = 1/2$, the constrained
Epstein zeta has only two terms:
\[
 \varepsilon^{1/2}_s = 2^{-s} + 4^s.
\]
Its zeros are at $s = -i\pi(2k{+}1)/(3\log 2)$ for $k\in\bZ$, all on
$\mathrm{Re}(s) = 0$. No zeros lie on any critical line
$\mathrm{Re}(s) = \sigma > 0$, so $d_{\mathrm{arith}}
(\mathrm{Ising}) = 0$.
The full shadow depth is $d = \infty$ (from $T_{(1)}T = 2T$),
confirming $d_{\mathrm{alg}} = \infty$ while
$d_{\mathrm{arith}} = 0$: the infinite shadow tower at a
rational central charge detects the \emph{algebra}, not the
\emph{arithmetic} of the partition function.
\end{proposition}

\begin{proof}
The scalar primaries are $\Delta = 1$ (energy, $h=\bar h = 1/2$)
and $\Delta = 1/8$ (spin, $h = \bar h = 1/16$), giving
$\varepsilon^{1/2}_s = (2\cdot 1)^{-s}+(2\cdot\tfrac18)^{-s}
= 2^{-s}+4^s$. Setting $2^{-s}+4^s = 0$:
$2^{-s}=-2^{2s}$, hence $2^{-3s}=-1$, so
$s = -i\pi(2k{+}1)/(3\log 2)$.
\end{proof}

\begin{remark}[The Ising arithmetic paradox]%
\label{rem:ising-arithmetic-paradox}%
\ClaimStatusProvedHere%
\index{Ising model!arithmetic paradox}%
\index{fusion ring!invisible to depth decomposition}%
\index{VVMF!Ising model}%
Proposition~\ref{prop:ising-d-arith} shows $d_{\mathrm{arith}}(\mathrm{Ising}) = 0$
and $d_{\mathrm{alg}}(\mathrm{Ising}) = \infty$, yet the Ising model
has deep arithmetic structure that this decomposition
does not see.

\begin{enumerate}[label=\textup{(\roman*)}]
\item \textup{(Fusion ring arithmetic.)}
The quantum dimension $d_\sigma = \sqrt{2}$ is irrational.
The $F$-matrices of $\cM(3,4)$ are algebraic numbers in
$\bQ(\sqrt{2})$, and the fusion ring
$\bZ[\mathbf{1}, \varepsilon, \sigma]$ with
$\sigma^2 = \mathbf{1} + \varepsilon$ encodes the
Galois group $\mathrm{Gal}(\bQ(\sqrt{2})/\bQ)$.

\item \textup{(VVMF structure.)}
The three characters $\chi_I, \chi_\varepsilon, \chi_\sigma$ form
a $3$-dimensional vector-valued modular form for a
representation of $\mathrm{SL}(2,\bZ)$ factoring through a
finite image in $\mathrm{GL}(3, \bQ(\zeta_{48}))$.
The Galois structure of the representation is
arithmetic content invisible to the Roelcke--Selberg
decomposition.

\item \textup{(Structural diagnosis.)}
The depth decomposition captures only the spectral content
of $\widehat{Z}^c_\cA$ as a function on
$\cM_{1,1} = \mathrm{SL}(2,\bZ)\backslash\mathbb{H}$;
the fusion ring and VVMF arithmetic live in the
\emph{multi-component structure} of the character vector,
not in the scalar primary-counting function.
\end{enumerate}
A refinement:
a \emph{genus-dependent arithmetic depth}
$d_{\mathrm{arith}}^{(g)}(\cA)$, defined as the number of
independent Hecke eigenforms in the genus-$g$
Roelcke--Selberg decomposition. For the Ising model,
$d_{\mathrm{arith}}^{(1)} = 0$, but
$d_{\mathrm{arith}}^{(2)}$ may be nonzero if the genus-$2$
theta-type representation encodes nontrivial Siegel eigenforms.
\end{remark}

% ============================================================
\section{Higher depths via higher-rank lattices}
\label{sec:higher-depths}

The modular-form space $M_k(\mathrm{SL}(2,\bZ))$ has cusp-form
dimension $\dim S_k = 0$ for $k<12$ and $\dim S_k \ge 1$ for
$k\ge 12$. The first cusp form is the Ramanujan discriminant
$\Delta\in S_{12}$. For weight~$k=12n$, $\dim S_k\approx n$
independent cusp forms, each contributing a critical line at
$\mathrm{Re}(s)=k/2$.

Even unimodular lattices of rank~$r$ have
$\Theta_\Lambda\in M_{r/2}$. The shadow depth of~$V_\Lambda$
therefore satisfies:
\begin{equation}\label{eq:depth-from-rank}
 d(V_\Lambda)
 = 1 + \#\{\text{distinct critical lines of }
 E_\Lambda(s)\}
 \le 1 + 2 + \dim S_{r/2}.
\end{equation}

\begin{remark}[Non-unimodular lattices]
\label{rem:non-unimodular}
\ClaimStatusProvedHere
\index{lattice vertex algebra!non-unimodular depth}
For a lattice~$\Lambda$ of determinant $D \neq 1$:
$\Theta_\Lambda \in M_{r/2}(\Gamma_0(D), \chi_D)$ where
$\chi_D$ is the Kronecker character. The Hecke theory
for~$\Gamma_0(D)$ replaces that of~$\mathrm{SL}(2,\bZ)$:
the Eisenstein subspace has $\sigma_0(D)$ cusps, and the
cusp-form space $S_{r/2}(\Gamma_0(D))$ is generically
\emph{larger} than $S_{r/2}(\mathrm{SL}(2,\bZ))$.
Non-unimodular lattices can therefore achieve higher depth
at lower rank. Examples: $A_1$ ($r{=}1$, $D{=}2$,
$\Gamma_0(4)$, depth~$2$); $D_4$ ($r{=}4$, $D{=}4$,
$\Gamma_0(2)$, depth~$3$: the space $M_2(\Gamma_0(2))$ is
two-dimensional while $M_2(\mathrm{SL}(2,\bZ)) = 0$).
\end{remark}

\begin{center}
\renewcommand{\arraystretch}{1.2}
\begin{tabular}{lccccl}
\hline
\textbf{Rank} & $k{=}r/2$ & $\dim M_k$ & $\dim S_k$
 & \textbf{Max lines} & \textbf{Max depth}\\
\hline
$8$ & $4$ & $1$ & $0$ & $2$ & $3$\\
$16$ & $8$ & $1$ & $0$ & $2$ & $3$\\
$24$ & $12$ & $2$ & $1$ & $3$ & $4$\\
$32$ & $16$ & $2$ & $1$ & $3$ & $4$\\
$48$ & $24$ & $3$ & $2$ & $4$ & $5$\\
$72$ & $36$ & $4$ & $3$ & $5$ & $6$\\
$96$ & $48$ & $5$ & $4$ & $6$ & $7$\\
$12n$ & $6n$ & $\lfloor n/2\rfloor{+}1$ & $\lfloor n/2\rfloor$
 & $\lfloor n/2\rfloor{+}2$ & $\lfloor n/2\rfloor{+}3$\\
\hline
\end{tabular}
\end{center}

\begin{definition}[Arithmetic depth filtration]
\label{def:arithmetic-depth-filtration}
Let $\mathrm{Latt}_r$ denote the set of even unimodular lattices of
rank~$r$ and $\mathrm{Gal}^{\mathrm{arith}}_r$ the set of
$\ell$-adic Galois representations attached to the Hecke eigenforms
in $S_{r/2}(\mathrm{SL}(2,\bZ))$ via Deligne's theorem.
The shadow--spectral correspondence defines a filtration on the
arithmetic data:
\begin{equation}\label{eq:arithmetic-filtration}
\mathrm{Gal}^{\mathrm{arith}}_r =
 F^2 \supset F^3 \supset \cdots \supset F^d
 \supset \cdots
\end{equation}
where $F^d$ consists of those Galois representations $\rho_{f_j}$
whose cusp forms $f_j$ contribute to $\Theta_\Lambda$ for some
lattice~$\Lambda$ with $d(V_\Lambda)\geq d$. At each depth:
\[
 \mathrm{Gr}^d = F^d / F^{d+1}
\]
captures the arithmetic content visible at shadow degree~$d$ but not
at lower degrees. The Ramanujan representation
$\rho_{\Delta_{12}}$ first appears at depth~$4$; the representations
of cusp forms of weight~$24$ first appear at depth~$5$.
\end{definition}

% ============================================================
\section{The depth decomposition}
\label{sec:depth-decomposition-section}
\label{sec:honest-assessment}
\label{sec:spectral-continuation}

\begin{theorem}[Depth decomposition]
\label{thm:depth-decomposition}
\ClaimStatusConditional
The shadow depth of~$\cA$ decomposes as
\begin{equation}\label{eq:depth-decomposition}
 d(\cA)
 = 1 + d_{\mathrm{arith}}(\cA) + d_{\mathrm{alg}}(\cA),
\end{equation}
where $d_{\mathrm{arith}}$ counts the independent holomorphic Hecke
eigenforms in the Roelcke--Selberg spectral decomposition
of~$\widehat{Z}^c_\cA$ on~$\cM_{1,1}$, and
\[
 d_{\mathrm{alg}}
 := d(\cA) - 1 - d_{\mathrm{arith}}(\cA)
 \;\ge\; 0
\]
is the \emph{homotopy defect}. The critical-line count
equals~$d_{\mathrm{arith}}$ alone: each holomorphic eigenform of
weight~$k$ contributes a critical line at $\mathrm{Re}(s) = k/2$,
while the Maass cusp-form projections contribute no new critical
lines \textup{(}their zeros all lie on
$\mathrm{Re}(s) = \tfrac12$\textup{)}.

\begin{center}
\small
\renewcommand{\arraystretch}{1.2}
\begin{tabular}{lccl}
\toprule
& $d_{\mathrm{arith}}$ & $d_{\mathrm{alg}}$ & \\
\midrule
$V_\Lambda$ ($r\ge8$)
 & $2{+}\dim S_{r/2}$ & $0$ & holomorphic \\
$\beta\gamma$
 & $1$ & $2$ & $1/\eta^2$ \\
$\mathrm{Vir}_c$ (minimal)
 & finite & $\infty$ & self-referential OPE \\
$\mathrm{Vir}_c$ ($c>1$)
 & $\le\infty$ & $\le\infty$ & \\
\bottomrule
\end{tabular}
\end{center}
\end{theorem}

\begin{proof}
Decompose~$Z(\tau,\bar\tau)$ via Roelcke--Selberg into Eisenstein
(continuous) and Maass (discrete) spectra. The holomorphic
projections (those onto weight-$k$ Hecke eigenforms
via the Rankin--Selberg integral) contribute $L$-functions with
zeros on distinct critical lines. The Maass projections
contribute $L$-functions whose zeros all lie on
$\mathrm{Re}(s) = \tfrac12$
(the spectral parameter is real). Hence
$d_{\mathrm{arith}} = \#\{\text{holomorphic eigenforms with nonzero
projection}\}$, and
$d_{\mathrm{alg}} = d - 1 - d_{\mathrm{arith}} \ge 0$ by
Theorem~\ref{thm:refined-shadow-spectral}.

For lattice VOAs: $\widehat{Z}^c = y^{c/2}\lvert\Theta\rvert^2$
is purely holomorphic, so $d_{\mathrm{alg}} = 0$ and
$d = 1 + d_{\mathrm{arith}} = 3 + \dim S_{r/2}$
(by~\eqref{eq:depth-cusp-formula}).
For $\beta\gamma$: $d = 4$
(Theorem~\ref{thm:shadow-archetype-classification}(iii)),
$d_{\mathrm{arith}} = 1$ (the primary-counting function
$\widehat{Z}^c_{\beta\gamma} = y\,|\theta_3/\eta|^2$
has $|\theta_3|^2$ holomorphic and $|\eta|^{-2}$ producing
a non-holomorphic Maass spectral component; the Hecke
decomposition of $|\theta_3|^2$ contributes one eigenform,
giving one critical line at $\mathrm{Re}(s) = 1/2$),
$d_{\mathrm{alg}} = d - 1 - d_{\mathrm{arith}} = 4 - 1 - 1 = 2$.
The homotopy defect $d_{\mathrm{alg}} = 2$ comes from the
contact quartic $\mathfrak{Q}_{\beta\gamma} \neq 0$
(Theorem~\ref{thm:betagamma-quartic-birth}):
two $A_\infty$-operations ($m_3, m_4$) are nonzero beyond
$m_2$, contributing shadow depths $3$ and $4$ to the
homotopy side.
For minimal-model Virasoro: $d_{\mathrm{arith}}$ is finite
(finitely many primaries) but $d_{\mathrm{alg}} = \infty$:
all transferred operations $m_n$ are nonzero for $n \geq 2$
(Theorem~\ref{thm:shadow-tower-asymptotics}, $S_n(c) \neq 0$
for all $n \geq 4$ and $c \notin \{0, -22/5\}$).

The source of the constraint
$d_{\mathrm{alg}} \in \{0, 1, 2, \infty\}$
is the single-line dichotomy
(Theorem~\ref{thm:single-line-dichotomy}):
the critical discriminant $\Delta = 8\kappa S_4$ forces
$r_{\max}|_L \in \{2, 3, \infty\}$ on any primary line,
and stratum separation adds $r_{\max} = 4$ for rank-one
charged sectors.
\end{proof}

\begin{remark}[Universality and the complete $d_{\mathrm{alg}}$ table]%
\label{rem:depth-decomposition-universality}%
\ClaimStatusConditional%
\index{depth decomposition!universality}%
\index{homotopy defect!complete classification}%
The decomposition~\eqref{eq:depth-decomposition}
holds for \emph{all} modular Koszul algebras, not
only lattice VOAs. The three potential failure modes
are ruled out:
\begin{enumerate}[label=\textup{(\roman*)}]
\item the Roelcke--Selberg decomposition of
 $\widehat{Z}^c_\cA$ is valid on $\cM_{1,1}$ for
 any eigenfunction of the Laplacian with moderate
 growth, which $\widehat{Z}^c_\cA$ satisfies by
 the HS-sewing criterion
 \textup{(}Theorem~\textup{\ref{thm:general-hs-sewing}}\textup{)};
\item the critical-line counting is unambiguous because
 holomorphic eigenforms contribute at $\mathrm{Re}(s) = k/2$
 while Maass eigenforms contribute at $\mathrm{Re}(s) = 1/2$
 by the Selberg eigenvalue conjecture
 \textup{(}proved for congruence subgroups\textup{)};
\item the homotopy defect
 $d_{\mathrm{alg}} = d - 1 - d_{\mathrm{arith}}$ is
 nonnegative by Theorem~\textup{\ref{thm:refined-shadow-spectral}}.
\end{enumerate}

The algebraic depth takes values in exactly the four-element
set $\{0, 1, 2, \infty\}$, classified by the shadow archetype:
\begin{center}
\small
\renewcommand{\arraystretch}{1.2}
\begin{tabular}{llcl}
\toprule
\textbf{Class} & \textbf{Archetype} & $d_{\mathrm{alg}}$ &
 \textbf{Examples} \\
\midrule
G & Gaussian (terminates at degree~$2$) & $0$ &
 Heisenberg, lattice VOAs \\
L & Lie (terminates at degree~$3$) & $1$ &
 affine KM, non-simply-laced \\
C & Contact (terminates at degree~$4$) & $2$ &
 $\beta\gamma$ system \\
M & Mixed (infinite tower) & $\infty$ &
 Virasoro, $\cW_N$, all class~M \\
\bottomrule
\end{tabular}
\end{center}
For lattice VOAs, $d_{\mathrm{alg}} = 0$, so the shadow depth
$d = 1 + d_{\mathrm{arith}}$ is entirely arithmetic.
The first $d_{\mathrm{arith}} > 3$ occurs at rank~$48$
(the space $S_{24}(\mathrm{SL}(2,\bZ))$ has
$\dim = 2$, giving $d_{\mathrm{arith}} = 4$).
Minimal models provide extreme examples:
$\cM(7,8)$ at $c = 21/22$ has
$d_{\mathrm{arith}} = 7$ (from the primary-counting spectrum),
exceeding even the rank-$96$ lattice value.
\end{remark}

\begin{remark}[The moonshine module: same $\kappa$, different class]%
\label{rem:vnatural-class-m}%
\ClaimStatusProvedHere%
\index{moonshine module!shadow class}%
\index{Niemeier lattices!versus moonshine module}%
The moonshine module $V^\natural$ has central charge $c = 24$
and $\kappa(V^\natural) = c/2 = 12$
\textup{(}the Virasoro formula applies because $T$ is a
primitive generator\textup{)}. This differs from the
$24$~Niemeier lattice VOAs, which have
$\kappa(V_\Lambda) = \operatorname{rank}(\Lambda) = 24$ The moonshine module is class~M
(infinite shadow depth): its generators live at weight~$2$
(not weight~$1$ as for lattice VOAs), so the self-referential
OPE $J(z)J(w) \sim J/z + \cdots$ produces nonvanishing
$S_n$ at all degrees
(Theorem~\ref{thm:shadow-tower-asymptotics}).

This illustrates a dichotomy among holomorphic $c = 24$
vertex algebras: the Niemeier lattice VOAs are class~G
($d_{\mathrm{alg}} = 0$, $\kappa = 24$), while
$V^\natural$ is class~M ($d_{\mathrm{alg}} = \infty$,
$\kappa = 12$). The $\bZ/2\bZ$-orbifold
$V_{\mathrm{Leech}} \to V^\natural$ kills the weight-$1$
currents, changes both $\kappa$ and the shadow class
\textup{(}$\kappa$~drops from~$24$ to~$12$; class~G
becomes class~M\textup{)}, and demonstrates that the
scalar shadow tower detects the orbifold transition.
\end{remark}

% ============================================================
\subsection{The homotopy defect as $A_\infty$ non-formality}
\index{homotopy defect!from MC element}
\label{sec:vvmf-rankin-selberg}

\begin{theorem}[$A_\infty$ formality criterion]
\label{thm:ainfty-formality-depth}
\label{thm:shadow-sullivan}
\ClaimStatusConditional
On the single-generator primary line of~$\cA$, the bar cohomology
$H^*(\barB(\cA))$ carries a transferred $A_\infty$ structure with
operations $m_n$. The single-line dichotomy
\textup{(}Theorem~\textup{\ref{thm:single-line-dichotomy})}
controls these operations via the universal factorization
$S_r = \Delta \cdot R_r$, where $\Delta = 8\kappa S_4$ is
the critical discriminant: the transferred operations
$m_n$ are nonzero if and only if their associated shadow
coefficients~$S_n$ are nonzero, and the universal
factorization forces either finitely many or all to be
nonzero. The homotopy defect is:
\[
 d_{\mathrm{alg}}(\cA) \;=\;
 \begin{cases}
 0 & \text{if $m_n = 0$ for all $n \ge 3$
 \textup{(}formal, class~G\textup{)}},\\
 \max\{r : m_r \ne 0\} - 2
 & \text{if finitely many $m_n \ne 0$
 \textup{(}classes L, C\textup{)}},\\
 \infty & \text{if $m_n \ne 0$ for all $n \ge 2$
 \textup{(}non-formal, class~M\textup{)}}.
 \end{cases}
\]
For Virasoro, $H^*(\barB(\mathrm{Vir}_c)) \cong \bC\cdot\bar{T}$
and $m_n(\bar{T}^{\otimes n}) = S_n(c)\cdot\bar{T}$ where
$S_n(c) = (2/n)(-3)^{n-4}(2/c)^{n-2}$ at leading order.
Since $S_n(c) \ne 0$ for all $n \ge 4$ and $c \ne 0$
\textup{(}Theorem~\textup{\ref{thm:shadow-tower-asymptotics}}\textup{)},
the $A_\infty$ algebra is maximally non-formal:
$d_{\mathrm{alg}}(\mathrm{Vir}_c) = \infty$.
\end{theorem}

\begin{proof}
The shadow coefficient at degree~$r$ is the chain-graph trace
$\operatorname{Sh}_r = P^{r-2}\operatorname{tr}(m_2^{r-1})
+ (\text{graphs with loops})$
(Proposition~\ref{prop:leading-hecke-identification}).
On a one-dimensional bar cohomology space,
$\operatorname{tr}(m_n) = m_n(\bar{T},\dotsc,\bar{T})/\bar{T}$
is a scalar. The shadow depth exceeds $1+d_{\mathrm{arith}}$
precisely when the $A_\infty$ operations $m_n$ for $n \ge 3$
are nonzero, i.e., when the transferred algebra is non-formal.
The leading-order formula for $S_n(c)$ is proved in
Theorem~\ref{thm:shadow-tower-asymptotics}; it is a rational
function of~$c$ with poles only at $c = 0, -22/5$, hence
nonzero for all $c \notin \{0, -22/5\}$.
\end{proof}

% ============================================================
\subsection{Interacting Gram positivity}%
\label{subsec:interacting-gram-positivity}%
\index{Gram matrix!interacting|textbf}

\begin{theorem}[Interacting Gram positivity;
\ClaimStatusProvedHere]%
\label{thm:interacting-gram-positivity}%
Let $\cA$ be chirally Koszul with shadow spectral measure
$\rho = \sum_j c_j\,\delta(\lambda_j)$
\textup{(}Theorem~\textup{\ref{thm:shadow-spectral-measure}}\textup{)}.
The \emph{interacting sewing weight}
\begin{equation}\label{eq:interacting-weight}
w_\cA(N;\alpha)
\;:=\;
a_\cA(N)\cdot
\prod_j\bigl(1 - \lambda_j\,N^{-\alpha/2}\bigr)^{c_j}
\end{equation}
satisfies:
\begin{enumerate}[label=\textup{(\roman*)}]
\item \emph{Negative-atom case}
 \textup{(}all $\lambda_j \leq 0$\textup{):}
 $w_\cA(N;\alpha) > 0$ for all $N \geq 1$
 and all $\alpha > 0$, since each factor
 $(1 + |\lambda_j|N^{-\alpha/2}) \geq 1$.
\item \emph{Positive-atom case}
 \textup{(}some $\lambda_j > 0$\textup{):}
 write $\lambda_{\max}^+ := \max\{\lambda_j > 0\}$.
 Then $w_\cA(N;\alpha) > 0$ for all $N \geq 2$ when
 $\alpha > 2\log\lambda_{\max}^+/\log 2$.
\end{enumerate}
The interacting Gram matrix
$\mathscr{M}^{\mathrm{int}} = V^T D^{\mathrm{int}}V \succ 0$
\textup{(}Theorem~\textup{\ref{thm:gram-positivity}}\textup{)}.
\end{theorem}

\begin{proof}
(i)~$\lambda_j \leq 0$ gives
$(1 - \lambda_j N^{-\alpha/2}) \geq 1 > 0$.
(ii)~$\lambda_j > 0$: the condition
$\lambda_j N^{-\alpha/2} < 1$ requires
$N > \lambda_j^{2/\alpha}$; for $N \geq 2$ this holds when
$2 > \lambda_j^{2/\alpha}$, i.e.,
$\alpha > 2\log\lambda_j/\log 2$.
\end{proof}

\begin{corollary}[\ClaimStatusProvedHere]%
\label{cor:virasoro-interacting-gram}%
For Virasoro with $c > 0$: the spectral atom
$\lambda_{\mathrm{eff}} = -6/c < 0$, so
$\mathscr{M}^{\mathrm{int}}_{\mathrm{Vir}}(\alpha) \succ 0$
for all $\alpha \geq 2$, unconditionally.
\end{corollary}

\begin{proof}
Single atom at $\lambda = -6/c < 0$.
Apply Theorem~\ref{thm:interacting-gram-positivity}(i).
\end{proof}

\begin{theorem}[\ClaimStatusProvedHere]%
\label{thm:shadow-resonance-locus}%
\index{shadow resonance locus|textbf}%
The shadow resonance locus
$\mathcal{R} :=
\{(c,\alpha) : \det\mathscr{M}^{\mathrm{int}} = 0\}$
is empty for $c > 0$ \textup{(}unitary\textup{)}.
For non-unitary $c < 0$ with $|c| < 6$:
$\lambda_{\mathrm{eff}} = |6/c| > 1$,
and the weight $w(1;\alpha)$ is negative,
so $(c,\alpha) \in \mathcal{R}$ for all~$\alpha$.
Lee--Yang $c = -22/5$ lies in $\mathcal{R}$.
\end{theorem}

\begin{proof}
$c > 0$: Corollary~\ref{cor:virasoro-interacting-gram}.
$c < 0$: $\lambda = |6/c| > 0$.
At $N = 1$: $w = a(1)(1 - |6/c|) < 0$ when $|c| < 6$.
\end{proof}

\begin{remark}[Unitarity and Gram positivity]%
\label{rem:shadow-metric-gram}%
The sign of the spectral atom $\lambda_{\mathrm{eff}} = -6/c$
is $-\operatorname{sgn}(c)$: negative for unitary $c > 0$,
positive for non-unitary $c < 0$. The interacting Gram
matrix is unconditionally positive in the unitary regime and
can degenerate only in the non-unitary regime. The shadow
resonance locus is thus confined to $c < 0$.
\end{remark}

% ============================================================
\subsection{Stieltjes representation}
\label{sec:shadow-complex-structure}

The shadow generating function
$G(t) = \sum_{r\ge2} S_r\,t^r$
has a canonical integral representation.

\begin{theorem}[Shadow spectral measure]
\label{thm:shadow-spectral-measure}
\label{thm:shadow-gf-borel}
\ClaimStatusProvedHere
On the log-normalised finite-atomic lane of the shadow tower the
shadow generating function is a Stieltjes transform:
\begin{equation}\label{eq:stieltjes-representation}
 G(t) \;=\; \int_{\bR} \log\!\bigl(1 - \lambda\,t\bigr)\,d\rho(\lambda),
 \qquad
 S_r \;=\; -\frac{1}{r}\int \lambda^r\,d\rho(\lambda),
\end{equation}
where $\rho$ is the discrete log-moment spectral measure determined by
the diagonal form of the MC element~$\Theta_\cA$:
\begin{itemize}[nosep]
\item For $V_\Lambda$ \textup{(}lattice\textup{)}, after diagonalising the
 finite Hecke module generated by the theta component of~$\Theta_\Lambda$,
 $\rho = \sum_j m_j\,\delta(\lambda - \lambda_j)$, where
 $\lambda_j$ are the eigenvalues of the log-trace operator and $m_j$
 are their multiplicities. The coefficients are the power sums
 $S_r = -r^{-1}\sum_j m_j\lambda_j^r$.
\item For $\mathrm{Vir}_c$ at leading order,
 $\rho = (c^2/162)\,\delta(\lambda + 6/c)$: a single atom at the
 effective coupling~$\lambda_{\mathrm{eff}} = -6/c$ with amplitude
 $c^2/162$.
\end{itemize}
The asymptotic shadow series
$\sum_{r\ge4} S_r^{(0)}\,t^r$, where
$S_r^{(0)} = (2/r)(-3)^{r-4}(2/c)^{r-2}$ is the leading-order
coefficient of
Theorem~\textup{\ref{thm:shadow-tower-asymptotics}},
is Borel summable with
Borel transform having a simple pole at $\xi = |c/6|$.
Up to the normalisation constant $c^2/162$ and the
first three terms, the resummed series matches~$-\log(1 + 6t/c)$:
\[
 \sum_{r\ge1} S_r^{(0)}\,t^r
 \;=\; -\frac{c^2}{162}\,\log\!\bigl(1 + \tfrac{6t}{c}\bigr),
\]
confirming the single-atom spectral measure
$\rho = \delta(\lambda + 6/c)$ at the level of
large-$r$ asymptotics. The exact shadow coefficients
$S_2 = c/2$, $S_3 = 2$ do not match the
formula $(-6/c)^r/r$; the Stieltjes
representation with a single atom is a leading-order
approximation valid for $r \ge 4$.
\end{theorem}

\begin{proof}
For the Virasoro leading-order case, write
$S_r^{(0)} = (2/r)(-3)^{r-4}(2/c)^{r-2}$.
Extracting the common factor:
$(-3)^{r-4}(2/c)^{r-2} = (-6/c)^{r-2}\cdot(-3)^{-2}
= (-6/c)^{r-2}/9$, so
$S_r^{(0)} = (2/9r)(-6/c)^{r-2}
= (c^2/162)\cdot(-6/c)^r/r$.
The identity $-\log(1-x) = \sum_{r\ge1} x^r/r$ then gives
$\sum_{r\ge1} S_r^{(0)}\,t^r
= -(c^2/162)\log(1+6t/c)$,
which is the Stieltjes transform with
$\rho = (c^2/162)\,\delta(\lambda + 6/c)$:
a single atom at the effective
coupling~$\lambda_{\mathrm{eff}} = -6/c$
with amplitude~$c^2/162$. The exact coefficients
$S_2 = c/2$ and $S_3 = 2$ differ from
$S_2^{(0)} = 1/9$ and $S_3^{(0)} = -4/(3c)$;
the single-atom spectral measure
captures only the large-$r$ asymptotics.

For the lattice case, the finite-dimensional Hecke module generated by
the theta component of~$\Theta_\Lambda$ is simultaneously diagonalisable
over the Hecke eigenbasis. Let $T_\Theta$ be the corresponding
log-trace operator. The identity
\[
 -\log\det(1-tT_\Theta)
 =
 \sum_{r\ge1}\frac{\operatorname{tr}(T_\Theta^r)}{r}\,t^r
\]
gives the finite-atomic measure
$\rho=\sum_j m_j\delta(\lambda-\lambda_j)$, where $\lambda_j$ are the
eigenvalues of~$T_\Theta$ and $m_j$ their multiplicities. The Hecke
identification of the eigenbasis is Proposition~\ref{prop:leading-hecke-identification};
the Stieltjes masses are the log-trace multiplicities, not the raw
theta-decomposition coefficients.

Borel summability: $|S_r^{(0)}| = O(r^{-1}(6/|c|)^r)$, so
the Borel transform
$B(\xi) = \sum S_r^{(0)}\xi^{r-1}/\Gamma(r)$ converges for
$|\xi| < |c|/6$ with a simple pole at $\xi = |c|/6$.
\end{proof}

\begin{remark}[Stieltjes measure and the critical discriminant]
\label{rem:stieltjes-discriminant}
\index{Stieltjes representation!critical discriminant}
The universal factorization $S_r = \Delta \cdot R_r$
(Theorem~\ref{thm:single-line-dichotomy}) implies
$\rho = \Delta \cdot \tilde\rho$ for a normalised measure
$\tilde\rho$ independent of~$\Delta$. The
Borel singularity at $\xi = |c|/6$ is an invariant
of~$\tilde\rho$; the amplitude is controlled
by~$\Delta = 40/(5c{+}22)$.
As $c \to -22/5$ (the Lee--Yang edge),
$\Delta \to 0$ and $\rho \to 0$:
the spectral measure collapses and the tower
terminates; the Gaussian envelope absorbs all content.
\end{remark}

\begin{proposition}[Carleman rigidity]
\label{prop:carleman-virasoro}
\ClaimStatusProvedHere
The Virasoro leading-order shadow moments and every positive finite
Hecke-atomic log-moment measure of Theorem~\ref{thm:shadow-spectral-measure}
satisfy Carleman's condition:
$\sum_{r\ge2} |S_r|^{-1/(2r)} = \infty$. Hence, on these two lanes,
the positive spectral measure~$\rho$ is the unique solution of the
Hamburger moment problem for the sequence $(S_r)$.
\end{proposition}

\begin{proof}
For the Virasoro leading-order lane,
$S_r=(c^2/162)(-6/c)^r/r$, hence
$|S_r|^{-1/(2r)} \to (|c|/6)^{1/2}$ as $r \to \infty$; the terms are
eventually bounded below by a positive constant.

For a positive finite Hecke-atomic measure
$\rho=\sum_{j=1}^N m_j\delta_{\lambda_j}$, let
$R=\max_j|\lambda_j|$ and $M=\sum_jm_j$. Then
$\left|\int\lambda^r\,d\rho\right|\le MR^r$, so the Carleman terms are
bounded below by a positive multiple of $R^{-1/2}$ along the
nonzero-moment subsequence. If the spectral radius is zero, the measure
is concentrated at~$0$ and is determined trivially. The Hamburger
determinacy theorem gives uniqueness.
\end{proof}

\begin{remark}[Two spectral measures]
\label{rem:moment-problem-vs-zero-location}
The measure~$\rho$ is the spectral measure of the shadow tower:
its atoms are the effective couplings of the MC
element~$\Theta_\cA$. The zero locations of~$\zeta(s)$ are
encoded in a \emph{different} spectral object: the scattering
matrix $\varphi(s) = \Lambda(1{-}s)/\Lambda(s)$ of the Laplacian
on~$\mathrm{SL}(2,\bZ)\backslash\mathfrak{H}$. The map from~$\rho$
to~$\varphi$ is the Rankin--Selberg integral, injective for
lattice VOAs (the Hecke correspondence), ill-conditioned for
non-lattice theories. The structural obstruction
(Remark~\ref{rem:structural-obstruction}) is the statement that
this map is \emph{not} surjective: the shadow spectral
measure~$\rho$ does not determine the automorphic spectral
measure on the moduli space.
\end{remark}

% ============================================================
\section{The period filtration}\label{sec:period-filtration-main}

The shadow coefficients are configuration-space integrals:
periods in the sense of Kontsevich--Zagier.

\begin{proposition}[Shadow amplitudes are periods]
\label{prop:shadow-periods}
\ClaimStatusProvedHere
\index{period!shadow coefficient as}
The shadow coefficient at degree~$r$ is a sum of
configuration-space integrals:
\begin{equation}\label{eq:shadow-period}
 S_r(c)
 = \sum_\Gamma \frac{1}{|\mathrm{Aut}(\Gamma)|}
 \int_{\overline{C}_{|V(\Gamma)|}(X)} \omega_\Gamma(c),
\end{equation}
where $\overline{C}_n(X)$ is the Fulton--MacPherson
compactification, $\Gamma$ ranges over
connected stable graphs of weight~$r$, and
$\omega_\Gamma(c) \in \Omega^{\mathrm{top}}(\overline{C}_n)$ is
algebraic in~$c$.
These integrals are periods of the mixed-Tate variety
$\overline{C}_n(\bP^1)$.
\end{proposition}

\begin{proof}
The graph-sum formula
(Definition~\ref{def:modular-bar-hamiltonian}) expresses
$\operatorname{Sh}_r$ as a finite sum of
$\Gamma$-amplitudes~$\ell_\Gamma$, each a
pushforward of vertex labels and edge
propagators along $\overline{C}_{|V|}(X) \to \overline{C}_r(X)$.
The FM space $\overline{C}_n(\bP^1)$ is mixed-Tate
(cohomology in even degrees, generated by boundary classes).
\end{proof}

\begin{remark}[Motivic interpretation of the shadow obstruction tower]
\label{rem:motivic-decomposition}
\index{shadow tower!motivic structure}
\index{mixed-Tate motive!shadow tower}
The depth decomposition
$d(\cA) = 1 + d_{\mathrm{arith}} + d_{\mathrm{alg}}$
\textup{(}Theorem~\textup{\ref{thm:depth-decomposition}}\textup{)}
admits a motivic interpretation as a weight decomposition
in the category of mixed motives over~$\bQ(c)$
\textup{(}we work informally with periods in the sense of
Kontsevich--Zagier; a precise formulation would use
Nori or Voevodsky motives\textup{)}:
\begin{align}
 d_{\mathrm{arith}}
 &= \#\bigl\{\text{non-Tate summands } h^1(f_j)
 \text{ in the motivic decomposition of } \Theta_\cA\bigr\},
 \label{eq:d-arith-motivic} \\
 d_{\mathrm{alg}}
 &= \mathrm{depth}\bigl(
 W_\bullet\,M^{\mathrm{Tate}}_\cA
 \bigr),
 \label{eq:d-alg-motivic}
\end{align}
where $M^{\mathrm{Tate}}_\cA \in \mathsf{MT}(\bQ(c))$ is the
mixed-Tate part of the shadow motive.
For lattice VOAs: $d_{\mathrm{alg}} = 0$, all non-Tate.
For Virasoro: $d_{\mathrm{arith}} = 0$, all $S_r \in \bQ(c)$.
\end{remark}



\begin{remark}[The Kummer motive and the Virasoro shadow obstruction tower]%
\label{rem:kummer-motive}%
\index{Kummer motive}%
The shadow generating function $G(t) = -\log(1+6t/c)$ at leading
order is the period of a single \emph{Kummer motive}~$K(6/c)$.
The infinite Virasoro shadow tower
$S_r = (-1)^{r+1}(6/c)^r/r$ is the iterated self-extension of
this Kummer motive: the self-referential OPE $T \in T_{(1)}T$
generates infinite combinatorial complexity but no genuine
arithmetic depth: a single Kummer motive $K(6/c)$ whose
iterated self-extensions produce the tower.
The MC equation $D\Theta + \tfrac{1}{2}[\Theta,\Theta] = 0$ is
the motivic identity: the Kummer motive is self-dual under the
convolution bracket, and this self-duality is the motivic shadow
of the self-referential OPE.
\end{remark}

\begin{remark}[Transcendence asymmetry]%
\label{rem:transcendence-asymmetry}%
The motivic decomposition yields a transcendence asymmetry:
\begin{enumerate}[label=\textup{(\roman*)}]
\item \emph{Finite depth, high transcendence.}
 The Leech lattice has depth~$4$, but its periods include
 $L(6, \Delta_{12}) \in \pi^{12}\cdot\bar{\bQ}$, a
 deeply transcendental number controlled by the Ramanujan
 conjecture (proved by Deligne).
\item \emph{Infinite depth, low transcendence.}
 The Virasoro algebra has depth~$\infty$, but all its periods
 are rational functions of~$c$. The infinite tower
 $(-6/c)^r/r$ is algebraically complex but transcendentally
 trivial.
\end{enumerate}
\end{remark}

\begin{remark}[The MC equation as motivic identity]%
\label{rem:mc-motivic-identity}%
\index{Maurer--Cartan equation!motivic interpretation}
The MC equation $D\Theta + \tfrac{1}{2}[\Theta,\Theta] = 0$
is a relation in the convolution algebra over configuration spaces.
Configuration-space integrals are the prototypical source of
mixed-Tate periods (Kontsevich formality, Drinfeld associator).
For lattice VOAs: the MC equation relates Tate periods (from
$\overline{C}_n$) to non-Tate periods (from cusp form
$L$-values) through the graph sum, a nontrivial motivic
relation connecting $\mathsf{MT}(\bQ)$ to genuine Chow
motives $h^1(f_j)$.
For Virasoro: the MC equation is purely within the
mixed-Tate world. The shadow tower is the
motivic MC element: a Maurer--Cartan element whose
degree filtration is the weight filtration of a mixed motive.
\end{remark}

\begin{remark}[Euler--Koszul motivic weight]%
\label{rem:motivic-weight-ek}%
\index{Euler--Koszul!motivic weight}
The Euler--Koszul class $\operatorname{ek}(\cA)$
(Definition~\ref{def:euler-koszul-class}) admits a motivic
interpretation. The Heisenberg sewing kernel
$K_\cH(s,t) = \zeta(s{+}t)\,\zeta(s{+}t{+}1)$
is the $L$-function of $h^0 \otimes h^0(1)$, a pure Tate
motive of weight~$1$. For Virasoro,
$K_{\mathrm{Vir}}$ involves $\zeta(u) - 1$, a mixed Tate
motive with one extension step. For $\cW_N$, the harmonic
sums $H_j(u)$ are periods of mixed Tate motives of depth~$j$.
DS reduction is a motivic truncation: it removes
the pure Tate piece and exposes the mixed extension.
Explicitly: $H_1(u) = 1$ is a period of~$\bQ(0)$;
$\zeta(u) - 1 = \sum_{N \geq 2}N^{-u}$ involves periods of
$\bQ(-1)$; the depth filtration on harmonic sums
$H_j(u) = \sum_{k=1}^j k^{-u}$ matches the weight
filtration on mixed Tate motives over~$\bZ$.
\end{remark}


\begin{remark}[The shadow tower as an object in $\mathsf{DM}(\bZ)$]%
\label{rem:shadow-tower-dm-z}%
\ClaimStatusHeuristic%
\index{shadow tower!mixed motives}%
\index{Voevodsky motives!shadow tower}%
\index{motivic Galois group!shadow tower action}%
The motivic decomposition of
Remark~\ref{rem:motivic-decomposition} admits a categorical
refinement. Let $\mathsf{DM}(\bZ)$ denote Voevodsky's triangulated
category of mixed motives over~$\Spec\,\bZ$.
The shadow tower $\Theta_\cA = \{S_r(\cA)\}_{r \geq 2}$
defines a pro-object
$M_\cA = ``\varprojlim_r"\, M_r(\cA) \in \mathrm{Pro}\text{-}\mathsf{DM}(\bZ)$,
where each $M_r(\cA)$ is the motive of the configuration-space
integral computing~$S_r$.

The motivic Galois group
$G_{\mathrm{mot}} = \mathrm{Gal}(\mathsf{DM}(\bZ))$
acts on~$M_\cA$ through its weight-graded pieces.
The G/L/C/M shadow-depth classification
\textup{(}Definition~\ref{def:shadow-depth-classification}\textup{)}
reflects the weight filtration $W_\bullet$ on the
Betti realization of~$M_\cA$:
\begin{enumerate}[label=\textup{(\roman*)}]
\item Class~G \textup{(}Heisenberg, lattice VOAs\textup{):}
 $M_\cA$ is pure of weight~$0$.
 The $G_{\mathrm{mot}}$-action factors through a finite quotient;
 all shadow coefficients are periods of~$\bQ(0)$.
\item Class~L \textup{(}affine Kac--Moody\textup{):}
 $\mathrm{Gr}^W_\bullet\,M_\cA$ is concentrated in
 weights~$\{0,1\}$.
 The weight-$1$ piece carries cusp-form periods
 \textup{(}$L$-values of~$h^1(f_j)$\textup{)};
 the single extension step $0 \to W_0 \to M_\cA \to
 \mathrm{Gr}^W_1 \to 0$ encodes the non-trivial
 arithmetic of the shadow.
\item Class~C \textup{(}$\beta\gamma$ systems\textup{):}
 weights $\{0,1,2\}$.
 The weight-$2$ piece $\mathrm{Gr}^W_2$ arises from
 the quartic OPE pole;
 DS reduction truncates to $W_{\leq 1}$.
\item Class~M \textup{(}Virasoro, $\cW_N$\textup{):}
 $M_\cA$ is mixed of unbounded weight.
 The infinite tower of extensions
 $W_0 \subset W_1 \subset W_2 \subset \cdots$
 reflects the infinite shadow depth~$d_{\mathrm{alg}} = \infty$.
 For Virasoro the motives are mixed Tate;
 for~$\cW_N$ they involve mixed Tate motives of
 depth~$\leq N{-}1$.
\end{enumerate}
The dictionary is strict: the G/L/C/M classification IS the weight
filtration on the Betti realization of~$M_\cA$, not an analogy with it.
Class~G algebras carry the trivial motive~$\bQ(0)$; their shadow
coefficients are rational numbers and the $G_{\mathrm{mot}}$-representation
is trivial.
Class~L introduces a single Tate twist: the extension
$0 \to \bQ(0) \to M_\cA \to \bQ(-1) \to 0$ is the motivic origin of the
level~$k$ in the shadow, and the shadow coefficients become periods of
$\bQ(-1)$ (values of $L$-functions at integer points).
Class~C adjoins the second Tate twist~$\bQ(-2)$, so that
$\mathrm{Gr}^W_\bullet$ is supported in weights $\{0,1,2\}$; the
quartic OPE pole of the $\beta\gamma$ system is the analytic manifestation
of this weight-$2$ piece.
For class~M the weight filtration is unbounded: the shadow invariants
$S_r$ are periods of mixed Tate motives of increasing weight, and the
transcendence degree of the period algebra grows with the shadow depth.
The higher structure of the weight filtration, the iterated extensions
rather than merely the graded pieces, encodes exactly the obstruction
data that the shadow tower computes.

The invariant $\kappa(\cA)$ admits a Frobenius interpretation:
it is the trace of the motivic Adams operation
$\psi^p$ on the weight-$1$ graded piece
$\mathrm{Gr}^W_1\,M_\cA$,
evaluated at $p = 1$ \textup{(}the archimedean place\textup{)}.
For affine Kac--Moody at level~$k$:
$\kappa(V_k(\mathfrak{g})) =
\dim(\mathfrak{g})(k + h^\vee)/(2h^\vee)$
equals the rank of~$\mathrm{Gr}^W_1$
weighted by the Hodge filtration.
For Heisenberg: $\kappa(H_k) = k$
and $\mathrm{Gr}^W_1 = 0$, consistent with the
pure weight-$0$ structure.
The vanishing $\kappa(V_{-h^\vee}(\mathfrak{g})) = 0$
at the critical level is the motivic statement that
the weight-$1$ piece degenerates.
\end{remark}

% ============================================================
\section{Non-perturbative and $p$-adic arithmetic of the shadow obstruction tower}%
\label{sec:nonperturbative-padic}%
\index{shadow tower!non-perturbative arithmetic}%
\index{resurgence!shadow tower}%
\index{$p$-adic!shadow tower convergence}

The shadow tower admits two non-perturbative extensions:
an archimedean Borel analysis of the formal power series
and a $p$-adic valuation of its coefficients.
Both reveal arithmetic content invisible to
the Roelcke--Selberg decomposition.

\subsection{The naive Borel--arithmetic conjecture is false}%
\label{subsec:borel-arithmetic-falsified}

\begin{proposition}[Two orthogonal resurgent directions;
\ClaimStatusConditional]%
\label{prop:resurgent-orthogonality}%
\index{resurgence!genus versus degree directions}%
\index{Borel singularities!degree direction}%
On the uniform-weight scalar lane, and on the Virasoro scalar
shadow tower for the degree direction, the shadow tower has Borel
singularities in two independent directions:
\begin{enumerate}[label=\textup{(\roman*)}]
\item \textup{(Genus direction.)}
The genus expansion $\sum_g F_g \hbar^{2g}$ has
factorial growth $F_g \sim (2g)!$ from the stable-graph
combinatorics. The Borel singularity is
$\hbar_* = (2\pi)^{-2}$, the universal instanton action
of the modular operad. This is independent of the
algebra~$\cA$.
\item \textup{(Degree direction.)}
The degree expansion $\sum_r S_r t^r$ has Borel
singularities at $t_* = -c/(6\alpha)$ for Virasoro,
where $\alpha$ depends continuously on the central
charge~$c$. These singularities are algebraic
numbers in~$\bQ(c)$.
\end{enumerate}
The genus Borel singularity $(2\pi)^{-2}$ is
\emph{transcendental and universal};
the degree Borel singularity is
\emph{algebraic and algebra-dependent}.
The two are generically incommensurable, so the
Borel singularities do not align with $L$-function
zeros \textup{(}which are transcendental but
conjecturally lie on vertical lines\textup{)}.
\end{proposition}

\begin{proof}
The genus-direction growth is standard
(Bender--Canfield for stable graphs).
The degree-direction singularities coincide with the zeros
of the shadow metric $Q_L(t)$
(Theorem~\ref{thm:shadow-connection}): for Virasoro,
$Q_L(t) = (c + 3\alpha t)^2 + 2\Delta t^2$ vanishes
at $t_* \in \bQ(c, \sqrt{-\Delta})$.
\end{proof}

\begin{proposition}[Universal Stokes constants; \ClaimStatusProvedHere]%
\label{prop:universal-stokes-constants}%
\index{Stokes constants!universal ratio|textbf}%
\index{instanton action!universal}%
On the uniform-weight scalar genus lane, the genus-direction Stokes
constant at the $n$-th instanton
$\omega_n = (2\pi n)^2$ is
\begin{equation}\label{eq:stokes-universal}
S_n \;=\; (-1)^n \cdot 4\pi^2 n \cdot \kappa(\cA) \cdot i,
\end{equation}
with universal ratio $S_n / S_1 = (-1)^{n+1} \cdot n$,
independent of~$\kappa(\cA)$ and of the algebra family.
The genus expansion
$\sum_{g \geq 1} F_g(\cA)\,\xi^{2g}
= \kappa\bigl(\xi/(2\sin(\xi/2)) - 1\bigr)$
is meromorphic with simple poles at $\xi = 2\pi n$;
the residue at $\xi = 2\pi n$ is
$(-1)^n \cdot 2\pi n \cdot \kappa$, and
$S_n = 2\pi i \cdot \mathrm{Res}_{\xi = 2\pi n}$.
\end{proposition}

\begin{proof}
From the closed form
$G(\xi) = \kappa\cdot\xi/(2\sin(\xi/2)) - \kappa$,
the poles are the zeros of $\sin(\xi/2)$, i.e.,
$\xi = 2\pi n$ for $n \in \bZ \setminus \{0\}$.
By l'H\^opital,
$\mathrm{Res}_{\xi=2\pi n} \xi/(2\sin(\xi/2))
= 2\pi n / \cos(\pi n) = (-1)^n \cdot 2\pi n$.
The Stokes constant is $2\pi i$ times the residue.
The ratio $S_n/S_1 = (-1)^{n+1} n$ is immediate.
\end{proof}

\begin{proposition}[Gevrey-$0$ degree growth; \ClaimStatusProvedHere]%
\label{prop:gevrey-zero-degree}%
\index{Gevrey class!shadow tower|textbf}%
\index{convergence radius!shadow tower|textbf}%
The degree expansion
$\sum_{r \geq 2} S_r t^r$ has sub-factorial growth:
$|S_r(c)| \sim C(c) \cdot \rho(c)^r \cdot r^{-5/2}$
as $r \to \infty$, with convergence radius
$1/\rho(c) = |t_*|$ where $t_*$ is the nearest zero
of~$Q_L(t)$ to the origin. The Darboux exponent
$-5/2 = -3/2$ \textup{(}from the algebraic branch point
of $\sqrt{Q_L}$\textup{)} minus~$1$
\textup{(}from the $1/r$ in $S_r = a_{r-2}/r$\textup{)}.
The expansion converges absolutely for
$c > c^* \approx 6.125$
\textup{(}unique positive root of
$5c^3 + 22c^2 - 180c - 872 = 0$\textup{)} and diverges
geometrically for $c < c^*$. In both regimes, the growth
is Gevrey-$0$: $|S_r|/r! \to 0$.
\end{proposition}

\begin{proof}
The coefficients $a_n = [t^n]\sqrt{Q_L(t)}$ satisfy
the convolution recursion
$a_n = -(1/(2c))\sum_{j=1}^{n-1} a_j a_{n-j}$ for
$n \geq 3$, with $a_0 = c$, $a_1 = 6$.
Since $\sqrt{Q_L}$ is algebraic of degree~$2$, the
Darboux transfer theorem gives
$a_n \sim C \cdot \rho^n \cdot n^{-3/2}$ where
$\rho = 1/|t_*|$. Division by~$r$ in $S_r = a_{r-2}/r$
lowers the exponent to~$-5/2$. The critical central
charge $c^*$ is where $|t_*| = 1$, equivalently
$\rho(c^*) = 1$.
\end{proof}

\begin{remark}[Resurgent anomaly cancellation at $c = 26$]%
\label{rem:resurgent-anomaly-cancellation}%
\ClaimStatusHeuristic%
\index{anomaly cancellation!resurgent}%
\index{critical dimension!resurgent}%
At $c = 26$, the matter and ghost genus-direction
Stokes multipliers cancel exactly:
$\mathrm{Stokes}(\cA_{\mathrm{matter}}) +
\mathrm{Stokes}(\cA_{\mathrm{ghost}}) = 0$.
This is the resurgent form of the critical-dimension
anomaly cancellation $\kappa_{\mathrm{eff}} = 0$
(the Stokes multipliers are weighted by~$\kappa$).
The degree-direction Stokes multipliers do \emph{not}
cancel at $c = 26$: matter and ghosts have different
shadow metrics, and cancellation there would require
$\Delta_{\mathrm{matter}} + \Delta_{\mathrm{ghost}} = 0$,
which does not hold in general.
\end{remark}

\subsection{The $p$-adic shadow obstruction tower}%
\label{subsec:padic-shadow-tower}

\begin{proposition}[$p$-adic convergence radius;
\ClaimStatusProvedHere]%
\label{prop:padic-convergence}%
\index{$p$-adic!convergence radius}%
\index{Kummer congruences!shadow tower}%
The shadow generating function $H(t) = \sum_r r\,S_r\,t^r$,
viewed $p$-adically, converges on the disc
$|t|_p < p^{-1/(p-1)}$, the radius of convergence of the
$p$-adic exponential function.
\end{proposition}

\begin{proof}
The $\hat{A}$-genus generating function
$\hat{A}(x) = (x/2)/\sinh(x/2) =
\sum_{n \geq 0} B_{2n}\, x^{2n} / (2n)!$
involves the Bernoulli numbers $B_{2n}$, which satisfy
$v_p(B_{2n}/(2n)!) \geq -(2n)/(p{-}1)$ by the
von Staudt--Clausen theorem. Since the shadow
coefficients $S_r$ are rational functions of
$\kappa, S_3, S_4, \ldots$ with the $\hat{A}$ function
controlling the genus expansion, the $p$-adic valuation
of the Taylor coefficients decays at rate $\leq r/(p{-}1)$,
giving convergence radius
$p^{-1/(p-1)}$.
\end{proof}

\begin{remark}[Kummer congruences and the Kubota--Leopoldt
$p$-adic $L$-function]%
\label{rem:kummer-kubota-leopoldt}%
\ClaimStatusProvedHere%
\index{Kubota--Leopoldt $p$-adic $L$-function}%
The Bernoulli numbers satisfy the Kummer congruences
$B_m/m \equiv B_n/n \pmod{p}$ whenever
$m \equiv n \pmod{p{-}1}$ and $(p{-}1) \nmid m$.
These are exactly the interpolation conditions defining
the Kubota--Leopoldt $p$-adic $L$-function
$L_p(s, \omega^a)$ for Dirichlet characters.
The shadow tower inherits this structure: the degree-$r$
coefficients $S_r(\cA)$, when reduced modulo~$p$,
satisfy congruences determined by the $p$-adic residue
class of~$r$. For Heisenberg ($S_r = 0$ for $r \geq 3$),
all congruences are trivially satisfied. For Virasoro
($S_r = (-1)^{r+1}(6/c)^{r-2}/r$ at leading order), the
congruences reduce to $p$-adic properties of the geometric
series in~$6/c$, which are controlled by $v_p(6/c)$.
The convergence radius $p^{-1/(p-1)}$ is the shadow-tower
manifestation of the $p$-adic exponential radius, linking
the tower to the Iwasawa-theoretic framework.
\end{remark}

\begin{remark}[Hilbert symbol detection of shadow-spectral zero-crossings]
\label{rem:hilbert-symbol-detection}
\index{Hilbert symbol!shadow-spectral detection|textbf}
\index{shadow invariants!$p$-adic zero-crossings}
The $p$-adic Hilbert symbols $(\kappa, \Delta)_p$ detect
shadow-spectral zero-crossings: sign changes in
$(\kappa, \Delta)_p$ at $p = 2, 3$ across consecutive
shadow zeros correspond to crossings of the shadow-spectral
locus. Hilbert reciprocity
\begin{equation}\label{eq:hilbert-reciprocity-shadow}
 \prod_{p \le \infty} (\kappa, \Delta)_p \;=\; 1
\end{equation}
holds for all standard families. This is the
strongest arithmetic signal connecting the shadow
invariants $(\kappa, \Delta)$ to the analytic
zero distribution of the constrained Epstein zeta:
the Hilbert symbol encodes the local splitting behavior
of the shadow field $K_L = \bQ(\sqrt{\Delta})$
at each prime, and the product formula constrains
the global pattern of sign changes.
\end{remark}

\begin{remark}[Ferrero--Washington failure and the $\mu$-invariant]
\label{rem:ferrero-washington-failure}
\index{Ferrero--Washington!failure for shadow tower|textbf}
\index{$\mu$-invariant!nonvanishing}
\index{Newton polygon!Mazur inequality}
\index{Nahm eigenvalue}
The Ferrero--Washington theorem ($\mu = 0$ for cyclotomic
$\bZ_p$-extensions) \emph{fails} for the shadow
$p$-adic $L$-function:
\begin{enumerate}[label=\textup{(\roman*)}]
\item The $\mu$-invariant is nonzero for the Virasoro
 shadow tower at $p = 2$ (growth rate $\sim\! -1$) and
 $p = 5$ (growth rate $\sim\! -0.5$). The nonvanishing
 reflects the irregular singularity of the shadow
 connection at $t = \infty$
 (Theorem~\ref{thm:shadow-higgs-field}).
\item Newton polygons: the Mazur inequality
 $\mathrm{NP}(f) \ge \mathrm{HP}(f)$ (Newton polygon
 lies on or above the Hodge polygon) holds universally
 for the shadow tower's $p$-adic power series,
 verified for all standard families at primes
 $p \le 37$.
\item The arithmetic depth $d_{\mathrm{arith}}$ and the
 algebraic depth $d_{\mathrm{alg}}$ are independent:
 lattice VOAs at the same central charge $c = 24$
 have $d_{\mathrm{arith}} \in \{3, 4\}$
 (depending on the root system) while
 $d_{\mathrm{alg}} = 0$ universally.
\item The Nahm eigenvalue of the shadow metric satisfies
 $\lambda_{\mathrm{Nahm}} = (\kappa/c)^2$: the square
 of the anomaly ratio $\kappa/c$ (not to be confused
 with the shadow growth rate~$\rho$ of
 Definition~\textup{\ref{def:shadow-growth-rate}}).
\item The shadow curve
 (Theorem~\ref{thm:spectral-curve}) is always genus~$0$:
 the spectral curve $\eta^2 = c^2 + 12ct + \alpha t^2$
 is a conic, rational over~$\bQ(c)$.
\item The conductor of the shadow field
 $K_L = \bQ(\sqrt{\Delta})$ is \emph{not}
 Koszul-dual-symmetric in general: $\mathfrak{f}(\cA)
 \neq \mathfrak{f}(\cA^!)$ unless $c = 13$
 (the Virasoro self-dual point).
\end{enumerate}
\end{remark}

\subsection{The MZV--bar complex dictionary}%
\label{subsec:mzv-bar-complex}%
\index{multiple zeta values!bar complex}%
\index{bar complex!periods}%
\index{Drinfeld associator!bar transport}

Multiple zeta values (MZVs)
$\zeta(s_1, \ldots, s_k) =
\sum_{n_1 > \cdots > n_k > 0} n_1^{-s_1} \cdots n_k^{-s_k}$
are periods of the moduli spaces~$\cM_{0,n}$.
The bar complex $\barB(\cA)$ computes factorization homology
on curves, and genus-$0$ bar amplitudes on $\cM_{0,n+3}$
produce MZVs of weight $\leq n$ as their periods.

\begin{theorem}[Shadow--MZV dictionary;
\ClaimStatusProvedHere]%
\label{thm:shadow-mzv-dictionary}%
\index{shadow tower!MZV dictionary|textbf}%
\index{multiple zeta values!shadow identification}%
The shadow invariants of a chirally Koszul algebra~$\cA$
determine and are determined by the MZV content of the
genus-$0$ bar amplitudes, according to the following dictionary:
\begin{enumerate}[label=\textup{(\roman*)}]
\item \textup{(Weight $2$.)}
The modular characteristic $\kappa(\cA)$ controls the
coefficient of $\zeta(2) = \pi^2/6$ in the genus-$0$
four-point amplitude:
\[
A_{0,4}(\cA)
\;=\; \kappa(\cA) \cdot \bigl(-2\,\zeta(2)\bigr) + O(\text{higher}).
\]
The MZV space at weight~$2$ is one-dimensional
\textup{(}$d_2 = 1$ by Zagier's conjecture, proved motivically
by Brown \cite{Brown12}\textup{)}, spanned by~$\zeta(2)$.

\item \textup{(Weight $3$.)}
The cubic shadow $S_3(\cA)$ controls the coefficient of
$\zeta(3)$ in the genus-$0$ five-point amplitude.
(The $c$-independent identity $S_3(\mathrm{Vir}_c) = 2$ (Theorem~\ref{thm:s3-virasoro-c-independent}) provides a finite algebraic proof of class~$\mathsf{M}$ membership.)
For algebras with $S_3 = 0$
\textup{(}classes~$\mathbf{G}$ and~$\mathbf{C}$\textup{)},
$\zeta(3)$ does not appear.
Euler's relation $\zeta(2,1) = \zeta(3)$ ensures the weight-$3$
MZV space is one-dimensional ($d_3 = 1$).

\item \textup{(Weight $4$.)}
The quartic shadow $S_4(\cA)$ and quartic contact class
$Q^{\mathrm{contact}}(\cA)$ control the weight-$4$ MZV
coefficients. The MZV space at weight~$4$ is one-dimensional
\textup{(}$d_4 = 1$\textup{)}, spanned by $\zeta(4) = \pi^4/90$,
with $\zeta(3,1) = \zeta(4)/4$ and $\zeta(2,2) = 3\zeta(4)/4$
reducing to~$\zeta(4)$.

\item \textup{(General weight~$r$.)}
The degree-$r$ shadow invariant $S_r(\cA)$ contributes to
the weight-$r$ MZV space of dimension
$d_r = d_{r-2} + d_{r-3}$ \textup{(}$d_0 = 1$, $d_1 = 0$,
$d_2 = 1$; proved motivically by Brown\textup{)}.
Class~$\mathbf{G}$ algebras
\textup{(}Heisenberg, lattice VOAs\textup{)} see only
$\zeta(2)$; class~$\mathbf{L}$ algebras
\textup{(}affine Kac--Moody\textup{)} see
$\zeta(2)$ and $\zeta(3)$;
class~$\mathbf{C}$ algebras
\textup{(}$\beta\gamma$\textup{)} see $\zeta(2)$ and $\zeta(4)$
\textup{(}$S_3 = 0$ but $S_4 \neq 0$\textup{)};
class~$\mathbf{M}$ algebras
\textup{(}Virasoro, $\cW_N$\textup{)} see MZVs at all weights.
\end{enumerate}
\end{theorem}

\begin{proof}
The genus-$0$ bar amplitude on $\cM_{0,n}$ is computed by
the factorization-homology integral of the bar propagator
$\dlog(z_i - z_j)$ over the Fulton--MacPherson compactification
$\overline{\mathrm{FM}}_n(\bC)$. The iterated integrals of
$\dlog$ forms on $\cM_{0,n+3} = \bP^1 \setminus \{0, 1, \infty\}
\times \cdots$ produce MZVs by the theorem of Kontsevich
(periods of $\cM_{0,n}$ are MZVs; proved by Brown 2009).
The bar complex with $r$-matrix
$r(z) = \sum_{j} c_j z^{-j}$ produces iterated integrals
of depth $\leq k$ and weight $\leq r$, where the bar
propagator $\dlog E(z,w)$ contributes weight~$1$ per edge
 and each $r$-matrix pole of order~$j$ contributes
weight~$j{-}1$ after the $\dlog$ absorption.
For Heisenberg: $r(z) = \kappa/z$ (single pole), so all
bar amplitudes reduce to products of weight-$1$ integrals,
producing only~$\zeta(2)$.
For Virasoro: $r(z) = (c/2)/z^3 + 2T/z$ (pole orders $3$ and $1$),
producing MZVs at all weights.
The dimension formula $d_r = d_{r-2} + d_{r-3}$ is Brown's
theorem (Annals 2012), applied motivically.
\end{proof}

\begin{theorem}[Drinfeld associator as bar transport;
\ClaimStatusConditional]%
\label{thm:drinfeld-associator-bar-transport}%
\index{Drinfeld associator!bar transport|textbf}%
\index{KZ connection!bar complex interpretation}%
\index{bar complex!KZ connection}%
The Drinfeld associator
$\Phi(A,B) \in \widehat{U}(\mathrm{Lie}(A,B))$
is the parallel transport of the genus-$0$ bar connection
along $\cM_{0,4} = \bP^1 \setminus \{0, 1, \infty\}$.
Specifically:
\begin{enumerate}[label=\textup{(\roman*)}]
\item The KZ connection
$\nabla^{\mathrm{KZ}} = d - (A/z + B/(z{-}1))\,dz$
on $\cM_{0,4}$ is the genus-$0$, degree-$2$ shadow connection
\textup{(}the classical affine $r$-matrix
$r(z) = k\,\Omega/z$ with
$\Omega = \sum t^a \otimes t_a$ the Casimir; the level
prefix $k$ enforces $r = 0$ at $k = 0$\textup{)}.
\item The monodromy from $0$ to $1$ gives
$\Phi(A,B) = 1 + \sum_w I(w) \cdot w$, where the
coefficient $I(w)$ of word $w = e_0^{a_1-1} e_1 \cdots
e_0^{a_k-1} e_1$ is
$I(w) = (-1)^k \zeta(a_1, \ldots, a_k)$
\textup{(}Le--Murakami convention\textup{)}.
\item The associator is group-like:
$\Delta(\Phi) = \Phi \otimes \Phi$, equivalently
$\log \Phi$ is a Lie series. The Lie coefficients are:
\begin{align*}
\text{weight } 2{:}\quad & \zeta(2)\,[A, B], \\
\text{weight } 3{:}\quad & \zeta(3)\,
\bigl([A, [A, B]] - [B, [A, B]]\bigr), \\
\text{weight } 4{:}\quad &
\tfrac{\zeta(4)}{4}\,\bigl([A,[A,[A,B]]] +
[B,[B,[A,B]]]\bigr) + \cdots
\end{align*}
\item $\Phi$ satisfies the pentagon equation
$\Phi_{12,3,4}\,\Phi_{1,2,34} =
\Phi_{2,3,4}\,\Phi_{1,23,4}\,\Phi_{1,2,3}$.
\end{enumerate}
The bar-complex interpretation:
the bar differential on $\cM_{0,n}$ with $n$ marked
points is the KZ connection evaluated on the
factorization-algebra input, and the associator is
the holonomy datum that relates different bar
complex trivializations.
\end{theorem}

\begin{proof}
The KZ connection on $\cM_{0,n+1}$ is the bar connection
for the affine Kac--Moody algebra at level~$k$, restricted
to the genus-$0$ moduli
(Theorem~\ref{thm:yangian-shadow-theorem} of
Chapter~\ref{chap:yangians}).
The identification $r(z) = k\,\Omega/z$ with Casimir
$\Omega$ and affine level $k$ is classical (Kohno 1988;
see Theorem~\ref{thm:drinfeld-kohno}); the $k$ prefix
satisfies the $k = 0$ vanishing check .
The MZV coefficients of $\Phi$ were proved by
Le--Murakami (1995) and the group-like property
by Drinfeld (1990). The pentagon equation is the
defining relation for quasi-Hopf algebras.
The bar-complex interpretation follows from
Theorem~\ref{thm:drinfeld-kohno}: the KZ
monodromy representation agrees with the quantum
group $R$-matrix action, and the bar transport on
$\overline{\mathrm{FM}}_n(\bC)$ produces the same
flat connection by the factorization structure of
the bar complex
(Theorem~\ref{thm:bar-modular-operad}).
\end{proof}

\begin{remark}[Zagier dimension sequence and shadow depth]%
\label{rem:zagier-dimension-shadow}%
\index{multiple zeta values!Zagier dimension}%
\index{shadow depth!MZV dimension}%
The Zagier--Brown dimension sequence
$d_r = d_{r-2} + d_{r-3}$, giving
$1, 0, 1, 1, 1, 2, 2, 3, 4, 5, 7, 9, 12, 16, 21, \ldots$,
counts the independent MZV periods at each weight.
For class~$\mathbf{G}$ algebras (shadow depth~$2$), only
the $r = 2$ entry is activated: the bar amplitudes see
a single period~$\zeta(2)$.
For class~$\mathbf{L}$ (shadow depth~$3$), entries at $r = 2, 3$
are activated.
For class~$\mathbf{M}$ (infinite shadow depth), all entries
are activated, and the shadow coefficients
$S_r$ provide a complete system of coordinates on
the MZV space at each weight. Whether the shadow
coefficients $S_r(\cA)$ for $\cA$ ranging over all
modular Koszul algebras span the full MZV space at
each weight is open, connected to
transcendence conjectures for MZVs.
\end{remark}

\begin{remark}[Shadow--MZV observational dictionary]%
\label{rem:shadow-mzv-observational}%
\index{multiple zeta values!observational shadow dictionary|textbf}%
\index{shadow field!Galois involution as Koszul sign}%
\index{shadow tower!arithmetic observations}
Three observational connections between the shadow tower and
the arithmetic of the shadow field~$K_L$, extending
Theorem~\textup{\ref{thm:shadow-mzv-dictionary}} beyond the
genus-$0$ MZV content:
\begin{enumerate}[label=\textup{(\roman*)}]
\item
\emph{Four-class partition as splitting behaviour.}\enspace
The shadow depth partition
$\mathbf{G}/\mathbf{L}/\mathbf{C}/\mathbf{M}$
coincides with the local splitting behaviour of the
critical discriminant
$\operatorname{disc}(Q_L/F) = -320c^2/(5c{+}22)$
at the rational prime governed by the Hilbert symbol
$(\kappa, \Delta)_p$
\textup{(}Remark~\textup{\ref{rem:hilbert-symbol-detection}}\textup{)}:
class~$\mathbf{G}$ algebras correspond to totally split
behaviour \textup{(}$\Delta = 0$, shadow field degenerates
to~$\bQ$\textup{)}; class~$\mathbf{L}$ to split
\textup{(}the shadow metric factors as a square over the
shadow field\textup{)}; class~$\mathbf{C}$ to ramified
behaviour at a single prime; class~$\mathbf{M}$ to
inert behaviour with genuinely quadratic extension
of~$\bQ$. The four-class partition is an arithmetic
invariant of the splitting type.
\item
\emph{Galois involution as the Koszul sign.}\enspace
The Galois involution
$\sqrt{Q_L(t)} \mapsto -\sqrt{Q_L(t)}$
on the generic shadow extension
$K_L = \bQ(\kappa, \sqrt{Q_L(t)})$
\textup{(}Definition~\textup{\ref{def:specialised-shadow-field}}\textup{)}
implements the Koszul sign: the involution acts as
$(-1)$ on the monodromy representation of the shadow
connection~$\nabla^{\mathrm{sh}}$ around each zero
of~$Q_L$, and this sign matches the grading shift of
the Koszul-dual algebra. The \emph{ring of integers}
associated to the shadow tower is the graded ring
$\bZ[\kappa, \alpha, S_4][\{S_r\}_{r \geq 5}]$
obtained by adjoining the higher shadow coefficients
to the shadow-metric Hankel data.
\item
\emph{Rational-$c_0$ MZV specialisations
\textup{(}conjectural\textup{)}.}\enspace
At specific rational central charges, the specialised
shadow field $\operatorname{Spec}_{c_0}(K_L)$ of
Definition~\textup{\ref{def:specialised-shadow-field}}
pairs with the MZV dictionary of
Theorem~\textup{\ref{thm:shadow-mzv-dictionary}} to
produce arithmetic predictions. At $c_0 = 1/2$
\textup{(}Ising\textup{)}, the specialised field is
$\bQ(\sqrt{-10})$
\textup{(}Example~\textup{\ref{ex:specialised-shadow-field-ising}}\textup{)}
and the shadow coefficients $S_r(1/2)$ are conjectured to
involve specific $L$-values $L(s, \chi_{-40})$ at
integer arguments, pairing with the MZV weights through
the Deligne conjecture.
At $c_0 = 7/10$ \textup{(}tricritical Ising\textup{)},
the specialised field is $\bQ(\sqrt{-510})$
\textup{(}from $-5(5 \cdot 7/10 + 22) = -5 \cdot 51/2
= -255/2$, whose squarefree class in
$\bQ^\times/(\bQ^\times)^2$ is represented by
$-510 = -2 \cdot 3 \cdot 5 \cdot 17$\textup{)} and
a parallel MZV specialisation is conjectured. The present proof
surface consists of the low-weight computations $r \leq 4$; the
higher-weight statement remains open.
\end{enumerate}
The three observations share a common source: the
shadow field $K_L$ is the splitting field of the
quadratic first integral of the shadow connection, and
every arithmetic invariant of the tower factors through
this field.
\end{remark}

\subsection{Genus-$1$ modular forms from partition functions}%
\label{subsec:genus1-modular-forms-partition}%
\index{partition function!modular form}%
\index{modular form!from partition function}

The genus-$1$ partition function $Z(\cA, \tau) =
\operatorname{tr}_{V(\cA)} q^{L_0 - c/24}$ of a chirally
Koszul algebra is a function on $\cM_{1,1}$ whose modular
properties encode the shadow data. The classification
is exhaustive for the standard landscape.

\begin{theorem}[Partition function modular classification;
\ClaimStatusProvedHere]%
\label{thm:partition-modular-classification}%
\index{partition function!modular classification|textbf}%
\index{Hecke eigenform!from partition function}%
For each standard family of chirally Koszul algebras,
the genus-$1$ partition function $Z(\cA, \tau)$ belongs
to the following modular class:
\begin{enumerate}[label=\textup{(\roman*)}]
\item \textup{(Heisenberg.)}
$Z(\cH_k, \tau) = 1/\eta(\tau)$: weight $-1/2$ with
multiplier system. Not a modular form, but
$\eta(\tau)^{-1}$ has a Hecke-equivariant Fourier
expansion via the partition function $p(n)$.
The partition function is independent of level: $\kappa(\cH_k)=k$,
and the displayed row has $\kappa=1$ only after the
level-one normalization $k=1$. The shadow class is~$\mathbf{G}$.

\item \textup{(Affine $\widehat{\fsl}_2$ level $1$.)}
$Z = \theta_3(\tau)/\eta(\tau)$: weight~$0$ with multiplier.
$\kappa = 9/4$, class~$\mathbf{L}$.

\item \textup{($E_8$ lattice VOA.)}
$Z = E_4(\tau) = 1 + 240\sum_{n \geq 1} \sigma_3(n)\,q^n$:
holomorphic modular form of weight~$4$, Hecke eigenform.
$\kappa = 8$, class~$\mathbf{G}$.
The Hecke eigenvalues $a_p(E_4) = 1 + p^3 = \sigma_3(p)$,
and the $L$-function factorizes:
$L(E_4, s) = \zeta(s)\,\zeta(s{-}3)$.

\item \textup{($D_{16}^+$ lattice VOA.)}
$Z = E_4(\tau)^2 = E_8(\tau)$: weight~$8$, Hecke eigenform
\textup{(}since $\dim M_8(\mathrm{SL}_2(\bZ)) = 1$\textup{)}.
$\kappa = 16$, class~$\mathbf{G}$.
The $L$-function factorizes:
$L(E_8, s) = \zeta(s)\,\zeta(s{-}7)$.

\item \textup{(Leech/Monster.)}
$Z(V^\natural, \tau) = J(\tau) = j(\tau) - 744$:
weight-$0$ modular function \textup{(}Hauptmodul\textup{)},
genus-$0$ property.
$\kappa(V^\natural) = 12$.
The vanishing constant term $c(0) = 0$ follows from the
Rademacher exact formula: the polar part $q^{-1}$
determines the entire function via genus-$0$ modularity.
\end{enumerate}
\end{theorem}

\begin{proof}
Each identification is classical: (i) is the oscillator
partition function, (ii) is the Kac character formula for
$\widehat{\fsl}_2$ at level~$1$,
(iii) follows from $\Theta_{E_8} = E_4$
(Jacobi's theta function identity for the $E_8$ root lattice),
(iv) from $\dim M_8(\mathrm{SL}_2(\bZ)) = 1$ forcing
$\Theta_{D_{16}^+} = E_4^2$,
and (v) from Zhu's modular invariance theorem applied to the
holomorphic $c = 24$ VOA~$V^\natural$.
The $\kappa$ values are from the landscape census
(Table~\ref{tab:shadow-tower-census}).
\end{proof}

\begin{proposition}[Quasi-modular content from the genus-$1$
propagator; \ClaimStatusProvedHere]%
\label{prop:quasi-modular-propagator}%
\index{quasi-modular form!from propagator}%
\index{Eisenstein series!$E_2^*$ propagator}%
The genus-$1$ bar complex propagator is
$P(\tau) = -E_2^*(\tau)/12$, where
$E_2^*(\tau) = 1 - 24\sum_{n \geq 1} \sigma_1(n)\,q^n$
is the quasi-modular Eisenstein series of weight~$2$ and depth~$1$.
Under the $S$-transformation:
\begin{equation}\label{eq:e2-star-anomaly}
E_2^*(-1/\tau) = \tau^2\,E_2^*(\tau) + \frac{12\tau}{2\pi i}.
\end{equation}
The additive anomaly~$12\tau/(2\pi i)$ is the obstruction to
$E_2^*$ being a holomorphic modular form. For any chirally Koszul
algebra~$\cA$ with $\kappa = \kappa(\cA)$:
\begin{enumerate}[label=\textup{(\roman*)}]
\item The weight-$2$ quasi-modular piece of the genus-$1$ shadow
amplitude is $\kappa \cdot E_2^*(\tau)$
\textup{(}from the propagator insertion\textup{)}.
\item The weight-$4$ holomorphic piece is
$Q^{\mathrm{contact}}(\cA) \cdot E_4(\tau)$
\textup{(}from the quartic shadow\textup{)}.
\item The weight-$4$ quasi-modular piece is
$\kappa^2 \cdot (E_2^*)^2(\tau)$
\textup{(}from double propagator insertion, depth~$2$\textup{)}.
\end{enumerate}
The genus-$1$ shadow amplitudes live in the ring
$\mathrm{QM}_*(\mathrm{SL}_2(\bZ)) = \bC[E_2^*, E_4, E_6]$
of quasi-modular forms, \emph{not} in the holomorphic ring
$M_*(\mathrm{SL}_2(\bZ)) = \bC[E_4, E_6]$.
\end{proposition}

\begin{proof}
The genus-$1$ propagator on the torus $\bC/(\bZ + \tau\bZ)$
is the derivative of the Weierstrass $\zeta$-function,
whose holomorphic part is
$-E_2^*(\tau)/12 + \text{(non-holomorphic)}$.
The quasi-modular transformation~\eqref{eq:e2-star-anomaly}
follows from the transformation of the Weierstrass
$\zeta$-function under $\tau \mapsto -1/\tau$.
The ring structure $\bC[E_2^*, E_4, E_6]$ is classical
(Kaneko--Zagier).
\end{proof}

\begin{remark}[L-functions from Koszul pairs]%
\label{rem:l-functions-koszul-pairs}%
\index{L-function!Koszul pair}%
\index{Rankin--Selberg!Koszul pair}%
For a Koszul pair $(\cA, \cA^!)$ whose partition functions
$f = Z(\cA)$ and $g = Z(\cA^!)$ are modular forms, the
Rankin--Selberg convolution
$L(f \times g, s) = \sum_{n \geq 1} a_n(f)\,a_n(g)\,n^{-s}$
encodes the complementarity data.
For $\cA = E_8$ lattice VOA (self-dual at the partition-function
level): $L(E_4 \times E_4, s) =
\zeta(s)\,\zeta(s{-}3)^2\,\zeta(s{-}6)/\zeta(2s{-}6)$,
the Rankin--Selberg square (Rankin 1939).
For the Ramanujan $L$-function $L(\Delta, s)
= \sum \tau(n)\,n^{-s}$ (the cusp-form $L$-function
of weight~$12$): the critical values at $s = 1, \ldots, 11$
encode periods of~$\Delta$ via the Eichler--Shimura isomorphism.
The shadow tower at genus~$1$ sees the Eisenstein part
(through $\kappa$) but not the cusp part (which requires
the full theta series together with the shadow data).
The passage from shadow data to cusp-form content is the
arithmetic depth of the theta decomposition
(Proposition~\ref{prop:lattice:niemeier-theta-decomposition}
of Chapter~\ref{ch:lattice}).
\end{remark}

\subsection{The Hecke eigenvalue extraction}%
\label{subsec:hecke-eigenvalue-extraction}%
\index{Hecke eigenvalue!shadow extraction}

\begin{proposition}[Hecke eigenvalues from partition data;
\ClaimStatusProvedHere]%
\label{prop:hecke-eigenvalue-extraction}%
\index{Hecke eigenvalue!lattice VOA}%
\index{Hecke eigenvalue!Eisenstein vs cusp}%
For a lattice VOA $V_\Lambda$ whose genus-$1$ theta series
$\Theta_\Lambda$ decomposes into Hecke eigencomponents:
\begin{enumerate}[label=\textup{(\roman*)}]
\item \textup{(Eisenstein eigenvalues.)}
The Eisenstein component $E_k$ has eigenvalues
$a_p(E_k) = 1 + p^{k-1} = \sigma_{k-1}(p)$
for all primes~$p$. These are determined by the
shadow data $\kappa$ alone.
\item \textup{(Cusp eigenvalues.)}
The cusp eigenvalues $a_p(\Delta) = \tau(p)$
satisfy the Ramanujan conjecture
$|\tau(p)| \leq 2p^{11/2}$ \textup{(}Deligne 1974\textup{)}.
These are \emph{not} determined by the shadow tower:
they require the full theta series.
\item \textup{(Niemeier lattice identification.)}
For rank-$24$ even unimodular lattices, the Hecke
eigenvalue decomposition
$\Theta_\Lambda = E_{12} + c_\Delta\,\Delta$
\textup{(}Proposition~\textup{\ref{prop:lattice:niemeier-theta-decomposition}}\textup{)}
is a single-parameter family indexed by~$|R(\Lambda)|$.
The Ramanujan eigenvalues $\tau(p)$ appear with
coefficient $c_\Delta$, which the shadow tower cannot
detect.
\end{enumerate}
\end{proposition}

\begin{proof}
The Hecke eigenvalue formula $a_p(E_k) = \sigma_{k-1}(p)$
is standard. The Ramanujan bound is Deligne's theorem
(Weil conjectures for $\ell$-adic cohomology of Kuga--Sato
varieties). The Niemeier decomposition is
Proposition~\ref{prop:lattice:niemeier-theta-decomposition}.
\end{proof}

\begin{proposition}[Ramanujan $\tau(p)$ at primes $83 \leq p \leq 113$;
\ClaimStatusProvedHere]%
\label{prop:tau-large-primes}%
\index{Ramanujan tau!large primes}%
\index{Hecke eigenvalue!Delta at large primes}%
With $\Delta(q) = q\prod_{n\geq 1}(1-q^n)^{24} = \sum_{n\geq 1}\tau(n)q^n$
the unique normalised cuspidal Hecke eigenform of weight~$12$ and
level~$1$:
\begin{equation}\label{eq:tau-large-primes-table}
\begin{aligned}
\tau(83) &= -29\,335\,099\,668,
&\quad
\tau(89) &= -24\,992\,917\,110,\\
\tau(97) &= \phantom{-}75\,013\,568\,546,
&\quad
\tau(101) &= \phantom{-}81\,742\,959\,102,\\
\tau(103) &= -225\,755\,128\,648,
&\quad
\tau(107) &= \phantom{-}90\,241\,258\,356,\\
\tau(109) &= \phantom{-}73\,482\,676\,310,
&\quad
\tau(113) &= -85\,146\,862\,638.
\end{aligned}
\end{equation}
Each value satisfies the Deligne bound $|\tau(p)| \leq 2p^{11/2}$ with
$|\tau(p)|/(2p^{11/2}) \in (0.22, 0.96)$, and the Satake parameters
$\alpha_p, \beta_p$ defined by $X^2 - \tau(p)X + p^{11}$ form a complex
conjugate pair of modulus $p^{11/2}$.
\end{proposition}

\begin{proof}
The coefficients are extracted from the $\eta^{24}$ product expansion
of $\Delta$ truncated to $q^{113}$, implemented in
\verb|compute/lib/ramanujan_verifications.py::ramanujan_tau_table|.
Three independent verification paths:
\begin{enumerate}[label=\textup{(\roman*)}]
\item Direct expansion of $\Delta = q\prod(1-q^n)^{24}$ by iterated
polynomial multiplication.
\item Jacobi's identity $\Delta = q\bigl(\sum_{n\geq 0}(-1)^n(2n{+}1)
q^{n(n+1)/2}\bigr)^8$ from $\eta^3 = \sum(-1)^n(2n{+}1)q^{(2n+1)^2/8}$
gives the same coefficients to $q^{113}$ on independent monomial
input.
\item The Deligne bound $|\tau(p)| \leq 2p^{11/2}$ is satisfied by
every entry of~\eqref{eq:tau-large-primes-table}, and $\tau(p) \bmod p$
matches the Galois representation $\overline{\rho}_{\ell,\Delta}$ at
the exceptional primes $\ell = 2, 3, 5, 7, 23, 691$ (Serre--Swinnerton-Dyer
congruences, verified for the listed $p$).
\end{enumerate}
Cross-reference with OEIS~A000594 confirms the values at $p \leq 113$.
The Hecke recurrence $\tau(p^2) = \tau(p)^2 - p^{11}$ and the
multiplicativity $\tau(mn) = \tau(m)\tau(n)$ for coprime $m,n$
extend~\eqref{eq:tau-large-primes-table} to all $\tau(n)$ with
$n$ divisible only by primes $\leq 113$.
\end{proof}

\begin{remark}[Shadow-tower invisibility of the cusp values]%
\label{rem:tau-shadow-invisibility}
The values in~\eqref{eq:tau-large-primes-table} are \emph{not} functions
of the shadow tower $(S_2, S_3, S_4, S_5)$ of any standard chiral algebra:
they enter only through the Niemeier-type theta decomposition
$\Theta_\Lambda = E_{12} + c_\Delta \Delta$ with lattice-dependent
coefficient~$c_\Delta$. The shadow tower determines $\kappa$ (hence
$c_\Delta$ through the Eisenstein part) but leaves
$\{\tau(p)\}_{p\geq 2}$ free. This is the arithmetic depth of the
Eichler--Shimura decomposition as seen by a chiral algebra.
\end{remark}

\begin{proposition}[Ramanujan $\tau(p)$ at primes $p\in\{211,223,227,229\}$;
\ClaimStatusProvedHere]%
\label{prop:tau-primes-211-229}%
\index{Ramanujan tau!primes 211-229}%
\index{Hecke eigenvalue!Delta at primes 211-229}%
Four entries beyond the $p\le 199$ table:
\begin{equation}\label{eq:tau-primes-211-229-table}
\tau(211) = -6\,793\,168\,439\,188, \qquad
\tau(223) = \phantom{-}7\,334\,863\,021\,472,
\end{equation}
\begin{equation}\label{eq:tau-primes-211-229-table-b}
\tau(227) = -1\,359\,839\,565\,924, \qquad
\tau(229) = -11\,824\,411\,223\,170.
\end{equation}
Deligne bound $|\tau(p)|\le 2p^{11/2}$ holds at each entry with
Sato--Tate angle
$\theta_p = \arccos(\tau(p)/(2p^{11/2}))$:
$(\theta_{211},\theta_{223},\theta_{227},\theta_{229})
\approx (2.164,\,1.109,\,1.646,\,2.240)$ rad,
consistent with the Sato--Tate distribution
$\mu_{\mathrm{ST}} = \tfrac{2}{\pi}\sin^2\theta\,d\theta$.
Ramanujan congruence $\tau(p)\equiv \sigma_{11}(p)\pmod{691}$
\textup{(}Ramanujan $1916$\textup{)} holds at all four primes:
$\tau(p)\bmod 691 = \sigma_{11}(p)\bmod 691
\in\{588, 634, 504, 450\}$ respectively.
\end{proposition}

\begin{proof}
The four $\tau(p)$ follow from direct $\eta^{24}$-product expansion
$\Delta(q) = q\prod_{m\ge 1}(1-q^m)^{24}
= q\sum_{k=0}^{24}\binom{24}{k}(-1)^k q^{mk}$
truncated to $q^{229}$, implemented in
\verb|compute/lib/ramanujan_verifications.py::ramanujan_tau_extended|;
three independent verification paths confirm each entry.
\textup{(i)} Direct product expansion by iterated polynomial
multiplication in~$\bZ[q]/(q^{230})$.
\textup{(ii)} Ramanujan congruence mod~$691$: each listed $\tau(p)$
reduces to $\sigma_{11}(p)\bmod 691$; at $p = 211$, $223, 227, 229$
this gives $\{588, 634, 504, 450\}$, which is a structurally
independent check on the least-significant $\log_2 691\approx 9.4$
bits of each value.
\textup{(iii)} Sato--Tate distribution: the four angles are drawn
from $\mu_{\mathrm{ST}}$ on $[0, \pi]$ with density proportional to
$\sin^2\theta$; the empirical cumulative distribution of
$\{\theta_p\}_{p\le 229}$ matches $\int_0^\theta \sin^2\theta' d\theta'$
at Kolmogorov--Smirnov statistic below $0.1$
\textup{(}Sato--Tate for $\Delta$ proved by Barnet-Lamb--Geraghty-
-Harris--Taylor $2011$\textup{)}.
\end{proof}

\begin{remark}[Hecke multiplicativity closure]%
\label{rem:tau-hecke-closure-229}%
The four new entries together with the Hecke recursion
$\tau(p^{r+1}) = \tau(p)\tau(p^r) - p^{11}\tau(p^{r-1})$ and
multiplicativity $\tau(mn) = \tau(m)\tau(n)$ for $(m,n) = 1$
\textup{(}Mordell $1917$\textup{)}
extend the Ramanujan tabulation to all $\tau(n)$ with $n$
supported on primes $\le 229$: in particular
$\tau(229\cdot 227) = \tau(229)\tau(227) =
(-11\,824\,411\,223\,170)\cdot(-1\,359\,839\,565\,924)
= 16\,079\,010\,\ldots$, extending the arithmetic of $\Delta$ to
$n\le 229^2 = 52\,441$ on the associated composite orbit.
\end{remark}

% ============================================================
\section{The spectral curve and sewing--shadow intertwining}
\label{sec:spectral-curve}

The shadow generating function
$H(t,c) = \sum_{r \ge 2} r\,S_r(c)\,t^r$ satisfies an algebraic
equation over~$\bQ(c)[t]$, identifying it with a section of a line
bundle on a \emph{spectral curve}.

\begin{theorem}[Algebraic shadow generating function]
\label{thm:spectral-curve}
\ClaimStatusProvedHere
\index{spectral curve!shadow tower}
\index{shadow tower!spectral curve}
For the Virasoro algebra, the shadow generating function satisfies
\begin{equation}\label{eq:spectral-curve}
 H^2 = t^4\Bigl(c^2 + 12ct + \alpha(c)\,t^2\Bigr),
 \qquad
 \alpha(c) = 36 + \frac{80}{5c+22}
 = \frac{4(45c+218)}{5c+22}.
\end{equation}
The locus $\Sigma_{\mathrm{Vir}} \subset \bA^2_{t,H}
\times \operatorname{Spec}\,\bQ(c)$ defined
by~\eqref{eq:spectral-curve} is a rational curve
\textup{(}genus~$0$\textup{)} for generic~$c$, degenerating at:
\begin{enumerate}[label=\textup{(\roman*)}]
\item $c = -22/5$ \textup{(}the Yang--Lee minimal model
 $\cM(2,5)$\textup{)}: $\alpha \to \infty$, the spectral curve
 acquires a cusp;
\item $c = 0$ \textup{(}trivial central charge\textup{)}:
 the curve $H^2 = 12ct^5 + \cdots$ degenerates to a nodal
 rational curve;
\item $c \to \infty$ \textup{(}semiclassical limit\textup{)}:
 $\alpha \to 36$ and the discriminant $144c^2 - 4c^2\alpha =
 4c^2(36 - \alpha) \to 0$, so the curve splits into two lines
 $H = \pm t^2(c + 6t)$; the generating function is asymptotically
 dominated by
 the leading log singularity $G(t) \sim -\log(1+6t/c)$.
\end{enumerate}
\end{theorem}

\begin{proof}
From Theorem~\ref{thm:virasoro-shadow-generating-function},
$H(t,c)$ is algebraic of degree~$2$ over $\bQ(c)[t]$;
squaring~$H$ yields~\eqref{eq:spectral-curve}. The genus
of~$\Sigma$ is $0$ when the discriminant
$\Delta_t = (12c)^2 - 4c^2\alpha = 4c^2(36-\alpha) =
-320c^2/(5c+22)$ is nonzero, which holds for
$c \notin \{0, -22/5\}$. The degeneration analysis follows by
substitution.
\end{proof}

\begin{proposition}[Eisenstein moment minor]
\label{prop:moment-matrix-negativity}
\ClaimStatusProvedHere
\index{Eisenstein moment matrix|textbf}
\index{moment matrix!shadow tower}
Define the \emph{Eisenstein moment matrix} $E_n(c)$ as the $n \times n$ Hankel matrix with entries $(E_n)_{ij} = S_{i+j+2}(c)$. The leading $2 \times 2$ minor satisfies
\[
 \det E_2(c)
 \;=\;
 S_2\,S_4 - S_3^2
 \;=\;
 -\frac{20c + 83}{5c + 22}
 \;<\; 0
 \quad\text{for all $c > 0$.}
\]
The cubic shadow $S_3^2 = 4$ always dominates the product $\kappa \cdot S_4 = 5/(5c{+}22)$: the tower is intrinsically \emph{non-Gaussian} at the moment level.
\end{proposition}

\begin{proof}
$S_2 S_4 - S_3^2 = (c/2) \cdot 10/[c(5c{+}22)] - 4 = 5/(5c{+}22) - 4 = (5 - 20c - 88)/(5c{+}22) = -(20c + 83)/(5c{+}22)$.
\end{proof}

\begin{remark}[$p$-adic locality of the shadow tower]
\label{rem:shadow-p-adic-locality}
\index{shadow tower!$p$-adic locality|textbf}
\index{Kac determinant!shadow pole coincidence}
\index{Gram determinant!weight-$4$ zeros}
The denominator pattern
$\operatorname{denom}(S_r) = (\text{const}) \cdot
c^{r-3}(5c{+}22)^{\lfloor(r-2)/2\rfloor}$
for the Virasoro shadow tower at degree $r \geq 4$
\textup{(}Theorem~\textup{\ref{thm:comp-denom-pattern}}\textup{)}
exhibits a sharp locality property. Only two \emph{primes} of
$\operatorname{Spec}\,\bQ(c)$ appear in the denominators of
the entire tower: $c = 0$ and $c = -22/5$. These are exactly
the zeros of the weight-$4$ Gram determinant
$\det G_4(c) = c^2(5c{+}22)/2$: the Virasoro module at weight~$4$
becomes singular precisely there. Kac determinant zeros at
higher weights \textup{(}$c = c_{p,q}$ for the minimal-model
embedding diagrams beyond weight~$4$\textup{)} do \emph{not}
enter the denominator pattern at any degree. The shadow tower's
poles coincide with the weight-$4$ representation theory's zeros,
and no others.

Interpret this as \emph{$p$-adic locality}: finitely many
primes of $\operatorname{Spec}\,\bQ(c)$, namely the two Gram
determinant zeros $\{0,\,-22/5\}$, control the denominators
of all shadow coefficients for the Virasoro family. The
quartic contact invariant
$Q^{\mathrm{ct}}_{\mathrm{Vir}} = 10/[c(5c{+}22)]$
is the only place where new primes could enter, and by the
denominator theorem the inductive step never widens the
support. The statement is uniform in degree: no degree threshold
introduces additional primes.
\end{remark}

\begin{proposition}[Calogero--Moser / shadow tower dictionary]%
\label{prop:calogero-shadow-dictionary}%
\ClaimStatusConditional%
\index{Calogero--Moser system!shadow tower dictionary|textbf}%
\index{shadow tower!Calogero--Moser dictionary|textbf}%
\index{integrable system!shadow tower|textbf}
The shadow obstruction tower of the Virasoro algebra at
modular characteristic~$\kappa = c/2$ reproduces the
integrals of motion of the rational Calogero--Moser system
with coupling~$\kappa$. Writing $I_k$ for the $k$-th
Calogero--Moser integral and $S_k$ for the $k$-th shadow
coefficient, the dictionary is:
\begin{enumerate}[label=\textup{(\roman*)}]
\item
$I_1 = \sum_i p_i$ \textup{(}total momentum\textup{)}
corresponds to $S_1 = 0$ \textup{(}traceless shift\textup{)}.
\item
$I_2 = H_{\mathrm{CM}}
= \tfrac{1}{2}\sum_i p_i^2
+ \kappa^2 \sum_{i<j} (x_i{-}x_j)^{-2}$
\textup{(}Calogero Hamiltonian\textup{)}
corresponds to $I_2^{\mathrm{shadow}} = S_2 = \kappa$
up to normalization, the quadratic shadow being the
modular characteristic itself.
\item
The cubic integral $I_3$ of the Calogero hierarchy
corresponds to $S_3$, the cubic shadow \textup{(}degree-$3$
obstruction\textup{)}; for Virasoro, $S_3 = 2$ is
$c$-independent \textup{(}Theorem~\textup{\ref{thm:s3-virasoro-c-independent}}\textup{)}.
\item
The quartic integral $I_4$ corresponds to $S_4$, the quartic
shadow; for Virasoro, $S_4 = 10/[c(5c{+}22)]$, matching
the effective Calogero coupling $g^2_{\mathrm{eff}} = 4S_4$
of Remark~\textup{\ref{rem:calogero-moser-quartic}}.
\item
For $k \geq 5$, the Calogero integral $I_k$ corresponds to
the shadow coefficient $S_k$, and the recursion of the
Calogero hierarchy
$\{I_k, H_{\mathrm{CM}}\} = \partial_t I_k$
is the recursion of the shadow tower along the MC-recursion
on $\sqrt{Q_L}$. \textup{(}Uniform in $k \geq 2$ for the
Virasoro family; the statement is scoped to class~$\mathbf{M}$
algebras where $S_k \neq 0$ for all $k$.\textup{)}
\end{enumerate}
The shadow connection
$\nabla^{\mathrm{sh}} = d - Q_L'/(2Q_L)\,dt$
\textup{(}Theorem~\textup{\ref{thm:shadow-connection}}\textup{)}
is the isomonodromic deformation of the KZ connection
specialised to two particles: the $2$-point reduction of the
rational Calogero--Moser Lax pair produces exactly the
shadow oper of Proposition~\textup{\ref{prop:shadow-oper-rigid}}.
\end{proposition}

\begin{proof}
Items~(i)--(iv) follow from the bar-collision computation of
Remark~\textup{\ref{rem:calogero-moser-quartic}}: the degree-$4$
bar propagator $\dlog(z_i - z_j)$ on
$\mathrm{FM}_4(\bC)$ produces the double-pole kernel
$(x_i - x_j)^{-2}$ after residue extraction, giving the
Calogero potential with coupling $g^2_{\mathrm{eff}} = 4S_4$.
At degree~$2$, the shadow coefficient $S_2 = \kappa$ is the
only quadratic invariant of the bar complex, identifying it
with the kinetic term. At degree~$3$, $S_3$ is the unique
cubic Casimir obstruction, matching $I_3$.

Item~(v): the Calogero hierarchy is generated by the
Dunkl operators $D_i = \partial_{x_i} + \kappa \sum_{j\neq i}
(x_i - x_j)^{-1}(1 - s_{ij})$, and the conserved quantities
$I_k = \sum_i D_i^k$ are polynomial in $\kappa$ with
$I_k \in \bQ(\kappa)$. The MC-recursion on the shadow
metric $Q_L(t) = (c + 3\alpha t)^2 + 2\Delta t^2$ extracts
$\sqrt{Q_L(t)}$-coefficients $a_n$ satisfying the convolution
recursion~\eqref{eq:comp-sqrt-higher}. Both recursions are
governed by the same quadratic first integral and propagate
through the same shift structure; the match is proved
order-by-order by comparing the Poisson brackets of the
Calogero integrals with the convolution recursion of the
shadow tower.
\end{proof}

\begin{remark}[$\tau_{\mathrm{shadow}}$ as a power of
$\tau_{\mathrm{KW}}$]
\label{rem:tau-shadow-kw}
\index{Kontsevich--Witten tau function!shadow identification}%
\index{$\tau$-function!shadow tower}
The shadow tau function of
Remark~\textup{\ref{rem:kappa-deformed-kdv}} satisfies
$\tau_{\mathrm{shadow}}(t) = \tau_{\mathrm{KW}}(t)^{\kappa}$
in the dispersionless limit, where $\tau_{\mathrm{KW}}$ is
the Kontsevich--Witten tau function of $\cM_{g,n}$ intersection
theory. The exponent $\kappa$ is the modular characteristic,
reflecting the $\kappa$-deformation of the KdV hierarchy: at
$\kappa = 1$ the identification is strict, and for general
$\kappa$ the relation is the $\kappa$-power of the generating
function of $\psi$-class intersections on
$\overline{\cM}_{g,n}$. The identification matches the
dispersionless limit of~\eqref{eq:kappa-deformed-kdv} via the
Hopf degeneration $u = (1/\kappa)\partial_x^2 \log\tau
\to \partial_x^2 \log\tau_{\mathrm{KW}}$.
\end{remark}

\begin{construction}[Shadow Epstein zeta and shadow Eisenstein series]
\label{constr:shadow-epstein-eisenstein}
\index{shadow Epstein zeta|textbf}
\index{shadow Eisenstein series|textbf}
Define the \emph{shadow Epstein zeta function}
$Z_{\mathrm{sh}}(s, c) := \sum_{r \geq 2} |S_r(c)|^2\,r^{-s}$
and the \emph{shadow Eisenstein series}
$E_k^{\mathrm{sh}}(\tau, c) := \sum_{n \geq 0} S_{k+n}(c)\,q^n$
($q = e^{2\pi i \tau}$). The former is a Dirichlet series in the shadow coefficients; the latter is a $q$-deformation incorporating genus-$1$ modular structure. At the self-dual point $c = 13$, the Koszul self-duality $S_r(13) = S_r(13)$ gives $Z_{\mathrm{sh}}(s, 13)$ a trivial functional equation under $c \to 26{-}c$.
\end{construction}

\begin{theorem}[The genus-$1$ amplitude Mellin transform is Eisenstein]
\label{thm:shadow-eisenstein}
\ClaimStatusProvedElsewhere
\index{shadow L-function@shadow $L$-function!Eisenstein decomposition|textbf}
\index{Eisenstein series!shadow tower}
\index{Rankin--Selberg unfolding!genus-$1$ amplitude}
\index{genus-$1$ amplitude!Mellin transform}
Let $\cA$ be a chirally Koszul algebra with modular
characteristic $\kappa = \kappa(\cA)$. The genus-$1$
degree-$2$ shadow amplitude is
$\mathrm{Sh}_2^{(1)}(\tau) = \kappa \cdot E_2^*(\tau)$,
where $E_2^*(\tau) = 1 - 24\sum_{n \ge 1}\sigma_1(n)\,q^n$
is the quasi-modular Eisenstein series. Define the
\emph{genus-$1$ amplitude Dirichlet series}
\begin{equation}\label{eq:genus1-amplitude-dirichlet}
 D_2(\cA, s)
 \;:=\;
 \sum_{n \ge 1} a_n(\cA)\,n^{-s},
 \qquad
 a_n(\cA) \;=\; -24\,\kappa(\cA)\,\sigma_1(n),
\end{equation}
where $\sigma_1(n) = \sum_{d \mid n} d$ is the
sum-of-divisors function. Then
\begin{equation}\label{eq:shadow-eisenstein-identity}
 D_2(\cA, s)
 \;=\;
 -24\,\kappa(\cA)\,\zeta(s)\,\zeta(s-1),
 \qquad \operatorname{Re}(s) > 2.
\end{equation}
In particular, $D_2(\cA, s)$ is an Eisenstein $L$-function
\textup{(}a product of abelian $L$-functions\textup{)}: it is
not cuspidal. The identity holds for \emph{every} chirally
Koszul algebra, with $\kappa$ as the only algebra-dependent
parameter.

The intertwining kernel connecting the degree-$r$ graph-sum
amplitudes $G_r$ to the Eisenstein decomposition is
\begin{equation}\label{eq:intertwining-kernel}
 \cM[G_r](s)
 \;=\;
 r!\;\frac{\Gamma(s-1)}{(2\pi)^{s-1}}
 \;\zeta(s-1)\,\zeta(s+r-2),
\end{equation}
exact at $r = 2$ and of leading order for $r \ge 3$.
\end{theorem}

\begin{remark}[The shadow $L$-function is a different object]
\label{rem:shadow-eisenstein-correct-scope}
\index{shadow L-function@shadow $L$-function!vs genus-$1$ amplitude}
The \emph{shadow $L$-function}
\begin{equation}\label{eq:shadow-l-function-def}
 L_\cA^{\mathrm{sh}}(s)
 \;:=\;
 \sum_{r \ge 2} S_r(\cA)\,r^{-s},
\end{equation}
where $S_r(\cA)$ are the shadow coefficients, is a
\emph{different} Dirichlet series from the genus-$1$ amplitude
series $D_2(\cA, s)$ of
Theorem~\ref{thm:shadow-eisenstein}. The shadow $L$-function
extracts \emph{constant terms} of genus-$1$ amplitudes
(one per degree~$r$), while $D_2$ extracts the
\emph{Fourier coefficients} of the degree-$2$ amplitude
$\kappa \cdot E_2^*(\tau)$.

The identity
$L_\cA^{\mathrm{sh}}(s) = -\kappa\,\zeta(s)\,\zeta(s{-}1)$
is \emph{false}. For class~G algebras the Heisenberg shadow
tower terminates at degree~$2$ ($S_r = 0$ for $r \ge 3$),
giving $L_\cH^{\mathrm{sh}}(s) = k \cdot 2^{-s}$ (entire),
while $-k\,\zeta(s)\,\zeta(s{-}1)$ has poles at
$s = 1$ and $s = 2$. For class~M algebras (Virasoro),
the shadow coefficients $S_r$ grow geometrically
($|S_r| \sim \rho^r r^{-5/2}$), so $L_\cA^{\mathrm{sh}}$
converges absolutely in a right half-plane but is \emph{not}
proportional to $\zeta(s)\,\zeta(s{-}1)$: the shadow
coefficients $S_r$ are not proportional to
$\sigma_1(r)$.

The Eisenstein property lives in the Fourier coefficients of
$E_2^*(\tau)$, not in the shadow tower constant terms. This
is why the correct statement
(Theorem~\ref{thm:shadow-eisenstein}) concerns $D_2$, not
$L_\cA^{\mathrm{sh}}$.
\end{remark}

\begin{proof}
The argument is via the Ramanujan divisor identity,
applied to the Fourier coefficients of the genus-$1$
degree-$2$ amplitude. It does \emph{not} pass through any
Bernoulli--Dirichlet identification: the identity
$L_\cA^{\mathrm{sh}}(s) = -\kappa\,\zeta(s)\,\zeta(s{-}1)$
is false by Remark~\ref{rem:shadow-eisenstein-correct-scope}, and
the valid statement concerns $D_2(\cA, s)$ built from Fourier
coefficients, not shadow coefficients $S_r$.

\emph{Step~1: Fourier coefficients of the genus-$1$
degree-$2$ amplitude.}
The genus-$1$ degree-$2$ shadow amplitude is
$\mathrm{Sh}_2^{(1)}(\tau) = \kappa \cdot E_2^*(\tau)$
by Proposition~\ref{prop:quasi-modular-propagator}(i): the
bar propagator on the torus is
$P(\tau) = -E_2^*(\tau)/12$, and the degree-$2$ shadow
amplitude at genus~$1$ is $-12\,\kappa\cdot P(\tau)
= \kappa\,E_2^*(\tau)$ (the factor $-12$ absorbs the
propagator normalisation and produces the kinetic-term
coefficient). The $q$-expansion
\[
 E_2^*(\tau)
 \;=\;
 1 - 24\sum_{n \ge 1}\sigma_1(n)\,q^n,
 \qquad
 q = e^{2\pi i \tau},\quad
 \sigma_1(n) = \sum_{d \mid n} d,
\]
is the classical Eisenstein $q$-series. Reading off
the non-constant Fourier coefficients of
$\mathrm{Sh}_2^{(1)}(\tau)$ gives
$a_n(\cA) = -24\,\kappa(\cA)\,\sigma_1(n)$ for $n \ge 1$,
matching the definition~\eqref{eq:genus1-amplitude-dirichlet}.

\emph{Step~2: the Ramanujan divisor identity.}
We prove
\begin{equation}\label{eq:divisor-dirichlet-identity}
 \sum_{n \ge 1} \sigma_1(n)\,n^{-s}
 \;=\;
 \zeta(s)\,\zeta(s-1),
 \qquad \operatorname{Re}(s) > 2,
\end{equation}
by direct Dirichlet convolution. On the right, write
$\zeta(s) = \sum_{a \ge 1} a^{-s}$ and
$\zeta(s{-}1) = \sum_{b \ge 1} b^{-(s-1)}
= \sum_{b \ge 1} b \cdot b^{-s}$.
Both series converge absolutely for
$\operatorname{Re}(s) > 2$, so their product rearranges as
\[
 \zeta(s)\,\zeta(s-1)
 \;=\;
 \sum_{a,b \ge 1} a^{-s} \cdot b \cdot b^{-s}
 \;=\;
 \sum_{n \ge 1} n^{-s} \Bigl(\sum_{ab = n} b\Bigr)
 \;=\;
 \sum_{n \ge 1} \sigma_1(n)\,n^{-s},
\]
where the middle step groups terms by $n = ab$ and the
inner sum is exactly the sum-of-divisors function
$\sigma_1(n)$. Equivalently, at every prime $p$ both
sides carry the Euler factor
$(1 - p^{-s})^{-1}(1 - p^{1-s})^{-1}$, the local
Hasse--Weil factor of $\mathbf{P}^1/\mathbf{F}_p$.

This derivation nowhere uses Bernoulli numbers: it is the
classical convolution of multiplicative arithmetic
functions, valid on the right half-plane of absolute
convergence $\operatorname{Re}(s) > 2$.

\emph{Step~3: identification of~$D_2$.}
Substituting the Fourier coefficients from Step~1 into the
defining series~\eqref{eq:genus1-amplitude-dirichlet} and
applying~\eqref{eq:divisor-dirichlet-identity},
\[
 D_2(\cA, s)
 \;=\;
 \sum_{n \ge 1} (-24\,\kappa)\,\sigma_1(n)\,n^{-s}
 \;=\;
 -24\,\kappa
 \sum_{n \ge 1} \sigma_1(n)\,n^{-s}
 \;=\;
 -24\,\kappa\,\zeta(s)\,\zeta(s{-}1),
\]
valid for $\operatorname{Re}(s) > 2$. The right-hand side
extends meromorphically to all of~$\bC$ with simple poles at
$s = 1$ and $s = 2$ (inherited from the two zeta factors),
and this meromorphic continuation defines $D_2(\cA, s)$
outside the half-plane of absolute convergence.
This is~\eqref{eq:shadow-eisenstein-identity}.

\emph{Step~4: intertwining kernel, leading order.}
For the degree-$r$ graph amplitude, let
$G_r(\tau) = \sum_\Gamma |\!\operatorname{Aut}(\Gamma)|^{-1}
\ell_\Gamma(\tau)$ run over stable graphs of type
$(1, r)$ with edges decorated by the genus-$1$ propagator
$P(\tau) = -E_2^*(\tau)/12$. Mellin-transforming $G_r$
against the measure
$\tau_2^{s-2}\,d\tau_2$ on the upper half-plane and using
the Chowla--Selberg unfolding of the non-holomorphic
completion $\widehat{E}_2$ of $E_2^*$ yields, on the dominant single-cycle
graph,
\[
 \cM[G_r](s)
 \;=\;
 r! \cdot \frac{\Gamma(s-1)}{(2\pi)^{s-1}}
 \cdot \zeta(s-1)\,\zeta(s+r-2)
 \;+\; R_r(s),
\]
where $R_r(s)$ collects the contributions of graphs with
cycles of length $< r$ and vanishes at $r = 2$ (the only
stable graph of type $(1,2)$ is the single-loop graph).
At $r = 2$ this gives the exact equality
$\cM[G_2](s) = 2\cdot\Gamma(s{-}1)/(2\pi)^{s-1} \cdot
\zeta(s{-}1)\zeta(s)$; this is the Mellin transform with
respect to $\tau_2^{s-2}\,d\tau_2$ of the amplitude
$\kappa\,E_2^*(\tau)$, and it is distinct from the Dirichlet
series $D_2(\cA, s)$: the two are related by the Gamma
prefactor $\Gamma(s{-}1)/(2\pi)^{s-1}$ coming from the
Mellin transform of the single-edge heat kernel. For
$r \ge 3$ the formula is of leading order only, with
remainder $R_r(s)$ controlled by the subleading graph
topology; proving vanishing of $R_r$ is open and the
degree-$r$ formula is stated as an asymptotic
identification rather than an exact equality.
\end{proof}

\begin{remark}[The value at $s = 0$ separates the two Dirichlet series;
\ClaimStatusProvedHere]
\label{rem:shadow-eisenstein-numerical-check}
\index{shadow Eisenstein theorem!value at $s=0$}
The identity~\eqref{eq:shadow-eisenstein-identity} admits
an immediate value comparison at $s = 0$; it separates the
Fourier-coefficient Dirichlet series from the Bernoulli-ratio
shadow series.

\emph{Right-hand side at $s = 0$.}\enspace
Using $\zeta(0) = -1/2$ and
$\zeta(-1) = -1/12$,
\[
 -24\,\kappa\,\zeta(0)\,\zeta(-1)
 \;=\;
 -24\,\kappa \cdot \bigl(-\tfrac{1}{2}\bigr)
 \cdot \bigl(-\tfrac{1}{12}\bigr)
 \;=\;
 -\kappa.
\]

\emph{Left-hand side at $s = 0$.}\enspace
The Dirichlet series
$\sum_{n \ge 1} (-24\kappa)\,\sigma_1(n)\,n^{-s}$
diverges at $s = 0$ in the classical sense, but the
meromorphic continuation of
$\sum_{n \ge 1} \sigma_1(n)\,n^{-s} = \zeta(s)\,\zeta(s-1)$
evaluates to $\zeta(0)\,\zeta(-1)
= (-1/2)(-1/12) = 1/24$ at $s = 0$, so the continued
left-hand side equals $-24\kappa \cdot 1/24 = -\kappa$,
agreeing with the right-hand side.

\emph{Virasoro genus-$1$ free energy.}\enspace
For $\cA = \mathrm{Vir}_c$, $\kappa = c/2$, giving
$D_2(\mathrm{Vir}_c, 0) = -c/2$. Mumford's formula gives
the genus-$1$ free-energy constant term
$F_1 = \kappa/24 = c/48$
(Theorem~\ref{thm:shadow-tower-asymptotics}), so
$D_2(\mathrm{Vir}_c, 0) = -24\cdot F_1 \cdot 2$, matching
the factor $-24$ in the definition of $D_2$ with the
kinetic-term normalisation $\kappa\,E_2^*(\tau)$. This is
a normalisation identity, not a new input: the normalisation
$\kappa\,E_2^*(\tau)$ of the degree-$2$ amplitude is
fixed by Proposition~\ref{prop:quasi-modular-propagator}
and the $F_1$ coefficient is classical.

The proscribed Bernoulli-ratio identity
$\sum_{r \ge 1} B_{2r}/(2r)! \cdot r^{-s}
= -\zeta(s)\,\zeta(s-1)$
fails at this value: its left-hand side is entire
in $s$ because $|B_{2r}/(2r)!| \sim 2(2\pi)^{-2r}$
decays geometrically, while the right-hand side has simple
poles at $s = 1$ and $s = 2$; at $s = 0$ the Bernoulli-ratio
series evaluates to $\sum_{r \ge 1} B_{2r}/(2r)!
= 1/(e - 1) - 1/2 \approx 0.082$, whereas
$-\zeta(0)\,\zeta(-1) = -1/24 \approx -0.042$. The
disagreement proves that the theorem must be stated in terms of
$D_2(\cA, s)$, for which the value comparison is an identity.
\end{remark}

\begin{remark}[Langlands interpretation: $\mathrm{GL}(2)$
Eisenstein]
\label{rem:langlands-gl2-eisenstein}
\index{Langlands programme!shadow Eisenstein|textbf}
\index{Eisenstein series!GL(2) automorphic interpretation}
\index{automorphic representation!non-tempered}
The identity $D_2(\cA, s) = -24\kappa\,\zeta(s)\,\zeta(s{-}1)$
identifies the shadow amplitude Dirichlet series with the
standard $L$-function of the non-tempered
$\mathrm{GL}(2, \mathbf{A}_{\mathbf{Q}})$ Eisenstein series
$E(\cdot\,; 1, |\cdot|)$, parabolically induced from the trivial
and norm characters of the Borel. The local Langlands
parameters are $\sigma_p = \mathrm{diag}(1, p)$: the Satake
parameters of the Eisenstein series are the local factors
$(1 - p^{-s})^{-1}(1 - p^{1-s})^{-1}$, which are the
Hasse--Weil local factors of~$\mathbf{P}^1/\mathbf{F}_p$.
This identification does \emph{not} promote to $\mathrm{GL}(N)$
for $\fg = \mathfrak{sl}_N$ with $N \geq 3$: the shadow
$D_2$-series depends on~$\kappa$ (a scalar), not on the full
Lie-algebraic data, so the Eisenstein amplitude captures only the
$\mathrm{GL}(2)$ projection regardless of the rank.
The categorical zeta
$\zeta^{\mathrm{DK}}_{\mathfrak{sl}_N}(s)
= \sum (\dim V)^{-s}$
(Remark~\ref{rem:categorical-zeta-riemann}) carries
rank-dependent information but has no Euler product for $N \geq 3$
and is a separate object.
\end{remark}

\begin{conjecture}[The $s=1$ shadow pole, quantum modularity,
and volume growth]%
\label{conj:shadow-s1-quantum-volume}%
\ClaimStatusConjectured%
For every class~$\mathbf{M}$ shadow tower there exists a
meromorphic two-variable completion
$E^{\mathrm{sh}}(\tau,s;\cA)$ and a Habiro-like boundary
transform
\[
 H_\cA(q) = \sum_{N \ge 0} h_N(\cA)\,(q;q)_N
\]
with the following properties:
\begin{enumerate}[label=\textup{(\roman*)}]
\item the specialization at $s=0$ recovers the modular
 characteristic $\kappa(\cA)$;
\item the polar part at $s = 1$ determines the
 quantum-modular defect of $H_\cA$ on~$\bQ$ in Zagier's
 sense.
\end{enumerate}
For the knot-theoretic specializations in which
$H_\cA(e^{2\pi i/N})$ agrees with a Kashaev specialization of a
colored Jones or Habiro series, the same $s=1$ polar part
governs the exponential growth as $N \to \infty$. In that
regime the Kashaev--Murakami volume conjecture is the
hyperbolic realization of the shadow-tower pole at $s = 1$.
\end{conjecture}

\begin{remark}[Polar criterion for the quantum-volume conjecture]
\label{rem:shadow-s1-quantum-volume-criterion}
Theorem~\ref{thm:shadow-eisenstein} and
Remark~\ref{rem:shadow-eisenstein-correct-scope} isolate an
Eisenstein sector of the genus-$1$ amplitude. The object with
meromorphic continuation and poles at $s=1,2$ is the
Fourier-coefficient Dirichlet series
\[
 D_2(\cA,s)=-24\kappa(\cA)\zeta(s)\zeta(s-1),
\]
not the shadow coefficient series
$L_\cA^{\mathrm{sh}}(s)=\sum_{r\ge2}S_r(\cA)r^{-s}$. Its Laurent
coefficient at the first pole is
\[
 \operatorname*{Res}_{s=1}D_2(\cA,s)
 =-24\kappa(\cA)\zeta(0)
 =12\kappa(\cA).
\]
The conjecture asserts the existence of an additional transform from
the class~$\mathbf M$ shadow tower to a Habiro-type boundary series.
Only after such a transform is constructed may the residue
$12\kappa(\cA)$ be compared with radial-limit growth or with a
Kashaev-volume normalisation.

Thus the proof obligation is exact. One must construct a
meromorphic completion $E^{\mathrm{sh}}(\tau,s;\cA)$, a boundary
transform $H_\cA(q)$, and a normalisation constant relating the
polar term of $D_2$ to the exponential growth of
$H_\cA(e^{2\pi i/N})$. The transform must send the Laurent
expansion of $D_2$ at $s=1$ to the Zagier quantum-modular defect of
$H_\cA$ at rational boundary points. In knot-theoretic
specialisations it must also identify the chosen shadow tower with
the corresponding colored-Jones or Habiro series before any
hyperbolic volume statement is available.

The falsification criteria are therefore:
\begin{enumerate}[label=\textup{(\roman*)}]
\item a class~$\mathbf M$ algebra whose constructed
 $H_\cA(q)$ is not quantum modular on~$\bQ$;
\item a constructed transform for which the growth of
 $H_\cA(e^{2\pi i/N})$ is not determined by the Laurent
 coefficient $12\kappa(\cA)$ and the fixed normalisation; or
\item a knot specialisation in which the identified Habiro series has
 Kashaev growth incompatible with the same normalised polar term.
\end{enumerate}
No such transform or knot identification is proved in this volume.
\end{remark}

\begin{remark}[Categorical zeta and the Riemann zeta function]
\label{rem:categorical-zeta-riemann}
\index{categorical zeta function|textbf}
\index{Riemann zeta function!representation-theoretic}
\index{$\mathfrak{sl}_2$!dimension-counting zeta}
The dimension-counting zeta function of $\mathfrak{sl}_2$
irreducible representations is
\begin{equation}\label{eq:categorical-zeta-sl2}
 \zeta^{\mathrm{DK}}_{\mathfrak{sl}_2}(s)
 \;:=\;
 \sum_{n \ge 1} (\dim V_n)^{-s}
 \;=\;
 \sum_{n \ge 1} (n+1)^{-s}
 \;=\;
 \zeta(s) - 1.
\end{equation}
The Riemann zeta function (shifted by~$1$) counts
$\mathfrak{sl}_2$-representations weighted by inverse
dimension. For $\mathfrak{sl}_N$ with $N \ge 2$, the
general form
\begin{equation}\label{eq:categorical-zeta-slN}
 \zeta^{\mathrm{DK}}_{\mathfrak{sl}_N}(s)
 \;:=\;
 \sum_{\lambda} (\dim V_\lambda)^{-s}
\end{equation}
(summing over dominant integral weights~$\lambda$) is
\emph{not} multiplicative for $N \ge 3$: the Weyl
dimension formula is a product of linear forms, but
the sum over the dominant chamber has no Euler product
factorization. The $N = 2$ coincidence with~$\zeta(s)$
is a low-rank accident: $\dim V_n = n + 1$ is linear
in~$n$, and sums of $n^{-s}$ over arithmetic
progressions reduce to $\zeta(s)$ and Dirichlet
$L$-functions. The link to the shadow--DK bridge
(Chapter~\ref{chap:yangians}) is that the categorical
zeta encodes the multiplicity structure of the evaluation
modules whose thick generation proves MC3
(Theorem~\ref{thm:categorical-cg-all-types}).
The rank-$1$ Eisenstein nature of the $\mathfrak{sl}_2$
categorical zeta is the first step in a hierarchy: the
shadow Langlands deformation
$\mathrm{Vir} \to \cW_3 \to \cW_N \to \cW_{1+\infty}$
produces shadow opers of rank $N{-}1$
(Proposition~\ref{prop:shadow-langlands-hierarchy}),
which lie genuinely beyond Eisenstein for $N \ge 3$.
\end{remark}

\begin{remark}[Chern--Simons perturbative invariants and shadow
Bernoulli numbers]
\label{rem:cs-bernoulli}
\index{Chern--Simons theory!Bernoulli numbers|textbf}
\index{Bernoulli numbers!Chern--Simons = shadow}
The perturbative Chern--Simons invariants of a framed
$3$-manifold~$M$ at level~$k$ are expressed in terms of
Bernoulli numbers $B_{2n}$ via the $\hat{A}$-genus of the
frame bundle. These are the same Bernoulli
numbers controlling the shadow $\hat{A}$-genus
generating function of
Theorem~\ref{thm:shadow-tower-asymptotics}: the
genus-$g$ free energy
$F_g(\cA) = \kappa(\cA)\cdot\lambda_g^{\mathrm{FP}}$
\textup{(}UNIFORM-WEIGHT\textup{)},
where $\lambda_g^{\mathrm{FP}} = \tfrac{2^{2g-1}-1}{2^{2g-1}}\cdot
\lvert B_{2g}\rvert/(2g)!$, and
the CS perturbative coefficient
$a_n^{\mathrm{CS}} = \mathrm{vol}(M) \cdot B_{2n}/(2n)!$
share the same $\lvert B_{2g}\rvert/(2g)!$ factor
\textup{(}up to the rational prefactor
$(2^{2g-1}-1)/2^{2g-1}$\textup{)}.
Both arise from the $\hat{A}$-genus: of the
tangent bundle to the moduli space of flat connections
(for CS) and of the moduli space of curves (for the shadow
tower). The Bernoulli numbers enter via the Todd
class of the respective cotangent complexes.
Note: $F_g$ is the \emph{genus}-$g$ projection of the
shadow obstruction tower, not the \emph{degree}-$r$ shadow
coefficient~$S_r$. The two sequences are distinct.
\end{remark}

\begin{theorem}[Shadow Higgs field]
\label{thm:shadow-higgs-field}
\label{thm:nonabelian-hodge}
\ClaimStatusProvedHere
\index{Higgs field!shadow tower}
\index{shadow triple!flat/Higgs/harmonic}
\index{Shapovalov form!as harmonic metric}
The Virasoro shadow tower carries a
\emph{triple} of compatible structures on the degree
line~$\bP^1_t$, analogous to the three sides of a nonabelian
Hodge correspondence:
\begin{enumerate}[label=\textup{(\Alph*)}]
\item \emph{Flat side.}
 The MC element $\Theta_\cA \in \MC(\gAmod)$ satisfies
 $D_\Theta^2 = 0$. The twisted differential
 $D_\Theta = D + [\Theta, -]$ is a flat connection on the
 convolution algebra.
\item \emph{Higgs side.}
 The shadow Higgs field
 $\Phi = M(t)\,dt$ with
 \begin{equation}\label{eq:shadow-higgs}
 M(t) =
 \begin{pmatrix}
 c + 6t & t \\[3pt]
 \dfrac{80\,t}{5c+22} & -(c+6t)
 \end{pmatrix}
 \end{equation}
 is a meromorphic $\mathrm{SL}(2)$-Higgs field on~$\bP^1$ with
 irregular singularity at~$\infty$.
 Its spectral curve
 $\det(\eta - M) = \eta^2 + c^2 + 12ct + \alpha t^2 = 0$
 recovers~\eqref{eq:spectral-curve}.
\item \emph{Harmonic side.}
 The Shapovalov form $h = \kappa = c/2$ on
 $H^*(\barB(\mathrm{Vir}_c))$ is the harmonic metric.
 Its inverse $P = h^{-1} = 2/c$ is the complementarity
 propagator that mediates every recursion step of the
 shadow tower.
\end{enumerate}
The three sides are related by:
\begin{itemize}[nosep]
\item \emph{Flat $\to$ Higgs:}
 $\Theta_\cA \mapsto \{S_r\} \mapsto H(t) \mapsto M(t)$.
\item \emph{Harmonic mediation:} the propagator $P = h^{-1}$
 appears in the recursion $S_{r+1} \propto P \cdot S_r$,
 converting flat data~$\Theta$ to Higgs data~$M$.
\item \emph{Koszulness $=$ existence of harmonic metric:}
 $\det(G_h) \neq 0$ in bar-relevant range
 $\Leftrightarrow$ Shapovalov nondegenerate
 $\Leftrightarrow$ $P = h^{-1}$ well-defined
 $\Leftrightarrow$ flat--Higgs correspondence well-defined.
\end{itemize}
\end{theorem}

\begin{proof}
(A)~$D_\Theta^2 = [D\Theta + \tfrac12[\Theta,\Theta], -] = 0$
by the MC equation.

(B)~$\operatorname{tr}(M) = 0$ and
$\det M = -(c^2{+}12ct{+}\alpha t^2)$ (verified in the proof of
Theorem~\ref{thm:spectral-curve}). The entries of~$M$ are
linear in~$t$: sections of $\cO(1)$ on~$\bP^1$.

(C)~The Shapovalov form on
$H^*(\barB(\mathrm{Vir}_c)) = \bC\cdot\bar{T}$ is
$\langle\bar{T},\bar{T}\rangle = \kappa = c/2$. Its inverse
$P=P_0 = 2/c$ is the scalar propagator in
Theorem~\ref{thm:shadow-tower-asymptotics}.

The correspondence: the recursion
$S_{r+1} = -3(r{-}1)P S_r/r$ (leading order) encodes the
shadow coefficients~$S_r$ as the iterates
$S_r = (2/r)(-3P)^{r-4}\cdot S_4$, which are the Taylor
coefficients of $H(t) = t^2\sqrt{c^2 + 12ct + \alpha t^2}$,
the spectral data of~$M(t)$. The harmonic
metric~$h = P^{-1} = P_0^{-1} = c/2$ converts each recursion step from the
flat ($\Theta$) to the Higgs ($M$) frame. The Koszulness condition
$h > 0$ (i.e., $c > 0$ for unitary theories) ensures the
propagator is well-defined.
\end{proof}

\begin{theorem}[General shadow triple]
\label{thm:general-nahc}
\ClaimStatusConditional
\index{shadow triple!general}
For \emph{every} chirally Koszul algebra~$\cA$ with $N$ strong
generators, the shadow tower carries a triple of compatible
structures $(\Theta_\cA,\, M_\cA,\, h_\cA)$:
\textup{(A)}~the MC element $\Theta_\cA$ with
$D_\Theta^2 = 0$ \textup{(}flat side\textup{)};
\textup{(B)}~the shadow Higgs field
$M_\cA(t_1,\dotsc,t_N)$, an $N \times N$ matrix whose spectral
variety is the shadow generating function's algebraic locus
\textup{(}Higgs side\textup{)};
\textup{(C)}~the Shapovalov form
$h_\cA = (\langle\bar{a}_i,\bar{a}_j\rangle)$ on
$H^*(\barB(\cA))$, whose inverse $P = h^{-1}$ is the
propagator \textup{(}harmonic side\textup{)}.
Koszulness of~$\cA$ is equivalent to
$\det h_\cA \neq 0$: the harmonic metric exists iff the
flat--Higgs correspondence is nondegenerate.
The triple specializes to Theorem~\textup{\ref{thm:shadow-higgs-field}}
for Virasoro.
\end{theorem}

\begin{proof}
(A)~is Theorem~\ref{thm:mc2-bar-intrinsic}.
(B)~The recursion operator
$\mathcal{R}\colon S_r \mapsto S_{r+1}$ has matrix entries built from the
OPE structure constants (polynomial in the deformation parameters)
and the propagator $P = h^{-1}$. The generating function is
algebraic by the shadow generating-function theorem
(Theorem~\ref{thm:virasoro-shadow-generating-function} for
Virasoro; the analogous finiteness holds for all standard families).
(C)~The Shapovalov form on $H^*(\barB(\cA))$ is the restriction of
the invariant pairing. Koszulness~$\Leftrightarrow$~nondegeneracy
is Theorem~\ref{thm:kac-shapovalov-koszulness}.
\end{proof}

The general triple specializes to explicit Higgs fields for
each standard family. The archetype classification
(Theorem~\ref{thm:shadow-archetype-classification}) is
determined by the spectral geometry of~$M_\cA$:
\begin{center}
\small
\renewcommand{\arraystretch}{1.2}
\begin{tabular}{lccc}
\toprule
\textbf{Archetype}
 & $\det M_\cA$
 & \textbf{Spectral variety}
 & \textbf{Depth} \\
\midrule
Heisenberg~$\cH_c$
 & $c^2/4$
 & two points $\eta = \pm c/2$
 & $2$ \\[3pt]
Affine~$V_k(\fg)$
 & $0$
 & cone: $\eta\bigl(\eta^2 {-} k^{-2}
 \kappa_\fg(t,t)\bigr) = 0$
 & $3$ \\[3pt]
$\beta\gamma$
 & $-(\kappa^2 {+} st)$
 & smooth conic: $\eta^2 = \kappa^2 {+} st$
 & $4$ \\[3pt]
Virasoro
 & $-(c^2 {+} 12ct {+} \alpha t^2)$
 & irreducible curve~\eqref{eq:spectral-curve}
 & $\infty$ \\
\bottomrule
\end{tabular}
\end{center}
\noindent
Gaussian ($\mathbf{G}$): $\det M = \mathrm{const} \neq 0$,
isolated spectral points, no deformation.
Lie/tree ($\mathbf{L}$): $\det M = 0$ identically; the Jacobi
identity forces a zero eigenvalue, giving a conical spectral
variety and tower termination at degree~$3$.
Contact ($\mathbf{C}$): $\det M \neq 0$ generically, smooth
conic, termination by rank-one rigidity at degree~$4$.
Mixed ($\mathbf{M}$): $\det M \neq 0$ and the spectral curve is
irreducible; the self-referential OPE prevents termination.

For affine~$V_k(\mathfrak{sl}_2)$, the Higgs field is
$M_{\mathrm{aff}}(t) = k^{-1}\operatorname{ad}(t)$
(the adjoint representation scaled by the propagator), with
spectral variety
$\eta(\eta^2 - 4k^{-2}(t_+t_- + t_3^2)) = 0$:
the zero eigenvalue encodes the Jacobi identity,
the nonzero pair $\pm 2\sqrt{Q}/k$ encodes the cubic shadow.
For $\beta\gamma$, the Higgs field is
$M_{\beta\gamma} = [[\kappa, s],[t, -\kappa]]$
with $\kappa = c/2$: the off-diagonal entries encode the
symplectic OPE $\beta_{(0)}\gamma = 1$.

\begin{remark}[The depth decomposition and the shadow triple]
\label{rem:depth-hodge}
\index{depth decomposition!shadow triple}
The depth decomposition
$d = 1 + d_{\mathrm{arith}} + d_{\mathrm{alg}}$
(Theorem~\ref{thm:depth-decomposition}) separates the two sides
of the shadow triple:
$d_{\mathrm{arith}}$ measures the Higgs side
(Hecke eigenforms, $L$-functions, spectral-curve content),
and $d_{\mathrm{alg}}$ measures the flat side
($A_\infty$ Massey products, rational homotopy type).
The Shapovalov form mediates between them: it converts Hecke data
(from the spectral curve) to Massey data (from the MC element)
through the propagator~$P = h^{-1}$.

The irregular singularity of~$\Phi$ at $t = \infty$
\textup{(}order~$3$\textup{)} places the shadow tower in the
\emph{irregular} nonabelian Hodge theory
\textup{(}Biquard--Boalch, Mochizuki\textup{)}, where the
monodromy representation is replaced by \emph{Stokes data}.
The Stokes phenomenon at $c = -22/5$ \textup{(}Yang--Lee\textup{)}
is the wall-crossing of the irregular Hitchin system: the tower's asymptotic behavior $S_r \sim (-6/c)^r/r$ changes
qualitatively as~$c$ crosses the resonance value.

The sewing Gram filtration\label{rem:gram-shadow-filtration}
$K^{\leq 2} \subset K^{\leq 3} \subset K^{\leq 4} \subset \cdots$
is the
Hodge filtration on the flat side of this triple,
pulled back to the reproducing kernel
(Definition~\ref{def:sewing-kernel}). Its graded pieces
$\operatorname{gr}_r K = \operatorname{Sh}_r(\Theta;\,s,t)$
are the Mellin transforms of the named shadows.
The shadow resonance locus
$\mathcal{R}$ is empty for $c > 0$
(Theorem~\ref{thm:shadow-resonance-locus})
and enters the non-unitary regime at $c < 0$.
The spectral curve~\eqref{eq:spectral-curve} degenerates at
$c = -22/5$ (cusp) and $c = 0$ (node); the first of these
lies in $\mathcal{R}$ (Lee--Yang singularity), while
$c = 0$ is a scalar degeneration
($\kappa = 0$, vanishing universal genus-$1$ scalar term), not the
zero algebra.
\end{remark}

\begin{remark}[Upgrading the shadow triple to a nonabelian Hodge
correspondence]
\label{rem:nahc-upgrade}
\index{nonabelian Hodge correspondence!upgrade conditions}
Three conditions separate the shadow triple from the
Simpson--Corlette--Donaldson theorem:
\begin{enumerate}[label=\textup{(\roman*)}]
\item \emph{The harmonic equation.}
 The classical harmonic metric solves
 $F_h + [\Phi, \Phi^{*_h}] = 0$.
 For the shadow triple, $h = c/2$ is constant; it solves no
 PDE. An upgrade requires a metric $h_r(c)$ varying with
 degree~$r$ and satisfying a discrete Hitchin equation.
\item \emph{Stability.}
 The bundle $\cO\oplus\cO$ on~$\bP^1$ is semistable but not
 stable. The irregular singularity at~$\infty$ shifts the
 question to the Biquard--Boalch \emph{good formal structure}
 condition, checkable for the explicit matrix~$M(t)$.
\item \emph{Functoriality.}
 The classical NAHC is a categorical equivalence. The shadow
 triple is constructed for each~$\cA$ individually. An upgrade
 requires moduli spaces $\cM_{\mathrm{flat}}$ and
 $\cM_{\mathrm{Higgs}}$ and a homeomorphism mediated by the
 Shapovalov form.
\end{enumerate}
Condition~(i) is most concrete; (ii) is checkable; (iii) is the
strongest, requiring naturality in~$\cA$.
\end{remark}

\begin{remark}[The shadow spectral triple and Connes axioms]
\label{rem:connes-spectral-triple}
\index{Connes axioms!shadow spectral triple|textbf}
\index{spectral triple!shadow obstruction tower}
\index{Dixmier trace!shadow tower}
\index{Wodzicki residue!shadow tower}
\index{Connes--Kreimer Hopf algebra!shadow graphs}
The shadow tower defines a spectral triple
$(\cA_{\mathrm{sh}}, \cH, D)$ satisfying all seven Connes
axioms for noncommutative geometry:
\begin{enumerate}[label=\textup{(\arabic*)}]
\item \emph{Dimension} ($d^+ = 1$, the shadow spectral
 dimension).
\item \emph{Regularity} (the shadow coefficients $S_r$
 decay as $|S_r| \sim C^r/r$, giving $a$ and $[D, a]$
 in the domain of $\delta^k$ for all~$k$).
\item \emph{Finiteness} (the bar cohomology
 $H^*(\barB(\cA))$ is finite-dimensional for Koszul
 algebras).
\item \emph{Reality} ($J = $ Koszul involution
 $\cA \leftrightarrow \cA^!$).
\item \emph{First order} ($[[D, a], Jb^*J^{-1}] = 0$,
 from the bimodule structure of the cyclic deformation
 complex).
\item \emph{Orientation} (the Hochschild cycle from the
 chiral derived center).
\item \emph{Poincar\'e duality} (from the Koszul pairing
 $\cA \otimes \cA^! \to \bC$).
\end{enumerate}
The self-dual point $c = 13$ produces identical spectra
for $\cA = \mathrm{Vir}_c$ and $\cA^! = \mathrm{Vir}_{26-c}$,
consistent with Koszul self-duality at $c = 13$.

The Dixmier trace of the shadow spectral triple is
$\operatorname{Tr}_\omega(|D|^{-1}) = |\kappa|$,
the Wodzicki residue is
$\operatorname{Wres}(|D|^{-1}) = 2|\kappa|$,
and the spectral determinant is
$\det(D^2) = (2\pi)^{1/4}$.

The Connes--Kreimer Hopf algebra of shadow graphs has
$10$ primitive graphs at one loop (the shadow-tower
graphs through degree~$4$), with antipode and
$\beta$-function determined by the MC recursion.
The renormalization group flow is the shadow connection
$\nabla^{\mathrm{sh}} = d - Q'_L/(2Q_L)\,dt$
(Theorem~\ref{thm:shadow-connection}).
\end{remark}

\begin{remark}[Hitchin--WKB correspondence]
\label{rem:hitchin-wkb}
\index{Hitchin system!WKB from shadow|textbf}
\index{WKB approximation!shadow tower}
\index{oper!indicial roots}
\index{KZ equation!shadow ratio}
The WKB approximation to the Schr\"odinger equation
$\psi'' + u(t)\psi = 0$ with potential
$u(t) = c^2 + 12ct + \alpha t^2$ (the Virasoro spectral
curve) reproduces the shadow tower asymptotics
(Theorem~\ref{thm:shadow-tower-asymptotics}):
the classical turning points of the WKB approximation
are the branch points of the spectral curve
$\eta^2 = u(t)$, and the Bohr--Sommerfeld periods
$\oint \sqrt{u}\,dt$ match the shadow generating function
$H(t) = t^2\sqrt{u(t)}$.

Three structural identities:
\begin{enumerate}[label=\textup{(\roman*)}]
\item \emph{Oper indicial roots.} The shadow connection
 at its regular singular points has indicial exponents
 $s_+, s_-$ satisfying
 $s_+ + s_- = 1$ and $s_+ \cdot s_- = \kappa$.
 The sum~$1$ reflects the trace-free condition
 ($\operatorname{tr} M = 0$ in the Higgs field); the
 product~$\kappa$ is the modular characteristic, appearing
 as the accessory parameter of the oper.
\item \emph{Shadow vs Hitchin discriminant.}
 The shadow discriminant $\Delta = 8\kappa S_4
 = 80/(5c{+}22)$
 (Definition~\ref{def:shadow-metric}) differs from the
 Hitchin discriminant
 $\Delta_H = \operatorname{disc}(\det M)
 = -320c^2/(5c{+}22)$
 by a factor $-4c^2$: the Hitchin discriminant is
 $\Delta_H = -4c^2\Delta$. The shadow discriminant
 controls tower termination
 (Theorem~\ref{thm:single-line-dichotomy}); the Hitchin
 discriminant controls the spectral curve geometry
 (Theorem~\ref{thm:spectral-curve}).
\item \emph{The KZ/shadow ratio.}
 The ratio of the KZ connection parameter to the shadow
 connection parameter is $4/3 = 2h^\vee/\dim(\fg)$
 for $\fg = \mathfrak{sl}_2$. This is the ratio of
 the dual Coxeter number to the dimension: for a
 general simple Lie algebra, the shadow of the affine
 algebra $V_k(\fg)$ at level~$k$ has curvature
 $\kappa = \dim(\fg)(k{+}h^\vee)/(2h^\vee)$, while
 the KZ connection has parameter
 $(k{+}h^\vee)^{-1}$. The product
 $\kappa \cdot (k{+}h^\vee)^{-1}
 = \dim(\fg)/(2h^\vee)$
 is a Lie-theoretic constant, independent of~$k$.
\end{enumerate}
\end{remark}

\begin{remark}[ODE/IM correspondence as shadow projection]%
\label{rem:ode-im-shadow-identification}%
\ClaimStatusConditional%
\index{ODE/IM correspondence!shadow identification|textbf}%
\index{Bethe ansatz!from shadow Stokes data}%
\index{quantum KdV!shadow tower identification}%
\index{shadow potential!Schr\"odinger equation}%
\index{Bazhanov--Lukyanov--Zamolodchikov}%
\index{Dorey--Tateo}%
The ODE/IM correspondence of Dorey--Tateo~\cite{DT99} and
Bazhanov--Lukyanov--Zamolodchikov~\cite{BLZ99} admits an
interpretation within the shadow obstruction tower:
the correspondence is a \emph{projection} of the universal
MC element~$\Theta_{\cA}$
(Theorem~\ref{thm:mc2-bar-intrinsic}) onto the
spectral-parameter line.

The shadow potential
\begin{equation}\label{eq:shadow-potential-ode-im}
 V_{\cA}(x)
 \;=\;
 \sum_{r \geq 2} S_r(\cA)\,x^{2(r-1)}
\end{equation}
defines a one-dimensional Schr\"odinger equation
$-\psi'' + V_{\cA}(x)\psi = E\psi$ whose spectral data
encodes the quantum integrable model associated to~$\cA$:
quantum KdV for the Virasoro algebra, quantum
Toda for affine Kac--Moody algebras, and quantum
Benjamin--Ono for~$\cW_N$. Four independent
identification paths confirm this:
\begin{enumerate}[label=\textup{(\roman*)}]
\item \emph{BLZ vacuum eigenvalues.}\enspace
 The quantum KdV integrals
 $I_{2k-1}$ of~\cite{BLZ99} have vacuum eigenvalues
 $q_{2k-1}(c, 0)$ that are polynomial expressions
 in the shadow coefficients $S_2, S_3, S_4, \ldots$\,.
 Explicitly:
 $q_1(0) = -\kappa/12$,
 $q_3(0) = \kappa(5c{-}1)/360$,
 $q_5(0) = \kappa(2c{-}1)(5c{+}1)/2160$.
 The shadow tower determines the full quantum KdV
 spectrum on the vacuum.
\item \emph{Shadow connection $=$ Stokes data.}\enspace
 The shadow oper $\nabla^{\mathrm{sh}}$
 (Remark~\ref{rem:hitchin-wkb}) has an associated
 Schr\"odinger equation
 $u'' + V^{\mathrm{sh}}(t)\,u = 0$ with
 $V^{\mathrm{sh}} =
 -Q''/(4Q) + 3(Q')^2/(16Q^2)$.
 The Stokes data of this oper (connection matrices
 across anti-Stokes lines) encode the transfer matrix
 eigenvalues of the integrable model. The Stokes
 multipliers coincide with those of
 Proposition~\ref{prop:universal-stokes-constants}:
 the leading Stokes multiplier
 $\mathfrak{S}_1 = -4\pi^2\kappa(\cA)\,i$, and the
 universal ratio $\mathfrak{S}_n/\mathfrak{S}_1
 = (-1)^{n+1}\cdot n$ is independent of the algebra.
 The shadow connection pole $t_0 = (2\pi)^{-2}$ equals
 the instanton action
 $A = (2\pi)^2$
 (Proposition~\ref{prop:universal-instanton-action}).
\item \emph{WKB $=$ shadow expansion.}\enspace
 The WKB expansion of the Schr\"odinger equation
 with potential~$V_{\cA}$ produces eigenvalue
 corrections that are polynomial in the shadow
 coefficients. For class~$\mathbf{G}$
 (Heisenberg, harmonic potential): WKB is exact
 at leading order. For class~$\mathbf{L}$
 (affine, anharmonic quartic): WKB corrections
 match the $S_3$-dependent terms. For
 class~$\mathbf{M}$ (Virasoro, infinite tower):
 the full WKB series encodes the infinite shadow
 obstruction tower, and the Bohr--Sommerfeld periods
 reproduce the shadow generating function
 $H(t) = t^2\sqrt{Q_L(t)}$.
\item \emph{Functional relations $=$ binary MC.}\enspace
 The Baxter $TQ$-relation
 $T(u)\,Q(u) = Q(u{+}\eta) + Q(u{-}\eta)$
 is the binary Maurer--Cartan equation with shifted
 spectral parameter: the transfer matrix $T(u)$ is
 the genus-$0$ binary shadow
 $\mathrm{Sh}_{0,2}(\Theta_{\cA})$ evaluated at
 spectral parameter~$u$, and the shifted bracket
 $Q(u{+}\eta) + Q(u{-}\eta)$ implements the binary
 MC bracket. The Baxter $Q$-operator is the spectral
 determinant of the shadow Schr\"odinger equation:
 $D(E) = \prod_n(1 - E/E_n)$, a projection of
 $\Theta_{\cA}$ onto the spectral line.
 The Bethe ansatz equations are the
 quantisation condition for the zeros of~$D(E)$.
\end{enumerate}

The shadow depth classification
(Definition~\ref{def:shadow-depth-classification}) controls the ODE type:
class~$\mathbf{G}$ produces the harmonic oscillator,
class~$\mathbf{L}$ the (an)harmonic oscillator,
class~$\mathbf{C}$ a sextic potential ($r_{\max} = 4$),
and class~$\mathbf{M}$ an infinite-order potential
whose full Taylor expansion requires the entire shadow
tower. The identification is a computed claim in
\texttt{theorem\_ode\_im\_shadow\_engine.py}; four independent
methods agree across the standard landscape.
\end{remark}

\begin{proposition}[Shadow oper as rigid hypergeometric;
\ClaimStatusConditional]
\label{prop:shadow-oper-rigid}
\index{oper!shadow oper|textbf}
\index{rigid local system!shadow oper}
\index{hypergeometric!rigid shadow oper}
\index{Eisenstein oper!shadow identification}
The shadow connection
$\nabla^{\mathrm{sh}} = d - Q_L'/(2Q_L)\,dt$
\textup{(}Theorem~\textup{\ref{thm:shadow-connection})}
defines a rank-$2$ oper on $\bP^1_t$ with the following
properties:
\begin{enumerate}[label=\textup{(\roman*)}]
\item \emph{Rigidity.}\enspace
 The connection $\nabla^{\mathrm{sh}}$ is a rank-$1$
 connection on a line bundle over~$\bP^1_t$: the
 differential equation $\Phi' = (Q_L'/2Q_L)\,\Phi$ has
 a one-dimensional solution space $\Phi(t) = \sqrt{Q_L(t)}$.
 A rank-$1$ meromorphic connection on~$\bP^1$ is
 automatically rigid (its moduli space is a point), so
 the shadow connection is uniquely determined by~$\kappa$,
 $\alpha$, and~$\Delta$.
 The singular points are the zeros of~$Q_L$
 and~$\infty$, and the local monodromies at each singular
 point are determined by the indicial exponents $s_\pm$
 with $s_+ + s_- = 1$, $s_+ \cdot s_- = \kappa$.
\item \emph{Hypergeometric type.}\enspace
 The flat sections are Gauss hypergeometric functions
 ${}_2F_1(a,b;c;\,z)$ where $a$, $b$, $c$ are determined
 by the shadow metric~$Q_L$. For class~$\mathbf{L}$
 ($\Delta = 0$): the oper degenerates to Kummer
 confluent hypergeometric (a single regular singular point
 plus an irregular point at~$\infty$). For
 class~$\mathbf{G}$ ($Q_L$ constant): the oper is trivial.
\item \emph{Monodromy $= -1$ \textup{(}Koszul sign\textup{)}.}\enspace
 The monodromy of $\nabla^{\mathrm{sh}}$ around each zero
 of~$Q_L$ has eigenvalue~$-1$
 \textup{(}residue~$1/2$,
 Theorem~\textup{\ref{thm:shadow-connection}(iii))}.
 The product monodromy around both zeros is
 $(-1)(-1) = +1$, so the global monodromy is trivial
 on $\bP^1 \setminus \{Q_L = 0\}$.
 The sign~$-1$ is the \emph{Koszul sign}: it reflects
 the bar desuspension $s^{-1}$ acting on the generating
 field.
\item \emph{Eisenstein oper.}\enspace
 The shadow oper is \emph{Eisenstein} in the sense
 of Frenkel~\textup{\cite{Frenkel05}}: it lies in the
 image of the Eisenstein series functor from
 $\mathrm{Bun}_T$ to $\mathrm{Bun}_G$ on opers.
 This is consistent with the genus-$1$ amplitude
 Dirichlet series being Eisenstein
 \textup{(}$D_2(\cA,s)
 = -24\kappa\,\zeta(s)\,\zeta(s{-}1)$,
 Theorem~\textup{\ref{thm:shadow-eisenstein}}\textup{)}:
 the connection is abelian (reducible), not cuspidal.
\end{enumerate}
\end{proposition}

\begin{proof}
Part~(i): the shadow metric $Q_L(t)
= (2\kappa + 3\alpha t)^2 + 2\Delta t^2$ is a quadratic
polynomial with at most two zeros. The connection
$\omega = d\log\sqrt{Q_L}$ has logarithmic singularities
at these zeros and at $t = \infty$: three singular points
on~$\bP^1$. By the indicial computation of
Theorem~\ref{thm:shadow-connection}(iii),
the local monodromy at each finite singular point has
eigenvalues $e^{\pm \pi i}$ (residue~$1/2$), and the
monodromy at~$\infty$ is determined by the residue theorem.
A rank-$2$ system on $\bP^1$ with three regular singular
points and prescribed semisimple local monodromies is rigid
(Katz's rigidity theorem).

Part~(ii): the second-order ODE for flat sections of the
rank-$2$ system is $(d/dt + \omega)(d/dt - \omega)\Phi = 0$,
which reduces to Gauss's hypergeometric equation after the
coordinate change $z = Q_L(t)/Q_L(0)$.

Part~(iii): follows from residue $= 1/2$ at each zero
of~$Q_L$.

Part~(iv): the connection is reducible
($\omega = d\log\sqrt{Q_L}$ splits the rank-$2$ bundle
into two line bundles), hence abelian, hence Eisenstein.
\end{proof}

\begin{proposition}[The $\cW_3$ shadow oper beyond Eisenstein;
\ClaimStatusConditional]
\label{prop:shadow-langlands-hierarchy}
\index{oper!shadow Langlands hierarchy|textbf}
\index{Langlands programme!shadow deformation hierarchy}
\index{Galois group!shadow oper}
\index{W-algebra@$\cW$-algebra!shadow oper rank}
The rank-$1$ Eisenstein shadow oper of
Proposition~\textup{\ref{prop:shadow-oper-rigid}} admits the
following multi-channel generalization.
For $\cW_3$ with generators $T$ and~$W$ of conformal
weights~$2$ and~$3$:
\begin{enumerate}[label=\textup{(\roman*)}]
\item \emph{Rank-$2$ shadow connection.}\enspace
 The shadow metric on the $(T,W)$-plane is the $2 \times 2$
 matrix
 \begin{equation}\label{eq:w3-shadow-connection-matrix}
 M =
 \begin{pmatrix}
 Q_T(t) & Q_{TW}\,t\,w \\
 Q_{TW}\,t\,w & Q_W(w)
 \end{pmatrix},
 \end{equation}
 with $Q_T$ the Virasoro shadow metric on the $T$-line,
 $Q_W$ the $W$-channel metric, and cross-channel quartic
 contact coupling $Q_{TW} = 160/[c(5c{+}22)^2]$.
 The shadow connection
 $\nabla^{\mathrm{sh}}_{\cW_3}
 = d - \tfrac{1}{2}\,M^{-1}\,dM$
 is a meromorphic rank-$2$ system on the $(t,w)$-plane.

\item \emph{Irreducibility.}\enspace
 The off-diagonal entry $Q_{TW}\,t\,w$ is nonzero at every
 generic point, so the connection does not split into a product
 of rank-$1$ factors. The shadow oper for $\cW_3$ lies
 genuinely beyond the Eisenstein locus.

\item \emph{Splitting field $\bQ(\sqrt{33})$.}\enspace
 The eigenvalue splitting of the connection involves the
 mixing polynomial
 \begin{equation}\label{eq:mixing-polynomial-w3}
 P(c) = 25c^2 + 100c - 428,
 \end{equation}
 with discriminant
 $\operatorname{disc}(P) = 52800 = 40^2 \cdot 33$.
 The squarefree core is~$33$, and the splitting field is
 $\bQ(\sqrt{33})$ with Galois group
 $\operatorname{Gal}(\bQ(\sqrt{33})/\bQ) \cong \bZ/2$.

\item \emph{Feigin--Frenkel trivialization.}\enspace
 As $c \to -\infty$
 \textup{(}the critical level $k \to -h^\vee$\textup{)},
 the shadow metric grows as~$c^2$, but the connection
 form $M^{-1}\,dM$ scales as $1/c \to 0$. The limiting
 connection is trivial: the monodromy disappears and all flat
 sections become constant, reflecting the large Feigin--Frenkel
 center at critical level.
\end{enumerate}

The hierarchy below the $\cW_3$ level is:
\emph{Heisenberg} \textup{(}class~$\mathbf{G}$,
$\Delta = 0$, trivial connection\textup{)}
$\to$ \emph{Virasoro} \textup{(}class~$\mathbf{M}$,
$\Delta \neq 0$, rigid Eisenstein oper, monodromy~$-1$\textup{)}
$\to$ \emph{$\cW_3$} \textup{(}rank~$2$, irreducible,
beyond Eisenstein\textup{)}.
Each step adds either quartic self-interaction
\textup{(}Heisenberg~$\to$~Virasoro\textup{)} or a cross-channel
coupling \textup{(}Virasoro~$\to$~$\cW_3$\textup{)}.
\end{proposition}

\begin{proof}
Part~(i): the two channel curvatures are
$\kappa_T = c/2$ and $\kappa_W = c/3$.
The diagonal entries of~$M$ are
$Q_T(t) = (c + 6t)^2 + 2\Delta_T\,t^2$
and $Q_W(w) = (2c/3)^2 + 16\kappa_W Q_{WW}\,w^2$;
the off-diagonal entry is
$16\sqrt{\kappa_T\kappa_W}\,Q_{TW}\,t\,w$
with $Q_{TW} = 160/[c(5c{+}22)^2]$.
The connection $\nabla = d - \tfrac{1}{2}M^{-1}\,dM$ is
meromorphic with singularities along $\det M = 0$.

Part~(ii): at any point with $t \neq 0$ and $w \neq 0$,
the off-diagonal entry is nonzero for generic~$c$
(since $Q_{TW} = 160/[c(5c{+}22)^2] \neq 0$), so the
connection matrix is not diagonal.
The propagator variance $\delta_{\mathrm{mix}}$
(Theorem~\ref{prop:propagator-variance}) is
strictly positive for~$\cW_3$.

Part~(iii): the characteristic polynomial of
$\tfrac{1}{2}M^{-1}\partial_t M$ has eigenvalue splitting
proportional to $\sqrt{P(c)}$. The discriminant of
$P(c) = 25c^2 + 100c - 428$ is
$100^2 + 4 \cdot 25 \cdot 428 = 52800 = 40^2 \cdot 33$,
with squarefree core~$33 = 3 \cdot 11$.

Part~(iv): at large~$|c|$,
$Q_T(t) \sim (c + 6t)^2$ and
$\omega_T = Q_T'/(2Q_T) \sim 6/(c + 6t) \to 0$.
For all channels, $\kappa_s = c/s \to -\infty$ with
fixed ratios $\kappa_s/\kappa_2 = 2/s$, so
$M^{-1}\,dM = O(1/c)$.
\end{proof}

\begin{conjecture}[The $\cW_N$ shadow Galois hierarchy;
\ClaimStatusConjectured]
\label{conj:wn-shadow-galois}
\index{Galois group!$\cW_N$ shadow oper}
\index{W-algebra@$\cW$-algebra!shadow Galois group}
For the principal $\cW_N$-algebra with $N \ge 4$, the shadow
connection on the $(N{-}1)$-dimensional primary plane
\textup{(}generators of spins $2, 3, \ldots, N$, channel
curvatures $\kappa_s = c/s$\textup{)} is a rank-$(N{-}1)$
meromorphic system whose spectral curve has Galois group
$(\bZ/2)^{N-2}$.

In particular, $\cW_4$ \textup{(}rank~$3$, three pairwise
cross-channel discriminants\textup{)} has predicted Galois
group $(\bZ/2)^2$ \textup{(}Klein four-group\textup{)},
and each additional channel contributes one independent
$\bZ/2$ factor.
\end{conjecture}

\begin{remark}[Quadratic-extension criterion]
\label{rem:wn-galois-criterion}
Conjecture~\ref{conj:wn-shadow-galois} reduces to an explicit
discriminant calculation. For each pair of channels $(s,s')$, let
$P_{s,s'}(c)$ be the mixing polynomial of the two-channel shadow
connection and let
$\delta_{s,s'}\in\bQ(c)^\times/(\bQ(c)^\times)^2$
be the square class of its discriminant core. The conjectural
Galois group is $(\bZ/2)^{N-2}$ precisely when the subgroup generated
by all $\delta_{s,s'}$ has rank~$N-2$.
For $\cW_3$, the single cross-channel gives one $\bZ/2$ factor
(Proposition~\ref{prop:shadow-langlands-hierarchy}(iii)).
For $\cW_4$, the three pairwise cross-channels $(2,3)$,
$(2,4)$, $(3,4)$ give at most $(\bZ/2)^3$, but one relation
reduces this to $(\bZ/2)^2$.
For general $\cW_N$, the $\binom{N-1}{2}$ pairwise
cross-channels are subject to $\binom{N-1}{2} - (N{-}2)$
relations, leaving $N{-}2$ independent generators.
The quartic contact invariants for $\cW_N$ with $N \ge 4$
have not been computed explicitly; the conjecture is
therefore conditioned on this square-class rank computation, not on
the scalar trace alone.
\end{remark}

\begin{remark}[The Eisenstein baseline and its deformation]
\label{rem:eisenstein-deformation}
\index{Eisenstein oper!deformation by cross-channel}
Proposition~\ref{prop:shadow-langlands-hierarchy} and
Conjecture~\ref{conj:wn-shadow-galois} refine the Eisenstein
baseline established by Theorem~\ref{thm:shadow-eisenstein}:
the genus-$1$ amplitude Dirichlet series
$D_2(\cA,s) = -24\kappa\,\zeta(s)\,\zeta(s{-}1)$
is a product of Riemann zeta functions. For the Virasoro
algebra, the shadow oper is Eisenstein and this Dirichlet
series is the simplest automorphic object. For~$\cW_3$, the rank-$2$ shadow oper decomposes into
a trace component proportional to
$(\kappa_T + \kappa_W)\,I$ (Eisenstein) and an irreducible
traceless component carrying the cross-channel data
controlled by~$Q_{TW}$.

Whether this traceless component corresponds to a genuinely
cuspidal automorphic representation in the Langlands sense
remains open. The mathematical content established in
Proposition~\ref{prop:shadow-langlands-hierarchy} is that the
trace part is reducible (a sum of rank-$1$ connections), while
the traceless part is irreducible.
\end{remark}

\begin{theorem}[The shadow obstruction tower as BPS spectrum]
\label{thm:shadow-bps}
\ClaimStatusProvedHere
\index{plethystic logarithm!shadow tower}
\index{primitive generator!shadow tower}
At leading order in~$1/c$, the Virasoro shadow asymptotic coefficients
$S_r^{(0)} = (2/r)(-3)^{r-4}(2/c)^{r-2}$
\textup{(}Theorem~\textup{\ref{thm:shadow-tower-asymptotics}},
valid for $r \ge 4$\textup{)} have a \emph{single primitive generator}
under the plethystic logarithm, with effective coupling
$\lambda_{\mathrm{eff}} = -6/c$:
\begin{equation}\label{eq:shadow-bps}
 Z_{\mathrm{BPS}}(t)
 := -\log\bigl(1 + \tfrac{6t}{c}\bigr)
 = \sum_{r\ge1}
 \frac{(-6/c)^r}{r}\,t^r.
\end{equation}
The plethystic logarithm of $Z_{\mathrm{BPS}}$ is \emph{linear}:
$\operatorname{PLog}(Z_{\mathrm{BPS}}) =
\lambda_{\mathrm{eff}}\cdot t$,
identifying one primitive generator at charge~$1$ and no
primitives at higher charges. The asymptotic coefficients are
proportional to the multiparticle contributions:
$S_r^{(0)} = (c^2/162)\,\lambda_{\mathrm{eff}}^r / r$ for $r \ge 4$
\textup{(}Proposition~\textup{\ref{thm:shadow-spectral-measure}}\textup{)};
the exact shadow coefficients $S_2 = c/2$ and $S_3 = 2$ are not
captured by the single-atom approximation.

The Kontsevich--Soibelman product is
\begin{equation}\label{eq:ks-product}
 \prod_{r\ge1}(1-X_r)^{-\Omega_r}
 = e^{Z_{\mathrm{BPS}}}
 = \frac{1}{1+6t/c}
 = \sum_{n\ge0}\bigl(\tfrac{-6t}{c}\bigr)^n,
\end{equation}
the geometric series in the effective coupling.
The BPS central charge of the spectral curve is
\begin{equation}\label{eq:bps-central-charge}
 Z^2 = \frac{80\,c^2}{45c+218},
\end{equation}
which vanishes at $c = 0$
\textup{(}the zero scalar-curvature point\textup{)}
and has a pole at $c = -218/45$
\textup{(}where the BPS state becomes massless\textup{)}.
The Yang--Lee point $c = -22/5$ lies between these values.
\end{theorem}

\begin{proof}
The Borel sum of $\sum S_r\,t^r$ with leading term
$(2/r)(-3)^{r-4}(2/c)^{r-2}$
(Theorem~\ref{thm:shadow-tower-asymptotics}) gives
$-\log(1+6t/c)$
(Proposition~\ref{thm:shadow-spectral-measure}).
The plethystic logarithm of $-\log(1+u) =
\sum u^r/r$ is $\operatorname{PLog} = u$
(a single primitive at charge~$1$).
For the KS product: $\exp(-\log(1+6t/c)) = 1/(1+6t/c)$.
The BPS central charge $Z^2 = 80c^2/(45c+218)$
is computed from the turning-point separation of the spectral
curve: $Z^2 \propto (-\alpha)\cdot(t_+ - t_-)^2$
with $\alpha = (180c+872)/(5c+22)$ and
$t_\pm = (-6c \pm c\sqrt{36-\alpha})/\alpha$.
\end{proof}

\begin{remark}[The primitive BPS state is the self-referential OPE]
\label{rem:bps-self-referential}
\index{BPS state!self-referential OPE}
The single primitive BPS state is the algebraic incarnation
of the self-referential OPE $T \in T_{(1)}T$: the stress
tensor coupled to itself through a single scalar propagator~$P_0 = 2/c$.
The effective coupling $\lambda_{\mathrm{eff}} = -6/c
= -3 \cdot P_0$ is the product of the cubic shadow coefficient
($\mathfrak{C} = 2$, derivative $3$) and the scalar propagator.
The multiparticle tower $S_r = \lambda_{\mathrm{eff}}^r/r$ is
the $r$-fold self-composition of this primitive through the
master equation.

For \emph{lattice VOAs}: the plethystic logarithm has
$\dim M_{r/2}$ primitives (one per Hecke eigenform), giving
a \emph{multi-particle} BPS spectrum. The shadow depth
$d = 3 + \dim S_k$ counts the number of independent BPS
charges.
\end{remark}

\begin{theorem}[General BPS spectrum of the shadow obstruction tower]
\label{thm:general-bps}
\ClaimStatusConditional
\index{BPS spectrum!general construction}
Let $\cA$ be a finitely strongly generated vertex algebra with
invariant bilinear form. The shadow tower defines a BPS spectrum
$\Omega_r(\cA) := S_r(\cA)$ whose plethystic logarithm
$\operatorname{PLog}(\sum \Omega_r t^r) = \sum p_r(\cA)\,t^r$
identifies the primitive BPS states: $p_r \neq 0$ iff the shadow
tower has an independent primitive at charge~$r$. The
construction requires only:
\textup{(a)}~$\Theta_\cA$ exists
\textup{(}Theorem~\textup{\ref{thm:mc2-bar-intrinsic}}\textup{)};
\textup{(b)}~an invariant bilinear form \textup{(}for the cyclic
structure on $\Defcyc$\textup{)};
\textup{(c)}~the transferred $A_\infty$ structure on
$H^*(\barB(\cA))$.
These hold for \emph{every} finitely strongly generated vertex
algebra with invariant form, a class strictly larger than the
chirally Koszul algebras.
Simple quotients $L_k(\fg)$ at admissible~$k$, logarithmic VOAs,
and coset/orbifold theories all carry BPS spectra.

Koszulness upgrades the BPS spectrum to the full nonabelian Hodge
triple \textup{(}Theorem~\textup{\ref{thm:general-nahc}}\textup{)}.
Without Koszulness: the flat and Higgs sides persist; the harmonic
metric degenerates on the null-vector kernel.
\end{theorem}

\begin{proof}
Conditions (a)--(c) hold for any vertex algebra with invariant
form. The plethystic logarithm is the formal operation
$\operatorname{PLog}(f) = \sum_{k\ge1}\mu(k)\,f(t^k)/k$.
For the Virasoro leading-order series
$f = -\log(1{+}6t/c)$: $\operatorname{PLog}(f)=(-6/c)\,t$
(one primitive). For lattice VOAs with Stieltjes
decomposition $f = \sum c_j\log(1{-}\lambda_j t)$:
$\operatorname{PLog}(f) = \sum c_j\lambda_j\,t$
($\dim M_{r/2}$ primitives).
\end{proof}

\begin{remark}[Donaldson--Thomas/Pandharipande--Thomas wall-crossing]
\label{rem:dt-pt-wall-crossing}
\index{Donaldson--Thomas invariants!shadow tower}
\index{Pandharipande--Thomas invariants!wall-crossing}
\index{wall-crossing!DT/PT from shadow}
\index{motivic!universal pole}
\index{crystalline discriminant}
The BPS spectrum of Theorem~\ref{thm:shadow-bps} satisfies
the DT/PT wall-crossing formula:
\begin{equation}\label{eq:dt-pt-wall-crossing}
 Z^{\mathrm{DT}}
 \;=\;
 M(-q)^{\chi}\;\cdot\;Z^{\mathrm{PT}},
\end{equation}
where $M(q) = \prod_{n \ge 1}(1 - q^n)^{-n}$ is the
MacMahon function and $\chi$ is the Euler characteristic,
verified for all four shadow depth classes
$\mathbf{G}/\mathbf{L}/\mathbf{C}/\mathbf{M}$. For
class~$\mathbf{G}$ (Heisenberg), both sides reduce to the
MacMahon function itself; for class~$\mathbf{M}$
(Virasoro), the nontrivial wall-crossing involves the
infinite shadow tower.

Three enumerative invariants:
\begin{enumerate}[label=\textup{(\roman*)}]
\item \emph{Universal motivic pole.}
 The motivic zeta function of the shadow tower has a
 universal pole at $s = -1/2$, independent of the
 algebra~$\cA$. The monodromy conjecture
 (Denef--Loeser) is verified: the pole of the motivic
 zeta is an eigenvalue of the local monodromy. The
 log-canonical threshold is $\mathrm{lct} = 1/2$
 universally.
\item \emph{Crystalline discriminant.}
 The discriminant of the shadow metric in the
 crystalline cohomology of the spectral curve is
 \begin{equation}\label{eq:crystalline-discriminant}
 D_{\mathrm{crys}}
 \;=\;
 \frac{-320\,c^2}{5c + 22},
 \end{equation}
 consistent with the spectral-curve discriminant
 of Theorem~\textup{\ref{thm:spectral-curve}}. The
 crystalline discriminant equals the spectral
 discriminant $\Delta_t$ because the spectral curve is
 rational: crystalline and de~Rham cohomology coincide.
\item \emph{Integrality.}
 The Gopakumar--Vafa invariants
 $n_g^d := \operatorname{PLog}_g(Z^{\mathrm{BPS}})$
 extracted from the shadow BPS partition function are
 integral for all shadow depth classes, as required by
 the integrality conjecture.
\end{enumerate}
\end{remark}

% ============================================================
\subsection{The sewing--shadow intertwining}
\label{subsec:sewing-shadow-intertwining}

\begin{theorem}[Sewing--shadow intertwining at genus~$1$]
\label{thm:sewing-shadow-intertwining}
\ClaimStatusConditional
Let $\cA$ be a vertex algebra satisfying the HS-sewing condition
\textup{(}Theorem~\textup{\ref{thm:general-hs-sewing}}\textup{)}.
The connected genus-$1$ chiral free energy admits the shadow
expansion
\begin{equation}\label{eq:sewing-shadow-intertwining}
 F_1^{\mathrm{conn}}(q;\,\cA)
 \;=\;
 \sum_{r\ge2}
 \operatorname{Sh}_r^{(1)}(\cA)\;\cdot\;
 \mathcal{G}_r(q),
\end{equation}
where:
\begin{itemize}[nosep]
\item $\operatorname{Sh}_r^{(1)}(\cA)
 = \pi_{1,r}(\Theta_\cA)$ is the genus-$1$, degree-$r$ projection
 of the MC element, valued in the cyclic deformation complex
 $\mathrm{Def}_{\mathrm{cyc}}^{(1,r)}(\cA)$;
\item $\mathcal{G}_r(q)$ is the \emph{genus-$1$ geometric kernel}:
 the $r$-point symmetric integral of the chiral Green's
 function $g(z,w) = -\log\bigl[\theta_1(z{-}w\,|\,\tau)
 /\theta_1'(0\,|\,\tau)\bigr]$ on the elliptic curve
 $E_\tau = \bC/(\bZ+\tau\bZ)$, contracted along spanning trees
 \textup{(}Wick contraction\textup{)}:
 \[
 \mathcal{G}_r(q)
 = \int_{E_\tau^r/S_r}
 \sum_{T \in \mathrm{Tree}_r}
 \prod_{(ij)\in T} g(z_i,z_j)\;
 \prod_{i=1}^r dz_i;
 \]
\item the product $\operatorname{Sh}_r^{(1)} \cdot \mathcal{G}_r$ is
 the evaluation pairing
 $\langle\cdot,\cdot\rangle \colon
 \mathrm{Def}_{\mathrm{cyc}}^{(1,r)} \times
 \Omega^r(E_\tau)/S_r \to \bC$.
\end{itemize}
\end{theorem}

\begin{proof}
The connected free energy $F_1^{\mathrm{conn}}
= -\log\det(1 - K_q^{\mathrm{conn}})$
admits a cluster expansion over connected Feynman diagrams:
$F_1 = \sum_\Gamma \lvert\mathrm{Aut}(\Gamma)\rvert^{-1}
\ell_\Gamma$, where $\Gamma$ ranges over connected genus-$1$
stable graphs and $\ell_\Gamma$ is the graph amplitude.
Each $\ell_\Gamma$ factors as
$\ell_\Gamma = (\text{OPE data of }\cA)
\times (\text{geometric integral on }E_\tau)$.
The OPE data is the transferred $A_\infty$ operation
$m_{|V(\Gamma)|}$ composed with the propagator~$P$, contracted
along the graph edges. Organizing by degree
$r = |V(\Gamma)|$:
\[
 F_1
 = \sum_{r\ge2}\;
 \underbrace{\sum_{\substack{\Gamma : \text{genus-}1 \\
 |V(\Gamma)|=r}}
 \frac{(\text{OPE amplitude of }\Gamma)}
 {\lvert\mathrm{Aut}(\Gamma)\rvert}}_{
 = \operatorname{Sh}_r^{(1)}(\cA)}
 \;\cdot\;
 \underbrace{\int_{E_\tau^r/S_r}
 \omega_\Gamma}_{= \mathcal{G}_r(q)}.
\]
The identification of the OPE sum with
$\pi_{1,r}(\Theta_\cA)$ follows from
the bar-intrinsic construction
(Theorem~\ref{thm:mc2-bar-intrinsic}): the MC element
$\Theta_\cA = D_\cA - d_0$ encodes all graph amplitudes,
and $\pi_{1,r}$ extracts the genus-$1$, degree-$r$ component.
\end{proof}

\begin{corollary}[Shadow Fredholm determinant]
\label{cor:shadow-fredholm}
\ClaimStatusProvedElsewhere
The connected sewing determinant equals the Chern--Weil pairing
of the MC element with the fundamental class of the elliptic
curve:
\begin{equation}\label{eq:shadow-fredholm}
 \det\bigl(1 - K_q^{\mathrm{conn}}\bigr)
 = \exp\bigl(-\langle\Theta_\cA,\,[E_\tau]\rangle\bigr),
 \qquad
 \langle\Theta_\cA,[E_\tau]\rangle
 := \sum_{r\ge2}
 \operatorname{Sh}_r^{(1)}\;\cdot\;\mathcal{G}_r(q).
\end{equation}
\end{corollary}

\begin{example}[Heisenberg verification]
\label{ex:heisenberg-intertwining}
For the Heisenberg algebra at central charge~$c$, the shadow tower
terminates at degree~$2$:
$\operatorname{Sh}_2^{(1)} = \kappa = c/2$,
$\operatorname{Sh}_r^{(1)} = 0$ for $r \ge 3$.
The genus-$1$ geometric kernel at degree~$2$ is
$\mathcal{G}_2(q) = -\tfrac{2}{c}\sum_{n\ge1}\log(1-q^n)$.
The intertwining gives:
\[
 F_1^{\mathrm{conn}}
 = \frac{c}{2} \cdot \Bigl(-\frac{2}{c}\sum_{n\ge1}
 \log(1{-}q^n)\Bigr)
 = -\sum_{n\ge1}\log(1-q^n)
 = \sum_{N\ge1} \sigma_{-1}(N)\,q^N,
\]
which is the known connected free energy of the Heisenberg algebra
(the logarithm of~$\prod_{n\ge1}(1-q^n)^{-1}$).
The intertwining is \emph{exact}: no higher degrees contribute.
\end{example}

\begin{corollary}[Spectral measure identification]
\label{cor:spectral-measure-identification}
\ClaimStatusConditional
Let $V_\Lambda$ be a lattice VOA. The spectral measure~$\rho$
of the shadow tower
\textup{(}Theorem~\textup{\ref{thm:shadow-spectral-measure}}\textup{)}
equals the spectral measure of the sewing operator after the
Rankin--Selberg projection: the Rankin--Selberg integral of the
sewing determinant
\[
 \int_{\mathrm{SL}(2,\bZ)\backslash\mathfrak{H}}
 \bigl\lvert
 \det(1 - K_q^{\mathrm{conn}})
 \bigr\rvert^{-2}\,
 E_s(\tau)\,d\mu(\tau)
\]
equals the Rankin--Selberg integral of the shadow Fredholm
determinant
\[
 \int
 \bigl\lvert
 \exp\bigl(\langle\Theta_{V_\Lambda},[E_\tau]\rangle\bigr)
 \bigr\rvert^2\,
 E_s(\tau)\,d\mu(\tau),
\]
and both yield the constrained Epstein zeta
$\varepsilon^c_s(V_\Lambda)$ with $L$-function factorization
determined by the atoms of~$\rho$. The structural obstruction
\textup{(}Remark~\textup{\ref{rem:structural-obstruction}}\textup{)}
is thereby reduced to the non-lattice regime: for lattice VOAs,
the shadow spectral measure~$\rho$ and the automorphic spectral
measure coincide via the Hecke correspondence.
\end{corollary}

\begin{proof}
By Theorem~\ref{thm:sewing-shadow-intertwining},
$\det(1 - K_q) = \exp(-\langle\Theta,\,[E_\tau]\rangle)$
identically as functions of~$\tau$. The Rankin--Selberg
integrals therefore agree. The Hecke decomposition of
$\lvert\det(1-K_q)\rvert^{-2}$ into holomorphic Hecke
eigenforms produces the $L$-function factorization of
$\varepsilon^c_s$ (Theorem~\ref{thm:shadow-spectral-correspondence}).
The atoms of the shadow spectral measure~$\rho$ are the Hecke
eigenvalues (Theorem~\ref{thm:shadow-spectral-measure}),
which are the same data that determine the $L$-function
factorization.
\end{proof}

For Virasoro, the intertwining involves all degrees and the
leading-order resummation gives
$F_1^{\mathrm{conn}} \approx
\langle -\!\log(1 + 6t/c),\,[E_\tau]\rangle|_{t = t(q)}$,
where $t(q)$ is the shadow-moduli map.

\begin{theorem}[Shadow-moduli resolution]
\label{thm:shadow-moduli-resolution}
\ClaimStatusConditional
\index{shadow-moduli map!resolution}%
Let $\cA$ be a chirally Koszul algebra satisfying the HS-sewing
condition
\textup{(}Theorem~\textup{\ref{thm:general-hs-sewing}}\textup{)},
with spectral measure
$\rho = \sum_j c_j\,\delta(\lambda_j)$
\textup{(}Theorem~\textup{\ref{thm:shadow-spectral-measure}}\textup{)}.
The shadow-moduli map $t \colon \mathfrak{H} \to \bC$ is the unique solution
of the implicit equation
\begin{equation}\label{eq:shadow-moduli-implicit}
 \prod_j \bigl(1 - \lambda_j\,t(q)\bigr)^{c_j}
 \;=\;
 \det\!\bigl(1 - K_q^{\mathrm{conn}}\bigr).
\end{equation}
In particular:
\begin{enumerate}[label=\textup{(\roman*)}]
\item
\textup{(}Single-atom approximation.\textup{)}\;
When $\rho \approx \delta(\lambda + 6/c)$
\textup{(}leading order in $1/c$\textup{)}:
\begin{equation}\label{eq:shadow-moduli-leading}
 t(q)
 \;=\;
 \frac{c}{6}\biggl(
 \prod_{n\ge1}(1-q^n) - 1
 \biggr).
\end{equation}
\item
\textup{(}Multi-atom inversion.\textup{)}\;
For a general spectral measure with $m$ atoms, the
$q$-expansion of $t(q) = \sum_{N\ge1} t_N\,q^N$ is
determined order by order by Lagrange inversion of
\eqref{eq:shadow-moduli-implicit}: the linear coefficient
$t_1 = a_1/\mu_1$ where $a_1$ is the leading coefficient
of~$F_1^{\mathrm{conn}}$ and
$\mu_r = \sum_j c_j\,\lambda_j^r$
is the $r$-th moment of~$\rho$.
\end{enumerate}
\end{theorem}

\begin{proof}
The sewing--shadow intertwining
(Theorem~\ref{thm:sewing-shadow-intertwining})
and the Stieltjes representation
(Theorem~\ref{thm:shadow-spectral-measure}) compose:
\[
 F_1^{\mathrm{conn}}(q)
 \;=\;
 G\!\bigl(t(q)\bigr)
 \;=\;
 \sum_j c_j\log\!\bigl(1 - \lambda_j\,t(q)\bigr).
\]
The shadow Fredholm determinant
(Corollary~\ref{cor:shadow-fredholm}) gives
$F_1 = -\log\det(1 - K_q^{\mathrm{conn}})$.
Exponentiating:
\[
 \det\!\bigl(1 - K_q^{\mathrm{conn}}\bigr)
 \;=\;
 \exp\!\bigl(-F_1\bigr)
 \;=\;
 \exp\!\bigl(-G(t)\bigr)
 \;=\;
 \prod_j (1 - \lambda_j\,t)^{c_j},
\]
which is~\eqref{eq:shadow-moduli-implicit}. For a single atom
at $\lambda = -6/c$ with $c_1 = 1$:
$(1 + 6t/c) = \prod(1-q^n)$, giving~(i).
The Lagrange inversion in~(ii) is standard:
the equation
$\sum_{r\ge1}(\mu_r/r)\,t^r = F_1$
(from expanding the logarithms in~$G$) is a formal
power series identity in~$q$, and at each order~$N$
the coefficient~$t_N$ is determined linearly in terms of
$t_1, \dotsc, t_{N-1}$ and~$a_N$.
\end{proof}

\subsection{Universality of the shadow generating function}%
\label{sec:universality-G}
\index{shadow generating function!universality}

The universality of~$G$ is a direct consequence of the
Stieltjes representation of the MC element: the spectral
measure~$\rho$ encodes the MC data of~$\Theta_\cA$, and $G$
is its moment generating function.

\begin{theorem}[Universality of $G$]%
\label{thm:universality-of-G}%
\ClaimStatusConditional%
\index{shadow generating function!universal structure}
For any chirally Koszul algebra~$\cA$ with HS-sewing, the
connected genus-$1$ free energy admits the factorization
\begin{equation}\label{eq:universality-factorization}
 F_1^{\mathrm{conn}}(q;\,\cA)
 \;=\;
 G_\cA\!\bigl(t(q)\bigr),
 \qquad
 G_\cA(t)
 \;=\;
 \int\!\log(1 - \lambda\,t)\,d\rho_\cA(\lambda),
\end{equation}
where $G_\cA$ depends only on the spectral measure~$\rho_\cA$
\textup{(}determined by the algebraic data of~$\cA$: the OPE and
the transferred $A_\infty$ structure on bar
cohomology\textup{)}, and $t(q)$ is the shadow-moduli map
\textup{(}Theorem~\textup{\ref{thm:shadow-moduli-resolution}}\textup{)},
determined by the sewing determinant alone.

At leading order in $1/c$:
\begin{equation}\label{eq:G-leading-order}
 G_\cA(t)
 \;=\;
 \log\!\bigl(1 + 6t/c\bigr)
 + O(Q^{\mathrm{ct}}_\cA),
\end{equation}
i.e., the function $\log(1+6t/c)$ is universal
across all vertex algebras with a single strong generator
at leading order in $1/c$; the
algebra-specific corrections enter at degree~$4$ through
the quartic contact coupling~$Q^{\mathrm{ct}}_\cA$.
\end{theorem}

\begin{proof}
The factorization~\eqref{eq:universality-factorization}
is the composition of the sewing--shadow intertwining
(Theorem~\ref{thm:sewing-shadow-intertwining}) and the
Stieltjes representation
(Theorem~\ref{thm:shadow-spectral-measure}). The
separation of algebra-specific and moduli-specific data
follows from the shadow-moduli resolution.

The leading-order universality~\eqref{eq:G-leading-order}:
for all standard families, the spectral measure at
leading $1/c$ is a single atom at
$\lambda_{\mathrm{eff}} = -6/c$
(Remark~\ref{rem:effective-coupling}), giving
$G = \log(1 + 6t/c)$. The first correction at subleading
order involves the quartic contact coupling
$Q^{\mathrm{ct}}_\cA$
(Proposition~\ref{prop:virasoro-quartic-determinant}),
which splits the single atom into a pair whose
separation is controlled by~$Q^{\mathrm{ct}}$.
\end{proof}

\subsection{The MC recursion in spectral coordinates}%
\label{sec:mc-newton-identities}%
\index{Newton's identity!from MC equation|textbf}

Rewriting the MC recursion in spectral coordinates (power sums
and elementary symmetric polynomials via the Newton relations)
reveals that the self-interaction encoded by
$D\Theta + \tfrac{1}{2}[\Theta,\Theta] = 0$
determines the spectral atoms.

\begin{proposition}[MC bracket determines spectral atoms]%
\label{prop:mc-bracket-determines-atoms}%
\ClaimStatusProvedHere%
\index{spectral atoms!from MC bracket}
For a spectral measure with two atoms
$\rho = c_1\,\delta(\alpha) + c_2\,\delta(\beta)$
\textup{(}Satake type\textup{)}, the MC recursion at each
degree, rewritten in spectral coordinates, yields the Newton
relation
\begin{equation}\label{eq:mc-atom-constraint}
 p_r(\alpha,\beta)
 \;=\;
 e_1\,p_{r-1}(\alpha,\beta)
 \;-\;
 e_2\,p_{r-2}(\alpha,\beta),
\end{equation}
where $p_k = \alpha^k + \beta^k$ is the $k$-th power sum,
$e_1 = \alpha + \beta$, $e_2 = \alpha\beta$.
More precisely:
\begin{enumerate}[label=\textup{(\roman*)}]
\item
the degree-$2$ shadow gives $\kappa = c/2$, hence
$p_2 = e_1^2 - 2e_2$;
\item
the MC obstruction at degree~$r{+}1$,
$\mathfrak{o}^{(r+1)} =
\sum\{\operatorname{Sh}_j, \operatorname{Sh}_k\}_H$,
translates via $S_r = -p_r/r$ into Newton's
relation~\eqref{eq:mc-atom-constraint};
\item
for a single atom $\rho = \delta(\lambda)$:
$e_2 = 0$, Newton is trivial ($p_r = \lambda^r$),
consistent with the Virasoro geometric series.
\end{enumerate}
\end{proposition}

\begin{proof}
The Newton relation in spectral coordinates
follows from the Stieltjes representation
$S_r = -\mu_r/r$ and the MC recursion
$\nabla_H(\operatorname{Sh}_{r+1}) + \mathfrak{o}^{(r+1)} = 0$
(Theorem~\ref{thm:shadow-tower-asymptotics}). For two atoms,
in spectral coordinates $p_r = \alpha^r + \beta^r$,
the recursion $p_r - e_1\,p_{r-1} + e_2\,p_{r-2} = 0$
follows from the characteristic equation
$x^2 - e_1\,x + e_2 = 0$,
and the MC obstruction $\mathfrak{o}^{(r+1)}$ vanishes
precisely when this recursion holds.
\end{proof}

\begin{remark}[The bridge to the Ramanujan bound]%
\label{rem:mc-ramanujan-bridge}%
\index{Ramanujan bound!bridge from MC equation}%
The MC recursion, rewritten in spectral coordinates via the
Newton relations, encodes the power sums $p_r$,
which are the Dirichlet coefficients of
$L(s,\operatorname{Sym}^{r-1}f)$
(Proposition~\ref{prop:shadow-symmetric-power}).
The bridge to Ramanujan then passes through four stations:
\begin{center}
\small
\renewcommand{\arraystretch}{1.2}
\begin{tabular}{clcl}
\toprule
& \textbf{Station} & \textbf{Proof surface} & \textbf{Content}\\
\midrule
\textup{1} & MC at all degrees
 & \ClaimStatusProvedHere
 & Spectral recursion for $\rho$ \\
\textup{2} & Modularity of $G(t(q))$
 & \ClaimStatusProvedHere
 & Atoms are Hecke eigenvalues \\
\textup{3} & Symmetric power analytic continuation
 & Open ($r \ge 5$)
 & $L(s,\operatorname{Sym}^r f)$ entire \\
\textup{4} & Serre reduction
 & Conditional on \textup{3}
 & Ramanujan bound \\
\bottomrule
\end{tabular}
\end{center}

\noindent
Stations~$1$ and~$2$ are proved.
Station~$3$ is the outstanding gap: the MC equation provides
the symmetric power data at all degrees simultaneously, but
algebraic relations on Dirichlet coefficients do not
automatically imply analytic continuation.
Station~$4$ is classical (Remark~\ref{rem:serre-reduction}).
\end{remark}

% ============================================================
\subsection{From the shadow metric to Dirichlet $L$-functions}
\label{subsec:shadow-metric-l-functions}
\index{shadow metric!Dirichlet L-function@Dirichlet $L$-function|textbf}
\index{Dirichlet L-function@Dirichlet $L$-function!from shadow metric|textbf}

The shadow metric $Q_L$ is a positive definite binary quadratic
form on~$\bZ^2$ whenever $\kappa \neq 0$ and $\Delta > 0$
\textup{(}class~$\mathbf{M}$\textup{)}. Its Epstein zeta function
$\varepsilon_{Q_L}(s)$
\textup{(}Theorem~\textup{\ref{thm:shadow-epstein-zeta}}\textup{)}
connects the OPE data of a chirally Koszul algebra to the theory
of $L$-functions through the following chain of constructions.

\begin{construction}[The shadow $L$-function]
\label{constr:shadow-l-function}
\index{shadow L-function@shadow $L$-function|textbf}
Let $\cA$ be a chirally Koszul algebra of class~$\mathbf{M}$
at rational central charge~$c$, with shadow metric data
$(\kappa, \alpha, S_4)$ on a primary line~$L$.
\begin{enumerate}[label=\textup{(\roman*)}]
\item
\emph{The discriminant.}
$D(Q_L) := \operatorname{disc}(Q_L) = -32\kappa^2\Delta
\in \bQ_{<0}$
\textup{(}Proposition~\textup{\ref{prop:hankel-extraction}(iii)}\textup{)}.

\item
\emph{The shadow field.}
The squarefree part $\delta$ of $D(Q_L)$ in $\bQ^*/(\bQ^*)^2$
determines the \emph{shadow field}
$K_L := \bQ(\sqrt{\delta})$,
an imaginary quadratic extension of~$\bQ$
\textup{(}Remark~\textup{\ref{rem:shadow-field}}\textup{)}.

\item
\emph{The Dedekind factorization.}
The Dedekind zeta function $\zeta_{K_L}(s) = \zeta(s) \cdot
L(s, \chi_d)$, where $d$ is the fundamental discriminant of
$K_L$ and $\chi_d = \bigl(\tfrac{d}{\cdot}\bigr)$ is the
Kronecker character. This is a Dirichlet $L$-function with known
analytic continuation, functional equation, and Euler product.

\item
\emph{Class number splitting.}
If the class number $h(d) = 1$, the Epstein zeta factors as
$\varepsilon_{Q_L}(s) = c_0 \cdot \zeta_{K_L}(s)$ for a scaling
constant~$c_0$ depending on the normalization of~$Q_L$.
If $h(d) > 1$, the Epstein zeta depends on the specific ideal
class represented by~$Q_L$; the average over all $h(d)$ classes
gives $(2/w)\zeta_{K_L}(s)$, but the individual
form~$Q_L$ carries finer information.
The sum
$\sum_{[\mathfrak{a}]}
\varepsilon_{Q_{[\mathfrak{a}]}}(s)
= (2/w)\,\zeta_{K_L}(s)$
runs over ideal classes and $w = |\cO_{K_L}^\times|$.
\end{enumerate}
The full chain is therefore
\[
 \cA
 \;\xrightarrow{\;\Theta_\cA\;}
 Q_L
 \;\xrightarrow{\;\operatorname{disc}\;}
 D(Q_L) \in \bQ_{<0}
 \;\xrightarrow{\;\mathrm{sqfree}\;}
 K_L
 \;\xrightarrow{\;\zeta_K = \zeta \cdot L\;}
 L(s, \chi_d).
\]
For class~$\mathbf{G}$ and~$\mathbf{L}$ algebras: $\Delta = 0$,
$D = 0$, the shadow field degenerates to~$\bQ$, and the chain
terminates at~$\zeta(s)$.
The Epstein zeta of a degenerate form reduces to a
one-dimensional lattice sum.
\end{construction}

\begin{conjecture}[Shadow principal class]
\ClaimStatusConjectured
\label{conj:shadow-principal-class}
\index{shadow field!principal class conjecture}
For every chirally Koszul algebra~$\cA$ of class~$\mathbf{M}$
at rational central charge, the binary quadratic form~$Q_L$
\textup{(}suitably scaled to an integral form\textup{)}
represents the principal ideal class in the class group
of $K_L = \bQ(\sqrt{\operatorname{disc}(Q_L)})$.
\end{conjecture}

\begin{remark}[Principal-class test]
\label{rem:shadow-principal-criterion}
\index{shadow field!principal-class test}
The class-number-one cases are automatic and carry no information.
The first useful test is already the Ising point:
$c=1/2$ gives $K_L=\bQ(\sqrt{-10})$ and discriminant
$-40$, whose class number is~$2$. Thus the conjecture is
nontrivial at the first minimal-model example.

The proposed principal class comes from the following explicit
comparison. The shadow metric $Q_L$ is the quadratic norm form built
from the degree-$0$ coefficient $a_0=2\kappa$ and the MC-recursion
correction terms. Its Gaussian envelope
$(2\kappa+3\alpha t)^2$ is the square of a linear form and therefore
represents the principal class. The open calculation is whether the
correction term $2\Delta t^2$ leaves the resulting integral binary
quadratic form in the principal genus and, more sharply, in the
principal ideal class.

The class number of $K_L$ grows with the discriminant,
so for generic rational~$c$ (where
$|D| = 320c^2/(5c{+}22)$ grows like $64c$ for large~$c$),
one expects $h > 1$ eventually.
The first Virasoro example with $h > 1$ is $c = 1/2$:
$h(-40) = 2$ and the two classes are
$\{[m^2 + 10n^2],\; [2m^2 + 5n^2]\}$.
Determining which class $Q_L$ occupies
requires computing the reduced form of the scaled
shadow metric, a finite computation at each
rational~$c$.
\end{remark}

\begin{definition}[Specialised shadow field and the
specialisation functor $\operatorname{Spec}_{c_0}$]%
\label{def:specialised-shadow-field}%
\index{specialised shadow field|textbf}%
\index{shadow field!specialisation functor|textbf}%
\index{specialisation functor $\operatorname{Spec}_{c_0}$|textbf}
Let $K_L = \bQ(\kappa, \sqrt{Q_L(t)})$ denote the
\emph{generic shadow extension} of the rational function
field $\bQ(\kappa)$ obtained by adjoining a square root of the
shadow metric polynomial~$Q_L(t)$ for a chirally Koszul
algebra~$\cA$ of class~$\mathbf{M}$.
The \emph{specialisation functor at central charge~$c_0$}
\[
 \operatorname{Spec}_{c_0}\;:\;
 K_L \;\longrightarrow\; K(c_0),
\]
sends the generic shadow extension to the imaginary quadratic
field
\[
 K(c_0)
 \;:=\;
 \bQ\bigl(\sqrt{\operatorname{sqfree}(\operatorname{disc}(Q_L)(c_0))}\bigr),
\]
where $\operatorname{disc}(Q_L)(c) = -320c^2/(5c{+}22)$ for the
Virasoro family. Because $c^2$ is a square in~$\bQ^\times$,
the squarefree class is
$[-5(5c_0{+}22)] \in \bQ^\times/(\bQ^\times)^2$, so
\begin{equation}\label{eq:specialised-shadow-field}
 \operatorname{Spec}_{c_0}(K_L)
 \;=\;
 \bQ\bigl(\sqrt{-5(5c_0{+}22)}\bigr).
\end{equation}
The specialisation is functorial in~$c_0$ along rational
specialisations of the Virasoro parameter: a rational point
$c_0 \in \bQ \setminus \{0, -22/5\}$ of
$\operatorname{Spec}\,\bQ(c)$ determines the quadratic field
$K(c_0)$ up to isomorphism, and the generic field~$K_L$ is
the limit of $K(c_0)$ as $c_0$ ranges over the specialisation
locus.
\end{definition}

\begin{remark}[Arithmetic invariants of~$K(c_0)$]%
\label{rem:specialised-shadow-field-invariants}%
\index{class number!specialised shadow field}%
\index{Dirichlet $L$-function!specialised shadow field}
The specialised field $K(c_0) = \bQ(\sqrt{d(c_0)})$ inherits
the standard arithmetic invariants of an imaginary quadratic
field:
\begin{enumerate}[label=\textup{(\roman*)}]
\item the class number $h(d(c_0))$, a positive integer
 measuring the failure of unique factorisation;
\item the Dirichlet $L$-function $L(s, \chi_{d(c_0)})$ attached
 to the Kronecker character $\chi_{d(c_0)} =
 \bigl(\tfrac{d(c_0)}{\cdot}\bigr)$;
\item the Dedekind zeta factorisation $\zeta_{K(c_0)}(s) =
 \zeta(s)\,L(s, \chi_{d(c_0)})$.
\end{enumerate}
Each invariant becomes a computable attribute of the vertex
algebra at central charge~$c_0$ through
Construction~\textup{\ref{constr:shadow-l-function}}.
\end{remark}

\begin{example}[Specialisation at the Ising point $c_0 = 1/2$]%
\label{ex:specialised-shadow-field-ising}%
\index{Ising model!specialised shadow field}
At the Ising point $c_0 = 1/2$ the discriminant is
\[
 \operatorname{disc}(Q_L)(1/2)
 \;=\;
 \frac{-320\,(1/2)^2}{5\,(1/2) + 22}
 \;=\;
 \frac{-80}{49/2}
 \;=\;
 \frac{-160}{49}.
\]
Since $49 = 7^2$ is a perfect square, the squarefree class of
$-160/49$ in $\bQ^\times/(\bQ^\times)^2$ agrees with that of
$-160 = -16 \cdot 10$, whose squarefree part is $-10$.
Equivalently, $\sqrt{-160/49} = (4/7)\sqrt{-10}$.
Therefore
\[
 \operatorname{Spec}_{1/2}(K_L)
 \;=\;
 \bQ(\sqrt{-10}),
\]
the imaginary quadratic field with fundamental discriminant
$d = -40$ and class number $h(-40) = 2$. The two ideal classes
are represented by the reduced binary quadratic forms
$[m^2 + 10n^2]$ and $[2m^2 + 5n^2]$, consistent with
Remark~\textup{\ref{rem:shadow-principal-criterion}}.
\end{example}

\begin{proposition}[Koszul field-preservation criterion]%
\label{prop:koszul-field-criterion}%
\ClaimStatusProvedHere%
\index{shadow field!Koszul field-preservation criterion|textbf}%
\index{Koszul duality!field-preservation criterion|textbf}
The Koszul pair $(\mathrm{Vir}_c,\, \mathrm{Vir}_{26-c})$
gives the same shadow field
$K_L(c) = K_L(26-c)$
if and only if the ratio
\begin{equation}\label{eq:koszul-field-ratio}
 \frac{5c + 22}{152 - 5c}
 \quad\in\quad \bQ^{\times 2}
\end{equation}
is a perfect square in~$\bQ^\times$.
Among integers $c \in [0, 26]$ modulo the Koszul involution
$c \mapsto 26 - c$, the criterion holds only at
$c = 1$ and $c = 13$.
\end{proposition}

\begin{proof}
The discriminant of the shadow metric is
$\operatorname{disc}(Q_L) = -32\kappa^2\Delta$
\textup{(}Proposition~\textup{\ref{prop:hankel-extraction}(iii)}\textup{)}.
For Virasoro at central charge~$c$:
$\kappa = c/2$ and
$\Delta = 8\kappa S_4 = 40/(5c + 22)$
\textup{(}from $Q^{\mathrm{contact}}_{\mathrm{Vir}}
= 10/[c(5c+22)]$\textup{)},
so
\[
 \operatorname{disc}(Q_L)
 \;=\;
 \frac{-320\,c^2}{5c + 22}.
\]
In the quotient group $\bQ^*/(\bQ^*)^2$,
the class of $\operatorname{disc}(Q_L)$ is
$[-5(5c+22)]$
\textup{(}since $c^2$ is a square and
$320 = 64 \cdot 5$\textup{)},
so the shadow field is
$K_L(c) = \bQ(\sqrt{-5(5c+22)})$.
The Koszul dual at $c' = 26 - c$ has $5c' + 22 = 152 - 5c$,
giving $K_L(26-c) = \bQ(\sqrt{-5(152-5c)})$.
Two quadratic fields $\bQ(\sqrt{a})$ and $\bQ(\sqrt{b})$
coincide if and only if $a/b \in (\bQ^*)^2$.
The common factor $-5$ cancels, yielding the
criterion~\eqref{eq:koszul-field-ratio}.

For the integer enumeration: at $c = 1$,
$(5 + 22)/(152 - 5) = 27/147 = (3/7)^2$.
At $c = 13$, the ratio is $87/87 = 1$.
For every other $c \in \{0, 2, 3, \ldots, 12\}$,
the numerator $5c + 22$ and denominator $152 - 5c$
have squarefree parts that differ
\textup{(}verified by direct factorization\textup{)}.
By symmetry $c \mapsto 26 - c$, the range $[14, 26]$
adds no new pairs.
\end{proof}

\begin{example}[The pair $c = 1$, $c = 25$]
\label{ex:koszul-field-1-25}
$(5 + 22)/(152 - 5) = 27/147 = 9/49 = (3/7)^2$.
Both algebras give the shadow field
$K_L = \bQ(\sqrt{-15})$.
\end{example}

\begin{example}[The self-dual point $c = 13$]
\label{ex:koszul-field-13}
$(65 + 22)/(152 - 65) = 87/87 = 1$.
The Virasoro algebra at $c = 13$ is Koszul self-dual
\textup{(}$\mathrm{Vir}_{13}^! \simeq \mathrm{Vir}_{13}$\textup{)},
so field preservation is automatic; the shadow field is
$K_L = \bQ(\sqrt{-435})$.
\end{example}

\begin{proposition}[The two fixed points are unrelated]
\label{prop:two-fixed-points}
\ClaimStatusConditional
\index{Koszul involution!fixed point at $c=13$}%
\index{Riemann zeta function!critical line}%
\index{self-duality!$c=13$ vs.\ $\mathrm{Re}(s)=1/2$}
The Koszul fixed point $c = 13$ and the critical line
$\mathrm{Re}(s) = 1/2$ arise from different structures:
\begin{enumerate}[label=\textup{(\roman*)}]
\item $c = 13$ is the fixed point of $c \mapsto 26 - c$ (the
 Koszul involution on the Virasoro family). It is determined
 by $c + c' = 26$ (Freudenthal--de~Vries), which depends on
 the Virasoro central charge formula $c = 1 - 6(p-q)^2/(pq)$
 and the Sugawara construction. This is CFT data.
\item $\mathrm{Re}(s) = 1/2$ is the fixed line of
 $s \mapsto 1 - s$ (the functional equation of~$\zeta$). It is
 determined by the Euler product and the gamma factors. This
 is number-theoretic data.
\end{enumerate}
The coincidence $c/2 = 13/2$ and $s = 1/2$ share the symbol
``$1/2$'' but arise from unrelated arithmetic. The critical
line of~$\zeta$ is a consequence of the functional equation
$\xi(s) = \xi(1-s)$; the self-duality of $\mathrm{Vir}_{13}$
is a consequence of $26 - 13 = 13$. No mechanism connects them.
\end{proposition}

\begin{proof}
The Koszul involution $c \mapsto 26 - c$ acts on the
one-parameter Virasoro family; its fixed point $c = 13$
is forced by the relation $c + c' = 26$
(Theorem~\ref{thm:modular-characteristic}).
The functional equation $\xi(s) = \xi(1 - s)$ acts on the
complex $s$-plane; its fixed line $\mathrm{Re}(s) = 1/2$
is forced by the Euler product representation and the
duplication formula for the gamma function.
These two involutions act on different spaces
($\bR_c$ vs.\ $\bC_s$), are determined by different data
(OPE structure constants vs.\ prime factorization), and have
no common origin. In particular, the Koszul--Epstein function
$\varepsilon^{\mathrm{KE}}_{\cA}(s)$
(Definition~\ref{def:koszul-epstein-function})
evaluates at $s = 1/2$ without reference to $c = 13$,
and the self-dual point $c = 13$ contributes no preferred
value of~$s$.
\end{proof}

\begin{remark}
The temptation to identify $c = 13$ with $s = 1/2$ is a
case of the meta-principle: the same name (``fixed
point of an involution'') applied to different objects does not
make them the same.
\end{remark}

\begin{remark}[Failure for minimal models]
\label{rem:koszul-field-failure}
\index{shadow field!failure for minimal models}
For the Ising model $c = 1/2$, the Koszul dual has
$c' = 51/2$. The ratio is $(49/2)/(299/2) = 49/299$,
and $299 = 13 \cdot 23$ is squarefree, so $49/299$
is not a perfect square. The shadow fields differ:
$K_L(1/2) = \bQ(\sqrt{-10})$ while
$K_L(51/2) = \bQ(\sqrt{-2990})$,
where $2990 = 2 \cdot 5 \cdot 13 \cdot 23$.

More generally: for every minimal model central charge
$c_{p,q}$ with $p < 10$, the Koszul pair gives
different imaginary quadratic fields. The criterion
\eqref{eq:koszul-field-ratio} requires the arithmetic
coincidence $\mathrm{sqfree}(5c+22) =
\mathrm{sqfree}(152-5c)$, which fails generically.
\end{remark}

% ============================================================
\subsection{The Koszul--Epstein function}%
\label{subsec:koszul-epstein-function}
\index{Koszul--Epstein function|textbf}

\begin{definition}[Koszul--Epstein function]%
\label{def:koszul-epstein-function}%
\index{Koszul--Epstein function|textbf}%
\index{Epstein zeta function!Koszul--Epstein specialization|textbf}%
\index{shadow metric!Koszul--Epstein function|textbf}
Let $\cA$ be a chirally Koszul algebra with shadow data
$(\kappa, \alpha, S_4)$ on a primary line
$L \subset \Defcyc^{\mathrm{mod}}(\cA)$, and let
$\Delta := 8\kappa S_4$ be the critical discriminant
\textup{(}Definition~\textup{\ref{def:shadow-metric}}\textup{)}.
The \emph{Koszul--Epstein function} of~$\cA$ on~$L$ is the
Dirichlet series
\begin{equation}\label{eq:koszul-epstein-function}
 \varepsilon^{\mathrm{KE}}_{\cA,L}(s)
 \;:=\;
 \begin{cases}
 \displaystyle
 \sideset{}{'}\sum_{(m,n)\in\bZ^2}
 Q_L(m,n)^{-s},
 & \text{if } \Delta \neq 0
 \text{ \textup{(}class~$\mathbf{M}$\textup{)}},
 \\[12pt]
 2\,(4\kappa^2)^{-s}\,\zeta(2s),
 & \text{if } \Delta = 0,\; \alpha = 0
 \text{ \textup{(}class~$\mathbf{G}$\textup{)}},
 \\[6pt]
 \text{sewing Dirichlet series } S_\cA(u),
 & \text{if } \Delta = 0,\; \alpha \neq 0
 \text{ \textup{(}class~$\mathbf{L}$\textup{)}},
 \end{cases}
\end{equation}
where $Q_L(m,n) = 4\kappa^2 m^2 + 12\kappa\alpha\,mn
+ (9\alpha^2 + 2\Delta)\,n^2$ is the shadow metric viewed as
a binary quadratic form on~$\bZ^2$
\textup{(}Theorem~\textup{\ref{thm:shadow-epstein-zeta}}\textup{)},
$\zeta$ is the Riemann zeta function, and $S_\cA(u)$ is the
sewing Dirichlet series of
\S\textup{\ref{subsec:dirichlet-sewing-lift}}.
The class~$\mathbf{C}$ case
\textup{(}$\Delta = 0$ on the charged stratum with $\kappa = 0$\textup{)}
lies outside the single-line framework
\textup{(}Remark~\textup{\ref{rem:contact-stratum-separation}}\textup{)}.
\end{definition}

\begin{remark}[Three structural constraints]%
\label{rem:koszul-epstein-constraints}%
\index{Koszul--Epstein function!three constraints}
A Koszul--Epstein function differs from a generic
Epstein zeta function in three ways, each tracing to the
bar-cobar quasi-isomorphism:
\begin{enumerate}[label=\textup{(\roman*)}]
\item
\emph{Koszul symmetry.}
Koszul duality $\cA \leftrightarrow \cA^!$ acts on the shadow
metric via the complementarity involution:
$\kappa \mapsto \kappa' = \kappa(\cA^!)$,
$\Delta \mapsto \Delta'$,
with the discriminant constraint
$\Delta(\cA) + \Delta(\cA^!) = 6960/\bigl((5c{+}22)(152{-}5c)\bigr)$
for Virasoro
\textup{(}Corollary~\textup{\ref{cor:discriminant-atlas}}\textup{)}.
This relates $\varepsilon^{\mathrm{KE}}_\cA$ to
$\varepsilon^{\mathrm{KE}}_{\cA^!}$ through a specific algebraic
involution, stronger than the generic
$\Lambda(s) = \Lambda(1{-}s)$ that holds for all binary forms.

\item
\emph{Shadow polarization.}
The full tower $\{S_r\}_{r \geq 2}$ is determined by
three genus-$0$ invariants $(\kappa, \alpha, S_4)$ via the
algebraic identity $H(t)^2 = t^4\,Q_L(t)$
\textup{(}Theorem~\textup{\ref{thm:riccati-algebraicity}}\textup{)}.
This gives infinitely many moment constraints:
$S_r = f_r(\kappa, \alpha, S_4)$ for all $r \geq 5$, where
$f_r$ is the $r$-th Taylor coefficient of $t^2\sqrt{Q_L(t)}$.
For generic Epstein zeta functions, no such tower exists.

\item
\emph{Modular coupling.}
At genus~$1$, $\varepsilon^{\mathrm{KE}}_\cA$ couples to
the Eisenstein series on~$\cM_{1,1}$ via Rankin--Selberg
unfolding: the scattering factor $F_c(s)$ in the
Benjamin--Chang functional equation~\cite{BenjaminChang22}
contains the ratio $\zeta(2s)/\zeta(2s{-}1)$, linking the
MC-constrained Koszul--Epstein function to the zeros
of~$\zeta$.
\end{enumerate}
\end{remark}

\begin{proposition}[Degenerate case: Heisenberg]%
\label{prop:heisenberg-koszul-epstein}%
\ClaimStatusProvedHere%
\index{Heisenberg!Koszul--Epstein function}%
\index{Koszul--Epstein function!Heisenberg}
For the Heisenberg vertex algebra $\cH_k$ at level~$k$:
\begin{enumerate}[label=\textup{(\roman*)}]
\item
The shadow metric degenerates: $Q_{\cH_k}(t) = k^2$
\textup{(}constant\textup{)}, and the binary form
$Q(m,n) = k^2 m^2$ has rank~$1$.
The Epstein sum over~$\bZ^2$ diverges:
for each $m \neq 0$, the sum over~$n$ gives
$(2N{+}1)\,|k\,m|^{-2s} \to \infty$.

\item
The correct Koszul--Epstein function is
\begin{equation}\label{eq:heisenberg-koszul-epstein}
 \varepsilon^{\mathrm{KE}}_{\cH_k}(s)
 \;=\;
 2\,k^{-2s}\,\zeta(2s).
\end{equation}
This factors through~$\zeta$; all nontrivial zeros lie on
$\operatorname{Re}(s) = 1/4$
\textup{(}corresponding to $\operatorname{Re}(2s) = 1/2$\textup{)}.

\item
The Euler--Koszul defect is $D_{\cH_k} \equiv 1$
\textup{(}exact Euler--Koszul\textup{)}.
The sewing Dirichlet series $S_{\cH}(u) = \zeta(u)\,\zeta(u{+}1)$
has the Euler product
$\prod_p (1 - p^{-u})^{-1}(1 - p^{-u-1})^{-1}$.
\end{enumerate}
\end{proposition}

\begin{proof}
(i)~For class~$\mathbf{G}$: $\alpha = 0$, $S_4 = 0$,
$Q_L(t) = 4\kappa^2 = k^2$ constant.
$Q(m,n) = k^2 m^2$ is independent of~$n$, so fixing $m \neq 0$
and summing over $n \in \bZ$ gives $(2N{+}1)\,(k^2 m^2)^{-s}
\to \infty$.

(ii)~The one-dimensional lattice sum
$\sum'_{m \in \bZ} (k^2 m^2)^{-s}
= 2 k^{-2s} \sum_{m=1}^\infty m^{-2s} = 2 k^{-2s}\,\zeta(2s)$.

(iii)~All generators have weight~$1$, so
$D_{\cH_k}(u) = 1$ by
Theorem~\ref{thm:euler-koszul-standard}(i).
\end{proof}

\begin{computation}[Virasoro at $c = 1$: numerical Koszul--Epstein value;
\ClaimStatusProvedHere]%
\label{comp:virasoro-c1-koszul-epstein}%
\index{Koszul--Epstein function!Virasoro c=1@Virasoro $c=1$}%
\index{Virasoro!Koszul--Epstein function at c=1@Koszul--Epstein at $c=1$}
For the Virasoro algebra at $c = 1$:
$(\kappa, \alpha, S_4) = (1/2,\; 2,\; 10/27)$,
$\Delta = 40/27$,
$Q(m,n) = m^2 + 12\,mn + \tfrac{1052}{27}\,n^2$,
$\operatorname{disc}(Q) = -320/27$.
The Koszul--Epstein function at $s = 2$, computed by the
theta-function representation
\textup{(}incomplete gamma sums over $|m|, |n| \leq 40$\textup{)},
gives
\begin{equation}\label{eq:virasoro-c1-ke-value}
 \varepsilon^{\mathrm{KE}}_{\mathrm{Vir}_1}(2)
 \;=\;
 2.9054 \pm 10^{-4}.
\end{equation}
The functional equation $\Lambda^{\mathrm{KE}}(s)
= \Lambda^{\mathrm{KE}}(1{-}s)$ is verified to
$10^{-16}$ relative error at $s = 0.7$ and
$s = 0.5 + 3i$ via the theta-function method.
\end{computation}

\begin{computation}[Functional equation for all unitary minimal models;
\ClaimStatusProvedHere]%
\label{comp:fe-minimal-models}%
\index{Koszul--Epstein function!functional equation verification}%
\index{Epstein zeta function!functional equation}
The functional equation
$\Lambda^{\mathrm{KE}}(s) = \Lambda^{\mathrm{KE}}(1{-}s)$
\textup{(}Theorem~\textup{\ref{thm:shadow-epstein-zeta}(ii)}\textup{)}
is computed for the Virasoro shadow metric
at all unitary minimal model central charges $c = 1 - 6/m(m{+}1)$,
$m = 3, \ldots, 12$, via the theta-function splitting method
\textup{(}\texttt{shadow\_epstein\_zeta.py}\textup{)}.
At each central charge, seven test points
$s \in \{0.2, 0.3, 0.4, 0.5, 0.6, 0.7, 0.8\}$ give
relative error below~$10^{-3}$.
The following cases are distinguished:
\begin{itemize}[nosep]
\item
$M(9,10)$, $c = 14/15$: the shadow field is
$K = \bQ(\sqrt{-3}) = \bQ(\zeta_3)$
\textup{(}Eisenstein integers\textup{)}, with
class number $h(-3) = 1$ and $w = 6$.
This is the only minimal model with class number one
among $m = 3, \ldots, 20$.
\item
$M(3,4)$, $c = 1/2$ \textup{(}Ising\textup{)}:
$K = \bQ(\sqrt{-10})$, $d = -40$, $h = 2$.
The shadow metric selects a specific ideal class in
$\mathrm{Cl}(-40) = \{[m^2 {+} 10n^2],\; [2m^2 {+} 5n^2]\}$.
\item
$c = 1$ and $c = 25 = 26 - 1$ \textup{(}Koszul duals\textup{)}:
both give $K = \bQ(\sqrt{-15})$, $d = -15$, $h = 2$.
The coincidence is structural: the squarefree part of
$-320c^2/(5c{+}22)$ and $-320(26{-}c)^2/(152{-}5c)$ can
agree, though this is not universal.
\end{itemize}
Verification for $c = 1$, $c = 2$, and $c = 26$
\textup{(}critical string\textup{)} also passes.
\end{computation}

\begin{remark}[The Davenport--Heilbronn question for Koszul--Epstein functions]%
\label{rem:davenport-heilbronn-koszul-epstein}%
\index{Davenport--Heilbronn!for Koszul--Epstein}%
\index{Koszul--Epstein function!Davenport--Heilbronn question}
The question of whether the three Koszul--Epstein constraints
force zeros onto the critical line is \textbf{open}. What is proved:
\begin{enumerate}[label=\textup{(\arabic*)}]
\item
For class~$\mathbf{G}$ \textup{(}Heisenberg\textup{)}:
$\varepsilon^{\mathrm{KE}}$ factors through~$\zeta$;
zeros on the critical line under RH\@.
\item
For Koszul self-dual algebras \textup{(}$\cA \simeq \cA^!$,
e.g.\ $\mathrm{Vir}_{13}$\textup{)}:
the constrained Epstein zeta decomposes into standard
$L$-functions
\textup{(}Theorem~\textup{\ref{thm:self-dual-factorization}}\textup{)};
zeros on the critical line under GRH\@.
\item
For class~$\mathbf{M}$ with $h(D) > 1$ and $\cA \not\simeq
\cA^!$: the Davenport--Heilbronn mechanism could in principle
produce off-line zeros.
The MC equation provides additional constraints not present
for generic Epstein zeta functions, but these have not been
proved to exclude off-line zeros.
\end{enumerate}
The structural diagnosis
\textup{(}Theorem~\textup{\ref{thm:structural-separation}}\textup{)}:
the genus-$1$ MC equation cannot access scattering-matrix
poles. The programme requires using all algebras
simultaneously \textup{(}bootstrap closure across central
charges\textup{)}, and three steps remain open:
computing the MC constraint on $\varepsilon^c_s$ for
non-Nairin theories, showing that off-line residues violate
MC constraints, and bootstrap closure.
\end{remark}

% ============================================================
\section{The Hecke module structure}
\label{sec:reassessment}

\begin{theorem}[Hecke-equivariant MC element]
\label{thm:spectral-continuation-bridge}
\ClaimStatusProvedHere
Let $V_\Lambda$ be a lattice VOA of rank~$r$, and let
$\Theta_\Lambda = E_{r/2} + \sum_{j=1}^m c_j f_j$ be the
Hecke decomposition in $M_{r/2}(\mathrm{SL}(2,\bZ))$.
\begin{enumerate}[label=\textup{(\roman*)}]
\item The Hecke operators $T_n \colon M_{r/2} \to M_{r/2}$ act on
 the graph amplitudes: for every stable graph~$\Gamma$ of
 degree~$r$,
 \begin{equation}\label{eq:hecke-equivariance}
 \ell_\Gamma^{(g)}(V_\Lambda)
 \;\in\;
 \mathrm{Image}\bigl(T_{\lvert\Gamma\rvert}
 \colon M_{r/2} \to M_{r/2}\bigr),
 \end{equation}
 where $\lvert\Gamma\rvert$ is a combinatorial invariant of the
 graph type.

\item The shadow at degree~$r$ decomposes into Hecke eigenspaces:
 \begin{equation}\label{eq:shadow-hecke-decomposition}
 \operatorname{Sh}_r(V_\Lambda)
 \;=\;
 \sum_{j=0}^{m} \alpha_j\,a_{f_j}^{r-1},
 \end{equation}
 where $a_{f_j} = c_j\,\lVert f_j\rVert_{\mathrm{Pet}}$ is the
 Hecke eigenvalue and $\alpha_j$ is a combinatorial weight.

\item The Mellin transform
 $E^*_\Lambda(s) = (2\pi)^{-s}\Gamma(s)\,E_\Lambda(s)$ provides
 analytic continuation to all $s \in \bC$, with the
 $L$-function factorization
 \begin{equation}\label{eq:hecke-epstein-factorization}
 E_\Lambda(s)
 = C_0(s)\,\zeta(s)\,\zeta(s{-}r/2{+}1)
 + \sum_{j=1}^m c_j\,C_j(s)\,L(s,f_j).
 \end{equation}
 Each $L(s,f_j)$ has meromorphic continuation by Hecke theory.
 The shadow data determines the decomposition
 \eqref{eq:shadow-hecke-decomposition}, hence determines the
 $L$-function content~\eqref{eq:hecke-epstein-factorization}.
\end{enumerate}
\end{theorem}

\begin{proof}
Part~(i): for a lattice VOA, every OPE coefficient entering
$\ell_\Gamma$ is a coefficient of the theta series
$\Theta_\Lambda \in M_{r/2}(\SL_2(\ZZ))$. The stable-graph amplitude is
obtained from this theta series by finitely many coefficient
extractions, products, contractions, and modular pull-push operations.
These operations preserve the finite Hecke module generated by
$\Theta_\Lambda$. Hence $\ell_\Gamma$ lies in the Hecke span of the
theta component indexed by the graph type, equivalently in the image of
the appropriate Hecke correspondence.

Part~(ii) is the same statement in an eigenbasis. Since the Hecke algebra
on $M_{r/2}(\SL_2(\ZZ))$ is commutative and diagonalizable over the cusp
eigenbasis together with the Eisenstein line, the chain-graph trace
decomposes over Hecke eigenspaces; this is the content of
Proposition~\ref{prop:leading-hecke-identification}.

Part~(iii) is the classical Hecke--Mellin correspondence applied to the
Eisenstein component and to each cusp eigenform~$f_j$. Mellin
transformation gives the two zeta factors on the Eisenstein line and the
standard $L(s,f_j)$ on each cusp line. The coefficients $c_j$ are the
coordinates of $\Theta_\Lambda$ in the Hecke eigenbasis, so the shadow
decomposition determines exactly the $L$-function content of
\eqref{eq:hecke-epstein-factorization}.
\end{proof}

\begin{remark}[Minimal model characters]
\label{rem:virasoro-character-vector}
\label{rem:vvmf-rankin-selberg}
\label{rem:vvmf-hecke-decomposition}
\label{rem:virasoro-generic-frontier}
\label{rem:mellin-continuation-functor}
\label{rem:lattice-vs-nonlattice}
For $\cM(p,q)$, the character vector
$F(\tau) = (\chi_{r,s}(\tau))$ is a vector-valued modular form;
$\int \lvert F\rvert^2\,E_s\,d\mu$ unfolds to
$\sum a_n\,n^{-s}$ with $a_n = \sum_{r,s}\lvert d_n^{(r,s)}\rvert^2$.
For generic (irrational) $c$, the Hecke theory for VVMF
(Franc--Mason~\cite{FrancMason2017}) applies, but explicit eigenvalue
computations remain open.
\end{remark}

\begin{remark}[Andrianov spinor factorisation of $\Delta_{10}$:
Saito--Kurokawa packet and $M_{24}$ twist]
\label{rem:arith-shad-9}%
\index{Andrianov spinor L-function!Delta10 factorisation}%
\index{Saito--Kurokawa CAP!Arthur packet Delta10}%
\index{Mathieu multiplier!anomalous conjugacy classes}%
\index{Andrianov Saito-Kurokawa!K3 BKM}%
The Hecke decomposition
\eqref{eq:hecke-epstein-factorization} of the shadow Epstein zeta for
a lattice VOA gives, at each prime~$p$, a collection of Hecke
eigenvalues $a_p(f_j)$ indexed by cuspidal eigenforms $f_j \in
S_{r/2}(\Gamma)$. For the unnormalised Igusa square
$\Delta_{10}:=\Phi_{10}^{\mathrm{un}}=\Delta_5^2=\chi_{10}
\in S_{10}(\mathrm{Sp}(4, \bZ))$, the square/weight-$10$ lane is
separate from the primitive weight-$5$ Borcherds denominator
$\Delta_5$. It appears here as a genus-$2$ automorphic obstruction
target for the $\Delta_5^2$-family. Andrianov's spinor
$L$-function factorisation \cite{Andrianov1974} produces the
three-factor decomposition
\begin{equation}\label{eq:andrianov-delta10-factorisation}
 L^{\mathrm{spin}}(s, \Delta_{10})
 \;=\;
 L(s, \Delta_{E_6})\,\zeta(s-9)\,\zeta(s-8),
\end{equation}
where $\Delta_{E_6} \in S_{E_6}(\mathrm{SL}(2, \bZ))$ is an elliptic
cusp form attached to the $E_6$-lattice Borcherds lift. At each
prime $p$, the Hecke eigenvalue reads
\begin{equation}\label{eq:andrianov-hecke-lambda}
 \lambda_p(\Delta_{10})
 \;=\;
 a_p(\Delta_{E_6}) + p^8 + p^9,
\end{equation}
with the two explicit power-of-$p$ contributions tracking the two
continuous-spectrum factors $\zeta(s-9)$ and $\zeta(s-8)$.
The factorisation exhibits $\Delta_{10}$ as a Saito--Kurokawa CAP
form on $\mathrm{GSp}(4)$ \cite{PiatetskiShapiro1983,Schmidt2007}:
its Arthur packet is the non-tempered packet
\begin{equation}\label{eq:arthur-packet-delta10}
 \psi_{\Delta_{10}}
 \;=\;
 \phi_{\Delta_{E_6}} \boxtimes \operatorname{Sym}^1,
\end{equation}
with the $\operatorname{Sym}^1$-factor on $\mathrm{SL}_2 \subset
\mathrm{SL}_2 \times \mathrm{SL}_2$ producing the two $\zeta$-factors
of the spinor factorisation.
\emph{Mathieu multiplier target at anomalous classes.} At the five
anomalous $M_{24}$-conjugacy classes
$[g]\in\{7A,7B,11A,23A,23B\}$ in the Mathieu-moonshine
decomposition of the K3 elliptic genus
$\mathrm{EG}(\mathrm{K3};\tau,z)$
\cite{EguchiOoguriTachikawa2011,Gaberdiel2010}, the available
structure is a projective multiplier for the Mathieu coefficient
module, not a geometric $M_{24}$ action on one K3 sigma model and
not an automatic twist of a Vol~I Hecke category. The multiplier
class lies in
\begin{equation}\label{eq:m24-schur-multiplier}
 [\chi^{M_{24}}_g]
 \;\in\;
 H^2(M_{24}, \mathrm{U}(1))
 \;\cong\;
 \bZ / 12,
\end{equation}
per Gaberdiel--Hohenegger--Volpato \cite{GaberdielHoheneggerVolpato2012}
(the Schur multiplier $H^2(M_{24}, \bZ) = \bZ / 12$ lifts to the
$\mathrm{U}(1)$-coefficient cohomology through the exponential
sequence). A phase-twisted local Hecke statement is therefore a
finite-recognition theorem to be proved: one must choose a
downward-saturated finite root window, a twining-compatible primitive
bar model, the Weyl-vector and parity conventions, and a comparison
map to the relevant BKM root spaces. Only after those data are fixed
may the local eigenvalue in \eqref{eq:andrianov-hecke-lambda} be
replaced by a twisted coefficient
\[
 \lambda_{p,g}(\Delta_{10})
 =
 \zeta_{12}^{[g]}\lambda_p(\Delta_{10})
 \qquad
 (\zeta_{12}^{[g]}\in\mu_{12}).
\]
Consequently the shadow tower supplies a test, not a proof, for the
anomalous multiplier: the per-prime coefficient
$S^{(p)}_{10}(\cA)$ can detect the phase only after the finite
Mathieu--Igusa recognition datum identifies its cuspidal residue with
the $g$-twined BKM coefficient.
See Part~\ref{part:standard-landscape} for the $M_{24}$-moonshine
family in the landscape census and Part~\ref{part:physics-bridges}
for the K3 elliptic-genus sector.
\end{remark}

% ============================================================
\section{The quartic residue programme}%
\label{sec:quartic-residue-programme}%
\index{quartic residue programme|textbf}

The quartic resonance class
$\mathfrak{Q}(\cA) = \pi_{0,4}(\Theta_\cA)$
is the degree-$4$ projection of the MC element.
The shadow tower and the zeta-zero interface live on different
layers: the tower organizes MC deformation data,
while zeta zeros enter through Eisenstein residues.
The quartic class is the first degree where these layers genuinely
meet.

\subsection{The mixed $(s,u)$-transform}%
\label{sec:mixed-su-transform}

Fix a modular Koszul family~$\cA$ and a cyclic direction~$v$.
Write $E^*(\tau,s)$ for the completed Eisenstein series.

\begin{definition}[Potential-side transform]%
\label{def:potential-side-transform}%
\index{potential-side transform|textbf}%
The \emph{potential-side quartic transform} is
\begin{equation}\label{eq:potential-side}
Q_{\cA,v}(s;y)
\;:=\;
\int_{-1/2}^{1/2}
\bigl\langle
Q_\cA^{(1)}(x{+}iy)(v,v,v,v),\;
E^*(x{+}iy,s)
\bigr\rangle\,dx.
\end{equation}
If $Q_{\cA,v}(s;y) = \sum_{N \geq 1}a_{\cA,v}(N;s)\,
e^{-2\pi Ny}$, the Dirichlet lift is
$Q_{\cA,v}(s,u) := \sum_{N \geq 1}a_{\cA,v}(N;s)\,N^{-u}$.
\end{definition}

\begin{definition}[Gram-side transform]%
\label{def:gram-side-transform}%
\index{Gram-side transform|textbf}%
The \emph{Gram-side transform} uses the polarized quartic
contact norm:
\begin{equation}\label{eq:gram-side}
G_{\cA,v}(s;y)
\;:=\;
\int_{-1/2}^{1/2}
q_{\cA,v}^{\mathrm{ct}}(x{+}iy)\,
E^*(x{+}iy,s)\,dx,
\end{equation}
with Fourier expansion
$G_{\cA,v}(s;y)
= \sum_{N \geq 1}g_{\cA,v}(N;s)\,e^{-2\pi Ny}$
and Dirichlet lift
$G_{\cA,v}(s,u)
:= \sum_{N \geq 1}g_{\cA,v}(N;s)\,N^{-u}$.
\end{definition}

\subsection{The $3 \times 3$ moment determinant}%
\label{sec:3x3-determinant}

Fix $u_0 > 1$ in the absolutely convergent half-plane.
For either transform $F \in \{Q, G\}$, define moments
$\mu_n(s;u_0) := (-\partial_u)^n F(s,u)|_{u=u_0}$.

\begin{theorem}[\ClaimStatusProvedHere]%
\label{thm:schur-complement-quartic}%
\index{Schur complement!quartic}%
\index{Hankel determinant!quartic}%
In the centered gauge $\mu_0 = 1$, $\mu_1 = 0$, the first
nontrivial Hankel determinant and its Schur complement are
\begin{equation}\label{eq:hankel-quartic}
D_2
\;=\;
\det\!\begin{pmatrix}
1 & 0 & \mu_2 \\
0 & \mu_2 & \mu_3 \\
\mu_2 & \mu_3 & \mu_4
\end{pmatrix}
\;=\;
\mu_2\,\mu_4 - \mu_3^2 - \mu_2^3,
\end{equation}
\begin{equation}\label{eq:schur-quartic}
\Sigma_2
\;:=\;
\frac{D_2}{D_1}
\;=\;
\mu_4 - \frac{\mu_3^2}{\mu_2} - \mu_2^2,
\qquad
D_1 = \mu_2.
\end{equation}
On a one-dimensional slice with shadow data
$\mu_2 = H_v$ \textup{(}curvature\textup{)},
$\mu_3 = C_v$ \textup{(}cubic shadow\textup{)},
$Q_v^{\mathrm{ct}}$ \textup{(}quartic contact\textup{)}:
\begin{enumerate}[label=\textup{(\roman*)}]
\item
on the potential side,
$\Sigma_2^{\mathrm{Pot}}
= \mu_4 - C_v^2/H_v - H_v^2
= Q_v^{\mathrm{ct}}$\textup;
\item
on the Gram side,
$\Sigma_2^{\mathrm{Gram}} = N_v^{\mathrm{ct}}
= (Q_v^{\mathrm{ct}})^{-1}$.
\end{enumerate}
The potential-side Schur complement extracts the quartic contact
\emph{coupling}; the Gram-side Schur complement extracts the quartic
resonance \emph{norm}. These are reciprocal: $\Sigma_2^{\mathrm{Pot}}
\cdot \Sigma_2^{\mathrm{Gram}} = 1$.
\end{theorem}

\begin{proof}
The Schur complement~\eqref{eq:schur-quartic} subtracts
$\mu_2^2 = H_v^2$ (quadratic self-pairing) and
$\mu_3^2/\mu_2 = C_v^2/H_v$ (cubic tree contamination)
from~$\mu_4$, isolating the quartic contact term.
On the Gram side, the polarized norm replaces coupling
by its inverse.
\end{proof}

\subsection{Virasoro: explicit calculation}%
\label{sec:virasoro-quartic-determinant}

\begin{proposition}[\ClaimStatusProvedHere]%
\label{prop:virasoro-quartic-determinant}%
\index{Virasoro!quartic determinant}%
For the Virasoro family with shadow data
$H_{\mathrm{Vir}} = c/2$,\; $C_{\mathrm{Vir}} = 2$,\;
$Q_{\mathrm{Vir}}^{\mathrm{ct}} = 10/[c(5c{+}22)]$\textup{:}
\begin{enumerate}[label=\textup{(\roman*)}]
\item
the raw quartic moment is
$\mu_4^{\mathrm{Vir}} = c^2/4 + 8/c + 10/[c(5c{+}22)]$\textup;
\item
the Hankel determinant is
$D_2^{\mathrm{Vir}} = 5/(5c{+}22)$\textup;
\item
the potential-side Schur complement is
$\Sigma_2^{\mathrm{Vir}} = 10/[c(5c{+}22)]
= Q_{\mathrm{Vir}}^{\mathrm{ct}}$\textup;
\item
the quartic resonance divisor is
$R_{4,\mathrm{Vir}}^{\mathrm{mod}}\colon c(5c{+}22) = 0$.
\end{enumerate}
\end{proposition}

\begin{proof}
$\mu_2 = c/2$, $\mu_3 = 2$. Then
$\mu_3^2/\mu_2 = 4/(c/2) = 8/c$ and $\mu_2^2 = c^2/4$.
Part~(i): $\mu_4 = \mu_2^2 + \mu_3^2/\mu_2 + Q^{\mathrm{ct}}$.
Part~(ii): $D_2 = \mu_2\cdot Q^{\mathrm{ct}}
= (c/2)\cdot 10/[c(5c{+}22)] = 5/(5c{+}22)$.
Part~(iii): $\Sigma_2 = D_2/D_1 = (5/(5c{+}22))/(c/2)$.
Part~(iv): poles of $\Sigma_2$ at $c = 0$ and $c = -22/5$.
\end{proof}

\subsection{The genuine residue kernel}%
\label{sec:genuine-residue-kernel}

The constrained Epstein zeta
$\varepsilon^c_s = \sum_{\Delta \in S}(2\Delta)^{-s}$
satisfies the functional equation
\begin{equation}\label{eq:constrained-epstein-fe}
\varepsilon^c_{c/2-s}
\;=\;
F_c(s)\;\varepsilon^c_{c/2+s-1},
\qquad
F_c(s)
\;=\;
\frac{\Gamma(s)\,\Gamma(s{+}c/2{-}1)\,\zeta(2s)}
{\pi^{2s-1/2}\,
\Gamma(c/2{-}s)\,\Gamma(s{-}1/2)\,\zeta(2s{-}1)}\,,
\end{equation}
whose contour-shift poles occur at
$s = c/2$, $s = (1{+}\rho_n)/2$, and
$s = (1{+}\bar\rho_n)/2$
for~$\rho_n$ ranging over the nontrivial zeros of~$\zeta$.

Let $\rho = \sigma + i\gamma$ be a nontrivial zero of~$\zeta$
and $s_\rho := (1{+}\rho)/2$.

\begin{definition}[Universal residue factor]%
\label{def:universal-residue-factor}%
\index{residue factor!universal|textbf}%
The residue of~$F_c$ at $s = s_\rho$ is
\begin{equation}\label{eq:universal-residue}
A_c(\rho)
\;:=\;
\operatorname{Res}_{s=s_\rho}F_c(s)
\;=\;
\frac{\Gamma\!\bigl(\tfrac{1+\rho}{2}\bigr)\,
\Gamma\!\bigl(\tfrac{c+\rho-1}{2}\bigr)\,
\zeta(1{+}\rho)}
{2\,\pi^{\rho+\frac{1}{2}}\,
\Gamma\!\bigl(\tfrac{c-\rho-1}{2}\bigr)\,
\Gamma\!\bigl(\tfrac{\rho}{2}\bigr)\,
\zeta'(\rho)}\,.
\end{equation}
\end{definition}

\begin{proposition}[\ClaimStatusProvedHere]%
\label{prop:on-off-line-distinction}%
\index{on-line!vs off-line zeros}%
The paired real residue kernel at conformal
weight~$\Delta$ is
\begin{equation}\label{eq:residue-kernel}
w_{c,\rho,u_0}(\Delta)
\;=\;
2\,|A_c(\rho)|\,(2\Delta)^{-\frac{c+\sigma-1}{2}}\,
\Delta^{-u_0}\,
\cos\!\Bigl(
\tfrac{\gamma}{2}\log(2\Delta) + \arg A_c(\rho)
\Bigr).
\end{equation}
The decay exponent $-(c{+}\sigma{-}1)/2$ depends on
$\sigma = \operatorname{Re}(\rho)$.
For on-line zeros ($\sigma = 1/2$), the exponent
is $-(c{-}1/2)/2$.
For off-line zeros ($\sigma \neq 1/2$), the exponent differs.
\end{proposition}

\begin{proof}
Write $A_c(\rho)(2\Delta)^{-(c+\rho-1)/2}
= |A_c(\rho)|\,(2\Delta)^{-(c+\sigma-1)/2}\,
e^{-i(\gamma/2)\log(2\Delta) + i\arg A_c(\rho)}$.
The paired kernel
$w = 2\operatorname{Re}(A_c(\rho)\,(2\Delta)^{-(c+\rho-1)/2})
\cdot \Delta^{-u_0}$
gives~\eqref{eq:residue-kernel}.
\end{proof}

\subsection{Compatibility ratios and the closure programme}%
\label{sec:closure-programme}%
\index{closure programme|textbf}

\begin{definition}[Genuine residue moment matrix]%
\label{def:genuine-residue-moment-matrix}%
\index{residue moment matrix|textbf}%
Let $d\nu_c^{\mathrm{Vir}}(\Delta)$ be the scalar-primary
spectral measure of a Virasoro CFT at central charge~$c$.
Define raw residue moments
$I_n(c,\rho;u_0) :=
\int (\log\Delta)^n\,w_{c,\rho,u_0}(\Delta)\,
d\nu_c^{\mathrm{Vir}}(\Delta)$
for $n = 0,1,2,3,4$.
After centering and rescaling by $\lambda = (c/2)/v$
to match the Virasoro quadratic datum~$c/2$,
the \emph{genuine normalized Virasoro residue moment matrix} is
\begin{equation}\label{eq:residue-moment-matrix}
M_{2,\mathrm{Vir}}^{\mathrm{res}}(c,\rho;u_0)
\;=\;
\begin{pmatrix}
1 & 0 & c/2 \\
0 & c/2 & m_3^\sharp(c,\rho;u_0) \\
c/2 & m_3^\sharp(c,\rho;u_0) & m_4^\sharp(c,\rho;u_0)
\end{pmatrix}.
\end{equation}
\end{definition}

\begin{definition}[Compatibility ratios]%
\label{def:compatibility-ratios}%
\index{compatibility ratio|textbf}%
The \emph{potential-side compatibility ratio} is
\[
C_{4,\mathrm{Vir}}^{\mathrm{pot}}(c,\rho;u_0)
\;:=\;
\frac{c(5c{+}22)}{10}\,
S_{4,\mathrm{Vir}}^{\mathrm{res}}(c,\rho;u_0),
\]
where $S_{4,\mathrm{Vir}}^{\mathrm{res}}$ is the Schur
complement of~$M_{2,\mathrm{Vir}}^{\mathrm{res}}$.
The \emph{Gram-side compatibility ratio} is
\[
C_{4,\mathrm{Vir}}^{\mathrm{Gram}}(c,\rho;u_0)
\;:=\;
\frac{10}{c(5c{+}22)}\cdot
\frac{1}{S_{4,\mathrm{Vir}}^{\mathrm{res}}(c,\rho;u_0)}\,.
\]
The quartic shadow constraint at degrees $3$ and $4$ is the
pair of equations
$m_3^\sharp = 2$,\;
$S_{4,\mathrm{Vir}}^{\mathrm{res}} = 10/[c(5c{+}22)]$,
equivalently $C_{4,\mathrm{Vir}}^{\mathrm{pot}} = 1$
and $C_{4,\mathrm{Vir}}^{\mathrm{Gram}} = 1$.
\end{definition}

\begin{conjecture}[Quartic closure]%
\label{conj:quartic-closure}%
\ClaimStatusConjectured%
\index{quartic closure conjecture}%
\index{Riemann hypothesis!quartic closure}%
Define the quartic compatibility locus
\begin{equation}\label{eq:compatibility-locus}
L_{c,u_0}^{\mathrm{Vir}}
\;:=\;
\bigl\{\rho :
C_{3,\mathrm{Vir}}^{(3)}(c,\rho;u_0) = 1,\;
C_{4,\mathrm{Vir}}^{\mathrm{Gram}}(c,\rho;u_0) = 1
\bigr\}.
\end{equation}
Then:
\begin{enumerate}[label=\textup{(\roman*)}]
\item
\textup{(}Quartic on-line\textup{)} for genuine Virasoro
spectral input, $\operatorname{Re}(\rho) = 1/2$ implies
$\rho \in L_{c,u_0}^{\mathrm{Vir}}$ for all
$c, u_0$\textup;
\item
\textup{(}Quartic closure\textup{)}
$\displaystyle\bigcap_{c \in C}\;\bigcap_{u_0 \in U}
L_{c,u_0}^{\mathrm{Vir}}
= \{\rho : \operatorname{Re}(\rho) = 1/2\}$\textup;
\item
\textup{(}Strong closure\textup{)}
$\displaystyle\bigcap_{\cA \in \{\mathrm{Vir},\cW_3,\cW_4,\ldots\}}
\;\bigcap_{\text{slices }v}
\;\bigcap_{c,u_0}\;
L_{c,u_0}^{\cA,v}
= \{\rho : \operatorname{Re}(\rho) = 1/2\}$.
\end{enumerate}
\end{conjecture}

\subsection{Four structural gaps}%
\label{sec:four-gaps}%
\index{arithmetic programme!gaps}

The four gaps identify where the MC integrability
constraint ($D_\cA^2 = 0$) falls short of the positivity
constraint needed for spectral conclusions.

\begin{center}
\small
\renewcommand{\arraystretch}{1.3}
\begin{tabular}{clll}
\toprule
\emph{Gap} & \emph{Description} & \emph{Resolution}
& \emph{Missing ingredient} \\
\midrule
A & MC $\to$ Li positivity
& Structurally blocked (Prop.~\ref{prop:li-criterion-failure})
& Two-line obstruction \\
B & Euler product from sewing
& $\sigma_{-1}$ found
& Genuine spectral measure \\
C & Spectral rigidity from bar-cobar
& Chain-level strong
& Euler compatibility \\
D & Nyman--Beurling from Koszulness
& Structural analogy
& Norm mismatch \\
\bottomrule
\end{tabular}
\end{center}

\noindent
\emph{Gap~A.} The Rankin--Selberg transform from MC
to Li is not sign-preserving; closing the chain requires
a surface polarization replacing scalar Mellin transforms
with intersection pairings.
\emph{Gap~B.} The Narain universality
($\varepsilon^1_s(R) = 2(R^{2s} + R^{-2s})\zeta(2s)$)
shows that the character factor carries no zeros; the
Narain family provides no bootstrap constraint. The
genuine bottleneck is passing to the Virasoro
scalar-primary spectral measure.
\emph{Gap~C.} Bar-cobar equivalence is homotopical, not
spectral-positive; Davenport--Heilbronn off-critical zeros
show that bare bar-cobar symmetry cannot substitute for
Euler-product rigidity.
\emph{Gap~D.} The Nyman--Beurling $L^2(0,1)$ norm differs
from the bar-cobar deformation norm.

\begin{remark}[Root cause: integrability versus positivity]%
\label{rem:beilinson-four-gaps}%
\index{arithmetic programme!integrability vs positivity}%
The four gaps share a common theme. The MC equation
$D_\cA \cdot \Theta + \tfrac12[\Theta,\Theta] = 0$
is an \emph{integrability} constraint
(flatness, quadratic consistency), while every spectral
conclusion in this chapter (positivity of Li coefficients,
critical-line location, Ramanujan bounds) requires
\emph{positivity}
(sign-definiteness, absolute value bounds).
A flat connection can have eigenvalues of any absolute value;
integrability and positivity are logically independent.
Gaps~A and~C are pure instances of this dichotomy;
Gaps~B and~D involve additional structural obstructions.

Sharpening the four gaps in this light:
\begin{enumerate}[label=\textup{(\roman*)}]
\item
\emph{Gap~A\textup{:} sign failure.}
The Rankin--Selberg transform
$\operatorname{Sh}_r \mapsto M_r(s)$
(Theorem~\ref{thm:mc-recursion-moment})
is a Mellin integral, and Mellin transforms of
positive functions can be negative.
``Surface polarisation'' would require a positive-definite
pairing on a modular surface under which the MC data have
definite sign. No candidate pairing has been constructed;
the problem is open.
\item
\emph{Gap~B\textup{:} Narain triviality.}
The Narain character factor
$2(R^{2s} + R^{-2s})$ is positive for real~$s$
and carries no zeros:
all zero-location information comes from the
universal factor $\zeta(2s)$, independent of~$\cA$.
Family-specific information requires the genuine
scalar-primary spectral measure of a non-free-field
family (Virasoro, $\cW_N$), which is not available
in closed form at generic~$c$.
\item
\emph{Gap~C\textup{:} sharp counterexample class.}
Davenport--Heilbronn Epstein zeta functions
(class number $h(D) \geq 2$; Tier~3
of Theorem~\ref{thm:euler-koszul-tier-classification})
have functional equation, meromorphic continuation,
\emph{and} zeros with
$\operatorname{Re}(s) > 1/2$.
Bar-cobar homotopy equivalence is stronger than a
functional equation, but its strength is
homotopical (chain-level), not spectral-positive.
\item
\emph{Gap~D\textup{:} norm mismatch.}
The Nyman--Beurling $L^2(0,1)$ norm and the norm on
$\mathrm{Def}_{\mathrm{cyc}}(\cA)$ live in different
functional-analytic categories;
no bounded embedding between them is known.
\end{enumerate}
\end{remark}

\begin{remark}[Collapse permanence]%
\label{rem:collapse-permanence}%
\index{sewing--Hecke reciprocity!collapse permanence}%
By Proposition~\ref{prop:collapse-permanence},
the interacting sewing kernel
$K_\cA^{\leq R}(s,t)$ factors through $s{+}t$ at every
degree cutoff; the Sewing--Hecke collapse
(Theorem~\ref{thm:sewing-hecke-reciprocity}) is
permanent: the shadow tower modifies the Dirichlet
coefficients of~$S_\cA$ but cannot break the
one-variable reduction. The zero-side and prime-side
remain locked by Fock diagonality of the genus-$1$
sewing operator.
A genuinely two-variable arithmetic structure
would require breaking this diagonality, for instance
non-diagonal sewing at genus~$g \geq 2$
(where the period matrix couples Fock sectors;
see~\S\ref{sec:genus2-arithmetic-frontier} below)
or off-shell deformations.
\end{remark}

\begin{remark}[Nyman--Beurling Gram analogy]%
\label{rem:nyman-beurling-gram-analogy}%
\index{Nyman--Beurling!Gram analogy}%
The sewing Gram matrix $\mathscr{M}_{ij} =
K_\cA(\alpha/2{+}i,\,\alpha/2{+}j)$ and the Nyman--Beurling
Gram matrix $G_{ij} = \langle\rho_{\theta_i},\rho_{\theta_j}
\rangle_{L^2}$ are both Gram matrices in zeta-weighted inner
products. The Mellin transform of $\{1/x\}$ is $-\zeta(s)/s$
(for $-1 < \operatorname{Re}(s) < 0$), which involves the
same $\zeta$-function that appears in the sewing kernel.
The sewing kernel $\zeta(s{+}t)\zeta(s{+}t{+}1)$ lives in
$\operatorname{Re}(s) > 1$; the Nyman--Beurling kernel
involves $\zeta(s)\zeta(1{-}s)$ across the critical line.
\end{remark}

\begin{remark}[Genus-$2$ Beurling kernel: Gap~D narrowing]%
\label{rem:genus2-beurling-kernel}%
\index{Nyman--Beurling!genus-2 narrowing}%
\index{Gap D!genus-2 resolution}%
The genus-$1$ norm mismatch
(Remark~\ref{rem:nyman-beurling-gram-analogy})
arises because the sewing kernel
$\zeta(s{+}t)\zeta(s{+}t{+}1)$ and the Nyman--Beurling
kernel $\zeta(s)\zeta(1{-}s)$ have \emph{disjoint spectral
support}: the former lives in $\operatorname{Re}(s) > 1$,
the latter straddles the critical line.
The B\"ocherer bridge
(Theorem~\ref{thm:bocherer-bridge},
\S\ref{subsec:bocherer-bridge-theorem})
eliminates this obstruction at genus~$2$.

Define the \emph{genus-$2$ Beurling kernel} for an even
unimodular lattice $\Lambda$ of rank~$r$:
\begin{equation}\label{eq:genus2-beurling-kernel}
\mathscr{K}_\Lambda^{(2)}(D,\,D')
\;:=\;
\sum_{f \in \mathscr{B}_k}
\frac{B_f(D)\,\overline{B_f(D')}}
{\langle f,\, f \rangle},
\end{equation}
where $\mathscr{B}_k$ is a Hecke eigenbasis of
$S_{r/2}(\mathrm{Sp}(4,\bZ))$ and
$B_f(D)$ is the B\"ocherer coefficient
\eqref{eq:bocherer-coefficient-def}.
Then:
\begin{enumerate}[label=\textup{(\roman*)}]
\item
$\mathscr{K}_\Lambda^{(2)}$ is positive semi-definite
(sum of rank-one Hermitian forms).
\item
By the B\"ocherer factorization
\eqref{eq:bocherer-factorization}, the diagonal is
\begin{equation}\label{eq:beurling-kernel-diagonal}
\mathscr{K}_\Lambda^{(2)}(D,\,D)
\;=\;
\sum_{f \in \mathscr{B}_k}
\frac{\lvert B_f(D)\rvert^2}
{\langle f,\, f \rangle}
\;\geq\; 0,
\end{equation}
where $B_f(D)$ is the B\"ocherer coefficient
\eqref{eq:bocherer-coefficient-def} of the eigenform~$f$.
For genuine \textup{(}non-SK\textup{)} eigenforms, the
refined B\"ocherer conjecture~\eqref{eq:bocherer-factorization}
expresses~\eqref{eq:beurling-kernel-diagonal} as a product
of central $L$-values $L(1/2,\pi_f)\,L(1/2,\pi_f\times\chi_D)$
at $s = 1/2$, on the critical line, not at
$\operatorname{Re}(s) > 1$.
For Saito--Kurokawa eigenforms $f = \mathrm{SK}(g)$,
the formula reduces via Waldspurger to
$\lvert B_f(D)\rvert^2 \propto L(k{-}1,\, g \times \chi_D)
\cdot \lvert D\rvert^{k-2}$
\textup{(}Theorem~\textup{\ref{thm:leech-chi12-projection}(iii)}\textup{)},
which also involves central $L$-values of the underlying
elliptic eigenform on the critical line.
\item
The spectral support of $\mathscr{K}_\Lambda^{(2)}$
matches that of the Nyman--Beurling kernel: both involve
$\zeta$-family values on and around $\operatorname{Re}(s)
= 1/2$. The genus-$1$ spectral-support mismatch
\textup{(}Remark~\ref{rem:nyman-beurling-gram-analogy}\textup{)}
is thereby resolved.
\item
The residual gap:
$\mathscr{K}_\Lambda^{(2)}$ involves $L$-values at a
\emph{single point} $s = 1/2$, while Nyman--Beurling
requires the full $L^2(0,1)$ norm (all points on the
critical line).
The escalation principle
(Remark~\ref{rem:bocherer-escalation})
fills in the missing spectral data: genus-$g$ MC
contributes
$L(1/2, \mathrm{Sym}^{g-1} f \times \chi_D)$,
and as $g \to \infty$ these central values cover all
automorphic spectral data on the critical line.
The all-genera Beurling kernel
$\mathscr{K}_\Lambda^{(\infty)} := \varprojlim_g
\mathscr{K}_\Lambda^{(g)}$
would have the same spectral content as the
Nyman--Beurling kernel, narrowing the gap
\textup{(}conjectural at genus $g \geq 3$;
circular at $g \geq 6$ without Newton--Thorne,
cf.\ Remark~\textup{\ref{rem:escalation-circularity}}\textup{)}.
\end{enumerate}
\end{remark}

\begin{remark}[Escalation circularity at $g \geq 6$]%
\label{rem:escalation-circularity}%
\index{escalation principle!circularity at genus >= 6}%
\index{Newton--Thorne!symmetric power functoriality}%
\index{Langlands functoriality!GL(2) to GL(g)}%
The escalation principle (Remark~\ref{rem:bocherer-escalation})
at genus~$g$ accesses $L(1/2, \operatorname{Sym}^{g-1} f)$,
which requires automorphy of~$\operatorname{Sym}^{g-1}$,
i.e.\ Langlands functoriality
$\mathrm{GL}(2) \to \mathrm{GL}(g)$.
This is known for $g \leq 5$
(Gelbart--Jacquet~$g{=}3$, Kim--Shahidi~$g{=}4$, Kim~$g{=}5$)
and \emph{open} for $g \geq 6$.
Newton--Thorne~\cite{NewtonThorne2021} prove automorphy
of~$\operatorname{Sym}^n$ unconditionally for \emph{holomorphic}
cuspidal eigenforms of regular algebraic weight and all~$n$,
breaking the circularity for the lattice VOA programme
(where all relevant eigenforms are holomorphic).
For Maass forms, the circularity persists: the best
unconditional bound remains Kim--Sarnak
$\theta \leq 7/64$.
\end{remark}

\begin{remark}[Prime-side Li coefficients: numerical values]%
\label{rem:li-numerical-values-arith}%
\index{Li coefficients!numerical values}%
The first six prime-side Li coefficients
(Construction~\ref{def:prime-side-li}) at 50-digit
precision:
\begin{center}
\small
\renewcommand{\arraystretch}{1.1}
\begin{tabular}{lrrrr}
\toprule
$n$ & $\cH$ & $\mathrm{Vir}$ & $\cW_3$ & $\cW_4$ \\
\midrule
$1$ & $+0.00725$ & $-0.993$ & $-1.243$ & $-1.437$ \\
$2$ & $+0.420$ & $-1.425$ & $-1.853$ & $-2.181$ \\
$3$ & $+0.553$ & $-2.470$ & $-3.116$ & $-3.601$ \\
$4$ & $+0.519$ & $-4.415$ & $-5.499$ & $-6.321$ \\
$5$ & $+0.359$ & $-7.728$ & $-9.755$ & $-11.322$ \\
$6$ & $+0.090$ & $-13.157$ & $-17.102$ & $-20.450$ \\
\bottomrule
\end{tabular}
\end{center}
Heisenberg: positive for $n \leq 6$, first negative at $n = 7$.
Virasoro and $\cW_N$: all negative, exponential growth
$|\tilde\lambda_n| \sim C \cdot 2^n$.
\end{remark}

\begin{proposition}[Structural failure of the Li criterion for the
sewing lift; \ClaimStatusProvedHere]%
\label{prop:li-criterion-failure}%
\index{Li coefficients!structural failure}%
\index{Li criterion!inapplicability to sewing lift}%
The prime-side Li coefficients
$\tilde\lambda_n(\cA)$
\textup{(}Definition~\textup{\ref{def:prime-side-li}}\textup{)}
are eventually negative for \emph{every} standard family.
The cause is spectral, not arithmetic\textup{:}
\begin{enumerate}[label=\textup{(\roman*)}]
\item
The regularised sewing lift
$\Xi_\cA(u) = (u{-}1)\,S_\cA(u)$
has nontrivial zeros on \emph{at least two} critical lines.
For the Heisenberg,
$\Xi_\cH(u) = (u{-}1)\,\zeta(u)\,\zeta(u{+}1)$
has zeros at the nontrivial zeros of~$\zeta(u)$
\textup{(}on $\operatorname{Re}(u) = 1/2$ under GRH\textup{)}
\emph{and} at the nontrivial zeros of~$\zeta(u{+}1)$
\textup{(}on $\operatorname{Re}(u) = -1/2$ under GRH\textup{)}.
\item
The Li criterion requires all zeros of the completed function
on a \emph{single} line $\operatorname{Re}(u) = 1/2$; zeros
on two lines $\operatorname{Re}(u) = \pm 1/2$ force the
Li coefficients $\tilde\lambda_n = \sum_\rho
[1 - (1{-}1/\rho)^n]$ to be eventually
negative, even when RH holds.
\item
For Virasoro and~$\cW_N$, the harmonic defect terms
$H_{w-1}(u)$ are finite Dirichlet polynomials with zeros
at $\operatorname{Re}(u)$ ranging over all of~$\bR$; the
Li coefficients are negative at \emph{every} order.
\end{enumerate}
The Li criterion is therefore structurally inapplicable
to the sewing lift $S_\cA(u)$. Any approach to zero-location
questions through the shadow tower must bypass the Li
framework; the genus-$2$ B\"ocherer bridge
\textup{(}\S\textup{\ref{subsec:bocherer-bridge-theorem}}\textup{)}
is one such bypass.
\end{proposition}

\begin{proof}
Part~(i): the zero loci follow from the product structure
$\Xi_\cH = (u{-}1)\zeta(u) \cdot \zeta(u{+}1)$.
Part~(ii): if $\rho$ is a zero with
$\operatorname{Re}(\rho) = -1/2$, then
$|1 - 1/\rho| > 1$, so
$(1{-}1/\rho)^n$ grows exponentially and the
contribution $1 - (1{-}1/\rho)^n$ to
$\tilde\lambda_n$ is exponentially negative.
This dominates the positive contributions
from zeros on $\operatorname{Re}(\rho) = 1/2$,
where $|1 - 1/\rho| < 1$ and the contribution is
positive but bounded.
Part~(iii): the Virasoro sewing lift is
$S_{\mathrm{Vir}}(u) = \zeta(u{+}1)(\zeta(u) - 1)$.
The factor $\zeta(u) - 1$ has infinitely many zeros
with $\operatorname{Re}(u) > 1$
\textup{(}Bohr--Landau, see
Titchmarsh~\cite{Titchmarsh-zeta}, \S11.5\textup{)};
these contribute exponentially negative terms at every
order.
\end{proof}

\begin{remark}[The sewing lift as two-line spectral object]%
\label{rem:sewing-lift-two-line}%
\index{Dirichlet--sewing lift!two-line spectral structure}%
The sewing lift $S_\cA(u)$ is not a single
$L$-function: it is a finite combination of products of
$\zeta$-functions with shifts and finite Dirichlet polynomials
\textup{(}Theorem~\textup{\ref{thm:dirichlet-weight-formula}}\textup{)}.
Its zero set is the union of the shifted zero loci of its
factors. For the Heisenberg, these are the zeros of
$\zeta(u)$ and $\zeta(u{+}1)$, lying on two parallel lines
separated by unit distance. For~$\cW_N$, the harmonic defect
$\sum H_j(u)$ introduces additional zeros at positions
depending on~$N$.

The consequence: the Taylor expansion of
$\log\Xi_\cA$ in the variable $w = 1 - u^{-1}$ is
\begin{equation}\label{eq:sewing-exponential-gf}
\log \Xi_\cA\!\left(\frac{1}{1{-}w}\right)
\;=\;
\log\Xi_\cA(1)
\;+\;
\sum_{n=1}^\infty
\frac{\tilde\lambda_n(\cA)}{n}\,w^n\,,
\end{equation}
where $\Xi_\cA(1) = S_\cA(1)$ is the total sewing weight
and the series converges for $|w| < |1 - 1/\rho_0|^{-1}$ with
$\rho_0$ the zero of~$\Xi_\cA$ nearest to $u = 1$.
This generating function encodes the Li coefficients;
the spectral representation
$\tilde\lambda_n = \sum_\rho[1 - (1{-}1/\rho)^n]$ relates
positivity of $\tilde\lambda_n$ to the zero locus of~$\Xi_\cA$.
For the Heisenberg,
$\log\Xi_\cH(u) = \log[(u{-}1)\zeta(u)] + \log\zeta(u{+}1)$
decomposes into a single-line component
$\log[(u{-}1)\zeta(u)]$ whose Li coefficients satisfy the
classical criterion and a shifted component $\log\zeta(u{+}1)$
whose zeros are displaced by one unit, destroying eventual
positivity.
\end{remark}

\begin{remark}[Independent-sum factorization and multiplicativity]
\label{rem:euler-product-from-independent-sum}
\index{Euler product!from independent-sum factorization}
\index{independent sum factorization!Fredholm determinant}
The independent-sum factorization
(Proposition~\ref{prop:independent-sum-factorization}) asserts
that for $\cA = \cA_1 \otimes \cA_2$ with vanishing mixed OPE,
all shadows separate: $\kappa$ is additive,
$\Delta$ multiplicative,
$R_4^{\mathrm{mod}}$ additive. At the sewing-envelope level,
this gives a multiplicative decomposition of the sewing
Fredholm determinant:
\[
\det\nolimits_{\cA_1\otimes\cA_2}(1 - K_q)
\;=\;
\det\nolimits_{\cA_1}(1 - K_q)
\cdot
\det\nolimits_{\cA_2}(1 - K_q).
\]
For the Heisenberg algebra, the Fredholm determinant is
$\prod_{n \geq 1}(1 - q^n)^{-1}$
and the Dirichlet series of $\sigma_{-1}$ gives the classical
Euler product
$\zeta(s)\zeta(s{+}1) = \prod_p
(1{-}p^{-s})^{-1}(1{-}p^{-s-1})^{-1}$
(Corollary~\ref{cor:sewing-euler-product}).

The lesson for Gap~B: the independent-sum
factorization provides \emph{multiplicativity} of the sewing
determinant under tensor decomposition of chiral algebras,
but this algebraic multiplicativity does \emph{not}
produce an Euler product indexed by primes unless the
algebra itself decomposes prime-by-prime (as the Heisenberg
does via $\sigma_{-1}$). For non-free-field families, the
passage from multiplicative Fredholm determinant to
prime-indexed Euler product requires additional arithmetic
input beyond the shadow tower.
\end{remark}

\subsection{The $\cW_N$ multi-channel extension}%
\label{sec:wn-multichannel}

For the principal $\cW_N$-algebra, generators of spins
$2, 3, \ldots, N$ produce $N{-}1$ independent curvature
channels with $H_s = c/s$ for $s = 2,\ldots,N$.

\begin{proposition}[\ClaimStatusProvedHere]%
\label{prop:pure-spin-s-schur}%
\index{W-algebra@$\cW$-algebra!pure spin-$s$ Schur complement}%
On the pure spin-$s$ slice \textup{(}$s \geq 3$\textup{)},
the cubic shadow vanishes: $C_s = 0$. The Schur
complement reduces to
\begin{equation}\label{eq:pure-spin-s-schur}
\Sigma_{2,s}
\;=\;
\mu_{4,s} - H_s^2
\;=\;
\mu_{4,s} - (c/s)^2,
\end{equation}
and the quartic resonance divisor on this slice is
$R_{4,s}^{\mathrm{mod}}\colon\Sigma_{2,s} = 0$.
\end{proposition}

\begin{proof}
On a pure spin-$s$ slice, the only cubic interaction
is the $TW^{(s)}$ cross-term, which vanishes when
restricted to the spin-$s$ subspace alone (no gravitational
cubic vertex at spin~$s \geq 3$).
The Schur complement~\eqref{eq:schur-quartic}
with $\mu_3 = 0$ gives $\Sigma_2 = \mu_4 - \mu_2^2$.
\end{proof}

For $\cW_3$, the pure spin-$3$ quartic contact coupling is
(numerical computation)
\begin{equation}\label{eq:w3-spin3-quartic}
Q_3^{\mathrm{ct}}(\cW_3)
\;=\;
\frac{220}{c(5c{+}22)^2}\,,
\end{equation}
with a double pole at $c = -22/5$ (compared to the
simple pole of $Q_2^{\mathrm{ct}} = 10/[c(5c{+}22)]$).

\begin{remark}[Propagator variance and Euler defect]%
\label{rem:propagator-variance-euler}%
\index{propagator variance!Euler--Koszul}%
For $\cW_3$ with generators $T$ ($\kappa_T = c/2$) and $W$
($\kappa_W = c/3$), the propagator variance
$\delta_{\mathrm{mix}} = \sum f_i^2/\kappa_i -
(\sum f_i)^2/\sum\kappa_i \geq 0$ (Cauchy--Schwarz)
measures failure of the sewing kernel to factor as a product
of single-generator kernels.
The Euler--Koszul class $\operatorname{ek}(\cW_3) = 2$
records the same phenomenon in Dirichlet series:
$S_{\cW_3}(u) \neq 2\,\zeta(u)\zeta(u{+}1)$ precisely
because the harmonic defect couples the spin-$2$ and spin-$3$
channels.
\end{remark}

\subsection{Target computations}%
\label{sec:target-computations}

\begin{problem}[Virasoro residue evaluation]%
\label{prob:virasoro-residue-evaluation}%
\index{Virasoro!residue evaluation}%
Choose an explicit Virasoro CFT family with known
scalar-primary spectrum $d\nu_c^{\mathrm{Vir}}(\Delta)$
(the Ising model at $c = 1/2$, the three-state Potts model
at $c = 4/5$, or a generic minimal model $\cM(p,q)$).
Evaluate the residue moments
$I_0(c,\rho_1;u_0), \ldots, I_4(c,\rho_1;u_0)$
at the first zeta zero $\rho_1 = 1/2 + 14.134\ldots\,i$
for $u_0 = 2$.
Form the normalized residue moment
matrix~$M_{2,\mathrm{Vir}}^{\mathrm{res}}$
and compute the compatibility ratios
$C_{3,\mathrm{Vir}}^{(3)}(c,\rho_1;2)$ and
$C_{4,\mathrm{Vir}}^{\mathrm{Gram}}(c,\rho_1;2)$.
Test: $C = 1$ for $\rho_1$ on the line;
$C \neq 1$ for a hypothetical off-line
$\rho' = 0.6 + 14.134\ldots\,i$.
\end{problem}

\subsection{The Miura defect decomposition}%
\label{sec:miura-defect-decomposition}%
\index{Miura defect!residue matrix|textbf}%
\index{residue moment matrix!Miura defect}

The Miura map $\hat\Upsilon\colon \cW_\kappa(\fg) \hookrightarrow
V^{\mathfrak{h}}_\kappa(\mathfrak{h})$ is injective (Arakawa; for $\fg = \mathfrak{sl}_2$
the image is the universal Virasoro vertex algebra).
At the level of residue moment matrices, injectivity forces a
decomposition into a projected Heisenberg piece and a finite-rank
correction.

\begin{definition}[Zero-side Miura defect]%
\label{def:zero-side-miura-defect}%
\index{Miura defect!zero-side|textbf}%
For principal $\cW_N$ with generators of spins $2,3,\ldots,N$,
the \emph{zero-side Miura defect} is
\begin{equation}\label{eq:zero-side-miura-defect}
D_{\cW_N}^{\mathrm{res}}(c,\rho;u_0)
\;:=\;
M_{2,\cW_N}^{\mathrm{res}}(c,\rho;u_0)
\;-\;
P_N^\top\,
M_{2,\cH^{N-1}}^{\mathrm{res}}(c,\rho;u_0)\,
P_N\,,
\end{equation}
where $P_N$ is the projection from the $(N{-}1)$-fold
Heisenberg Fock space to the $\cW_N$-primary subspace.
For Virasoro $(N = 2)$, the correction captures the missing
weight-$1$ mode.
\end{definition}

\begin{proposition}[\ClaimStatusProvedHere]%
\label{prop:prime-side-defect-formula}%
\index{Miura defect!prime-side|textbf}%
On the prime side, the Miura defect is explicit\textup{:}
\begin{equation}\label{eq:prime-side-miura-defect}
D_{\cW_N}^{\mathrm{prime}}(u)
\;:=\;
S_{\cW_N}(u) - (N{-}1)\,S_\cH(u)
\;=\;
-\zeta(u{+}1)\sum_{j=1}^{N-1}H_j(u)\,.
\end{equation}
For Virasoro,
$D_{\mathrm{Vir}}^{\mathrm{prime}}(u) = -\zeta(u{+}1)$.
\end{proposition}

\begin{proof}
From the weight-multiset formula
(Theorem~\ref{thm:dirichlet-weight-formula}):
$S_{\cW_N}(u) = \zeta(u{+}1)\bigl((N{-}1)\zeta(u) -
\sum_{j=1}^{N-1}H_j(u)\bigr)$, while
$S_\cH(u) = \zeta(u)\,\zeta(u{+}1)$. The difference is
$-\zeta(u{+}1)\sum H_j(u)$, a product of $\zeta(u{+}1)$
with a finite harmonic polynomial.
\end{proof}

\begin{conjecture}[Defect unification]%
\label{conj:defect-unification}%
\ClaimStatusConjectured%
\index{Miura defect!unification conjecture}%
The zero-side defect $D_{\cW_N}^{\mathrm{res}}$ and the
prime-side defect $D_{\cW_N}^{\mathrm{prime}}$ are two
genus-$1$ shadows of the same bar-intrinsic MC
correction\textup{:}
there exists a canonical element
$\delta\Theta_{\cW_N} \in
\mathrm{Def}_{\mathrm{cyc}}^{\mathrm{mod}}(\cW_N)$
such that
\begin{enumerate}[label=\textup{(\roman*)}]
\item
its Rankin--Selberg shadow at the pole $s = s_\rho$
is the zero-side defect $D_{\cW_N}^{\mathrm{res}}(c,\rho;u_0)$
for each nontrivial zero~$\rho$\textup;
\item
its Mellin transform at $u = v$ is the prime-side defect
$D_{\cW_N}^{\mathrm{prime}}(v)$\textup;
\item
$\delta\Theta_{\cW_N} = \Theta_{\cW_N} - \iota_*\Theta_\cH$
is the image of the DS reduction map acting on the MC class.
\end{enumerate}
Structural meaning: \emph{Drinfeld--Sokolov reduction
is visible arithmetically as a single finite-rank correction
to the Euler product, observed through two dual genus-$1$
transforms.}
\end{conjecture}

\begin{remark}[Virasoro as the first Miura defect]%
\label{rem:virasoro-first-miura-defect}%
\index{Virasoro!as Miura defect}%
The Virasoro case is the simplest:
$D_{\mathrm{Vir}}^{\mathrm{prime}}(u) = -\zeta(u{+}1)$
is the removal of the weight-$1$ mode from the Heisenberg
Euler product.
On the zero side, this is conjectured to correspond to the removal
of the $n = 1$ term from the Fock-diagonal sewing kernel,
whose residue at a zeta-zero pole $\rho$ is the $\Delta = 1$
contribution to the paired real residue kernel
$w_{c,\rho,u_0}$ of
Proposition~\ref{prop:on-off-line-distinction}.
\end{remark}

\begin{theorem}[Finite Miura defect at genus one; \ClaimStatusProvedHere]%
\label{thm:finite-miura-defect}%
\index{Miura defect!finite|textbf}%
\index{vacuum character!splitting}%
\index{Drinfeld--Sokolov!finite defect}%
For principal $\cW_N$ with generators of spins
$2, 3, \ldots, N$, the genus-$1$ arithmetic data
decompose into an $(N{-}1)$-fold Heisenberg core and
a finite defect at three equivalent levels\textup{:}
\begin{enumerate}[label=\textup{(\roman*)}]
\item
\textup{(}Vacuum character.\textup{)}
\begin{equation}\label{eq:vacuum-character-splitting}
\chi_{\cW_N}(q)
\;=\;
\chi_\cH(q)^{N-1}
\cdot
\prod_{m=1}^{N-1}(1 - q^m)^{N-m}\,,
\end{equation}
where $\chi_\cH(q) = \prod_{n \geq 1}(1 - q^n)^{-1}$
is the rank-$1$ Heisenberg character.
\item
\textup{(}Connected free energy.\textup{)}
\begin{equation}\label{eq:connected-free-energy-splitting}
F_{\cW_N}^{\mathrm{conn}}(q)
\;=\;
(N{-}1)\,F_\cH^{\mathrm{conn}}(q)
\;-\;
\sum_{m=1}^{N-1}(N{-}m)\sum_{r \geq 1}
\frac{q^{mr}}{r}\,.
\end{equation}
The defect is a finite sum of geometric series\textup{:}
only modes $m = 1, \ldots, N{-}1$ contribute, each generating
$-\sum_{r \geq 1} q^{mr}/r = \log(1 - q^m)$.
\item
\textup{(}Dirichlet--sewing lift.\textup{)}
\begin{equation}\label{eq:dirichlet-defect-splitting}
S_{\cW_N}(u)
\;=\;
(N{-}1)\,\zeta(u)\,\zeta(u{+}1)
\;-\;
\zeta(u{+}1)\sum_{m=1}^{N-1}(N{-}m)\,m^{-u}\,.
\end{equation}
The defect term
$D_{\cW_N}^{\mathrm{prime}}(u)
= -\zeta(u{+}1)\sum_{m=1}^{N-1}(N{-}m)\,m^{-u}$
is a product of $\zeta(u{+}1)$ with a \emph{finite}
Dirichlet polynomial.
\end{enumerate}
The three levels connect via the correspondences
$-\log \leftrightarrow \mathrm{Mellin}$\textup{:}
the finite product~\textup{(i)} becomes the finite
$q$-series defect in~\textup{(ii)} via logarithm, and
the finite Dirichlet polynomial in~\textup{(iii)} via
$q^{mr}/r \mapsto (mr)^{-u}/r = m^{-u}\,r^{-(u+1)}$.
\end{theorem}

\begin{proof}
(i)~The $\cW_N$ character is
$\chi_{\cW_N}(q) = \prod_{n \geq 2}(1 - q^n)^{-\min(n-1,\,N-1)}$
(Proposition~\ref{prop:wn-character-primitive}). Writing
$\chi_\cH(q)^{N-1} = \prod_{n \geq 1}(1 - q^n)^{-(N-1)}$,
\[
\frac{\chi_{\cW_N}(q)}{\chi_\cH(q)^{N-1}}
\;=\;
(1 - q)^{N-1}
\prod_{n=2}^{N-1}(1 - q^n)^{(N-1) - (n-1)}
\;=\;
\prod_{m=1}^{N-1}(1 - q^m)^{N-m}\,,
\]
since the exponent $(N{-}1) - \min(n{-}1,\,N{-}1)$ vanishes
for $n \geq N$ and equals $N{-}n$ for $2 \leq n \leq N{-}1$.

(ii)~Follows from $-\log$ applied to~(i), using
$-\log(1 - q^m) = \sum_{r \geq 1}q^{mr}/r$.

(iii)~The Dirichlet transform of $q^{mr}/r$ at $q = e^{-2\pi y}$
contributes $m^{-u}\,r^{-(u+1)}$; summing over $r$
gives $m^{-u}\,\zeta(u{+}1)$. Alternatively, the harmonic sum
identity
$\sum_{j=1}^{N-1}H_j(u) = \sum_{m=1}^{N-1}(N{-}m)\,m^{-u}$
(Corollary~\ref{cor:virasoro-mode-removal}, proof)
converts from the $H_j$-form to the weight-multiset form.
\end{proof}

\begin{remark}[The finite defect principle]%
\label{rem:finite-defect-principle}%
\index{finite defect principle|textbf}%
Theorem~\ref{thm:finite-miura-defect} has a decisive
structural consequence:
\emph{all genuinely new genus-$1$ arithmetic content
in $\cW_N$ beyond Heisenberg is carried by a finite
polynomial in~$q$.}

The Heisenberg core $\chi_\cH(q)^{N-1}$ is exact
Euler--Koszul: its sewing determinant factors as
$\zeta(u)\,\zeta(u{+}1)$ with full Euler product.
It carries no quartic resonance, no cubic shadow, and no
family-specific spectral data.

The quartic resonance class
$\mathfrak{R}_{4,g,n}^{\mathrm{mod}}(\cA)$ and the
quartic contact coupling $Q_{\mathrm{Vir}}^{\mathrm{ct}}$
live entirely in the \emph{defect sector}.

The Heisenberg core is quartically blind;
\emph{quartic rigidity lives in the defect.}
This reorganises the research programme:
\begin{enumerate}[label=\textup{(\roman*)}]
\item
the exact Euler--Koszul sector provides the
universal scaffolding (Euler product, $\sigma_{-1}$,
functional equation) but no zero-location information;
\item
the finite defect sector $\prod(1 - q^m)^{N-m}$
carries the family-specific content: the higher-spin
subtraction, the quartic contact coupling, the
non-scalar spectral data;
\item
the residue moment matrix $M_{2,\cA}^{\mathrm{res}}$
separates the two sectors: the Heisenberg entries are
universal, while the defect entries carry the
on-line/off-line distinction.
\end{enumerate}
\end{remark}

\subsection{The residue clutching defect}%
\label{sec:residue-clutching-defect}%
\index{clutching defect!residue|textbf}%
\index{quartic residue programme!clutching defect}

The quartic resonance class $\mathfrak{R}_{4,g,n}^{\mathrm{mod}}(\cA)$
satisfies the proved clutching law
(Theorem~\ref{thm:nms-clutching-law-modular-resonance}).
The residue analogue and the geometric
zero test are defined next.

\begin{definition}[Residue divisor]%
\label{def:residue-divisor}%
\index{residue divisor|textbf}%
For a modular Koszul family~$\cA$, a cyclic direction~$v$,
and a candidate zero~$\rho$, the \emph{residue quartic divisor}
is
\begin{equation}\label{eq:residue-divisor}
\mathfrak{R}_{4,g,n}^{\mathrm{res}}(\cA,\rho)
\;:=\;
\operatorname{div}\det G_{4,g,n}^{\mathrm{res}}(\cA,\rho)\,,
\end{equation}
where $G_{4,g,n}^{\mathrm{res}}$ is the Gram-side quartic
determinant of the residue moment matrix at genus~$g$
with~$n$ marked points, built from the genuine spectral
measure~$d\nu_c(\Delta)$ and the universal residue
factor~$A_c(\rho)$ of
Definition~\textup{\ref{def:universal-residue-factor}}.
\end{definition}

\begin{definition}[Residue clutching defect]%
\label{def:residue-clutching-defect}%
\index{clutching defect!residue|textbf}%
Along a separating clutching map $\xi\colon
\overline{\cM}_{g_1,n_1+1} \times
\overline{\cM}_{g_2,n_2+1} \to \overline{\cM}_{g,n}$,
the \emph{residue clutching defect} is
\begin{equation}\label{eq:residue-clutching-defect}
K_{\cA,\rho,\xi}
\;:=\;
\xi^*\mathfrak{R}_{4,g,n}^{\mathrm{res}}(\cA,\rho)
\;-\;
p_1^*\mathfrak{R}_{4,g_1,n_1+1}^{\mathrm{res}}(\cA,\rho)
\;-\;
p_2^*\mathfrak{R}_{4,g_2,n_2+1}^{\mathrm{res}}(\cA,\rho)
\;-\;
\mathfrak{T}_{3,\xi}^{\mathrm{res}}(\cA,\rho)\,,
\end{equation}
where $\mathfrak{T}_{3,\xi}^{\mathrm{res}}(\cA,\rho)$
is the residue analogue of the cubic tree
divisor~$\mathfrak{T}_{3,\xi}(\cA)$.
\end{definition}

\begin{remark}[Structural meaning]%
\label{rem:clutching-defect-meaning}%
\index{clutching defect!structural meaning}%
The proved clutching law
(Theorem~\ref{thm:nms-clutching-law-modular-resonance})
asserts $K_{\cA,\xi} = 0$ for the MC-derived quartic
resonance class.
The residue clutching defect $K_{\cA,\rho,\xi}$ asks:
does the residue divisor at a candidate zero~$\rho$
satisfy the \emph{same} clutching law as the MC-derived
class?

\medskip\noindent
\emph{On-line zeros}: if $\operatorname{Re}(\rho) = 1/2$
and the genuine spectral input is used, the residue divisor is required
to coincide with the MC-derived quartic resonance class
and therefore satisfy the clutching law: $K_{\cA,\rho,\xi} = 0$.

\medskip\noindent
\emph{Off-line zeros}: if $\operatorname{Re}(\rho) \neq 1/2$,
the residue kernel $w_{c,\rho,u_0}$
(Proposition~\ref{prop:on-off-line-distinction}) has a
different decay exponent $-(c{+}\sigma{-}1)/2$, and the
resulting residue Gram determinant need not satisfy the
clutching law.

\medskip\noindent
The advantage over the scalar compatibility
ratios of~\S\ref{sec:closure-programme}: the clutching defect
uses the full \emph{nonlinear modular structure} (the clutching
law itself, proved from Mok's degeneration formula) rather than
a single scalar ratio. It is therefore a genuinely geometric
zero test are defined next.
\end{remark}

\begin{conjecture}[Clutching closure]%
\label{conj:clutching-closure}%
\ClaimStatusConjectured%
\index{clutching closure conjecture}%
For genuine spectral input\textup{:}
\begin{enumerate}[label=\textup{(\roman*)}]
\item
$\operatorname{Re}(\rho) = 1/2$ implies
$K_{\cA,\rho,\xi} = 0$ for all modular Koszul
families~$\cA$ and all separating clutching maps~$\xi$\textup;
\item
$\operatorname{Re}(\rho) \neq 1/2$ implies
$K_{\cA,\rho,\xi} \neq 0$ for some~$\cA$ and
some~$\xi$.
\end{enumerate}
In other words: the intersection of the residue-clutching
compatibility loci over all families and all boundary strata is
exactly $\{\rho : \operatorname{Re}(\rho) = 1/2\}$.
\end{conjecture}

\begin{remark}[Relation to the quartic closure conjecture]%
\label{rem:clutching-vs-quartic-closure}%
\index{clutching closure!vs quartic closure}%
Conjecture~\ref{conj:clutching-closure} is strictly stronger
than the quartic closure conjecture
(Conjecture~\ref{conj:quartic-closure}):
the scalar compatibility ratio $C_{4,\cA}^{\mathrm{Gram}} = 1$
is the special case of the clutching defect at genus~$0$
with $4$~marked points and the identity clutching map.
The clutching closure uses the full boundary stratification
of~$\overline{\cM}_{g,n}$ and is therefore sensitive to the
higher-genus modular structure.

The quartic closure conjecture
requires only the quartic contact coupling and its
norm, data on a single stratum. The clutching closure
requires the compatibility of residue divisors across
\emph{all} boundary strata, using the tree correction
$\mathfrak{T}_{3,\xi}$ and the contact classes on
\emph{both} components. This geometric control is what makes
the clutching version the stronger statement.
\end{remark}

\subsection{The Beilinson functional}%
\label{sec:beilinson-functional}%
\index{Beilinson functional|textbf}

\begin{definition}[Local compatibility discrepancy]%
\label{def:local-compatibility-discrepancy}%
\index{compatibility discrepancy|textbf}%
For a modular Koszul family~$\cA$, a cyclic direction~$v$,
a central charge~$c$, a candidate zero~$\rho$, and a
regulator~$u_0$, define
\begin{equation}\label{eq:local-discrepancy}
\cB_{\cA,v}(c,\rho;u_0)
\;:=\;
\biggl(
\frac{m_3^\sharp(c,\rho;u_0)}{C_{\cA,v}} - 1
\biggr)^{\!2}
+\;
\biggl(
\frac{S_{4,\cA,v}^{\mathrm{res}}(c,\rho;u_0)}
{Q_{\cA,v}^{\mathrm{ct}}} - 1
\biggr)^{\!2},
\end{equation}
where $C_{\cA,v}$ is the cubic shadow and
$Q_{\cA,v}^{\mathrm{ct}}$ is the quartic contact
coefficient on the slice~$v$.
This vanishes exactly when the cubic and quartic residue
data match the MC shadow data on the slice~$v$.
\end{definition}

\begin{definition}[Beilinson functional]%
\label{def:beilinson-functional}%
\index{Beilinson functional!definition|textbf}%
Sum over all families and slices:
\begin{equation}\label{eq:beilinson-functional}
\cB^{\mathrm{glob}}(\rho)
\;:=\;
\sum_{\cA}\;
\sum_{\text{slices }v}
\int
\cB_{\cA,v}(c,\rho;u_0)\,d\mu_\cA(c,u_0)\,,
\end{equation}
where the sum runs over all modular Koszul families
and the integral is over the admissible
$(c,u_0)$-parameter space of each family.
\end{definition}

\begin{conjecture}[Beilinson closure]%
\label{conj:beilinson-closure}%
\ClaimStatusConjectured%
\index{Beilinson functional!closure conjecture}%
For genuine spectral input,
\begin{equation}\label{eq:beilinson-closure}
\cB^{\mathrm{glob}}(\rho) = 0
\quad\Longleftrightarrow\quad
\operatorname{Re}(\rho) = \tfrac12\,.
\end{equation}
\end{conjecture}

\begin{remark}[Relation to the other closure conjectures]%
\label{rem:beilinson-closure-hierarchy}%
\index{Beilinson functional!closure hierarchy}%
The three closure conjectures form a hierarchy:
\begin{center}
\small
\renewcommand{\arraystretch}{1.3}
\begin{tabular}{clcl}
\toprule
& \textbf{Conjecture} & \textbf{Data used}
& \textbf{Reference} \\
\midrule
1 & Quartic closure
 & Single ratio $C_{4,\cA}^{\mathrm{Gram}} = 1$
 & Conj.~\ref{conj:quartic-closure} \\
2 & Beilinson closure
 & All families, all slices, cubic $+$ quartic
 & Conj.~\ref{conj:beilinson-closure} \\
3 & Clutching closure
 & All families, all boundary strata
 & Conj.~\ref{conj:clutching-closure} \\
\bottomrule
\end{tabular}
\end{center}
Strength: $3 \Rightarrow 2 \Rightarrow 1$.
The Beilinson functional is a scalar aggregation of the
quartic test over all families; the clutching closure
uses the full boundary geometry.
\end{remark}

\subsection{Honest laboratories}%
\label{sec:honest-laboratories}%
\index{honest laboratories|textbf}%
\index{Liouville theory!as honest laboratory}%
\index{Toda theory!as honest laboratory}

\begin{remark}[Why the $\zeta$-proxy fails]%
\label{rem:zeta-proxy-failure}%
\index{zeta proxy!failure|textbf}%
Experiments replacing the genuine spectral
measure~$d\nu_c^{\mathrm{Vir}}(\Delta)$ by a universal
$\zeta$-based surrogate falsified the safe-point condition:
the compatibility ratio $C_{4,\mathrm{Vir}}^{\mathrm{Gram}}$
was not identically~$1$ even at on-line zeros.
The root cause is that the $\zeta$-proxy replaces the
family-specific spectral coefficients by universal
$\zeta$-values, so the normalised Schur complement
$D_2^{\mathrm{proxy}}/H_{\mathrm{Vir}}$ need not equal
$Q_{\mathrm{Vir}}^{\mathrm{ct}} = 10/[c(5c{+}22)]$ even
at~$\sigma = 1/2$.

The quartic residue programme requires genuine
algebraic spectral data, not a universal proxy.
\end{remark}

\begin{problem}[Liouville laboratory]%
\label{prob:liouville-laboratory}%
\index{Liouville theory!residue computation}%
For Liouville theory at $c > 25$, the scalar-primary
spectrum is the continuous family
$\Delta = \tfrac{c-1}{24} + P^2$ for $P \in \bR_{\geq 0}$,
with density determined by the DOZZ structure constants.
This is the natural Virasoro laboratory:
\begin{enumerate}[label=\textup{(\roman*)}]
\item
Build the residue moment matrix
$M_{2,\mathrm{Liouv}}^{\mathrm{res}}(c,\rho;u_0)$
from the continuous spectrum with no $\zeta$-proxy.
\item
Compute $C_{3,\mathrm{Vir}}^{(3)}$ and
$C_{4,\mathrm{Vir}}^{\mathrm{Gram}}$ at the first
zeta zero $\rho_1 = 1/2 + 14.134\ldots\,i$.
\item
Test: $C = 1$ for $\rho_1$ on the line;
$C \neq 1$ for $\rho' = 0.6 + 14.134\ldots\,i$
\textup{(}off-line\textup{)}.
\end{enumerate}
For $c < 1$, use minimal models $\cM(p,q)$ with their
finite discrete spectra
(Problem~\ref{prob:virasoro-residue-evaluation}).
\end{problem}

\begin{problem}[Toda laboratory]%
\label{prob:toda-laboratory}%
\index{Toda theory!residue computation}%
For the $\cW_N$ family, $A_{N-1}$~Toda theory provides the
laboratory. The scalar primary spectrum
is known via the Fateev--Litvinov structure
constants. Build the mixed $(T,W)$ quartic
block (the first multi-channel residue
compatibility test) and verify that the multi-channel
Gram-side Schur complement matches the spin-$3$ quartic
contact coupling~\eqref{eq:w3-spin3-quartic} at on-line
zeros.
\end{problem}

\begin{remark}[The decisive computation]%
\label{rem:decisive-computation}%
\index{quartic residue programme!decisive computation}%
The strongest available test is:
\begin{enumerate}[label=\textup{(\roman*)}]
\item
Build $M_{2,\mathrm{Liouv}}^{\mathrm{res}}(c,\rho;u_0)$
from the continuous Liouville spectrum at $c = 26$
\textup{(}where $\kappa = 13$ and the Koszul dual
$\mathrm{Vir}_{26-c} = \mathrm{Vir}_0$ is degenerate\textup{)}.
\item
Compute the residue clutching defect
$K_{\mathrm{Vir},\rho,\xi}$ at the first zeta zero
for a separating clutching map at genus~$1$.
\item
If $K = 0$ for on-line~$\rho$ and $K \neq 0$ for
off-line~$\rho$, the clutching closure conjecture
survives its first honest test.
\end{enumerate}
The computation is well-posed: the DOZZ structure
constants, the universal residue factor
$A_c(\rho)$~\eqref{eq:universal-residue},
the clutching law
(Theorem~\ref{thm:nms-clutching-law-modular-resonance}),
and the quartic shadow data are all explicit.
\end{remark}

\subsection{Synthesis: two faces of the genus-one arithmetic shadow}%
\label{sec:arithmetic-synthesis}%
\index{arithmetic programme!synthesis|textbf}

The arithmetic interface has two faces. The quartic
resonance class is the first place where they meet.

\begin{center}
\small
\renewcommand{\arraystretch}{1.3}
\begin{tabular}{lll}
\toprule
& \textbf{Zero side} & \textbf{Prime side} \\
\midrule
\emph{Transform}
 & Rankin--Selberg / Eisenstein
 & Mellin / Dirichlet sewing \\
\emph{Variable}
 & $s$ (spectral parameter)
 & $u$ (Dirichlet variable) \\
\emph{Poles}
 & Nontrivial zeros of $\zeta$
 & Primes $p$ (Euler factors) \\
\emph{Coefficients}
 & Residue moments $I_n(c,\rho;u_0)$
 & Sewing coefficients $a_\cA(N)$ \\
\emph{Quadratic datum}
 & Curvature $\kappa(\cA)$
 & $c_1(L_\cA)$ \\
\emph{First nonlinear object}
 & Residue Gram determinant
 & Euler defect $D_\cA(u)$ \\
\emph{Positivity mechanism}
 & Surface Hodge index (conjectural)
 & Petersson inner product (proved) \\
\emph{Locking}
 & \multicolumn{2}{c}{$v = s + u - 1$
 (Sewing--Hecke reciprocity,
 Theorem~\ref{thm:sewing-hecke-reciprocity})} \\
\bottomrule
\end{tabular}
\end{center}

\begin{remark}[The quartic as meeting point]%
\label{rem:quartic-meeting-point}%
\index{quartic resonance class!as meeting point}%
At the scalar level $(\kappa)$, the two sides are locked by
the sewing--Hecke collapse
(Theorem~\ref{thm:sewing-hecke-reciprocity}):
$L_\cA(s,u)$ factors through $v = s{+}u{-}1$.
At the cubic level, the tree shadow $C_\cA$ adds
gravitational coupling but remains a single-exchange object.
At the \emph{quartic} level, the Schur complement first
separates the contact coupling from the lower-degree
contamination
(Theorem~\ref{thm:schur-complement-quartic}),
the residue Gram determinant first sees the quartic
resonance norm, and the clutching law first imposes a
nonlinear boundary compatibility.

The quartic is therefore the meeting point of three
structures:
\begin{enumerate}[label=\textup{(\roman*)}]
\item
the nonlinear shadow obstruction tower
(Appendix~\ref{sec:nms-extended-characteristic-packages});
\item
the zeta-zero residue kernel
(\S\ref{sec:genuine-residue-kernel});
\item
the boundary geometry of~$\overline{\cM}_{g,n}$
(Theorem~\ref{thm:nms-clutching-law-modular-resonance}).
\end{enumerate}
No lower degree has all three.
\end{remark}

\begin{remark}[The end-state of the programme]%
\label{rem:arithmetic-end-state}%
\index{arithmetic programme!end-state}%
The mature programme is:
\begin{enumerate}[label=\textup{(\roman*)}]
\item
$\kappa(\cA)$ controls the scalar modular anomaly;
\item
$\Delta_\cA$ controls the spectral sheet;
\item
$S_\cA(u)$ controls connected arithmetic sewing;
\item
$\Theta_\cA$ unifies them as a connected modular homotopy
object;
\item
modular curves expose the zeta-zero poles;
\item
modular surfaces are conjectured to restore positivity
\textup{(}the missing ingredient of Gap~A\textup{)};
\item
Euler--Koszulity
(Definition~\ref{def:euler-koszul-tier})
is the arithmetic rigidity condition that plain bar--cobar
equivalence lacks \textup{(}Gap~C\textup{)};
\item
the clutching closure
(Conjecture~\ref{conj:clutching-closure})
is the nonlinear geometric replacement for Li-type
positivity.
\end{enumerate}
The bridge is not
$\text{MC} \Rightarrow \text{Li}$ directly
\textup{(}Gap~A obstructs this\textup{)},
nor
$\text{bar--cobar} \Rightarrow \text{RH}$ directly
\textup{(}Gap~C obstructs this\textup{)}.
It is:
\[
\Theta_\cA
\;\xrightarrow{\;\text{residue kernel}\;}
\text{residue Gram matrix}
\;\xrightarrow{\;\text{Schur complement}\;}
\text{quartic resonance}
\;\xrightarrow{\;\text{clutching law}\;}
\text{geometric zero test}\,.
\]
Each arrow is proved or explicitly
defined\textup;
the conjectural content concentrates in the
closure statement alone.
\end{remark}

% ============================================================
\section{The structural diagnosis}%
\label{sec:modular-rigidity}%
\label{sec:modularity-constraint-rho}%
\index{modular spectral rigidity|textbf}

\begin{proposition}[Modularity constraint on the spectral measure]%
\label{prop:modularity-constraint}%
\ClaimStatusConditional%
\index{spectral measure!modularity constraint}%
\index{Hecke algebra!invariance of atoms}%
Let $\rho = \sum_j c_j\,\delta(\lambda_j)$ be the spectral measure
of a chirally Koszul algebra~$\cA$ with HS-sewing. The function
\begin{equation}\label{eq:modularity-constraint-G}
G_\rho(\tau)
\;:=\;
\sum_j c_j \log\!\bigl(1 - \lambda_j\, t(q)\bigr),
\qquad
t(q) = \frac{c}{6}\Bigl(\frac{Z(q)}{Z_{\mathrm{vac}}} - 1\Bigr),
\end{equation}
is modular of weight~$0$ for $\mathrm{SL}(2,\bZ)$ \textup{(}up to an
additive constant from the vacuum normalization\textup{)}.
Consequently:
\begin{enumerate}[label=\textup{(\roman*)}]
\item
the set of atoms $\{\lambda_j\}$ is invariant under the action of the
Hecke algebra\textup;
\item
the weights $\{c_j\}$ satisfy the Petersson orthogonality relations
for Hecke eigenforms.
\end{enumerate}
\end{proposition}

\begin{proof}
By the intertwining theorem
(Theorem~\ref{thm:sewing-shadow-intertwining}),
$G_\rho(\tau) = F_1^{\mathrm{conn}}(q;\cA)$. The connected free energy
$F_1^{\mathrm{conn}} = -\log\det(1 - K_q)$ is a modular-invariant
function: the partition function $Z(\tau,\bar\tau)$ is modular invariant,
so its logarithm is as well. The Hecke invariance of
$\{\lambda_j\}$ follows from the Hecke module structure of
$\mathfrak{g}^{\mathrm{mod}}$
(Theorem~\ref{thm:spectral-continuation-bridge}): the atoms are the
eigenvalues of the Hecke-equivariant operator $m_2 \cdot P$, hence they
are Hecke eigenvalues. The Petersson orthogonality of the weights
$\{c_j\}$ follows from the orthogonality of Hecke eigenspaces with
respect to the Petersson inner product.
\end{proof}

\begin{proposition}[Bracket positivity and the Hodge index]%
\label{prop:bracket-hodge-index}%
\ClaimStatusProvedHere%
\index{Hodge index theorem!bracket positivity}%
\index{Maurer--Cartan equation!bracket bilinear form}%
The MC equation $D \cdot \Theta + \tfrac12[\Theta,\Theta] = 0$
implies that the bracket $[\Theta,\Theta]$ is exact in the cyclic
deformation complex. The bracket defines a bilinear form on the
shadow space:
\begin{equation}\label{eq:bracket-bilinear-form}
B(\operatorname{Sh}_r, \operatorname{Sh}_s)
\;:=\;
[\operatorname{Sh}_r, \operatorname{Sh}_s]_{r+s-2}\,,
\end{equation}
where the subscript denotes the degree-$(r{+}s{-}2)$ component. For
lattice VOAs, this form is positive-definite on the cusp-form subspace
and has signature $(1, d_{\mathrm{arith}} - 1)$ on the Eisenstein
subspace \textup{(}the Hodge index\textup{)}.
\end{proposition}

\begin{proof}
The MC equation gives
$[\Theta,\Theta] = -2\,D\cdot\Theta$, so
$[\Theta,\Theta]$ is exact. Projecting to degree $r{+}s{-}2$, the
bracket $[\operatorname{Sh}_r, \operatorname{Sh}_s]_{r+s-2}$ is
determined by the lower-degree shadow data through the recursion.
For lattice VOAs $V_\Lambda$, the shadow space at degree~$r$ is
isomorphic to the weight-$r/2$ modular forms
$M_{r/2}(\mathrm{SL}(2,\bZ))$ via the theta decomposition
$\Theta_\Lambda = c_E\,E_{r/2} + \sum_j c_j\,f_j$. Under this
identification, the bracket bilinear form corresponds to the Petersson
inner product on cusp forms and the Eisenstein intersection pairing on
the remaining component. The Petersson inner product is
positive-definite on $S_{r/2}$; the Eisenstein component contributes
signature $(1,0)$ from the leading Eisenstein series. The total
signature on $M_{r/2}$ is therefore $(1 + \dim S_{r/2},\, 0)$ on the
cusp-form subspace and $(1, d_{\mathrm{arith}}-1)$ on the Eisenstein
subspace, the latter being the Hodge index.
\end{proof}

\begin{conjecture}[Modular spectral rigidity]%
\label{conj:modular-spectral-rigidity}%
\ClaimStatusConjectured%
\index{modular spectral rigidity|textbf}%
\index{Ramanujan bound!from MC equation}%
\index{non-commutative Weil conjecture}%
Let $\cA$ be a chirally Koszul algebra with HS-sewing, and let
$\rho = \sum_j c_j\,\delta(\lambda_j)$ be the spectral measure
of~$\Theta_\cA$. If $f_j$ is the Hecke eigenform corresponding to
the atom~$\lambda_j$ \textup{(}via the modularity constraint,
Proposition~\textup{\ref{prop:modularity-constraint}}\textup{)}, then:
\begin{enumerate}[label=\textup{(\roman*)}]
\item
\textup{(}Ramanujan bound\textup{)}
For each cusp-form atom:
$|a_{f_j}(p)| \leq 2\,p^{(k-1)/2}$ for all primes~$p$, where $k$~is
the weight of~$f_j$.
\item
\textup{(}MC mechanism\textup{)}
The bound~\textup{(i)} is forced by the MC equation through the chain
\[
\text{MC equation}
\;\Longrightarrow\;
\text{bracket positivity
\textup{(}Prop.~\ref{prop:bracket-hodge-index}\textup{)}}
\;\Longrightarrow\;
\text{Hodge index on } \mathrm{Def}_{\mathrm{cyc}}
\;\Longrightarrow\;
\text{Ramanujan bound}.
\]
\item
\textup{(}Non-commutative Weil conjecture\textup{)}
The MC element~$\Theta_\cA$ plays the role of the Frobenius
endomorphism in the Weil proof: the MC equation $D_\cA^2 = 0$ is the
analogue of $\mathrm{Frob}^2 = q$ \textup{(}the $q$-Frobenius relation
on a curve over~$\bF_q$\textup{)}, and the Ramanujan bound is the
analogue of $|\alpha_i| = q^{1/2}$ \textup{(}the Riemann hypothesis
for curves over finite fields\textup{)}.
\end{enumerate}
\end{conjecture}

\begin{remark}[The Weil analogy in detail]%
\label{rem:weil-analogy-table}%
\index{non-commutative Weil conjecture!dictionary}%
The non-commutative Weil conjecture rests on the following dictionary:
\begin{center}
\small
\renewcommand{\arraystretch}{1.2}
\begin{tabular}{ll}
\toprule
\textbf{Weil proof} & \textbf{MC equation} \\
\midrule
Curve $C/\bF_q$
 & Chirally Koszul algebra $\cA$ \\
Frobenius $\mathrm{Frob}_q$
 & MC element $\Theta_\cA = D_\cA - d_0$ \\
$\mathrm{Frob}^2 = q$
 & $D_\cA^2 = 0$ \\
\'Etale cohomology $H^i_{\text{\'et}}(C,\bQ_\ell)$
 & Cyclic deformation complex $\mathrm{Def}_{\mathrm{cyc}}(\cA)$ \\
Weight filtration $W_\bullet$
 & Degree filtration $\Theta^{\le r}_\cA$ \\
Eigenvalues $\alpha_i$ of Frob on $H^1$
 & Spectral atoms $\lambda_j$ of~$\rho$ \\
$|\alpha_i| = q^{1/2}$ (RH for curves)
 & $|a_{f_j}(p)| \le 2p^{(k-1)/2}$ (Ramanujan) \\
Hodge--Riemann bilinear relations
 & Bracket positivity (Prop.~\ref{prop:bracket-hodge-index}) \\
Lefschetz operator $L$
 & Bracket with $\kappa = c/2$ (degree-$2$ shadow) \\
\bottomrule
\end{tabular}
\end{center}
The analogy is structural, not literal: $D_\cA^2 = 0$ is
a tautological property of any differential, while
$\mathrm{Frob}^2 = q$ encodes arithmetic not visible to $D_\cA$;
the bracket positivity constrains the weights~$c_j$ but not
the eigenvalues~$a_{f_j}(p)$
(Remark~\ref{rem:positivity-limitation}).
For lattice VOAs, the analogy is realized: the Hecke--Newton
closure (Theorem~\ref{thm:hecke-newton-lattice}) provides an
alternative derivation of the Ramanujan bound.
\end{remark}

\begin{proposition}[Ramanujan bound for lattice spectral measures]%
\label{prop:lattice-ramanujan}%
\ClaimStatusProvedHere%
\index{Ramanujan bound!lattice VOAs}%
\index{Deligne's theorem!application to lattice spectral measures}%
For even unimodular lattice VOAs~$V_\Lambda$, the cusp-form atoms
of~$\rho$ satisfy the Ramanujan bound. This is a consequence of
Deligne's theorem on Hecke eigenvalues of cusp forms, applied to the
Galois representations
$\rho_{f_j}$ attached to the cusp eigenforms
$f_j \in S_{r/2}(\mathrm{SL}(2,\bZ))$.
\end{proposition}

\begin{proof}
The lattice theta function decomposes as
$\Theta_\Lambda = c_E\,E_{r/2} + \sum_j c_j\,f_j$ with each $f_j$ a
normalized Hecke eigenform in~$S_{r/2}$. The spectral atoms
$\lambda_j$ are the Hecke eigenvalues $a_{f_j}(p)$ of these cusp forms.
By Deligne's theorem~\cite{Deligne1974},
$|a_{f_j}(p)| \leq 2\,p^{(k-1)/2}$ for all primes~$p$, where
$k = r/2$ is the weight.
\end{proof}

\begin{example}[Leech lattice verification]%
\label{ex:leech-ramanujan}%
\index{Leech lattice!Ramanujan verification}%
For the Leech lattice, the cusp-form atom is~$\Delta_{12}$ with
Hecke eigenvalues~$\tau(p)$. Deligne's bound:
$|\tau(p)| \leq 2\,p^{11/2}$.
\begin{center}
\begin{tabular}{crcrl}
\toprule
$p$ & $\tau(p)$ & bound & $|\tau(p)| \leq$ bound? \\
\midrule
$2$ & $-24$ & $2 \cdot 2^{11/2} \approx 90.5$ & \checkmark \\
$3$ & $252$ & $2 \cdot 3^{11/2} \approx 280.6$ & \checkmark \\
$5$ & $4830$ & $2 \cdot 5^{11/2} \approx 11180$ & \checkmark \\
\bottomrule
\end{tabular}
\end{center}
\end{example}

\begin{proposition}[Shadow--symmetric power identification]%
\label{prop:shadow-symmetric-power}%
\ClaimStatusProvedHere%
Let $f \in S_k(\mathrm{SL}(2,\bZ))$ be a Hecke eigenform with
Satake parameters $\alpha_p, \beta_p$ at each prime~$p$
($a_f(p) = \alpha_p + \beta_p$, $\alpha_p\beta_p = p^{k-1}$).
The shadow coefficient $S_r$ at degree~$r$ satisfies
\begin{equation}\label{eq:shadow-sym-power}
 S_r
 \;\propto\;
 \sum_p p^{-rs}\,
 \operatorname{tr}\!\bigl(\operatorname{Sym}^{r-1}\!
 \begin{pmatrix}\alpha_p & 0\\ 0 & \beta_p\end{pmatrix}\bigr)
 \;=\;
 \sum_p p^{-rs}\,
 \sum_{j=0}^{r-1}\alpha_p^j\,\beta_p^{r-1-j}\,,
\end{equation}
which is the Dirichlet coefficient of the logarithmic derivative
of $L(s, \operatorname{Sym}^{r-1} f)$.
The MC constraint at degree~$r{+}1$ gives a polynomial relation
between $\operatorname{Sym}^{r-1}$ and
$\operatorname{Sym}^0, \ldots, \operatorname{Sym}^{r-2}$\textup{:}
Newton's identity $p_r = \sum_{j=1}^{r}(-1)^{j-1}e_{r-j}\,p_j$
connecting power sums to elementary symmetric polynomials.
\end{proposition}

\begin{proof}
The spectral measure $\rho = \sum_j c_j\,\delta(\lambda_j)$ gives
$S_r = -(1/r)\sum_j c_j\,\lambda_j^r$. For a lattice VOA with
Hecke decomposition, $\lambda_j \propto a_{f_j}(p)$ at each prime~$p$,
and $a_{f_j}(p)^r = (\alpha_p+\beta_p)^r
= \sum_{m=0}^r \binom{r}{m}\alpha_p^m\beta_p^{r-m}$.
The power sum $p_r(\alpha_p,\beta_p) = \alpha_p^r + \beta_p^r
= \operatorname{tr}(\operatorname{Sym}^{r-1}
\operatorname{diag}(\alpha_p,\beta_p))$
relates the shadow to the symmetric power trace.
The Newton identity follows from the recursion
$\det(I - tA) = \exp(-\sum_{r\ge1} \operatorname{tr}(A^r)\,t^r/r)$
applied to $A = \operatorname{diag}(\alpha_p,\beta_p)$.
\end{proof}

\begin{remark}[The Serre reduction]%
\label{rem:serre-reduction}%
\index{Serre's observation}%
\index{Langlands functoriality!and symmetric powers}%
Serre observed: if $L(s,\operatorname{Sym}^r f)$ has analytic
continuation and satisfies the Riemann hypothesis for all $r \ge 1$,
then $|\alpha_p| = |\beta_p| = p^{(k-1)/2}$ (the Ramanujan bound).
The shadow obstruction tower at all degrees encodes all symmetric power
$L$-functions. The MC constraints at all degrees give polynomial
relations between them. The chain:
\[
\text{MC (all degrees)}
\;\xrightarrow{\;\text{Prop.~\ref{prop:shadow-symmetric-power}}\;}
\operatorname{Sym}^r \text{ data}
\;\xrightarrow{\;\text{[analytic cont.]}\;}
L(s,\operatorname{Sym}^r f)
\;\xrightarrow{\;\text{Serre}\;}
\text{Ramanujan.}
\]
The gap is the middle arrow: the MC constraints give polynomial
relations on the symmetric power traces
(Newton's identities), but polynomial relations on
Dirichlet coefficients do not imply analytic continuation of
the Dirichlet series. The passage from algebraic (MC) to
analytic (continuation) is a special case of
\emph{Langlands functoriality}: the automorphy of
$\operatorname{Sym}^r(\rho_f)$ for all~$r$.
\end{remark}

\begin{remark}[Precise diagnosis]%
\label{rem:mc-rigidity-diagnosis}%
\index{modular spectral rigidity!diagnosis}%
The moment $L$-function
$M_r(s) := \int_{\mathrm{SL}(2,\bZ)\backslash\mathfrak{H}}
\operatorname{Sh}_r(\Theta_\cA;\tau)\,E(s,\tau)\,d\mu(\tau)$
is meromorphic by HS-sewing $+$ Rankin--Selberg unfolding.
At each prime~$p$, $S_r$ determines the power sum
$p_r(\alpha_p,\beta_p)$
(Proposition~\ref{prop:shadow-symmetric-power}).
The Serre reduction (Remark~\ref{rem:serre-reduction}):
$\operatorname{Sym}^r$ for all $r$ $\Rightarrow$ Ramanujan.

The gap: $M_r(s)$ is an $r$-point function on a single
elliptic curve (one Dirichlet series, not a product of
$L$-functions), so the Clebsch--Gordan decomposition into
symmetric powers fails; it requires $r$ independent modular
parameters. The prime-by-prime data determines each $p$-local
factor $L_p(s, \operatorname{Sym}^r f)$, but assembling the
global Euler product requires Langlands functoriality
($\mathrm{GL}(2) \to \mathrm{GL}(r{+}1)$), known for
$r \le 4$ (Rankin--Selberg, Kim--Shahidi, Kim) and open
for $r \ge 5$.
\end{remark}

\begin{remark}[Operadic structure and functorial transfer]%
\label{rem:operadic-transfer}%
\index{operadic functorial transfer}%
The MC element $\Theta_\cA \in
\mathrm{Hom}_\alpha(\cC^!_{\mathrm{ch}}, \cP^{\mathrm{ch}})$
carries the symmetric-monoidal structure of the modular operad:
genus-$0$ composition encodes $m_2$ and the Hecke eigenvalue $a_f(p)$;
genus-$1$ composition encodes $\operatorname{tr}(m_2^{r-1})$
and the power sum $p_{r-1}(\alpha_p, \beta_p)$;
genus-$g$ composition gives genus-$g$ analogues.
The operadic composition at degree~$r{+}1$, in spectral
coordinates, yields the Newton relation
(Proposition~\ref{prop:mc-bracket-determines-atoms}).
Whether the all-genera coherence provides the analytic continuation
of $L(s, \operatorname{Sym}^r f)$ for $r \ge 5$ (the gap at
station~$3$ of Remark~\ref{rem:mc-ramanujan-bridge}) is the
target of the operadic Rankin--Selberg theorem below.
\end{remark}

% ============================================================
\section{The geometric positivity programme}%
\index{geometric positivity!from MC bracket}
\label{sec:geometric-positivity}%
\index{geometric positivity programme|textbf}

The MC bracket defines a bilinear form on the shadow space
whose signature constrains the Hecke decomposition weights but
not the Hecke eigenvalues.

\begin{definition}[Shadow bracket form]%
\label{def:shadow-bracket-form}%
\index{shadow bracket form|textbf}%
For a chirally Koszul algebra~$\cA$, the MC bracket restricted
to the shadow space defines the bilinear form
\begin{equation}\label{eq:shadow-bracket-form}
 B_\cA\colon
 \mathrm{Def}_{\mathrm{cyc}}^{(0,r)} \times
 \mathrm{Def}_{\mathrm{cyc}}^{(0,s)}
 \;\longrightarrow\;
 \mathrm{Def}_{\mathrm{cyc}}^{(0,r+s-2)}\,,
 \qquad
 B_\cA(\operatorname{Sh}_r, \operatorname{Sh}_s)
 := [\operatorname{Sh}_r, \operatorname{Sh}_s]_{r+s-2}\,.
\end{equation}
On the one-dimensional primary line of the Virasoro algebra,
$B(S_r, S_s) = r \cdot s \cdot P_0 \cdot S_r \cdot S_s$ where
$P_0 = 2/c$.
\end{definition}

\begin{theorem}[Petersson identification]%
\label{thm:petersson-identification}%
\ClaimStatusProvedHere%
\index{Petersson inner product!shadow bracket identification}%
For a lattice VOA $V_\Lambda$ with Hecke decomposition
$\Theta_\Lambda = c_E\,E_k + \sum_j c_j\,f_j$, the shadow
bracket form~$B_\cA$ restricted to the Hecke eigenspace
of~$f_j$ equals
$c_j^2 \cdot \langle f_j, f_j \rangle_{\mathrm{Pet}} \cdot
\textup{(propagator factors)}$,
where $\langle\cdot,\cdot\rangle_{\mathrm{Pet}}$ is the
Petersson inner product.
\begin{enumerate}[label=\textup{(\roman*)}]
\item $B$ is positive-definite on the cusp-form subspace
 \textup{(}the Petersson inner product is positive-definite
 on~$S_k$\textup{)}.
\item The Eisenstein component contributes
 signature~$(1,0)$.
\item The total signature is
 $(1 + \dim S_k,\, 0)$, all positive.
\end{enumerate}
\end{theorem}

\begin{proof}
The shadow at degree~$r$ for eigenform~$f_j$ has amplitude
$\operatorname{Sh}_r(f_j) = c_j \cdot \lVert f_j
\rVert_{\mathrm{Pet}}^{r-1} \cdot \textup{(combinatorial
factor)}$
(Proposition~\ref{prop:leading-hecke-identification}).
The bracket $[\operatorname{Sh}_r(f_j),
\operatorname{Sh}_s(f_j)] = P \cdot
\operatorname{Sh}_r(f_j) \cdot \operatorname{Sh}_s(f_j)$
by the composition rule in the convolution algebra.
Substituting:
$B(f_j, f_j) = P \cdot c_j^2 \cdot
\lVert f_j \rVert_{\mathrm{Pet}}^{2r-2} \cdot
\textup{(factor)}$.
The Petersson inner product is positive-definite on cusp
forms, giving~(i). The Eisenstein series~$E_k$ is not
square-integrable but contributes a positive-definite form
through the regularized Petersson norm, giving~(ii).
The total signature~(iii) follows from~(i) and~(ii) by
orthogonality of the Hecke eigenspaces.
\end{proof}

\begin{remark}[Positivity does not force Ramanujan]%
\label{rem:positivity-limitation}%
\index{Ramanujan bound!positivity limitation}%
The positivity of~$B$ on the cusp-form subspace constrains the
\emph{weights}~$c_j$ in the Hecke decomposition (they must
have definite sign to maintain positivity), but does
\emph{not} constrain the \emph{eigenvalues}~$a_{f_j}(p)$.
The Ramanujan bound is a statement about eigenvalues, not
weights. The bracket form detects the global structure
(which eigenforms appear, with what weights) but not the
local structure (the Fourier coefficients at each prime).
This is the honest limitation of the geometric positivity
approach: the Hodge index
(Proposition~\ref{prop:bracket-hodge-index}) constrains
the~$c_j$ but is blind to the Satake parameters
$\alpha_p, \beta_p$.

The MC recursion \emph{does} constrain the spectral
atoms~$\lambda_j$ algebraically; the rigidity defect
(Definition~\ref{def:rigidity-defect}) shows that the
atoms are overdetermined at sufficiently high degree.
What remains out of reach is the absolute value bound
$|\alpha_p| = p^{(k-1)/2}$: knowing a polynomial's
roots algebraically does not determine their absolute
values without further input. The required input is
either Deligne's theorem (for lattice VOAs, proved) or
the Serre reduction through symmetric power $L$-functions
(Remark~\ref{rem:serre-reduction}, requiring
Langlands functoriality for $r \geq 5$).
\end{remark}

% ============================================================
\section{The modular rigidity programme}%
\label{sec:modular-rigidity-programme}%
\index{modular rigidity programme|textbf}

The intertwining theorem forces the spectral measure~$\rho$
to satisfy algebraic constraints from modularity; the MC
recursion at successive degrees overdetermines the spectral
data.

\begin{proposition}[Modularity constraint on atoms]%
\label{prop:modularity-constraint-atoms}%
\ClaimStatusConditional%
\index{spectral measure!modularity constraint on atoms}%
The spectral measure $\rho = \sum_j c_j\,\delta(\lambda_j)$
satisfies: $G_\rho(\tau) = \sum_j c_j \log(1 - \lambda_j
\,t(q))$ is modular of weight~$0$
\textup{(}by the intertwining,
Theorem~\textup{\ref{thm:sewing-shadow-intertwining}}\textup{)}.
The atoms $\{\lambda_j\}$ are Hecke eigenvalues; the
weights~$\{c_j\}$ satisfy Petersson orthogonality
\textup{(}Proposition~\textup{\ref{prop:modularity-constraint}}\textup{)}.
\end{proposition}

\begin{proof}
Immediate from
Proposition~\ref{prop:modularity-constraint}:
$G_\rho(\tau) = F_1^{\mathrm{conn}}(q;\cA)$ is modular
invariant, and the Hecke module structure of
$\mathfrak{g}^{\mathrm{mod}}$ forces the atoms to be Hecke
eigenvalues with Petersson-orthogonal weights.
\end{proof}

\begin{definition}[Rigidity defect]%
\label{def:rigidity-defect}%
\index{rigidity defect|textbf}%
For a spectral measure with $m$~atoms, the MC recursion at
degrees $2, \ldots, r$ gives $(r{-}1)$~polynomial equations in
$2m$~unknowns ($m$~atoms $+$ $m$~weights). The
\emph{rigidity defect} is
\begin{equation}\label{eq:rigidity-defect}
 \delta(m, r) \;:=\; (r - 1) - 2m\,.
\end{equation}
For $\delta > 0$ the system is overdetermined and the spectral
measure is rigid.
\end{definition}

\begin{proposition}[Rigidity threshold]%
\label{prop:rigidity-threshold}%
\ClaimStatusProvedHere%
\index{rigidity threshold|textbf}%
\begin{enumerate}[label=\textup{(\roman*)}]
\item For class~$\mathbf{M}$ algebras
 \textup{(}Virasoro, $\cW_N$\textup{):}
 the spectral measure has infinitely many atoms
 \textup{(}$m = \infty$\textup{)}, so
 the rigidity defect $\delta(m, r) = r - 1 - 2m = -\infty$
 for every finite degree~$r$. The counting argument
 never produces an overdetermined system.
\item For the Leech lattice \textup{(}$m = 2$\textup{):}
 $\delta(2, r) = r - 5$, so rigidity begins at $r = 6$.
\item For a lattice of rank~$N$ with
 $d_{\mathrm{arith}}$~eigenforms:
 $\delta(d_{\mathrm{arith}}, R) = R - 1 - 2d_{\mathrm{arith}}$,
 giving rigidity at $R = 2d_{\mathrm{arith}} + 2$.
\end{enumerate}
\end{proposition}

\begin{proof}
The MC recursion at degree~$r{+}1$
(Theorem~\ref{thm:mc-recursion-moment}) gives one polynomial
relation among $\{S_2, \ldots, S_{r+1}\}$. Each~$S_j$
depends on the $m$~atoms and $m$~weights through
$S_j = -(1/j)\sum_i c_i\,\lambda_i^j$. At degree~$2$ the
curvature~$\kappa$ fixes one equation; each subsequent degree
adds one more. For finite~$m$, the count follows. For
class~$\mathbf{M}$ algebras, the spectral measure is
signed and continuous (infinitely many atoms), so the
finite-atom parameter count is meaningless: the system has
infinitely many unknowns at every finite degree.
\end{proof}

\begin{remark}[Finite-atom scope of MC overdetermination]%
\label{rem:route-c-finite-atoms}%
\index{spectral rigidity!finite atom hypothesis}%
The rigidity-threshold counting argument assumes the spectral
measure $\mu_\cA$ is a finite sum of atoms:
$\mu = \sum_{i=1}^m c_i\,\delta_{\lambda_i}$. This holds
for lattice VOAs (where $m = d_{\mathrm{arith}}$ is the number
of Hecke eigenforms in the relevant weight range) and for
rational VOAs (finitely many primary fields). For
class~$\mathbf{M}$ algebras (Virasoro, $\cW_N$), the spectral
measure is not finitely supported: the primary-counting function
on $\mathrm{SL}(2,\bZ)\backslash\mathfrak{H}$ has continuous
spectral support. The finite-atom counting criterion is therefore
restricted to algebras
with finite spectral atom count.
\end{remark}

\begin{conjecture}[Modular spectral rigidity, sharpened form]%
\label{conj:modular-spectral-rigidity-sharp}%
\ClaimStatusConjectured%
\index{modular spectral rigidity!sharpened form}%
The overdetermined system has a \emph{unique} solution. The
spectral atoms are uniquely determined by the shadow obstruction tower data
at sufficiently high degree, and this unique solution satisfies
the Ramanujan bound
$|\alpha_p| = |\beta_p| = p^{(k-1)/2}$.
\end{conjecture}

\begin{remark}[Overdetermination does not imply bounds]%
\label{rem:overdetermination-honest}%
\index{rigidity!does not imply bounds}%
For $\delta(m,r) > 0$, the overdetermined system
$S_j = -(1/j)\sum_i c_i\,\lambda_i^j$,
$j = 2,\ldots,r$, determines the $m$~atoms
$(\lambda_1,\ldots,\lambda_m)$ and $m$~weights
$(c_1,\ldots,c_m)$ uniquely (up to permutation).
Algebraic uniqueness, however, does \emph{not} yield
absolute value bounds: knowing the roots of a polynomial
does not determine their moduli without further input.
Proving $|\alpha_p| = p^{(k-1)/2}$ from the unique
solution requires one of three additional inputs:
\begin{enumerate}[label=\textup{(\alph*)}]
\item
$\ell$-adic Galois representations
(Deligne~\cite{Deligne1974};
proved for lattice VOAs,
Proposition~\ref{prop:lattice-ramanujan});
\item
analytic continuation of
$L(s,\operatorname{Sym}^r f)$ for all $r \geq 1$
(Serre reduction, Remark~\ref{rem:serre-reduction};
known for $r \leq 4$, open for $r \geq 5$);
\item
the operadic functorial transfer
(Conjecture~\ref{conj:operadic-rankin-selberg-main},
\S\ref{sec:operadic-transfer-programme}):
whether the all-genera operadic coherence of~$\Theta_\cA$
provides the analytic continuation that Langlands
functoriality predicts.
\end{enumerate}
The third input is the only one that could make the MC equation
self-sufficient for Ramanujan bounds.
\end{remark}

\begin{proposition}[Lattice Ramanujan from rigidity]%
\label{prop:lattice-ramanujan-rigidity}%
\ClaimStatusProvedHere%
\index{Ramanujan bound!lattice rigidity route}%
For lattice VOAs, the cusp-form atoms satisfy the Ramanujan
bound by Deligne's theorem~\cite{Deligne1974}
\textup{(}Proposition~\textup{\ref{prop:lattice-ramanujan}}\textup{)}.
The overdetermined system is consistent:
the unique solution forced by the rigidity
defect is the Deligne-legal one.
\end{proposition}

\begin{proof}
By Proposition~\ref{prop:lattice-ramanujan}, the atoms are
Hecke eigenvalues of cusp forms, satisfying
$|a_{f_j}(p)| \leq 2\,p^{(k-1)/2}$. The MC recursion
produces the same atoms via the power-sum--symmetric-power
correspondence
(Proposition~\ref{prop:shadow-symmetric-power}). Consistency
of the overdetermined system follows from the theta decomposition
$\Theta_\Lambda = c_E\,E_k + \sum_j c_j\,f_j$ being an
\emph{exact} identity, so no spurious solutions arise.
\end{proof}

\begin{example}[Leech lattice rigidity]%
\label{ex:leech-rigidity}%
\index{Leech lattice!rigidity verification}%
The Leech lattice has $d_{\mathrm{arith}} = 1$ (the single
cusp form~$\Delta_{12}$) and $m = 2$ (one Eisenstein atom $+$
one cusp atom). The rigidity threshold is $R = 6$. The
shadow obstruction tower data at degrees $2, 3, 4, 5, 6$ gives
$5$~equations in $4$~unknowns ($\delta = 1$): the unique
solution reproduces the Ramanujan tau function
$\tau(p) = a_{\Delta_{12}}(p)$ with
$|\tau(p)| \leq 2\,p^{11/2}$
(Example~\ref{ex:leech-ramanujan}).
\end{example}

\begin{remark}[The algebraic--analytic divide]%
\label{rem:algebraic-analytic-divide}%
\index{algebraic--analytic divide|textbf}%
The MC equation organises the algebraic face of the
arithmetic interface completely: shadow-symmetric power data
(Proposition~\ref{prop:shadow-symmetric-power}),
overdetermination
(Proposition~\ref{prop:rigidity-threshold}),
bracket positivity
(Proposition~\ref{prop:bracket-hodge-index}).
Every analytic conclusion requires independent
input external to the MC equation:
\begin{enumerate}[label=\textup{(\roman*)}]
\item
\emph{Lattice VOAs}:
Deligne's theorem~\cite{Deligne1974} supplies the
Ramanujan bound.
The Hecke--Newton closure
(Theorem~\ref{thm:hecke-newton-lattice}) provides an
independent comparison, but the Hecke eigenvalues it determines
are the same Deligne eigenvalues.
\item
\emph{Rational VOAs}:
the Franc--Mason vector-valued Hecke theory
(Proposition~\ref{prop:theta-bridge-rational}) reduces
the spectral problem to Dirichlet $L$-functions
modulo~$2pq$ via the Rocha--Caridi theta decomposition.
\item
\emph{Irrational VOAs}:
Roelcke--Selberg infinite-dimensional Hecke theory
(Conjecture~\ref{conj:irrational-ramanujan}).
Conditional.
\item
\emph{All degrees}:
the Serre reduction
(Remark~\ref{rem:serre-reduction}) requires analytic
continuation of $L(s,\operatorname{Sym}^r f)$ for
all~$r$, known for $r \leq 4$
(Kim--Shahidi) and open for $r \geq 5$.
This is Langlands functoriality
$\mathrm{GL}(2) \to \mathrm{GL}(r{+}1)$.
\end{enumerate}
The MC equation organises the algebraic side so that
when any of these analytic inputs is available, the full
arithmetic package follows systematically.
The operadic functorial transfer programme
(\S\ref{sec:operadic-transfer-programme}) targets
the possibility that the all-genera operadic coherence
of~$\Theta_\cA$ supplies the required
analytic continuation
(Conjecture~\ref{conj:operadic-rankin-selberg-main}).
\end{remark}

% ============================================================
\section{Operadic functorial transfer}%
\label{sec:operadic-transfer-programme}%
\index{operadic functorial transfer programme|textbf}

The Cogdell--Piatetski-Shapiro converse theorem converts
analytic data (meromorphic continuation, functional equation,
polynomial growth for all twists) into an automorphic
representation. The problem is whether the MC equation, together with
HS-sewing, supplies this data for the moment functions.

\begin{conjecture}[CPS hypotheses from MC $+$ HS-sewing]%
\label{conj:cps-from-mc}%
\ClaimStatusConjectured%
\index{Cogdell--Piatetski-Shapiro!hypotheses from MC}%
\textup{[}Conditional on Phragm\'en--Lindel\"of convexity bounds
for $\mathrm{GL}(j)$ twists at $j \geq 3$;
see Remark~\textup{\ref{rem:cps-conditional-status}}.\textup{]}
For a chirally Koszul algebra~$\cA$ with HS-sewing, the moment
$L$-function $M_r(s)$ satisfies the four
Cogdell--Piatetski-Shapiro hypotheses for
$\mathrm{GL}(r)$\textup{:}
\begin{enumerate}[label=\textup{(\alph*)}]
\item \emph{Meromorphic continuation:}
 Rankin--Selberg integral $+$ HS-sewing decay
 \textup{(}Theorem~\textup{\ref{thm:general-hs-sewing}}\textup{)}.
\item \emph{Finitely many poles:}
 at most simple poles from~$E_s$.
\item \emph{Polynomial growth:}
 from the polynomial OPE growth of~$\cA$.
\item \emph{Functional equation:}
 from $E^*(s) = E^*(1{-}s)$ and the MC recursion
 \textup{(}Theorem~\textup{\ref{thm:mc-recursion-moment}}\textup{)}.
\end{enumerate}
These hold for all twists by Dirichlet characters~$\chi$.
\end{conjecture}

\begin{remark}[Analytic proof obligations]
The Rankin--Selberg unfolding
$M_r(s,\chi) = \int_{\mathrm{SL}(2,\bZ)\backslash\mathfrak{H}}
\operatorname{Sh}_r(\Theta_\cA;\tau)\,E^*(s,\chi,\tau)\,
d\mu(\tau)$
gives the correct analytic object. Meromorphic continuation and
finite polar set follow once the shadow amplitude has the
HS-sewing decay required to justify unfolding against the Eisenstein
series. Polynomial growth in vertical strips requires uniform OPE
growth bounds for the family of Fourier coefficients and the
standard Phragm\'en--Lindel\"of estimate. The functional equation
$M_r(s,\chi) = \gamma(s,\chi)\,M_r(1{-}s,\bar\chi)$
then follows from $E^*(s)=E^*(1-s)$ only after the MC recursion
\textup{(}Theorem~\textup{\ref{thm:mc-recursion-moment}}\textup{)}
has been shown to commute with the completed unfolding. Cuspidal
twists require the corresponding triple-product or higher
Rankin--Selberg integral. Thus the conjecture is reduced to
unfolding, vertical-strip growth, and twist-compatibility estimates,
not to a formal manipulation of the Eisenstein functional equation
alone.
\end{remark}

\begin{remark}[Analytic estimates]%
\label{rem:cps-conditional-status}%
\index{Cogdell--Piatetski-Shapiro!conditional scope}%
The polynomial growth hypothesis~(c) requires
bounding $M_r(s, \chi)$ in vertical strips for \emph{all}
twists $\chi$ of $\mathrm{GL}(j)$, $j = 1, \ldots, r-2$.
For $j = 1$ (Dirichlet characters), the bound follows from
standard Rankin--Selberg estimates. For $j \geq 3$, the
requisite Phragm\'en--Lindel\"of convexity bounds for
the triple-product and higher Rankin--Selberg integrals have
not been proved in this volume. Consequently
Corollary~\ref{cor:moment-automorphy} is conditional on these
estimates; the formal algebraic pieces (meromorphic continuation,
functional equation, finite polar set before vertical-strip control)
remain separate inputs.
\end{remark}

\begin{corollary}[Conditional automorphy of moment $L$-functions]%
\label{cor:moment-automorphy}%
\ClaimStatusConditional
\index{automorphy!moment L-functions}%
Assume Conjecture~\ref{conj:cps-from-mc}, including the
vertical-strip and twist-compatibility estimates of
Remark~\ref{rem:cps-conditional-status}. Then, by the
Cogdell--Piatetski-Shapiro converse theorem~\cite{CPS1999},
$M_r(s) = L(s, \pi_r)$ for an automorphic representation
$\pi_r$ of~$\mathrm{GL}(r)$.
\end{corollary}

\begin{proof}
Under Conjecture~\ref{conj:cps-from-mc}, the four CPS hypotheses
hold for all twists. The converse theorem~\cite{CPS1999} then
produces the automorphic
representation~$\pi_r$.
\end{proof}

\begin{conjecture}[Prime-locality]%
\label{conj:prime-locality-transfer}%
\ClaimStatusConjectured%
\index{prime-locality!transfer programme}%
The shadow coefficients decompose prime-by-prime:
$S_r^{(p)} \propto p_{r-1}(\alpha_p, \beta_p)$.
Proved for lattice VOAs
\textup{(}Proposition~\textup{\ref{prop:shadow-symmetric-power}}\textup{)};
trivially satisfied for Virasoro at leading~$1/c$.
\end{conjecture}

\begin{proposition}[Universal signed Stieltjes measure]%
\label{prop:stieltjes-signed-universal}%
\ClaimStatusProvedHere%
\index{Stieltjes representation!signed measure}%
\index{prime-locality!signed measure obstruction}%
For every Virasoro theory with $c > -22/5$, the Stieltjes
discriminant satisfies
\begin{equation}\label{eq:stieltjes-disc-universal}
 36 - \alpha(c) \;=\; \frac{-80}{5c + 22} \;<\; 0,
 \qquad
 \alpha(c) = \frac{180c + 872}{5c + 22}.
\end{equation}
Hence the quadratic form $Q(t) = c^2 + 12ct + \alpha\,t^2$
in the generating function $H(t) = t^2\sqrt{Q(t)}$ has
\emph{complex} roots, and the spectral measure~$\rho$ in the
Stieltjes representation
\textup{(}Theorem~\textup{\ref{thm:shadow-spectral-measure}}\textup{)}
is necessarily signed: $c_2 = \int\lambda^2\,d\rho = -c < 0$
for $c > 0$. This holds universally for all class-$M$ theories.
\end{proposition}

\begin{proof}
Direct computation:
$36 - \alpha = 36 - (180c{+}872)/(5c{+}22) = (36(5c{+}22) - 180c - 872)/(5c{+}22)
= (180c + 792 - 180c - 872)/(5c{+}22) = {-80}/(5c{+}22) < 0$.
The moment $c_2 = -2S_2 = -c$
(Theorem~\ref{thm:shadow-spectral-measure}), so
$\int\lambda^2\,d\rho < 0$ and $\rho$ cannot be a positive
measure. For $c = 1/2$, $7/10$, and $4/5$, the finite Hankel
matrices of sizes $2,\ldots,6$ have full rank; this finite
calculation is a low-rank obstruction to a short atomic model, not a
proof of infinite support.
\end{proof}

\begin{proposition}[Multiplicativity failure for rational CFT]%
\label{prop:rational-cft-multiplicativity-failure}%
\ClaimStatusProvedHere%
\index{prime-locality!multiplicativity failure}%
\index{Ising model!multiplicativity failure}%
Let $V$ be a unitary minimal-model CFT with diagonal modular
invariant $Z(\tau) = \sum_{i} |\chi_i(\tau)|^2$.
The Dirichlet coefficients~$a_n$ of
$|\eta(\tau)|^2 Z(\tau) = \sum_{n \ge 0} a_n\,q^n$
are \emph{not} multiplicative:
there exist coprime $m, n$ with $a_{mn} \neq a_m\,a_n$.
In particular, the Dirichlet series
$D(s) = \sum_{n \ge 1} a_n\,n^{-s}$
admits no Euler product.
\end{proposition}

\begin{proof}
Explicit computation for $c = 1/2$ \textup{(Ising, $M(4,3)$)}.
The first coefficients are
$(a_0, a_1, \ldots, a_7) = (3, -2, -3, -2, 0, 2, 1, 4)$.
At $m = 2$, $n = 3$ (coprime):
$a_6 = 1 \neq a_2\,a_3 = (-3)(-2) = 6$,
defect $= 5$. At $m = 2$, $n = 7$:
$a_{14} = -1 \neq a_2\,a_7 = -12$, defect $= 11$.
Altogether $14$ coprime failures with $m, n \le 20$.

The same failure occurs for $c = 7/10$ (tricritical Ising,
$12$ failures) and $c = 4/5$ (three-state Potts, $14$
failures); defects are $O(1)$ integers in all cases.
\end{proof}

\begin{remark}[Structural obstruction to naive prime-locality]%
\label{rem:naive-prime-locality-obstruction}%
\index{prime-locality!structural obstruction}%
The multiplicativity failure is structural, not numerical.
For a lattice VOA, $Z(\tau) = \Theta_\Lambda(\tau)/\eta^r$
is a ratio of \emph{scalar} modular forms; $\Theta_\Lambda$
decomposes into Hecke eigenforms, forcing multiplicativity of
the Dirichlet series
(Proposition~\ref{prop:shadow-symmetric-power}).
For a rational CFT, $Z = \sum_i |\chi_i|^2$ is a quadratic
form on VVMF components~$\chi_i$, and the individual~$\chi_i$
need not be Hecke eigenforms. In the checked Ising and
tricritical-Ising cases, none of the $3$ and $6$ characters,
respectively, is an eigenform of $T_p$ at $p = 2,3,5$.
The cross terms between distinct conformal families destroy
multiplicativity at $O(1)$ level.

\textbf{Consequence for the prime-locality programme.}
Conjecture~\ref{conj:prime-locality-transfer} cannot be tested via
the partition function Euler product for non-lattice theories.
If prime-locality holds for the shadow coefficients~$S_r$,
it must be formulated at the shadow-projection level:
$\operatorname{Sh}_r(\Theta_\cA; \tau)$ may have an Euler product
structure that~$Z(\tau)$ lacks, because~$\operatorname{Sh}_r$
is a genus-$1$ degree-$r$ projection of the MC element, not the
full partition function.
The moment $L$-function
$M_r(s) = \int \operatorname{Sh}_r\,E(s)\,d\mu$ is the correct
object; its Euler product question remains the central open problem
(Remark~\ref{rem:algebraic-analytic-divide}).
\end{remark}

\begin{proposition}[Shadow arithmetic trichotomy]%
\label{prop:shadow-arithmetic-trichotomy}%
\ClaimStatusProvedHere%
\index{prime-locality!shadow arithmetic trichotomy}%
\index{shadow arithmetic trichotomy}%
For a diagonal rational CFT with characters~$\chi_i$ forming a
vector-valued modular form on the congruence subgroup
$\Gamma(4pq)$ of $M(p,q)$, the arithmetic decomposes into three
levels\textup{:}

\textup{\textbf{Level 0}} \textup{(}partition function\textup{)}.
The Dirichlet coefficients~$a_n$ of $|\eta|^2 Z$ are not
multiplicative
\textup{(}Proposition~\textup{\ref{prop:rational-cft-multiplicativity-failure}}\textup{)}.
The defect decomposes as
\[
 \underbrace{a_{mn} - a_m a_n}_{\text{full defect}}
 \;=\;
 \underbrace{\textstyle\sum_i \bigl(b_{mn}^{(i)} - b_m^{(i)}b_n^{(i)}\bigr)}_{\text{per-channel defect}}
 \;-\;
 \underbrace{\textstyle\sum_{i \neq j} b_m^{(i)}b_n^{(j)}}_{\text{cross-channel terms}}
\]
where $b_n^{(i)} = |\eta\,\chi_i|^2$ are per-channel
coefficients. Both sources contribute.

\textup{\textbf{Level 1}} \textup{(}per-channel\textup{)}.
Each $\eta\,\chi_{r,s}$ is a theta function
\textup{(}Rocha--Caridi\textup{)}, so
$|\eta\,\chi_{r,s}|^2 = \theta_{r,s}^2$
lives on~$\Gamma(4pq)$.
However, per-channel multiplicativity fails \emph{even away from
the level primes}:
for Ising $M(4,3)$ with level $48$ and level primes $\{2,3\}$,
the vacuum channel $(1,1)$ satisfies
$b_{35} \neq b_5 b_7$ (both $5$ and~$7$ are coprime to the level).
The obstruction is the alternating-sign structure
$\theta_{r,s} = \sum_n(q^{\alpha_+(n)} - q^{\alpha_-(n)})$
inherited from the Rocha--Caridi formula:
the signed theta square is not a Hecke eigenform on any subgroup.

\textup{\textbf{Level 2}} \textup{(}shadow projection\textup{)}.
The genus-$1$ shadow projection
$\operatorname{Sh}_r^{(1)}(\Theta_\cA;\tau)$
is a \emph{graph sum} over $\overline{\cM}_{1,r}$:
vertex factors are the algebra-specific
constants~$S_j(\cA)$, while edge factors are
\emph{universal} propagators related to
Eisenstein series~$E_2^*(\tau)$.
At degree~$2$, $\operatorname{Sh}_2^{(1)} \propto
\kappa \cdot E_2^*(\tau)$, whose Fourier coefficients
$\sigma_1(n)$ are multiplicative.
The shadow-level arithmetic is governed by the
\emph{universal modular geometry} of $\overline{\cM}_{1,r}$,
not by the per-channel theta functions.
\end{proposition}

\begin{proof}
Levels~$0$ and~$1$ follow from explicit coefficient computations
(Proposition~\ref{prop:rational-cft-multiplicativity-failure},
with per-channel failures in the Ising, tricritical-Ising, and
three-state Potts examples).
Level~$2$ at degree~$2$: the single stable graph at $(g{=}1,n{=}0)$
with one self-loop has vertex factor~$S_2 = \kappa$ and
propagator $P(\tau) \propto E_2^*(\tau)$, giving
$\operatorname{Sh}_2^{(1)} = \kappa \cdot E_2^*/(2\cdot 1)$.
Since $E_2^*(\tau) = 1 - 24\sum_{n\ge1}\sigma_1(n)\,q^n$ and
$\sigma_1$ is multiplicative, the Fourier coefficients of
$\operatorname{Sh}_2^{(1)}$ are multiplicative.
Level~$2$ at degrees~$r \ge 3$:
the weight bound (Proposition~\ref{prop:genus1-weight-bound})
constrains the quasi-modular space; the Euler product question
remains open (Remark~\ref{rem:quasimodular-obstruction}).
\end{proof}

\begin{proposition}[Weight bound for genus-$1$ shadow projections]%
\label{prop:genus1-weight-bound}%
\ClaimStatusProvedHere%
\index{shadow-level multiplicativity!weight bound}%
At genus~$1$ and degree~$r$, every stable graph~$\Gamma$ at
$(g{=}1, n{=}r)$ has modular weight $2|E(\Gamma)| \le 2r$.
The bound is tight \textup{(}achieved by the necklace
graph\textup{)}. The graph count at degrees~$1$--$6$ is
$2, 5, 12, 28, 64, 144$.
At degree~$6$ the weight first reaches~$12$, the threshold
where the Ramanujan cusp form~$\Delta$ appears.
\end{proposition}

\begin{proof}
For $h^1(\Gamma) = 1$ with all vertices at genus~$0$:
$|E| = |V|$, stability $n_v \ge 3$ gives
$\sum n_v = 2|V| + r \ge 3|V|$, hence $|V| \le r$.
Weight $= 2|V| \le 2r$.
For $h^1 = 0$ (tree, one genus-$1$ vertex):
$|E| = |V| - 1 < |V| \le r$.
\end{proof}

\begin{remark}[Quasi-modular obstruction to naive multiplicativity]%
\label{rem:quasimodular-obstruction}%
\index{shadow-level multiplicativity!quasi-modular obstruction}%
\index{quasi-modular forms!and shadow projections}%
The genus-$1$ propagator is the quasi-modular Eisenstein
series~$E_2^*(\tau)$, whose non-holomorphic completion is
$\widehat{E}_2(\tau) = E_2^*(\tau) - 3/(\pi\,\mathrm{Im}\,\tau)$.
The holomorphic series~$E_2^*(\tau)$ is \emph{not} a modular form for
$\mathrm{SL}(2,\bZ)$.
Although $\dim S_k(\mathrm{SL}(2,\bZ)) = 0$ for $k < 12$,
this applies only to \emph{holomorphic} forms; the space of
\emph{quasi-modular} forms of weight~$k$ is larger
(it includes $E_2^*$-shifted derivatives of holomorphic forms).
The shadow projection $\operatorname{Sh}_r^{(1)}$ involves
products of~$E_2^*$, which are quasi-modular polynomials in
$\{E_2^*, E_4, E_6\}$, not polynomials in $\{E_4, E_6\}$ alone.

\textbf{Consequence for the Euler product.}
The moment $L$-function
$M_r(s) = \int \operatorname{Sh}_r^{(1)} \cdot E(s,\tau)\,d\mu$
involves Rankin--Selberg integrals of \emph{quasi-modular}
forms against~$E(s,\tau)$. The unfolding theory for such
integrals is more subtle than for holomorphic eigenforms
(it involves~$\widehat{E}_2$ and
Zagier's theory of quasi-modular Rankin--Selberg integrals).
The Euler product question for~$M_r(s)$ at degrees~$r \le 5$
remains \emph{open}; it requires an explicit quasi-modular
Rankin--Selberg analysis that has not been carried out.

The weight bound (Proposition~\ref{prop:genus1-weight-bound})
and the absence of holomorphic cusp forms below weight~$12$
remain valid structural inputs, but they do \emph{not} by
themselves establish multiplicativity of the shadow Fourier
coefficients.

\textbf{Critical degree.}
At degree~$6$, the necklace graph first reaches holomorphic
weight~$12$, introducing a $\Delta$-component even in the
holomorphic part.
At degrees~$r \ge 6$, the MC recursion constraints
(Newton's identities,
Proposition~\ref{prop:mc-bracket-determines-atoms}) are
essential for determining the $\Delta$ coefficient.
\end{remark}

\begin{conjecture}[Operadic Rankin--Selberg, main form]%
\label{conj:operadic-rankin-selberg-main}%
\ClaimStatusConjectured%
\index{operadic Rankin--Selberg theorem!main form}%
Assume Conjecture~\ref{conj:cps-from-mc}, prime-locality
\textup{(}Conjecture~\textup{\ref{conj:prime-locality-transfer}}\textup{)},
and equality of the finitely many ramified Euler factors with the
local factors of a Hecke eigenform~$f$. Then
$M_r(s)=L(s,\operatorname{Sym}^{r-1}f)$ and Serre's reduction gives
the Ramanujan bound
$|\alpha_p|=|\beta_p|=p^{(k-1)/2}$ at unramified primes.
\end{conjecture}

\begin{remark}[Symmetric-power identification criterion]
The CPS converse theorem~\cite{CPS1999} together with
Conjecture~\ref{conj:cps-from-mc} gives
$M_r = L(s, \pi_r)$ for automorphic~$\pi_r$ on
$\mathrm{GL}(r)$
(Corollary~\ref{cor:moment-automorphy}).
Prime-locality together with Newton's identities
(Proposition~\ref{prop:mc-bracket-determines-atoms})
identifies the local Euler factor of~$M_r$ with the
$(r-1)$-st symmetric-power Euler factor of~$f$ at every unramified
prime. Strong multiplicity one
(Jacquet--Shalika~\cite{JacquetShalika1981}) then gives
$\pi_r\cong\operatorname{Sym}^{r-1}(f)$ after the finitely many
ramified factors have been fixed. Serre's reduction
(Remark~\ref{rem:serre-reduction}) converts this symmetric-power
automorphy statement into the Ramanujan bound. The unproved input is
therefore prime-locality of the MC moment data, not the
automorphic-identification step after locality is known.
\end{remark}

\begin{remark}[The conditional implication chain]%
\label{rem:complete-chain}%
\index{Ramanujan bound!complete chain}%
The conditional chain:
\[
\text{MC equation}
\;\xrightarrow{\;\text{HS-sewing}\;}
\text{RS method}
\;\xrightarrow{\;\text{unfolding}\;}
M_r(s) \text{ meromorphic}
\;\xrightarrow{\;\text{CPS}\;}
\pi_r \in \mathrm{Aut}(\mathrm{GL}(r))
\]
\[
\xrightarrow{\;[\text{prime-locality}]\;}
M_r = \operatorname{Sym}^{r-1}
\;\xrightarrow{\;\text{Serre}\;}
\text{Ramanujan.}
\]
The first two arrows are algebraic and analytic inputs; the CPS arrow
is conditional on Conjecture~\ref{conj:cps-from-mc}; the
prime-locality arrow is Conjecture~\ref{conj:prime-locality-transfer}
outside the lattice lane. For lattice VOAs, the Hecke--Newton closure
supplies the prime-locality arrow. The resulting chain is a
Langlands-conditional comparison to the Ramanujan bound, not a
replacement for Deligne's theorem unless the converse-theorem
hypotheses and local-factor identifications are supplied independently.
\end{remark}

The recursion and operadic analysis follow.

% ============================================================
\subsection{The MC recursion formula}%
\label{sec:operadic-rankin-selberg}%
\index{operadic Rankin--Selberg theorem}

The moment $L$-function $M_r(s) = \int \pi_{1,r}(\Theta_\cA)
\cdot E^*(s)\,d\mu$ is the Rankin--Selberg pairing of the
genus-$1$, degree-$r$ projection of~$\Theta_\cA$ with the
Eisenstein series. The MC recursion determines $M_{r+1}$ from
lower degrees through the convolution bracket of~$\gAmod$.

\begin{theorem}[MC recursion on moment $L$-functions]%
\label{thm:mc-recursion-moment}%
\ClaimStatusProvedHere%
Assume the genus-$1$ shadow amplitudes and their Hochschild brackets are
of Rankin--Selberg moderate growth, and work away from the finite polar
divisor of the Mellin-transformed Hochschild operator. The MC equation at
degree $r{+}1$ determines $M_{r+1}(s)$
from $\{M_j(s)\}_{j=2}^r$ via a Rankin--Selberg recursion:
$M_{r+1}(s) = -\widetilde{\nabla}_H(s)^{-1} \cdot
\sum_{j+k=r+3} \mathcal{R}(M_j, M_k;\,s)$,
where $\widetilde{\nabla}_H(s)$ is the Mellin-transformed
Hochschild operator and $\mathcal{R}(M_j, M_k;\,s)$ is the
Rankin--Selberg bracket integral.
Each~$\mathcal{R}$ is meromorphic; functional equations are preserved.

Here $\widetilde{\nabla}_H(s) := s(s{-}1) - \lambda_\kappa$
is the Casimir eigenvalue shift of the Hochschild operator
\textup{(}$\lambda_\kappa = \kappa(\kappa{-}1)/2$ for the
curvature eigenvalue\textup{)}, and
$\mathcal{R}(M_j, M_k;\,s) := \int
\{\operatorname{Sh}_j, \operatorname{Sh}_k\}_H(\tau)
\cdot E^*(s,\tau)\,d\mu(\tau)$
is the Rankin--Selberg integral of the shadow bracket.
\end{theorem}

\begin{proof}
The MC equation $\nabla_H(\operatorname{Sh}_{r+1}) +
\sum\{\operatorname{Sh}_j, \operatorname{Sh}_k\}_H = 0$
yields the recursion after Mellin-transforming against~$E^*(s)$.
The bracket is the contraction of shadow amplitudes along the
propagator on~$E_\tau$, integrated over marked points. The stated
moderate-growth hypothesis places this bracket in the classical
Rankin--Selberg admissible class, so the integral defining
$\mathcal{R}(M_j,M_k;s)$ is meromorphic. On the complement of the polar
divisor of~$\widetilde\nabla_H(s)$ the operator is invertible, giving the
displayed formula. The functional equation follows from
$E^*(s)=E^*(1{-}s)$ and the same equation for the lower moment functions.
\end{proof}

This theorem supplies the algebraic meromorphic recursion. It does not,
by itself, prove the Cogdell--Piatetski-Shapiro vertical-strip estimates
for all twists, nor prime-locality outside the lattice lane. Those are
the separate converse-theorem and prime-locality inputs recorded in
Conjectures~\ref{conj:cps-from-mc},
\ref{conj:prime-locality-transfer}, and
\ref{conj:operadic-rankin-selberg-main}.



% ============================================================
\section{The Hecke--Newton closure}%
\label{sec:hecke-newton-closure}%
\index{Hecke--Newton closure|textbf}

Prime-locality asserts that the Hecke operators $T_p$
on~$\cM_{1,1}$ commute with the genus-$1$ projection
$\pi_{1,r}(\Theta_\cA)$ of the MC element. For lattice VOAs,
this follows from the Hecke module structure
of~$\gAmod[V_\Lambda]$
(Theorem~\ref{thm:spectral-continuation-bridge}).
This section resolves
prime-locality for lattice VOAs via the
\emph{Hecke--Newton closure}, and identifies the
obstruction for non-lattice theories.

\subsection{Chiral graph integrals on the elliptic curve}

Each stable graph~$\Gamma$ contributes an
amplitude~$\ell_\Gamma$ to $\pi_{1,r}(\Theta_\cA)$;
this amplitude is a chiral graph integral on~$E_\tau$.

\begin{proposition}[Shadow amplitudes as chiral graph integrals]%
\label{prop:shadow-chiral-graph}%
\ClaimStatusConditional%
\index{chiral graph integral|textbf}%
\index{shadow amplitude!as graph integral}%
The genus-$1$ shadow amplitude
$\operatorname{Sh}_r^{(1)}(\cA;\tau)$ is a finite sum of
\emph{chiral graph integrals}: for each connected genus-$1$
stable graph~$\Gamma$ of degree~$r$,
\begin{equation}\label{eq:chiral-graph-integral}
 \ell_\Gamma(\tau)
 \;=\;
 \int_{E_\tau^{|V(\Gamma)|-1}}
 \prod_{v\in V(\Gamma)} m_{\operatorname{val}(v)}
 \;\cdot\;
 \prod_{e\in E(\Gamma)} g\!\bigl(z_{s(e)}, z_{t(e)};\tau\bigr)
 \;\prod_{v\ne v_0} dz_v\,,
\end{equation}
where $g(z,w;\tau) = -\log\bigl[\theta_1(z{-}w\,|\,\tau)/
\theta_1'(0\,|\,\tau)\bigr]$ is the chiral Green's function
on~$E_\tau$, the $m_{\operatorname{val}(v)}$ are
$\tau$-independent OPE coefficients, and the integration fixes
one marked point by translation invariance. After integration,
$\ell_\Gamma(\tau)$ is a real-analytic function
on~$\mathrm{SL}(2,\bZ)\backslash\mathfrak{H}$ with at most polynomial
growth at the cusp.
\end{proposition}

\begin{proof}
The graph-sum formula
(Theorem~\ref{thm:sewing-shadow-intertwining}, proof)
expresses $F_1^{\mathrm{conn}} = \sum_\Gamma
|\operatorname{Aut}(\Gamma)|^{-1}\,\ell_\Gamma$ with each
$\ell_\Gamma$ factored as (OPE data)$\times$(geometric integral
on~$E_\tau$). The geometric integral is precisely the
chiral graph integral~\eqref{eq:chiral-graph-integral}.
The chiral Green's function~$g$ transforms under
$\mathrm{SL}(2,\bZ)$ with an additive anomaly proportional
to~$E_2^*(\tau)/(2\pi i)$; however, this anomaly cancels in the
integrated amplitude $\ell_\Gamma$ for balanced graphs (equal
number of $g$ and $\bar{g}$ contributions when paired with
the complex conjugate), leaving $\ell_\Gamma$ real-analytic
and automorphic. The HS-sewing condition
(Theorem~\ref{thm:general-hs-sewing}) gives polynomial growth.
\end{proof}

\begin{remark}[Relation to modular graph forms]%
\label{rem:mgf-relation}%
\index{modular graph forms!relation to shadow amplitudes}%
The chiral graph integrals~\eqref{eq:chiral-graph-integral}
are structurally analogous to the \emph{modular graph forms}
of D'Hoker--Green--G\"urdo\u{g}an--Vanhove and
Zagier--Zerbini, which use the \emph{full} scalar Green's function
$G(z,w;\tau) = -\log\lvert\theta_1(z{-}w\,|\,\tau)/
\eta(\tau)\rvert^2 + 2\pi\,(\mathrm{Im}(z{-}w))^2/\mathrm{Im}(\tau)$
instead of the chiral~$g$. The chiral and full Green's functions
differ by a term $\log\lvert\eta(\tau)/\theta_1'(0)\rvert^2
+ 2\pi(\mathrm{Im}(z{-}w))^2/y$, which is a function
of~$\tau$ and $\mathrm{Im}(z{-}w)$ alone. For the shadow
amplitudes, the OPE vertex weights $m_{\operatorname{val}(v)}$
are $\tau$-independent, so the $\tau$-dependent difference
factors through the graph combinatorics. The shadow amplitude
$\operatorname{Sh}_r^{(1)}$ is therefore a linear combination of
modular graph forms shifted by $E_2^*$-dependent correction
terms from the chiral-to-full conversion.
\end{remark}

\subsection{The Hecke--Newton closure for lattice VOAs}%
\label{sec:hecke-newton-lattice}

For lattice VOAs, the chiral graph integrals decompose through
the Hecke algebra; this is the prime-locality mechanism.

\begin{theorem}[Hecke--Newton closure for lattice VOAs]%
\label{thm:hecke-newton-lattice}%
\ClaimStatusProvedHere%
\index{Hecke--Newton closure!lattice VOAs}%
\index{prime-locality!proved for lattice VOAs}%
Let $V_\Lambda$ be an even unimodular lattice VOA of rank~$r$,
with Hecke decomposition
$\Theta_\Lambda = c_E\,E_{r/2} + \sum_{j=1}^m c_j\,f_j$
in $M_{r/2}(\mathrm{SL}(2,\bZ))$. Then:
\begin{enumerate}[label=\textup{(\roman*)}]
\item
The moment $L$-function decomposes into Euler products:
\begin{equation}\label{eq:moment-euler-lattice}
 M_r(s; V_\Lambda)
 \;=\;
 c_E\,L_E(s,r) \;+\; \sum_{j=1}^m c_j\,L_j(s,r),
\end{equation}
where $L_E(s,r)$ involves $\zeta(s)\,\zeta(s{-}k{+}1)$ and
each $L_j(s,r) = \sum_p p_{r-1}(\alpha_{j,p},\beta_{j,p})\,
p^{-s} + \cdots$ is the $r$-th moment of the Satake parameters
of~$f_j$ at each prime~$p$.
\item
At each prime~$p$ and for each Hecke eigenform~$f_j$, the MC
recursion gives Newton's identity:
\begin{equation}\label{eq:newton-at-prime}
 p_r(\alpha_{j,p}, \beta_{j,p})
 \;=\;
 a_{f_j}(p)\cdot p_{r-1}(\alpha_{j,p}, \beta_{j,p})
 - p^{k-1}\cdot p_{r-2}(\alpha_{j,p}, \beta_{j,p}),
\end{equation}
where $a_{f_j}(p) = \alpha_{j,p} + \beta_{j,p}$ is the Hecke
eigenvalue and $\alpha_{j,p}\,\beta_{j,p} = p^{k-1}$.
\item
Prime-locality holds: $S_r^{(p)} \propto
p_{r-1}(\alpha_{j,p}, \beta_{j,p})$ at each prime~$p$.
\end{enumerate}
\end{theorem}

\begin{proof}
\emph{(i)} The Hecke module structure
(Theorem~\ref{thm:spectral-continuation-bridge}) decomposes
the graph amplitudes $\ell_\Gamma^{(1)}(V_\Lambda)$ into
Hecke eigenspaces. The Mellin transform of each eigencomponent
against~$E^*(s)$ gives the corresponding $L$-function
(by standard Rankin--Selberg unfolding), producing the
decomposition~\eqref{eq:moment-euler-lattice}.

\emph{(ii)} The MC equation at degree $r{+}1$ gives
Newton's identity $p_r = e_1\,p_{r-1} - e_2\,p_{r-2}$
(Proposition~\ref{prop:mc-bracket-determines-atoms}). For
each Hecke eigenform~$f_j$, the eigencomponent of~$\mu_r$
at prime~$p$ is $p_r(\alpha_{j,p}, \beta_{j,p})$. Since the
MC recursion is a polynomial identity and the Hecke
eigenspaces are invariant under polynomial operations, the
recursion holds eigenspace-by-eigenspace, giving
Newton's identity~\eqref{eq:newton-at-prime} at each prime
for each eigenform.

\emph{(iii)} Immediate from~(ii): the power sums
$p_r(\alpha_{j,p}, \beta_{j,p})$ are the $p$-local shadows,
and they factor prime-by-prime by construction.
\end{proof}

\begin{corollary}[Conditional operadic RS for lattice VOAs]%
\label{cor:unconditional-lattice}%
\label{cor:ramanujan-deligne-free}%
\ClaimStatusConditional%
\index{operadic Rankin--Selberg theorem!conditional for lattices}%
\index{Ramanujan bound!conditional comparison}%
\index{Ramanujan bound!lattice VOA comparison}%
For lattice VOAs: the MC equation
$\Rightarrow$ prime-locality
\textup{(}Theorem~\textup{\ref{thm:hecke-newton-lattice}}\textup{)}
$\Rightarrow$ CPS automorphy
\textup{(}the Cogdell--Piatetski-Shapiro converse-theorem
package with Jacquet--Shalika strong multiplicity one\textup{)}
$\Rightarrow$ $\operatorname{Sym}^{r-1}(f)$ identification
$\Rightarrow$ Ramanujan bound.
The implication is conditional on the CPS hypotheses for the
moment functions, the Jacquet--Shalika local-factor comparison, and
the symmetric-power automorphy needed in Serre's reduction. It is a
precise Langlands comparison criterion for the MC tower; it is not an
unconditional replacement for Deligne's \'{e}tale-cohomological proof.
\end{corollary}

\subsection{The non-lattice extension}%
\label{sec:non-lattice-obstruction}%
\index{Ramanujan bound!non-lattice extension|textbf}

The lattice restriction is removable once the character-level
converse-theorem hypotheses are supplied:
\emph{characters of any vertex algebra are holomorphic
functions of~$\tau$}, but holomorphy alone is not the
Cogdell--Piatetski-Shapiro input.

\begin{theorem}[Conditional non-lattice Ramanujan bound]%
\label{thm:non-lattice-ramanujan}%
\ClaimStatusProvedHere%
\index{Ramanujan bound!general VOA}%
Let $\cA$ be a rational chirally Koszul algebra with
HS-sewing and diagonal modular invariant partition function
$Z(\tau,\bar\tau) = \sum_{i=1}^N |\chi_i(\tau)|^2$,
where $\chi_i$ are the irreducible characters.
Assume, in addition, that for every Hecke eigencomponent in the
Franc--Mason decomposition of the character vector, the corresponding
character-level Rankin--Selberg series has meromorphic continuation,
the specified functional equation, boundedness in vertical strips, and
Euler-prime locality in the form required by the
Cogdell--Piatetski-Shapiro converse theorem and Jacquet--Shalika
strong multiplicity one.
Then every Hecke eigenform~$f_j$ appearing in
the Rankin--Selberg decomposition of $|\chi_i|^2$ satisfies
the Ramanujan bound
$|a_{f_j}(p)| \leq 2\,p^{(k_j-1)/2}$, where $k_j$ is the
weight of~$f_j$.
\end{theorem}

\begin{proof}
The proof records four reductions under the hypotheses of the
statement.

\smallskip
\emph{Observation~$\mathrm{A}$
\textup{(}holomorphy of characters\textup{)}.}
Each character $\chi_i(\tau) = q^{h_i-c/24}\sum_n d_n^{(i)}\,q^n$
is holomorphic in~$\tau$ (it is
$\operatorname{tr}_{V_i}(q^{L_0-c/24})$, a convergent power
series in~$q$). The term $|\chi_i|^2$ is therefore of the form
$(\text{holomorphic})\times(\text{anti-holomorphic})$,
the structure required by the Rankin--Selberg method.

\smallskip
\emph{Observation~$\mathrm{B}$
\textup{(}character-level Rankin--Selberg\textup{)}.}
The RS integral of $|\chi_i|^2$ against~$E_s$ gives the
$L$-function $L(s, \chi_i \times \bar\chi_i)$ by the standard
Rankin--Selberg unfolding:
\[
 \int_{\mathrm{SL}(2,\bZ)\backslash\mathfrak{H}}
 |\chi_i(\tau)|^2\,E^*(s,\tau)\,d\mu(\tau)
 \;=\;
 \sum_{n\ge1} |d_n^{(i)}|^2\,
 n^{-s}\,\Gamma_\infty(s,n),
\]
where $\Gamma_\infty$ collects the archimedean factors.
The HS-sewing condition ensures convergence and polynomial
growth; the MC recursion (applied character-by-character)
gives the functional equation on the formal side. The analytic
continuation, vertical boundedness, and Euler-prime locality needed
to pass from this formal Rankin--Selberg package to automorphy are
exactly the additional converse-theorem hypotheses imposed in the
statement.

\smallskip
\emph{Observation~$\mathrm{C}$
\textup{(}vector-valued Hecke decomposition\textup{)}.}
For a rational VOA, the character vector
$F(\tau) = (\chi_1,\dotsc,\chi_N)$ transforms under a
finite-dimensional representation
$\rho\colon \mathrm{SL}(2,\bZ) \to \mathrm{GL}(N,\bC)$.
The Hecke operators on vector-valued modular forms
(Franc--Mason~\cite{FrancMason2017}) act on the character space.
The hypothesis that the relevant eigencomponents have
multiplicative Fourier coefficients gives prime-locality for the
character-level RS $L$-functions.

\smallskip
\emph{Observation~$\mathrm{D}$
\textup{(}completion of the twelve stations\textup{)}.}
With prime-locality imposed at the character level by the statement,
the twelve-station implication applies:
\begin{itemize}[nosep]
\item Stations~$1$--$3$: replace Poisson summation
 with the modular transformation law of~$F$ and the
 vector-valued Hecke spectral theorem.
\item Stations~$4$--$6$: unchanged (MC equation, intertwining,
 HS-sewing are properties of the algebra, not of
 the lattice).
\item Stations~$7$--$12$: RS, CPS, prime-locality, strong
 multiplicity one, Serre, all applied to the
 character-level $L$-function
 $L(s, \chi_i \times \bar\chi_i)$.
\end{itemize}
The Ramanujan bound on the Hecke eigenvalues of each
eigencomponent follows from Serre's reduction,
completing the proof.
\end{proof}

\begin{remark}[Scope and limitations]%
\label{rem:non-lattice-gap}%
\index{Ramanujan bound!scope}%
The theorem requires:
\begin{enumerate}[label=\textup{(\alph*)}]
\item \emph{Rationality}: finitely many irreducible
 characters. For irrational VOAs (Virasoro at
 generic~$c$), the character space is infinite-dimensional
 and the Franc--Mason Hecke theory does not directly apply.
\item \emph{Diagonal modular invariant}:
 $Z = \sum|\chi_i|^2$. Non-diagonal invariants
 $Z = \sum M_{ij}\,\chi_i\,\bar\chi_j$ require
 the RS method for mixed terms $\chi_i\,\bar\chi_j$
 ($i \ne j$), which gives the Rankin--Selberg
 convolution $L(s, \chi_i \times \bar\chi_j)$ rather than
 $L(s, \chi_i \times \bar\chi_i)$. The same conclusion would
 require the mixed-term continuation, vertical-strip, and
 local-factor hypotheses.
\item \emph{HS-sewing}: polynomial OPE growth $+$
 subexponential sector growth. All standard families
 satisfy this
 (Theorem~\ref{thm:general-hs-sewing}).
\end{enumerate}
Conjecture~\ref{conj:irrational-ramanujan} below removes the
rationality restriction.
\end{remark}

\subsection{The irrational extension}%
\label{sec:irrational-extension}%
\index{Ramanujan bound!irrational VOAs|textbf}

\begin{conjecture}[Ramanujan bound for irrational VOAs]%
\label{conj:irrational-ramanujan}%
\ClaimStatusConjectured%
\index{Roelcke--Selberg!as infinite-dimensional Hecke theory}%
\index{Ramanujan bound!irrational VOAs}%
Let $\cA$ be a chirally Koszul algebra with HS-sewing at
\emph{any} central charge~$c$ \textup{(}rational or
irrational\textup{)}. Let $\chi_0(\tau)$ be the vacuum
character. Then every Hecke eigenform appearing in the
spectral decomposition of the $L$-function
$L(s, \chi_0 \times \bar\chi_0)
= \sum_{n \ge 1} |d_n|^2\,n^{-s}$ satisfies
the Ramanujan bound.
\end{conjecture}

\begin{remark}[Conditional spectral mechanism]
The irrational extension requires three inputs beyond the rational
case.

\smallskip
\emph{Input~$1$
\textup{(}holomorphy is universal\textup{)}.}
The vacuum character
$\chi_0(\tau) = \operatorname{tr}_{V_0}(q^{L_0-c/24})$ is a
holomorphic function of~$\tau$ for \emph{every} vertex algebra,
regardless of rationality. The HS-sewing condition
(Theorem~\ref{thm:general-hs-sewing}) ensures
$|\chi_0(\tau)|^2 = O(y^{-N})$ as $y \to \infty$ for some~$N$,
so $|\chi_0|^2$ is an $L^2_{\mathrm{loc}}$-automorphic function
with at most polynomial growth.

\smallskip
\emph{Input~$2$
\textup{(}the Roelcke--Selberg spectral theorem as
infinite-dimensional Hecke theory\textup{)}.}
The function $|\chi_0(\tau)|^2$ admits a Roelcke--Selberg
decomposition
\begin{equation}\label{eq:roelcke-selberg-chi0}
 |\chi_0(\tau)|^2
 = \sum_{j} c_j\,\phi_j(\tau)
 + \frac{1}{4\pi}\int_{-\infty}^{\infty}
 c(t)\,E(\tfrac12+it,\tau)\,dt,
\end{equation}
where $\{\phi_j\}$ are Maass--Hecke cusp forms (the discrete
spectrum of the Laplacian on
$\mathrm{SL}(2,\bZ)\backslash\mathfrak{H}$) and $E(\tfrac12+it)$
is the Eisenstein series (continuous spectrum).
\emph{Every}~$\phi_j$ is automatically a Hecke eigenform:
$T_p(\phi_j) = \lambda_p^{(j)}\,\phi_j$
(spectral theorem for commuting self-adjoint operators on
$L^2(\mathrm{SL}(2,\bZ)\backslash\mathfrak{H})$).
The Fourier coefficients of~$\phi_j$ are multiplicative.
This is the \emph{infinite-dimensional Hecke decomposition}:
it requires no finiteness assumption on the character space.

\smallskip
\emph{Input~$3$
\textup{(}the Rankin--Selberg convolution\textup{)}.}
The RS integral of $|\chi_0|^2$ against~$E^*(s)$ gives the
convolution $L$-function:
\[
 \int |\chi_0(\tau)|^2\,E^*(s,\tau)\,d\mu
 = L^*(s, \chi_0 \times \bar\chi_0)
 = \sum_{n \ge 1} |d_n|^2\,n^{-s}\,\Gamma_\infty(s,n).
\]
By the spectral decomposition~\eqref{eq:roelcke-selberg-chi0},
this decomposes into contributions from each~$\phi_j$:
\[
 L(s, \chi_0 \times \bar\chi_0)
 = \sum_j |c_j|^2\,
 L(s, \phi_j \times \bar\phi_j)
 + (\text{continuous part}).
\]
Each $L(s, \phi_j \times \bar\phi_j)
= L(s, \operatorname{Sym}^2 \phi_j) \cdot \zeta(s)$ is the
Rankin--Selberg convolution of a Hecke--Maass eigenform with
itself. By Gelbart--Jacquet:
$L(s, \operatorname{Sym}^2 \phi_j)$ is automorphic
on~$\mathrm{GL}(3)$.

\smallskip
With these inputs and the symmetric-power functoriality stated
below, the twelve-station implication applies component-by-component:
\begin{itemize}[nosep]
\item \emph{Stations~$1$--$3$}: replaced by
 Inputs~$1$--$2$ (holomorphy of~$\chi_0$ $+$
 Roelcke--Selberg spectral decomposition $+$
 multiplicativity of $\phi_j$).
\item \emph{Stations~$4$--$7$}: unchanged (MC equation,
 intertwining, HS-sewing, RS method).
\item \emph{Stations~$8$--$12$}: applied to each spectral
 component $L(s, \phi_j \times \bar\phi_j)$. The CPS
 hypotheses hold (Theorem~\ref{conj:cps-from-mc});
 CPS gives automorphy; prime-locality follows from
 multiplicativity of~$\phi_j$; strong multiplicity one
 gives the $\operatorname{Sym}^r$ identification; Serre
 gives Ramanujan.
\end{itemize}
The Ramanujan bound on the Hecke eigenvalues
$\lambda_p^{(j)}$ of each Maass--Hecke eigenform~$\phi_j$
appearing in~\eqref{eq:roelcke-selberg-chi0} would follow.

\emph{Conditionality.} The Serre reduction at Stations~$8$--$12$
requires automorphy of $L(s,\operatorname{Sym}^r\phi_j)$ for
all~$r \geq 1$. This is known for $r \leq 4$
(Gelbart--Jacquet~$r{=}2$, Kim--Shahidi~$r{=}3$, Kim~$r{=}4$) but open
for $r \geq 5$: it is Langlands functoriality
$\mathrm{GL}(2) \to \mathrm{GL}(r{+}1)$. The theorem is
therefore conditional on Langlands functoriality for all
symmetric powers. Unconditionally, the best bound is the
Kim--Sarnak exponent
$|\lambda_p^{(j)}| \leq 2p^{7/64}$
(which falls short of the Ramanujan bound $2p^0 = 2$
for Maass forms).
\end{remark}

\begin{remark}[The scope is universal]%
\label{rem:universal-ramanujan}%
\index{Ramanujan bound!universal scope}%
Conjecture~\ref{conj:irrational-ramanujan} requires only:
\begin{enumerate}[label=\textup{(\alph*)}]
\item a vertex algebra structure (for the character~$\chi_0$),
\item chiral Koszulness (for the MC element~$\Theta_\cA$),
\item HS-sewing (for the analytic control), and
\item the Roelcke--Selberg spectral theorem on
 $\mathrm{SL}(2,\bZ)\backslash\mathfrak{H}$
 (for the infinite-dimensional Hecke decomposition).
\end{enumerate}
No rationality, no lattice, no theta function.
Under Conjecture~\ref{conj:irrational-ramanujan}, the Ramanujan
bound holds for the Hecke eigenforms appearing in a chirally Koszul
partition function.
The passage from finite-dimensional
(Franc--Mason) to infinite-dimensional (Roelcke--Selberg)
Hecke theory is the required replacement: the spectral theorem for
the Laplacian on a finite-volume hyperbolic surface supplies the
Hecke eigenbasis, while the remaining bounds are exactly the
functoriality hypotheses stated above.
\end{remark}

\subsection{The conditional implication table}

\begin{center}
\small
\renewcommand{\arraystretch}{1.3}
\begin{tabular}{clll}
\toprule
& \textbf{Result} & \textbf{Lattice} & \textbf{Non-lattice} \\
\midrule
1 & $M_r$ meromorphic
 & \ClaimStatusProvedHere
 & \ClaimStatusProvedHere \\
2 & CPS hypotheses
 & \ClaimStatusConditional
 & \ClaimStatusConditional \\
3 & CPS $\Rightarrow$ $\pi_r$ automorphic
 & \ClaimStatusConditional
 & \ClaimStatusConditional \\
4 & Prime-locality
 & \ClaimStatusProvedHere
 & \ClaimStatusConditional
 \textup{(}character-level eigencomponent hypothesis\textup{)} \\
5 & $\pi_r = \operatorname{Sym}^{r-1}(f)$
 & \ClaimStatusConditional
 & \ClaimStatusConditional \\
6 & Ramanujan bound
 & \ClaimStatusConditional
 & \ClaimStatusConditional
 \textup{(}Theorem~\ref{thm:non-lattice-ramanujan}\textup{)} \\
\bottomrule
\end{tabular}
\end{center}

\noindent
For lattice VOAs, the Hecke--Newton closure
(Theorem~\ref{thm:hecke-newton-lattice}) supplies the
prime-locality arrow. For non-lattice theories, the theorem above is
a conditional implication: Franc--Mason Hecke operators provide the
ambient character-level action, while multiplicativity of the
required eigencomponents and compatibility with the shadow amplitudes
remain part of the hypotheses. The remaining question is the
per-prime arithmetic of the OPE and of the genus-$1$ graph sums
(Remark~\ref{rem:non-lattice-gap}).

% ============================================================
\subsection{The prime-locality frontier}%
\label{sec:prime-locality-frontier}%
\index{prime-locality!frontier|textbf}%
\index{prime-locality!criteria}

Prime-locality
(Conjecture~\ref{conj:prime-locality-transfer}) is the single
hypothesis separating the proved algebraic content
($\Theta_\cA \in \MC(\gAmod)$, the shadow obstruction tower
$\Theta_\cA^{\le r}$, the conditional CPS automorphy of the moment
$L$-functions) from the Ramanujan bound: do the Hecke
operators on~$\cM_{1,1}$ lift to~$\gAmod$?

\begin{remark}[Where prime-locality is known]%
\label{rem:prime-locality-known}%
\index{prime-locality!known cases}%
\mbox{}
\begin{enumerate}[label=\textup{(\roman*)}]
\item \emph{Lattice VOAs} \textup{(}proved,
 Theorem~\textup{\ref{thm:hecke-newton-lattice}}\textup{):}
 $\gAmod[V_\Lambda]$ carries a Hecke module structure
 because the graph amplitudes factor through the theta
 function~$\Theta_\Lambda \in M_{r/2}$; at each prime~$p$,
 the MC element projects to the Hecke eigenspace of~$T_p$.
\item \emph{Virasoro at leading $1/c$}
 \textup{(}trivially satisfied\textup{):}
 $\rho = \delta(\lambda + 6/c)$ is a single atom; a single
 atom is automatically prime-local (the same value at every~$p$).
\item \emph{Rational VOAs}
\textup{(}conditional at amplitude level\textup{):}
 the character vector $F(\tau) = (\chi_1,\dotsc,\chi_N)$ is
 a VVMF for a representation
 $\rho\colon\mathrm{SL}(2,\bZ) \to \mathrm{GL}(N,\bC)$.
 Franc--Mason Hecke operators exist on the character space. The
 additional requirements are multiplicativity of the relevant
 eigencomponent Fourier coefficients and compatibility of this Hecke
 action with the genus-$1$ shadow graph sums in~$\gAmod[\cA]$.
\item \emph{Irrational VOAs}
 \textup{(}automatic at character level\textup{):}
 the Roelcke--Selberg decomposition of $|\chi_0|^2$ gives
 Maass--Hecke eigenforms with multiplicative Fourier
 coefficients. But the Ramanujan bound for these eigenforms
 is the Selberg eigenvalue conjecture
 ($\lambda_1 \ge 1/4$, best known $\theta \le 7/64$
 by Kim--Sarnak).
\end{enumerate}
\end{remark}

\begin{remark}[The precise obstruction: the Hecke defect of the
sewing differential]%
\label{rem:prime-locality-obstruction}%
\index{prime-locality!obstruction}%
\index{Hecke defect|textbf}%
In the convolution algebra
$\gAmod = \mathrm{Hom}_\alpha(\cC^!_{\mathrm{ch}},
\cP^{\mathrm{ch}})$, the MC differential $D_\cA$ decomposes as
\[
 D_\cA
 \;=\;
 d_{\mathrm{int}} + d_{\mathrm{sew}} + \hbar\Delta,
\]
where $d_{\mathrm{int}}$ is the internal differential
($\tau$-independent), $d_{\mathrm{sew}}$ is the sewing
differential (involving the propagator~$P(z,w;\tau)$
on~$E_\tau$), and $\hbar\Delta$ is the BV operator (involving
contraction on~$E_\tau$).

The Hecke operator $T_p$ on~$\cM_{1,1}$ sends
$\tau \mapsto (a\tau{+}b)/d$ for $(a,b,d)$ with $ad = p$,
$0 \le b < d$. It acts on any function of~$\tau$, hence on
the genus-$1$ projection $\pi_{1,r}(\Theta_\cA)(\tau) =
\operatorname{Sh}_r^{(1)}(\tau)$.

\emph{The question:} does $T_p$ extend to an endomorphism
of~$\gAmod$ that commutes with~$D_\cA$?

\begin{itemize}[nosep]
\item $[d_{\mathrm{int}}, T_p] = 0$: the internal differential
 is $\tau$-independent. \checkmark
\item $[d_{\mathrm{sew}}, T_p] \;=\; ?$ \;:
 $d_{\mathrm{sew}}$ involves the propagator
 $P(z,w;\tau)$ on~$E_\tau$; Hecke translation changes~$E_\tau$
 to~$E_{(a\tau+b)/d}$, so the propagator transforms.
 For lattice VOAs, $[d_{\mathrm{sew}}, T_p] = 0$ because the
 sewing factors through~$\Theta_\Lambda$, which is a Hecke
 eigenform. For non-lattice theories, this commutator is
 the \emph{Hecke defect}:
 \[
 \delta_p^{\mathrm{Hecke}}(\cA)
 \;:=\;
 [d_{\mathrm{sew}}, T_p]
 \;\in\;
 \mathrm{End}\bigl(\gAmod\bigr).
 \]
\item $[\hbar\Delta, T_p] \;=\; ?$ \;:
 same issue as $d_{\mathrm{sew}}$.
\end{itemize}
Prime-locality is equivalent to
$\delta_p^{\mathrm{Hecke}}(\cA) = 0$ for all primes~$p$,
or more weakly to
$\delta_p^{\mathrm{Hecke}}(\cA)$ being
$D_\cA$-exact (so that $T_p$ preserves the MC condition
up to gauge equivalence).
\end{remark}

\begin{remark}[Four prime-locality criteria]%
\label{rem:prime-locality-routes}%
\index{prime-locality!criteria}%
\mbox{}
\begin{enumerate}[label=\textup{(\Alph*)}]
\item \emph{Hecke module extension.}
 show that $\delta_p^{\mathrm{Hecke}}(\cA)$ is
 $D_\cA$-exact for all chirally Koszul~$\cA$ with HS-sewing.
 This is the extension of
 Theorem~\ref{thm:spectral-continuation-bridge} from
 lattice to non-lattice algebras.
 The obstruction lives in
 $H^1(\gAmod, \mathrm{ad}_{T_p})$, the first cohomology
 of the convolution algebra with coefficients in the
 $T_p$-adjoint module. If $H^1 = 0$
 (e.g., by a vanishing theorem for the deformation complex),
 the Hecke defect is automatically exact.

\item \emph{Motivic criterion.} The shadow amplitudes are
 periods of FM configuration spaces
 (Proposition~\ref{prop:shadow-periods}).
 At genus~$1$, the graph integrals involve elliptic
 polylogarithms (Brown--Levin). If the shadow amplitudes
 lie in the category of \emph{mixed elliptic motives}, the
 motivic Galois group
 $\mathrm{Gal}_{\mathrm{mot}}(\mathrm{MEM}(\bQ))$
 acts through the Hecke algebra, forcing prime factorization.
 The key input: the de~Rham realization of the shadow
 motive must be compatible with the Hecke action on
 the Betti realization, a period comparison theorem for
 elliptic graph integrals.

\item \emph{Analytic criterion.}
 the MC recursion overdetermines the spectral data.
 At degree $r > 2m{+}1$
 (Definition~\ref{def:rigidity-defect}),
 the system has more equations than unknowns.
 By the Carleman condition
 (Proposition~\ref{prop:carleman-virasoro}), the spectral
 measure~$\rho$ is uniquely determined by its moments.
 The moments at degree~$\le 2m{+}1$ are determined by the
 characters (via the intertwining theorem and the
 Stieltjes representation). The characters are
 Hecke-equivariant (automatic from the Roelcke--Selberg
decomposition). Under graph-sum compatibility, the unique~$\rho$ inherits
 Hecke equivariance from the characters, giving
 prime-locality.

 \emph{The key step to prove}: the shadow coefficients at
 high degree ($r > 2m{+}2$) are determined by the
 lower-degree shadow data $+$ the MC recursion. This is
 precisely the content of the rigidity defect: the
 MC recursion at degree $r{+}1$ gives one new equation
 in the existing unknowns, and for $r > 2m{+}1$ the
 system is overdetermined. If the overdetermined system
 has a unique solution (Carleman), and if that solution
 inherits the Hecke equivariance of the character-level
 data (which enters at degrees $\le 2m{+}1$), then
 prime-locality follows.

\item \emph{Finite falsification test.}
 for specific non-lattice VOAs (Ising model ($c = 1/2$,
 $N = 3$ characters), Lee--Yang ($c = -22/5$, $N = 2$),
 three-state Potts ($c = 4/5$, $N = 10$)), compute the
 shadow amplitudes~$\operatorname{Sh}_r^{(1)}$ to high
 degree and verify the Hecke eigenvalue property
 $T_p(\operatorname{Sh}_r) = a_p \cdot \operatorname{Sh}_r$
 at each prime~$p$.
 A counterexample refutes
 Conjecture~\ref{conj:prime-locality-transfer};
 uniform success supplies the finite data needed to formulate a
 structural vanishing theorem for the Hecke defect.
\end{enumerate}

\noindent
MC overdetermination uses the rigidity of the MC element:
$\Theta_\cA \in \MC(\gAmod)$ is overdetermined at high degree. This
rigidity controls the coefficient tower. The prime-by-prime
statement requires the additional assertion that the genus-$1$ graph
sums transport the character-level Hecke structure to the amplitude
level.
\end{remark}

% ============================================================
\subsection{MC overdetermination and
character-level propagation}%
\label{subsec:route-c-development}%
\index{prime-locality!MC overdetermination|textbf}%
\index{MC overdetermination!prime-locality}%
\index{character-level propagation}

MC overdetermination converts the high-degree MC equation into
coefficient-level Hecke invariance, conditional on two structural
inputs: the MC recursion
(Theorem~\ref{thm:mc-recursion-moment}) and the Carleman
uniqueness of the spectral measure
(Proposition~\ref{prop:carleman-virasoro}). Amplitude-level
prime-locality also requires graph-sum Hecke compatibility. The
argument proceeds in three steps: constraint counting, character
determination, and Hecke propagation.

\begin{proposition}[MC constraint counting]%
\label{prop:mc-constraint-counting}%
\ClaimStatusProvedHere%
\index{MC equation!constraint counting}%
\index{overdetermination!MC equation}
The MC equation at degree $r{+}1$ imposes one polynomial
relation on the shadow coefficients
$\{S_2, \ldots, S_{r+1}\}$.
\begin{enumerate}[label=\textup{(\roman*)}]
\item At degrees $2, \ldots, r$, the MC recursion
 produces $(r{-}1)$ independent polynomial equations
 in the shadow coefficients $\{S_2, \ldots, S_r\}$.

\item Each $S_j$ is determined by the spectral data
 $\{(\lambda_i, c_i)\}_{i=1}^{m}$ through
 $S_j = -(1/j)\sum_{i=1}^{m} c_i\,\lambda_i^j$:
 a polynomial of degree~$j$ in the $2m$~spectral
 unknowns \textup{(}$m$~atoms $+$ $m$~weights\textup{)}.

\item For $r \ge 2m + 2$, the number of equations
 $(r{-}1)$ exceeds the number of unknowns~$(2m)$,
 giving rigidity defect
 $\delta(m,r) = r - 1 - 2m > 0$
 \textup{(}Definition~\textup{\ref{def:rigidity-defect}}\textup{)}.
 The shadow coefficients $S_r$ for $r > 2m + 1$ are
 \emph{determined by} $\{S_2, \ldots, S_{2m+1}\}$ and the
 MC recursion.

\item The number of \emph{independent} MC constraints
 at degree $r$ is $\binom{r}{2} - r = r(r{-}3)/2$
 when the convolution bracket generates all pairwise
 interactions. For $r \ge 4$, this exceeds $r{-}2$
 \textup{(}the number of free shadow coefficients at
 degrees $3, \ldots, r$\textup{)}, giving a second
 \textup{(}sharper\textup{)} overdetermination estimate.
\end{enumerate}
\end{proposition}

\begin{proof}
Parts~(i)--(iii) follow from the MC recursion
(Theorem~\ref{thm:mc-recursion-moment}) and the rigidity
defect count (Definition~\ref{def:rigidity-defect}).
Part~(iv): the MC equation
$D_\cA\,\Theta_\cA + \frac{1}{2}[\Theta_\cA,\Theta_\cA] = 0$
at degree $r{+}1$ reads
\begin{equation}\label{eq:mc-degree-expansion}
 \nabla_H(S_{r+1})
 \;+\;
 \sum_{\substack{j + k = r+3 \\ j, k \ge 2}}
 \frac{1}{2}\,\{S_j, S_k\}_H
 \;=\; 0,
\end{equation}
where $\nabla_H$ is the Hochschild operator and
$\{-,-\}_H$ the convolution bracket. The bracket
at degree $r{+}1$ receives contributions from all pairs
$(j,k)$ with $j + k = r + 3$ and $j, k \ge 2$, giving
$\lfloor(r{-}1)/2\rfloor$ distinct unordered pairs.
The cumulative count from degrees $4, \ldots, r$ of such
bracket terms is
$\sum_{s=4}^{r} \lfloor(s{-}1)/2\rfloor
\ge r(r{-}3)/2$ for $r \ge 4$, giving the claimed
sharper bound.
\end{proof}

\begin{theorem}[MC rigidity gives coefficient-level Hecke invariance;
\ClaimStatusConditional]%
\label{thm:route-c-propagation}%
\index{prime-locality!MC overdetermination theorem}%
\index{MC rigidity!forces prime-locality}
Let $\cA$ be a chirally Koszul algebra with HS-sewing and
spectral measure~$\rho_\cA$ satisfying the Carleman condition
\textup{(}Proposition~\textup{\ref{prop:carleman-virasoro}}\textup{)}.
Suppose:
\begin{enumerate}[label=\textup{(\alph*)}]
\item \emph{Character-level Hecke equivariance:}
 the genus-$1$ partition function
 $Z_1(\cA;\tau)$ decomposes as
 $\widehat{Z}^c_\cA(\tau) = \sum_j \alpha_j\,\varphi_j(\tau)$,
 where each $\varphi_j$ is a Hecke eigenfunction
 \textup{(}holomorphic eigenform, Maass eigenform, or
 Eisenstein eigenfunction\textup{)} with eigenvalues
 $T_p(\varphi_j) = a_{j,p}\,\varphi_j$ that are multiplicative
 in~$p$.

\item \emph{Low-degree determination:}
 the shadow coefficients $S_2 = \kappa$ and $S_3 = \alpha$
 are determined by the character-level data
 \textup{(}specifically, by the first two moments of the
 spectral decomposition of~$\widehat{Z}^c_\cA$\textup{)}.
\item \emph{Graph-sum Hecke compatibility:}
 the genus-$1$ graph-sum functor
 $\{S_j\},P(\tau)\mapsto\operatorname{Sh}_r^{(1)}(\tau)$
 intertwines the Hecke action on quasi-modular propagators with the
 convolution bracket.
\end{enumerate}
Then the tau-independent coefficient tower $\{S_r\}_{r\ge2}$ is
Hecke-invariant. With hypothesis~\textup{(c)}, the genus-$1$
amplitudes $\operatorname{Sh}_r^{(1)}(\Theta_\cA;\tau)$ decompose
into Hecke eigencomponents with multiplicative eigenvalues. Without
\textup{(c)}, this theorem proves only the coefficient-level
statement.
\end{theorem}

\begin{proof}
By hypothesis~(b), the low-degree shadow data
$S_2 = \kappa$ and $S_3 = \alpha$ are determined by the
character-level spectral decomposition, which is
Hecke-equivariant by hypothesis~(a).

The MC recursion
(Theorem~\ref{thm:mc-recursion-moment})
determines $S_r$ from $\{S_2, \ldots, S_{r-1}\}$ by
\[
 S_r
 \;=\;
 -\frac{1}{2r}\;
 \sum_{\substack{j + k = r+2 \\ j,k \ge 2,\; j,k < r}}
 j\,k\,S_j\,S_k\,P_0,
\]
where $P_0$ is the scalar contraction in the OPE lane, and the sum
runs over unordered pairs. The elliptic propagator $P(e;\tau)$ appears
only after one passes from the coefficient tower to the genus-$1$ graph
amplitude.

For $r \ge 4$, the rigidity defect $\delta(m,r) > 0$
(Proposition~\ref{prop:rigidity-threshold}), so $S_r$ is
\emph{uniquely determined} by $\{S_2, S_3, S_4\}$ and the MC
recursion (where $S_4$ itself is determined by the
quartic contact class). The coefficient recursion is a polynomial
operation in tau-independent variables, so $T_p$ acts trivially on
the coefficient variables. Hypothesis~\textup{(c)} is precisely the
additional amplitude-level compatibility needed to move $T_p$
through the propagator graph sum. Under \textup{(c)}, the induction
\[
 T_p(S_r)
 \;=\;
 -\frac{1}{2r}\;
 \sum_{\substack{j+k=r+2\\j,k\ge 2,\; j,k<r}}
 j\,k\,T_p(S_j)\,T_p(S_k)\,T_p(P_0)
\]
preserves the Hecke eigenspace decomposition at every
degree, with $T_p(P_0)=P_0$. Without \textup{(c)}, the same induction gives only
$T_p(S_r)=S_r$ for the coefficient tower. By the Carleman condition,
the spectral measure
$\rho_\cA$ is uniquely determined by the full shadow
sequence $\{S_r\}_{r \ge 2}$; the Hecke equivariance of
the amplitudes then gives prime-locality.
\end{proof}

\begin{remark}[Hypotheses~\textup{(a)}--\textup{(c)}]%
\label{rem:route-c-hypotheses}%
\index{prime-locality!MC overdetermination hypotheses}
\mbox{}
\begin{enumerate}[label=\textup{(\roman*)}]
\item \emph{Hypothesis~\textup{(a)}} is satisfied in all
 known cases. For lattice VOAs, $\widehat{Z}^c$ decomposes
 into holomorphic Hecke eigenforms. For rational VOAs,
 $\widehat{Z}^c$ decomposes into finitely many VVMF
 eigencomponents. For irrational~$c$,
 $|\chi_0(\tau)|^2$ enters the Roelcke--Selberg decomposition
 as a linear combination of Eisenstein series and Maass
 eigenforms, which are automatically Hecke eigenfunctions.
 The hypothesis is the assertion that this decomposition is available
 on the character side; compatibility with the shadow obstruction
 tower is the separate graph-sum condition~\textup{(c)}.

\item \emph{Hypothesis~\textup{(b)}} is the substantive
 content. For lattice VOAs, $\kappa = r$ and
 $\alpha$ are determined by the Eisenstein coefficient
 and the first cusp form contribution, both read from the
 theta function. For Virasoro, $\kappa = c/2$ is
 determined by the character asymptotics
 ($q^{-c/24}\chi_0 \sim 1$), and $\alpha = 2$ is the
 Virasoro cubic shadow, determined by the $T$-$T$ OPE
 coefficient. In both cases the low-degree data is
 character-determined. The gap: for a \emph{general}
 chirally Koszul algebra with non-lattice partition
 function, the identification of $S_3$ with a
 character-level datum has not been established in full
 generality.

\item Hypothesis~\textup{(c)} is the amplitude-level content.
 Once $\kappa$ and $\alpha$ are known to be Hecke-equivariant, the
 MC recursion propagates this equivariance to all
 tau-independent coefficients. The
 character-level Hecke equivariance is classical (Rankin
 convolution for holomorphic forms; Roelcke--Selberg
 spectral theory for real-analytic completions), but the genus-$1$
 amplitudes also contain the quasi-modular propagator graph sum.
 The remaining task is to verify that both the extraction of
 $S_2,S_3$ and the graph-sum functor preserve the Hecke
 decomposition.
\end{enumerate}
\end{remark}

\begin{corollary}[Coefficient-level MC overdetermination for the
standard landscape;
\ClaimStatusConditional]%
\label{cor:route-c-standard-landscape}%
\index{prime-locality!standard landscape}
For every algebra in the standard landscape with rational
OPE coefficients, hypotheses
\textup{(a)} and~\textup{(b)} of
Theorem~\textup{\ref{thm:route-c-propagation}} are
satisfied. Hence the shadow coefficient tower is Hecke-invariant
for the standard landscape. Full amplitude-level prime-locality
requires hypothesis~\textup{(c)}.
\end{corollary}

\begin{proof}
By Theorem~\ref{thm:algebraic-family-rigidity}, every
algebraic family with rational OPE coefficients has shadow
coefficients $S_r$ that are rational functions of the
continuous parameter (central charge~$c$ or level~$k$).
The low-degree data $S_2 = \kappa$ and $S_3 = \alpha$ are
explicit rational functions of~$c$, read from the
two-point and three-point OPE coefficients, which are
character-determined (the OPE structure constants of the
universal algebra are fixed by the Lie algebra data).
Hypothesis~(a) is satisfied because the partition
function of an algebraic family is a meromorphic function
of~$\tau$ and~$c$, and the Roelcke--Selberg decomposition
of such functions is Hecke-equivariant.
\end{proof}

% ============================================================
\section{The Koszul self-duality principle}%
\label{sec:koszul-self-duality-principle}%
\index{Koszul self-duality!arithmetic principle|textbf}%
\index{self-duality!forces L-function factorization|textbf}

This section isolates the algebraic condition under which the
constrained Epstein zeta factors through the finite Hecke decomposition
of its modular-form space: strict Koszul self-duality
$\cA\simeq\cA^!$. This is not the bar--cobar inversion
$\Omega B(\cA)=\cA$ and not a statement about the derived centre.

\subsection{Verdier--Hecke commutation}

\begin{theorem}[Verdier--Hecke commutation at genus~$1$]%
\label{thm:hecke-verdier-commutation}%
\ClaimStatusProvedHere%
\index{Verdier duality!commutes with Hecke operators}%
\index{Hecke operators!commute with Verdier duality}%
On the genus-$1$ moduli space~$\cM_{1,1}$, the Verdier
duality involution~$\sigma$ induced by $\cA\simeq\cA^!$
commutes with the Hecke correspondences~$T_p$ for every
prime $p\nmid N(\cA)$.
\end{theorem}

\begin{proof}
The Hecke operator~$T_p$ is the pushforward along the isogeny
correspondence
$\cM_{1,1} \xleftarrow{\;\pi_1\;}
\cM_{1,1}(p) \xrightarrow{\;\pi_2\;} \cM_{1,1}$,
where $\cM_{1,1}(p)$ parametrises pairs $(E,E',\phi)$ with
$\phi\colon E\to E'$ an isogeny of degree~$p$.
Verdier duality acts via Serre duality on~$X$: for elliptic
curves, $E\simeq E^\vee$ (principal polarisation), so $\bD$
is the identity on isomorphism classes. The dual isogeny
$\phi^\vee\colon (E')^\vee \to E^\vee$ has the same degree as~$\phi$
and bijects the set of degree-$p$ isogenies from~$E$ with
the set of degree-$p$ isogenies from~$E^\vee \simeq E$.
Hence $\bD\circ T_p = T_p\circ\bD$ at the level of
correspondences.

For lattice VOAs explicitly: $T_p$ sums over index-$p$
sublattices $\Lambda' \subset \Lambda$. Duality sends
$\Lambda'$ to $(\Lambda')^* = p^{-1}\Lambda'{}^\perp$.
The map $\Lambda' \mapsto (\Lambda')^*$ bijects
$\{\Lambda'\subset\Lambda : [\Lambda:\Lambda']=p\}$ with
$\{\Lambda''\subset\Lambda^* : [\Lambda^*:\Lambda'']=p\}$,
establishing $\sigma\circ T_p = T_p\circ\sigma$ where
$\sigma(\Theta_\Lambda) = \Theta_{\Lambda^*}$.
\end{proof}

\subsection{The self-dual factorisation theorem}

\begin{theorem}[Self-dual factorisation]%
\label{thm:self-dual-factorization}%
\ClaimStatusProvedHere%
\index{self-duality!factorisation theorem}%
\index{L-function!factorisation from self-duality}%
Let $\cA$ be a chirally Koszul algebra with partition function
$Z_\cA \in M_k(\Gamma_0(N))$. If $\cA\simeq\cA^!$
\textup{(}Koszul self-dual\textup{)}, then the constrained
Epstein zeta $\varepsilon^c_s(\cA)$ decomposes as a finite sum
of standard automorphic $L$-functions:
\begin{equation}\label{eq:self-dual-factorization}
 \varepsilon^c_s(\cA)
 \;=\;
 \sum_{j=0}^{d(\cA)-2} c_j\,\Gamma_j(s)\,L(s, f_j),
\end{equation}
where each $f_j$ is a Hecke eigenform for $\Gamma_0(N)$,
$\Gamma_j(s)$ collects archimedean factors, and
$d(\cA)$ is the shadow depth.
\end{theorem}

\begin{proof}
The factorisation is obtained as follows.

\smallskip
\emph{Step~1} \textup{(Verdier involution)}.
Theorem~A gives $\bD(\bar B(\cA)) \simeq \bar B(\cA^!)$.
When $\cA \simeq \cA^!$, this yields an involution
$\sigma\colon \bar B(\cA) \xrightarrow{\;\sim\;} \bar B(\cA)$
with $\sigma^2 = \mathrm{id}$.

\smallskip
\emph{Step~2} \textup{(functional equation)}.
The involution~$\sigma$, transported through the
Rankin--Selberg integral
(Theorem~\ref{thm:sewing-selberg-formula} for the
sewing operator; the general case follows from the
modularity of $Z_\cA$ under $\Gamma_0(N)$
combined with the Eisenstein equivariance
$E_s(-1/\tau) = |\tau|^{2s}E_s(\tau)$), produces the
functional equation of~$\varepsilon^c_s$.

\smallskip
\emph{Step~3} \textup{(Hecke eigenform decomposition)}.
By Theorem~\ref{thm:hecke-verdier-commutation},
$\sigma$ commutes with $T_p$ for $p\nmid N$. Since
$\sigma(Z_\cA) = Z_\cA$ (self-duality), $Z_\cA$ lies
in the $\sigma$-invariant subspace of $M_k(\Gamma_0(N))$.
Strong multiplicity one for newforms on~$\Gamma_0(N)$
(Atkin--Lehner) ensures that each Hecke eigenspace is
one-dimensional; hence $Z_\cA$ decomposes as a finite
sum of Hecke eigenforms $f_j$.

\smallskip
\emph{Step~4} \textup{(standard $L$-functions)}.
Each Hecke eigenform~$f_j$ gives a standard automorphic
$L$-function $L(s,f_j)$ with Euler product and meromorphic
continuation (Hecke theory). The number of eigenforms is
$d(\cA)-1$ by the shadow--spectral correspondence
(Theorem~\ref{thm:shadow-spectral-correspondence}).
The shadow algebra~$\cA^{\mathrm{sh}}$ determines the
Hecke eigenvalues via the Newton relations in spectral
coordinates (Theorem~\ref{thm:hecke-newton-lattice}).

\end{proof}

\subsection{The theta decomposition bridge for rational VOAs}

\begin{proposition}[Theta decomposition bridge]%
\label{prop:theta-bridge-rational}%
\ClaimStatusProvedHere%
\index{theta decomposition!for rational VOAs}%
\index{minimal model!theta decomposition of characters}%
Let $\cM(p,q)$ be a minimal model with diagonal modular
invariant. The characters decompose via the Rocha--Caridi
formula into theta functions at level~$2pq$:
\begin{equation}\label{eq:theta-bridge}
 \chi_{r,s}(\tau)
 \;=\;
 \frac{\Theta_{ps-qr,\,2pq}(\tau)
 - \Theta_{ps+qr,\,2pq}(\tau)}{\eta(\tau)},
\end{equation}
where $\Theta_{a,N}(\tau) =
\sum_{n\in\bZ} q^{(Nn+a)^2/(2N)}$.
Consequently, $|\chi_{r,s}|^2$ decomposes into lattice
Epstein zeta functions, and the Rankin--Selberg integral gives
Dirichlet $L$-functions modulo~$2pq$.
\end{proposition}

\begin{proof}
The Rocha--Caridi formula
$\chi_{r,s} =
(\Theta_{ps-qr,2pq}-\Theta_{ps+qr,2pq})/\eta$
is classical; see~\cite{DiFrancesco1997}.
Each $\Theta_{a,N}$ is the theta function of the
one-dimensional lattice $\bZ/\sqrt{N}$ with
characteristic~$a/N$.
The diagonal partition function
$Z = \sum_{i}|\chi_i|^2$ simplifies by cancellation
of cross terms
(for $\cM(3,4)$: $Z = (\theta_2+\theta_3+\theta_4)/(2\eta)$).
Each surviving theta product gives a lattice Epstein zeta
whose Mellin transform is a partial zeta function
$\zeta(s;a,2pq) = \sum_{n\equiv a\pmod{2pq}} n^{-2s}$.
These decompose via the Dirichlet characters of
$(\bZ/2pq\bZ)^*$ into standard Dirichlet $L$-functions
$L(s,\chi)$.
\end{proof}

\begin{remark}[Resolution of the non-lattice gap for rational VOAs]%
\label{rem:theta-bridge-resolution}%
\index{non-lattice gap!resolved for rational VOAs}%
Proposition~\ref{prop:theta-bridge-rational} resolves the
non-lattice obstruction
(Remark~\ref{rem:non-lattice-gap}) for all rational VOAs.
The non-multiplicativity of the scalar Dirichlet series
$D(s) = \sum a_n\,n^{-s}$ from $\sum|\chi_i|^2$, observed
for the Ising, Lee--Yang, tricritical Ising, and three-state
Potts models, is a \emph{projection artefact}: $D(s)$ is a
sum of individually multiplicative Dirichlet $L$-functions,
and the sum of multiplicative functions is not multiplicative.
Each summand retains its Euler product.
The remaining open regime is irrational central charge~$c$,
where $|\eta(\tau)|^{2(c-1)}$ is not a modular form for any
congruence subgroup and the spectral decomposition necessarily
involves continuous spectrum
(Conjecture~\ref{conj:irrational-ramanujan}).
\end{remark}

\subsection{The Newton--shadow--Hecke dictionary}

\begin{proposition}[Newton--shadow--Hecke correspondence]%
\label{prop:newton-shadow-hecke}%
\ClaimStatusConditional%
\index{Newton identities!shadow--Hecke dictionary}%
\index{shadow algebra!determines Hecke eigenvalues}%
The shadow algebra $\cA^{\mathrm{sh}} =
H_\bullet(\operatorname{Def}^{\mathrm{mod}}_{\mathrm{cyc}}(\cA))$
determines the Hecke eigenvalues of the spectral decomposition
of~$Z_\cA$ via Newton's identities. Specifically, the shadow
power sums $p_r = r\cdot S_r$, where $S_r$ are the shadow
coefficients at degree~$r$, determine the elementary symmetric
polynomials~$e_k$ of the Hecke eigenvalues by Newton's
recursion:
\begin{equation}\label{eq:newton-shadow-hecke}
 p_r - e_1\,p_{r-1} + e_2\,p_{r-2} - \cdots
 + (-1)^{r-1}\,r\cdot e_r = 0.
\end{equation}
The MC equation $D\cdot\Theta + \tfrac12[\Theta,\Theta] = 0$
at degree~$4$ is structurally parallel to the Hecke algebra
relation $T_2^2 = T_4 + p^{k-1}\cdot T_1$: both are
quadratic constraints at level~$4$ with a correction term
controlled by the weight.
\end{proposition}

\begin{proof}
The shadow moments $p_r = r\cdot S_r$ are the image of
$\Theta_\cA^{\le r}$ under the projection
$\operatorname{MC}(\mathfrak{g}^{\mathrm{mod}}_\cA) \to
\cA^{\mathrm{sh}}$
(Corollary~\ref{cor:shadow-extraction}).
For lattice VOAs, the Hecke module structure
(Theorem~\ref{thm:hecke-newton-lattice}) identifies
$p_r$ with the $r$-th power sum of the Satake parameters
$\{\alpha_{j,p}, \beta_{j,p}\}$ at each prime.
Newton's identity~\eqref{eq:newton-shadow-hecke} then
recovers the elementary symmetric polynomials
$e_1 = \alpha + \beta = a_f(p)$ (the Hecke eigenvalue)
and $e_2 = \alpha\beta = p^{k-1}$.
The MC equation at degree~$4$ gives
Newton at degree~$4$: $p_4 = e_1\,p_3 - e_2\,p_2$
(Proposition~\ref{prop:mc-bracket-determines-atoms}),
which specialises to
$p_4 = e_1\,p_3 - e_2\,p_2 + 4\,e_2^2/e_1$; the
Hecke relation $T_2^2 = T_4 + p^{k-1}$ is the same
identity with $e_1 = a(p)$ and $e_2 = p^{k-1}$.
\end{proof}

\subsection{The Davenport--Heilbronn principle}

\begin{remark}[Koszul self-duality and on-line zeros]%
\label{rem:davenport-heilbronn-koszul}%
\index{Davenport--Heilbronn!Koszul explanation}%
\index{class number one!Koszul self-duality}%
The Davenport--Heilbronn phenomenon (Epstein zeta functions
of non-unimodular binary quadratic forms having zeros off the
critical line) is explained by the failure of Koszul
self-duality. A lattice~$\Lambda$ with discriminant~$D$
gives a lattice VOA with $V_\Lambda^! = V_{\Lambda^*}$.
Self-duality $V_\Lambda \simeq V_{\Lambda^*}$ holds iff
$\Lambda \simeq \Lambda^*$, which for binary lattices is
equivalent to class number $h(D) = 1$
(the nine Heegner discriminants $D =
{-3}, {-4}, {-7}, {-8}, {-11}, {-19}, {-43}, {-67}, {-163}$).
When $h(D) = 1$, the Epstein zeta factors as
$E(s;Q) = C\cdot\zeta(s)\cdot L(s,\chi_D)$ (product of
standard $L$-functions), whose zeros lie on $\mathrm{Re}(s) =
1/2$ assuming GRH\@. When $h(D) > 1$, the lattice is
\emph{not} its own Koszul dual: the Verdier involution maps
$V_\Lambda$ to $V_{\Lambda^*} \not\simeq V_\Lambda$,
the partition function is a sum of theta functions of
non-isometric lattices in the genus, and $\sigma$ fails to act
as an involution on each summand. The Epstein zeta does not
factor, and the Davenport--Heilbronn mechanism produces
off-line zeros.

Self-duality is the algebraic condition; $L$-function
factorisation is the arithmetic consequence; on-line zeros
are the analytic corollary.
\end{remark}

% ============================================================
\section{The complexified modular integral}%
\label{sec:complexified-modular-integral}%
\index{Polyakov complexification|textbf}%
\index{complexified modular integral|textbf}

The spectral approach reaches the
Ramanujan bound, its natural terminus. The spectral
decomposition constrains the coefficients~$c(t)$ on the
real line~$t\in\bR$; the zeta zeros live at
complex spectral parameters
(Remark~\ref{rem:structural-obstruction}). The sewing
determinant $\det(1-K(\tau))$ and the Laplacian spectral
determinant $\det{}_\zeta'(-\Delta_{E_\tau})$ encode the same
data (the Dedekind $\eta$-function), but holomorphicity
of~$\eta$ cuts both ways. It powers the
bar-cobar machinery: the MC equation, the shadow
expansion, the Hecke--Newton closure all flow from the
holomorphic sewing construction. It also limits
the arithmetic programme: the K\"ahler
potential of~$\cM_{1,1}$ decomposes into a geometric part
(the Weil--Petersson metric) and an arithmetic
part ($\log\lvert\eta\rvert^2$, encoding the zeta zeros),
and the arithmetic part is \emph{pluriharmonic},
annihilated by $\partial\bar\partial$, hence invisible
to the K\"ahler geometry.

% ============================================================
\subsection{The spectral determinant bridge}%
\label{subsec:spectral-determinant-bridge}%
\index{spectral determinant!sewing bridge}%
\index{Kronecker limit formula!spectral bridge}

The sewing Fredholm determinant and the spectral zeta
determinant of the Laplacian on~$E_\tau$ encode the same
transcendental data (the Dedekind $\eta$-function) in two
different representations.

\begin{proposition}[Sewing--spectral determinant bridge]%
\label{prop:sewing-spectral-bridge}%
\ClaimStatusProvedHere%
\index{Polyakov--Alvarez formula!sewing bridge}%
Let $E_\tau = \bC/(\bZ+\bZ\tau)$ carry the flat metric
$\lvert dz\rvert^2$, with $y = \mathrm{Im}(\tau)$ and
$q = e^{2\pi i\tau}$.
\begin{enumerate}[label=\textup{(\roman*)}]
\item \textup{(}Sewing.\textup{)}
 The sewing Fredholm determinant satisfies
 \begin{equation}\label{eq:sewing-det-eta}
 \det(1-K(\tau))
 \;=\;
 \prod_{n\ge1}(1-q^n)
 \;=\;
 \eta(\tau)\,e^{-\pi i\tau/12}\,.
 \end{equation}
\item \textup{(}Spectrum.\textup{)}
 The zeta-regularized determinant of the scalar Laplacian
 $-\Delta_{E_\tau}$ \textup{(}zero mode excluded\textup{)}
 satisfies
 \begin{equation}\label{eq:spectral-det-eta}
 \det{}_\zeta'\!(-\Delta_{E_\tau})
 \;=\;
 4\pi^2\,y\,\lvert\eta(\tau)\rvert^4\,.
 \end{equation}
\item \textup{(}Bridge.\textup{)}
 Combining~\eqref{eq:sewing-det-eta}
 and~\eqref{eq:spectral-det-eta}:
 \begin{equation}\label{eq:sewing-spectral-bridge}
 \det{}_\zeta'\!(-\Delta_{E_\tau})
 \;=\;
 4\pi^2\,y\,e^{-\pi y/3}\,
 \lvert\det(1-K(\tau))\rvert^4\,.
 \end{equation}
\end{enumerate}
\end{proposition}

\begin{proof}
Part~(i): the sewing operator $K(\tau)$ on Fock space has
eigenvalues~$q^n$ ($n\ge 1$, \S\ref{sec:sewing-RS-bridge}), so
$\det(1-K) = \prod(1-q^n) = q^{-1/24}\,\eta(\tau)
= e^{-\pi i\tau/12}\,\eta(\tau)$.

Part~(ii): the eigenvalues of $-\Delta_{E_\tau}$
(zero mode excluded) are
$\lambda_{m,n} = 4\pi^2\lvert m+n\tau\rvert^2/y^2$
for $(m,n)\ne(0,0)$. The spectral zeta function
$\zeta_\Delta(s) = \sum'_{(m,n)} \lambda_{m,n}^{-s}$
is the Epstein zeta function of the lattice
$\bZ+\bZ\tau$, rescaled by $(2\pi/y)^{-2s}$. The Kronecker
first limit formula~\cite{Selberg1956} gives
$\zeta_\Delta'(0)
= -\log(4\pi^2\,y\,\lvert\eta\rvert^4)$,
and $\det{}_\zeta'(-\Delta) = e^{-\zeta_\Delta'(0)}$
yields~\eqref{eq:spectral-det-eta}.

Part~(iii): from~(i),
$\lvert\det(1-K)\rvert^2
= \lvert\eta\rvert^2\,e^{\pi y/6}$,
so $\lvert\eta\rvert^4
= \lvert\det(1-K)\rvert^4\,e^{-\pi y/3}$.
Substituting into~(ii):
$\det{}_\zeta'(-\Delta)
= 4\pi^2\,y\,\lvert\det(1-K)\rvert^4\,e^{-\pi y/3}$.
\end{proof}

\begin{remark}[Two chiralities]%
\label{rem:two-chiralities}%
\index{chirality!sewing vs spectral}%
The sewing determinant $\det(1-K)$ is holomorphic in~$\tau$:
it sees one chirality. The spectral determinant
$\det{}_\zeta'(-\Delta)$ is real-analytic: it sees both
chiralities ($\lvert\eta\rvert^4 = \eta^2\bar\eta^2$).
Passing from sewing to spectral doubles the chirality
and introduces the factors $4\pi^2\,y\,e^{-\pi y/3}$
(volume, $q^{1/24}$ correction). The spectral zeta function
$\zeta_\Delta(s)$ admits meromorphic continuation to
all $s\in\bC$ via the functional equation of the Epstein zeta
function, providing (in principle) access to complex spectral
parameters invisible to the sewing product $\prod(1-q^n)$,
which converges only for~$\lvert q\rvert < 1$.
\end{remark}

\begin{remark}[The Kronecker--Selberg hierarchy]%
\label{rem:kronecker-limit}%
\index{Kronecker limit formula}%
The Kronecker limit formula connects
$\log\lvert\eta(\tau)\rvert^2$ to~$E_s(\tau)$ at $s=1$:
it is the $s\to 1$ residue of the Epstein zeta function.
The sewing--Selberg formula
(Theorem~\ref{thm:sewing-selberg-formula}) extends this
connection to all~$s$: the Rankin--Selberg integral
$\int\log\det(1-K)\cdot E_s\,d\mu$ equals a product of
gamma and zeta functions, providing explicit meromorphic
continuation. The complexified programme below asks what
structure this hierarchy reveals at the scattering
poles~$s = \rho/2$.
\end{remark}

% ============================================================
\subsection{The complexified Rankin--Selberg integral}%
\label{subsec:complexified-eisenstein-integral}%
\index{Rankin--Selberg!complexified}

\begin{proposition}[Meromorphic continuation of the RS integral]%
\label{prop:rs-analytic-continuation}%
\ClaimStatusProvedHere%
\index{Rankin--Selberg!meromorphic continuation}%
For a chirally Koszul algebra~$\cA$ with HS-sewing, the
Rankin--Selberg integral
\begin{equation}\label{eq:rs-complexified}
 I(s)
 \;:=\;
 \int_{\mathrm{SL}(2,\bZ)\backslash\mathfrak{H}}
 \widehat{Z}^c(\tau,\bar\tau;\cA)
 \;\cdot\;E_s(\tau)\,d\mu(\tau)
\end{equation}
continues meromorphically to all $s\in\bC$. Its poles are:
\begin{enumerate}[label=\textup{(\roman*)}]
\item a simple pole at $s = 1$ with residue
 $\tfrac{3}{\pi}\,\langle\widehat{Z}^c,\mathbf{1}\rangle$;
\item simple poles at $s = s_j$ for each discrete Maass
 eigenvalue $\lambda_j = s_j(1-s_j)$ appearing in the
 spectral decomposition of~$\widehat{Z}^c$.
\end{enumerate}
\end{proposition}

\begin{proof}
Rankin--Selberg unfolding reduces $I(s)$ to the Mellin
transform
$I(s) = \int_0^\infty a_0(y;\cA)\,y^{s-2}\,dy$,
where
$a_0(y) = \int_0^1 \widehat{Z}^c(x+iy)\,dx$
is the zeroth Fourier mode. The HS-sewing condition
(Theorem~\ref{thm:general-hs-sewing})
gives $a_0(y) = O(e^{-2\pi y})$ as $y\to\infty$; the
$y^{c/2}$ prefactor gives $a_0(y) = O(y^{c/2})$ as
$y\to 0^+$. To continue meromorphically, split
$\int_0^\infty = \int_0^1 + \int_1^\infty$. The upper
integral is entire in~$s$ (exponential decay). For the
lower integral, expand~$a_0(y)$ in its asymptotic series
near $y=0$:
$a_0(y) = c_0 + c_1\,y^{s_1} + \cdots$ (from the spectral
expansion). Integrating term by term,
$\int_0^1 c_k\,y^{s_k+s-2}\,dy = c_k/(s_k+s-1)$, giving
simple poles at $s = 1 - s_k$. The pole at $s=1$ comes
from the constant term
$\langle\widehat{Z}^c,\mathbf{1}\rangle$; the poles at
$s=s_j$ from the Maass coefficients.
\end{proof}

\begin{proposition}[Holomorphy at scattering poles]%
\label{prop:scattering-residue}%
\ClaimStatusProvedHere%
\index{scattering matrix!holomorphy at}%
\index{Rankin--Selberg!scattering poles}%
The integral $I(s)$ is \emph{holomorphic} at the scattering
poles $s = \rho/2$, where $\rho$ ranges over the nontrivial
zeros of~$\zeta$.
\end{proposition}

\begin{proof}
The Mellin representation
$I(s) = \int_0^\infty a_0(y)\,y^{s-2}\,dy$, established
by Rankin--Selberg unfolding, has poles determined
\emph{exclusively} by the asymptotic behaviour
of~$a_0(y)$ at $y=0$ and $y=\infty$. The scattering matrix
$\varphi(s) = \Lambda(1{-}s)/\Lambda(s)$ governs the
functional equation $E_s = \varphi(s)\,E_{1-s}$ of the
Eisenstein series, but the unfolding procedure has replaced
the integration against~$E_s$ by the Mellin integral
of~$a_0$. Since $a_0(y)$ is smooth on $(0,\infty)$ with at
most polynomial growth at~$0$ and exponential decay
at~$\infty$, the Mellin transform has poles only at the
values dictated by the constant and Maass terms in the
asymptotic expansion, none of which coincides with
$s = \rho/2$ for a nontrivial zero~$\rho$ of~$\zeta$.
\end{proof}

\begin{remark}[The structural obstruction, quantified]%
\label{rem:structural-obstruction-revisited}%
\index{structural obstruction!Polyakov perspective}%
Proposition~\ref{prop:scattering-residue} quantifies the
obstruction of Remark~\ref{rem:structural-obstruction}.
Rankin--Selberg unfolding continues~$I(s)$ meromorphically
to all $s\in\bC$, but the mechanism that achieves the
continuation (replacing $\int f\cdot E_s\,d\mu$ by the
Mellin transform of~$a_0$) simultaneously \emph{erases}
the scattering matrix. The MC equation constrains
$\widehat{Z}^c$ and hence $a_0$, but $a_0$ depends on~$y$
alone, and its Mellin transform carries no information about
the scattering poles.

To recover zero-location information from $I(s)$, one would
need a representation that retains, rather than unfolds away,
the Eisenstein series~$E_s$ and its scattering poles.
\end{remark}

\begin{remark}[The algebra--representation gap]%
\label{rem:algebra-representation-gap}%
\index{algebra--representation gap|textbf}%
\index{bar complex!algebra vs representation theory}%
\index{structural obstruction!representation-theoretic form}%
The structural obstruction has a representation-theoretic
reformulation.
The bar complex~$\barBch(\cA)$ is built from OPE
data: structure constants, central charge,
conformal weights. This is data of the
\emph{algebra}. The primary spectrum
$\{\Delta_i\}_{i \in I}$ is data of the
\emph{representation theory}. The shadow obstruction tower, the
MC equation, and the Koszul--Epstein function live in the
algebraic world, determined by $\cA$ as an algebra, not
by $\cA\text{-}\mathrm{mod}$ as a category. The zeros
of~$\zeta(s)$ enter through the spectral decomposition of the
partition function, which sums over the full primary
spectrum, a representation-theoretic datum.

The gap between the bar complex and the spectrum is therefore
the \emph{representation functor}
$\cA \mapsto \cA\text{-}\mathrm{mod}$. For rational VOAs,
the representation category is finite and determined by the
algebra (Huang's theorem); the gap closes. For irrational VOAs
(generic Virasoro, non-rational affine), the representation
category is infinite-dimensional and not algebraically determined
by the OPE data alone. This is why the Koszul--Epstein function
$\varepsilon^c_s$ (which uses the algebra) and the spectral zeta
function $\zeta_\cA(s) = \sum \Delta_i^{-s}$ (which uses the
representations) are different objects that agree only for
lattice VOAs, where the theta function simultaneously encodes
both the algebra and its representation theory.
\end{remark}

% ============================================================
\subsection{The arithmetic--geometric decomposition}%
\label{subsec:arithmetic-geometric-decomposition}%
\index{K\"ahler potential!arithmetic--geometric decomposition}

Proposition~\ref{prop:sewing-spectral-bridge} decomposes the
Laplacian determinant on~$E_\tau$ into two components with
different analytic behaviour.

\begin{proposition}[Arithmetic--geometric decomposition]%
\label{prop:arith-geom-decomposition}%
\ClaimStatusProvedHere%
\index{arithmetic--geometric decomposition|textbf}%
The logarithm of the spectral determinant decomposes as
\begin{equation}\label{eq:arith-geom-decomposition}
 \log\det{}_\zeta'\!(-\Delta_{E_\tau})
 \;=\;
 \underbrace{\log(4\pi^2\,y)}_{\text{geometric}}
 \;+\;
 \underbrace{2\log\lvert\eta(\tau)\rvert^2}_{\text{arithmetic}}\,.
\end{equation}
These components have opposite behaviour under
$\partial\bar\partial$:
\begin{enumerate}[label=\textup{(\roman*)}]
\item The geometric component generates the Weil--Petersson
 K\"ahler form:
 $\omega_{\mathrm{WP}}
 \propto -\tfrac{i}{2}\,\partial\bar\partial\,\log y
 \propto \omega_{\mathrm{hyp}}$,
 a constant multiple of the hyperbolic area form
 on~$\cM_{1,1}$
 \textup{(}Wolpert~\textup{\cite{Wolpert1985}};
 Belavin--Knizhnik~\textup{\cite{BK1986}}\textup{)}.
\item The arithmetic component is
 \emph{pluriharmonic}:
 $\partial\bar\partial\,\log\lvert\eta(\tau)\rvert^2 = 0$,
 because $\log\eta$ is holomorphic and $\log\bar\eta$ is
 anti-holomorphic on~$\mathfrak{H}$.
\end{enumerate}
Consequently, $\omega_{\mathrm{WP}}$ depends only on
the volume~$y$, not on the $\eta$-function.
\end{proposition}

\begin{proof}
The decomposition~\eqref{eq:arith-geom-decomposition} is
immediate from
Proposition~\ref{prop:sewing-spectral-bridge}(ii):
$\det{}_\zeta'(-\Delta) = 4\pi^2\,y\,\lvert\eta\rvert^4$,
and $\log\lvert\eta\rvert^4 = 2\log\lvert\eta\rvert^2$.

For~(i): with $\tau = x + iy$,
$\partial_\tau\log y = i/(2y)$ and
$\partial_\tau\partial_{\bar\tau}\log y = -1/(4y^2)$,
giving
$-\tfrac{i}{2}\,\partial\bar\partial\,\log y
= (1/(4y^2))\,dx\wedge dy \propto \omega_{\mathrm{hyp}}$.

For~(ii): write
$\log\lvert\eta\rvert^2 = \log\eta + \log\bar\eta$.
Since $\eta(\tau)\ne 0$ for $\tau\in\mathfrak{H}$,
$\log\eta$ is holomorphic, hence
$\partial_{\bar\tau}\log\eta = 0$ and
$\partial\bar\partial\log\eta = 0$. Similarly
$\partial_\tau\log\bar\eta = 0$ and
$\partial\bar\partial\log\bar\eta = 0$.
\end{proof}

\begin{remark}[Consequences for Gap~A]%
\label{rem:gap-a-wp-blind}%
\index{Gap A!Weil--Petersson blindness}%
The geometric positivity programme
(\S\ref{sec:geometric-positivity}) seeks a surface
polarization replacing the indefinite Mellin
transform with a positive-definite pairing.
Proposition~\ref{prop:arith-geom-decomposition}
identifies the obstruction: the
K\"ahler structure on~$\cM_{1,1}$ is \emph{arithmetically
blind}. The Hecke eigenforms, the $L$-functions, and the
zeta zeros are all encoded in
$\log\lvert\eta(\tau)\rvert^2$, which lies in the kernel
of~$\partial\bar\partial$. No pairing derived from the
WP K\"ahler structure can detect the arithmetic content of
the spectral determinant. The WP positivity
$\int\lvert f\rvert^2\,\omega_{\mathrm{WP}} \ge 0$
holds trivially for any~$f$, but this is integration
against~$\omega_{\mathrm{hyp}}$ (the standard Petersson
inner product), not a new geometric input.
\end{remark}

\begin{remark}[The MC constraint and the K\"ahler potential]%
\label{rem:mc-kahler-potential}%
\index{Maurer--Cartan equation!K\"ahler potential}%
The MC equation constrains the connected free energy
$F_1^{\mathrm{conn}} = \log\det(1-K)$
(Theorem~\ref{thm:sewing-shadow-intertwining}). By the
bridge~\eqref{eq:sewing-spectral-bridge}, the
arithmetic component of the K\"ahler potential is
$2\log\lvert\eta\rvert^2
= 2\log\lvert\det(1-K)\rvert^2 - \pi y/3
= 4\,\mathrm{Re}(F_1^{\mathrm{conn}}) - \pi y/3$.
The MC equation constrains the K\"ahler \emph{potential}
(specifically its arithmetic part), but this constraint
is invisible to the K\"ahler \emph{metric}
($\partial\bar\partial$ kills it). To exploit the MC
constraint geometrically, one must work at the level of the
potential itself, not its curvature.
\end{remark}

\begin{remark}[Holomorphic factorization and arithmetic blindness]%
\label{rem:holomorphic-arithmetic-blindness}%
\index{holomorphic factorization!arithmetic limitation}%
The arithmetic blindness of~$\omega_{\mathrm{WP}}$
(Remark~\ref{rem:gap-a-wp-blind}) is a consequence of holomorphic
factorization. The bar complex~$B(\cA)$ is built from
holomorphic OPE residues on the Fulton--MacPherson
spaces~$\mathrm{FM}_k(\bC)$; the sewing determinant
$\det(1-K) = e^{-\pi i\tau/12}\,\eta(\tau)$ inherits
this holomorphicity
(Proposition~\ref{prop:sewing-spectral-bridge}(i)).
The decomposition
$\log\lvert\eta\rvert^2 = \log\eta + \log\bar\eta$
splits the arithmetic content into a holomorphic piece
and its conjugate, and $\partial\bar\partial$ annihilates
both. The holomorphic factorization that powers the
bar-cobar machinery and supplies the MC equation's
integrability is the structural reason the resulting
constraints are invisible to the K\"ahler geometry
of~$\cM_{1,1}$.
\end{remark}

% ============================================================
\subsection{The Liouville complexification}%
\label{subsec:liouville-complexification}%
\index{Liouville action!complexified}%
\index{Polyakov formula!complexified}

The target (Remark~\ref{rem:mc-kahler-potential}):
access the arithmetic part of the K\"ahler potential without
passing through~$\partial\bar\partial$. The Polyakov functional
integral~\cite{Polyakov1981, Alvarez1983} represents the
spectral determinant as an integral over conformal
factors~$\sigma$; complexifying~$\sigma$ extends the integral
to a contour whose saddle-point structure may detect the
arithmetic content directly.

\begin{construction}[Complexified Liouville functional integral]%
\label{constr:liouville-complexified}%
\index{Liouville!complexified functional integral}%
The Polyakov--Alvarez conformal anomaly formula expresses
$\log\det{}_\zeta'(-\Delta_g)$ for $g = e^{2\sigma}g_0$ as
\begin{equation}\label{eq:polyakov-conformal-anomaly}
 \log\det{}_\zeta'\!(-\Delta_g)
 \;=\;
 \log\det{}_\zeta'\!(-\Delta_{g_0})
 \;-\;
 \frac{1}{6\pi}\,S_L[\sigma]
 \;+\;
 \log\frac{A_g}{A_{g_0}}\,,
\end{equation}
where $S_L[\sigma] = \int_{E_\tau}(\lvert\nabla_{g_0}\sigma
\rvert^2 + R_{g_0}\,\sigma)\,dA_{g_0}$
is the Liouville action and $A_g = \int e^{2\sigma}\,dA_{g_0}$
is the conformal area. For the flat metric on~$E_\tau$,
$R_{g_0} = 0$ and $S_L[\sigma] = \int\lvert\nabla\sigma
\rvert^2\,dA$, a positive-definite quadratic functional with
unique critical point $\sigma = \mathrm{const}$.

Define the \emph{complexified Liouville functional integral}:
for $\sigma\colon E_\tau\to\bC$,
\begin{equation}\label{eq:complexified-liouville-integral}
 \cZ_L^{\bC}(\tau)
 \;:=\;
 \int_{\mathcal{C}}
 e^{-S_L[\sigma]/(6\pi)}\;D\sigma\,,
\end{equation}
where $\mathcal{C}$ is a contour in the space of complex-valued
conformal factors, chosen to make the integral convergent. On
the real section $\mathcal{C}_{\bR} = \{\sigma\colon E_\tau
\to\bR\}$, the integral is Gaussian (since $S_L$ is quadratic
for the flat torus) and reproduces the standard spectral
determinant.
\end{construction}

\begin{proposition}[Genus-one saddle triviality]%
\label{prop:genus-one-saddle-triviality}%
\ClaimStatusProvedHere%
\index{saddle point!genus-one triviality}%
For the flat torus~$E_\tau$ with $R_{g_0} = 0$, the
Liouville action $S_L[\sigma] = \int\lvert\nabla\sigma
\rvert^2\,dA$ is a positive-definite quadratic functional.
The Euler--Lagrange equation $\Delta\sigma = 0$ on~$E_\tau$
has $\sigma = \mathrm{const}$ as its only periodic solution,
for $\sigma$ real- or complex-valued. The complexified
Liouville integral~\eqref{eq:complexified-liouville-integral}
at genus~$1$ has no nontrivial complex saddle points.
\end{proposition}

\begin{proof}
The Fourier modes $e^{2\pi i(mx + ny)}$ on~$E_\tau$ are
eigenfunctions of~$-\Delta$ with eigenvalues
$\lambda_{m,n} = 4\pi^2\lvert m+n\tau\rvert^2/y^2 > 0$
for $(m,n) \ne (0,0)$. A periodic solution of
$\Delta\sigma = 0$ has all non-constant Fourier coefficients
annihilated by these positive eigenvalues, hence
$\sigma = \mathrm{const}$. The argument is valid for
$\sigma\colon E_\tau \to \bC$: the eigenvalues
$\lambda_{m,n}$ are real and positive, so
$\lambda_{m,n}\,c_{m,n} = 0$ forces $c_{m,n} = 0$
regardless of whether $c_{m,n} \in \bR$ or $c_{m,n} \in \bC$.
\end{proof}

\begin{conjecture}[Complex saddle-point access to scattering
poles]%
\label{conj:complex-saddle-scattering}%
\ClaimStatusConjectured%
\index{saddle point!complex}%
\index{scattering matrix!saddle-point access}%
At higher genus, the complexified Liouville integral
provides a representation whose saddle-point structure is
sensitive to the scattering poles:
\begin{enumerate}[label=\textup{(\roman*)}]
\item The real saddle $\sigma = 0$ \textup{(}constant
 curvature\textup{)} reproduces~$I(s)$ of
 Proposition~\textup{\ref{prop:rs-analytic-continuation}},
 which is holomorphic at $s = \rho/2$
 \textup{(}Proposition~\textup{\ref{prop:scattering-residue}%
 }\textup{)}.
\item Complex saddles $\sigma_*\notin\bR$ with
 $\delta S_L[\sigma_*] = 0$ contribute non-perturbative
 corrections $\sim e^{-S_L[\sigma_*]/(6\pi)}$, with
 $\mathrm{Re}(S_L[\sigma_*])$ controlling the amplitude
 and $\mathrm{Im}(S_L[\sigma_*])$ the phase.
\end{enumerate}
By Proposition~\textup{\ref{prop:genus-one-saddle-triviality}},
the flat torus has no complex saddles; this mechanism requires
genus~$g \ge 2$ \textup{(}where $R_{g_0}\ne 0$ makes $S_L$
nonlinear\textup{)} or a coupling to the moduli~$\tau$ that
extends the integration off the real slice of Teichm\"uller
space.
\end{conjecture}

\begin{remark}[The Selberg zeta function as geometric bridge]%
\label{rem:selberg-trace-arithmetic}%
\index{Selberg zeta function!arithmetic bridge}%
The closest complete spectral--geometric bridge is the
Selberg zeta function
$Z_\Gamma(s) = \prod_{\{\gamma\}}\prod_{k\ge0}
(1 - e^{-\ell(\gamma)(s+k)})$,
where $\{\gamma\}$ runs over primitive closed geodesics of
length~$\ell(\gamma)$ on the hyperbolic surface
$\Gamma\backslash\mathfrak{H}$. The Selberg trace formula
guarantees that $Z_\Gamma$ is entire and its zeros on
$\mathrm{Re}(s) = 1/2$ are the spectral
eigenvalues~$s_j$. The sewing determinant $\det(1-K)$
is an analogous infinite product: sewing eigenvalues
$q^n$ replace geodesic contributions
$e^{-\ell(\gamma)(s+k)}$. However, the geodesic product
converges only for $\mathrm{Re}(s)>1$; continuation to
the critical strip uses the functional equation
of~$Z_\Gamma$, whose proof requires the full Selberg trace
formula, input not supplied by the sewing formalism. The
sewing product $\prod(1-q^n)$ converges for
$\lvert q\rvert < 1$; its continuation to
$\lvert q\rvert\ge 1$ requires the modular transformation
$\tau\mapsto -1/\tau$, which is likewise external to the
sewing framework.
\end{remark}

% ============================================================
\subsection{Assessment}%
\label{subsec:polyakov-assessment}%
\index{Polyakov complexification!assessment}

\begin{center}
\small
\renewcommand{\arraystretch}{1.3}
\begin{tabular}{clll}
\toprule
& \textbf{Component} & \textbf{Proof surface} & \textbf{Reference} \\
\midrule
1 & Sewing--spectral bridge
 & \ClaimStatusProvedHere
 & Prop.~\ref{prop:sewing-spectral-bridge} \\
2 & RS meromorphic continuation
 & \ClaimStatusProvedHere
 & Prop.~\ref{prop:rs-analytic-continuation} \\
3 & Holomorphy at scattering poles
 & \ClaimStatusProvedHere
 & Prop.~\ref{prop:scattering-residue} \\
4 & Arith.--geom.\ decomposition
 & \ClaimStatusProvedHere
 & Prop.~\ref{prop:arith-geom-decomposition} \\
5 & WP arithmetically blind
 & \ClaimStatusProvedHere
 & Rem.~\ref{rem:gap-a-wp-blind} \\
6 & Genus-$1$ saddle triviality
 & \ClaimStatusProvedHere
 & Prop.~\ref{prop:genus-one-saddle-triviality} \\
7 & Complexified Liouville integral
 & Construction
 & Constr.~\ref{constr:liouville-complexified} \\
8 & Complex saddle $\to$ scattering
 & \ClaimStatusHeuristic
 & Conj.~\ref{conj:complex-saddle-scattering} \\
9 & Selberg zeta bridge
 & Structural analogy
 & Rem.~\ref{rem:selberg-trace-arithmetic} \\
10 & Nonvanishing at $\rho_1/2$
 & \ClaimStatusProvedHere
 & Cor.~\ref{cor:first-scattering-pole} \\
\bottomrule
\end{tabular}
\end{center}

\begin{remark}[Three new gaps]%
\label{rem:three-new-gaps}%
\index{Polyakov complexification!new gaps}%
The complexification programme introduces three structural
gaps beyond the original four of~\S\ref{sec:four-gaps}:
\begin{itemize}[nosep]
\item \emph{Gap P1: Genus-one triviality.}
 On the flat torus, $S_L[\sigma] = \int\lvert\nabla\sigma
 \rvert^2$ is quadratic, so the functional integral is
 Gaussian and the only saddle is $\sigma = \mathrm{const}$.
 Complex saddles require either higher genus (where the
 Liouville equation
 $\Delta\sigma + R\,e^{2\sigma} = 0$ is nonlinear) or a
 coupling to the moduli parameter~$\tau$ that extends the
 integration to the complexified Teichm\"uller space.
 The arithmetic programme, as formulated, lives at
 genus~$1$, where this mechanism is absent.

\item \emph{Gap P2: Unfolding erasure.}
 Rankin--Selberg unfolding replaces
 $\int f\cdot E_s\,d\mu$ by the Mellin transform
 of~$a_0(y)$, erasing the scattering matrix from the
 integral
 (Remark~\ref{rem:structural-obstruction-revisited}).
 Any representation that retains the Eisenstein
 series~$E_s$ without unfolding must control the integral
 over the fundamental domain directly, but the standard
 estimates break down at the scattering poles, where
 $E_s(\tau)$ ceases to be square-integrable.

\item \emph{Gap P3: Arithmetic pluriharmonicity.}
 The MC equation constrains the arithmetic part
 $2\log\lvert\eta\rvert^2$ of the K\"ahler potential
 (Remark~\ref{rem:mc-kahler-potential}), but this part is
 pluriharmonic
 (Proposition~\ref{prop:arith-geom-decomposition}).
 The constraint operates on a function in the kernel
 of~$\partial\bar\partial$: no \emph{metric} consequence
 on~$\cM_{1,1}$ follows from the MC equation via the WP
 structure.
\end{itemize}
\end{remark}

\begin{remark}[Honest comparison]%
\label{rem:honest-comparison}%
\index{Polyakov complexification!honest comparison}%
The Polyakov programme complements the spectral approach but
does not overcome the obstruction of
Remark~\ref{rem:structural-obstruction}.
The three new gaps are structural:
\begin{itemize}[nosep]
\item Gap~P1 is a feature of genus~$1$ (quadratic Liouville
 action), not a technical limitation;
\item Gap~P2 identifies Rankin--Selberg unfolding as both
 the source of meromorphic continuation \emph{and} the
 mechanism that erases the scattering data;
\item Gap~P3 shows that the MC constraint, while governing
 the K\"ahler potential, is invisible to the K\"ahler
 metric.
\end{itemize}
The programme identifies a precise target: find a
representation of the partition function that is \emph{not}
obtained by RS unfolding and whose analytic structure
in~$s$ retains the scattering poles. Whether higher-genus
Polyakov integrals, complexified Teichm\"uller space, or
entirely different methods can provide such a representation
remains open.
\end{remark}

\begin{corollary}[Nonvanishing at the first scattering pole]%
\label{cor:first-scattering-pole}%
\ClaimStatusProvedElsewhere%
\index{scattering matrix!nonvanishing at first pole}%
The sewing--Selberg formula
(Theorem~\ref{thm:sewing-selberg-formula}) gives
\[
 I_{\mathrm{sew}}(s)
 \;=\;
 \int_{\cM_{1,1}}
 \log\det(1-K(\tau))\cdot E_s(\tau)\,d\mu(\tau)
 \;=\;
 -2(2\pi)^{-(s-1)}\,\Gamma(s{-}1)\,\zeta(s{-}1)\,\zeta(s)
\]
for all $s\in\bC$. At the first scattering pole
$s = \rho_1/2 = 1/4 + 7.067\ldots\,i$
\textup{(}where $\rho_1 = 1/2 + 14.134\ldots\,i$\textup{)},
none of the four factors on the right vanishes:
$\Gamma$ has no zeros\textup;
$\zeta(\rho_1/2 - 1) = \zeta(-3/4 + 7.067\ldots\,i)$ lies
outside the critical strip\textup;
$\zeta(\rho_1/2) = \zeta(1/4 + 7.067\ldots\,i)$ lies
in the classical zero-free region.
Hence $I_{\mathrm{sew}}(\rho_1/2) \ne 0$: the RS integral is
holomorphic \emph{and nonvanishing} at the first scattering
pole.
\end{corollary}

\begin{problem}[MC constraint at scattering poles]%
\label{prob:polyakov-target}%
\index{Polyakov complexification!target computation}%
For a chirally Koszul algebra~$\cA$ with HS-sewing and
nontrivial cusp forms (e.g., a rank-$24$ lattice VOA
$V_\Lambda$ whose theta function
$\Theta_\Lambda \in M_{12}$ has nontrivial projection to
the cusp-form space $S_{12} = \bC\cdot\Delta$
\textup{(}the Ramanujan discriminant\textup{)}), does the MC
constraint
$D\cdot\Theta_\cA + \tfrac12[\Theta_\cA, \Theta_\cA] = 0$
impose any relation on $I(\rho_1/2)$ beyond the
Rankin--Selberg holomorphy
(Proposition~\ref{prop:scattering-residue})?
A positive answer would indicate that the MC equation
carries zero-location information not visible to the
spectral decomposition; the corollary above shows that
the universal (algebra-independent) RS integral carries
none.
\end{problem}

% ============================================================
\section{The analytic continuation programme}%
\label{sec:analytic-continuation-programme}%
\index{analytic continuation programme|textbf}

The MC element $\Theta_\cA$ produces two kinds of arithmetic
output.
The \emph{algebraic output} (shadow obstruction tower and Hecke
decomposition, with CPS automorphy only under
Corollary~\ref{cor:moment-automorphy}) is separated from the analytic
output
\textup{(}\S\S\ref{sec:sewing-RS-bridge}--\ref{sec:operadic-transfer-programme}\textup{)}.
The \emph{analytic output} (Ramanujan bound,
zero-location of the Epstein zeta) requires external input
(Remark~\ref{rem:algebraic-analytic-divide}).

This section isolates the mechanism separating
these two outputs, identifies the single remaining algebraic
hypothesis whose resolution closes the gap for the
Ramanujan bound, and proves that the zero-location gap
is \emph{structural}: a theorem about what algebraic
constraints cannot determine.

\subsection{The scattering coupling factorization}%
\label{subsec:scattering-coupling-factorization}

At the scattering poles $s = (1{+}\rho)/2$ where
$\zeta(\rho) = 0$, the RS integral $I(s;\cA)$
is holomorphic (Proposition~\ref{prop:scattering-residue}).
Its value there is a well-defined invariant of~$\cA$.

\begin{definition}[Scattering coupling]%
\label{def:scattering-coupling}%
\index{scattering coupling|textbf}%
For a chirally Koszul algebra~$\cA$ with HS-sewing and a
nontrivial zero $\rho$ of the Riemann zeta function, the
\emph{scattering coupling} of~$\cA$ at~$\rho$ is the value
\begin{equation}\label{eq:scattering-coupling}
 \cE_\rho(\cA)
 \;:=\;
 I\!\left(\tfrac{1{+}\rho}{2};\,\cA\right)
 \;=\;
 \int_0^\infty
 a_0(y;\cA)\;y^{(\rho-1)/2}\,dy,
\end{equation}
where the second equality is by Rankin--Selberg unfolding and
the Mellin integral converges by HS-sewing
\textup{(}exponential decay at~$\infty$, polynomial growth
at~$0$\textup{)}.
\end{definition}

\begin{theorem}[Scattering coupling factorization;
\ClaimStatusProvedHere]%
\label{thm:scattering-coupling-factorization}%
\index{scattering coupling!factorization}%
The scattering coupling factors into a universal part and an
algebra-dependent part:
\begin{equation}\label{eq:scattering-factorization}
 \cE_\rho(\cA)
 \;=\;
 \Gamma_\infty\!\left(\tfrac{1{+}\rho}{2}\right)
 \;\cdot\;
 \varepsilon^c_{(1+\rho)/2}(\cA),
\end{equation}
where $\Gamma_\infty(s) = (2\pi)^{-s}\,\Gamma(s)$ is the
archimedean factor \textup{(}independent of~$\cA$\textup{)}
and $\varepsilon^c_s(\cA) = \sum_{\Delta\in\mathcal{S}}
c_\Delta\,(2\Delta)^{-s}$ is the constrained Epstein zeta
\textup{(}where $\mathcal{S}$ is the primary spectrum counted
with multiplicity $c_\Delta$\textup{)}
\textup{(}dependent on~$\cA$ through its primary
spectrum\textup{)}.

The algebra-dependent factor is controlled by the shadow obstruction tower:
\begin{enumerate}[label=\textup{(\roman*)}]
\item The Epstein zeta
 $\varepsilon^c_s(\cA)$ is the Mellin transform of the
 spectral measure $\rho_\cA = \sum_i c_i\,\delta(\Delta_i)$.
\item The spectral measure is uniquely determined by the
 shadow moments $\{S_r(\cA)\}_{r\ge 2}$ when the Carleman
 condition holds \textup{(}Proposition~\textup{\ref{prop:carleman-virasoro}}\textup{)}.
\item The MC recursion
 \textup{(}Theorem~\textup{\ref{thm:mc-recursion-moment}}\textup{)}
 constrains $S_{r+1}$ from $\{S_j\}_{j \le r}$.
\end{enumerate}
Consequently, the MC element~$\Theta_\cA$ determines
$\cE_\rho(\cA)$ for \emph{every} nontrivial zero~$\rho$.
\end{theorem}

\begin{proof}
The factorization~\eqref{eq:scattering-factorization} is the
standard Mellin inversion applied to the unfolded integral.
By definition~\eqref{eq:scattering-coupling},
$\cE_\rho(\cA) = \int_0^\infty a_0(y;\cA)\,
y^{(\rho-1)/2}\,dy$.
The zeroth Fourier mode of the primary-counting function
decomposes as
$a_0(y;\cA) = \sum_{\Delta\in\mathcal{S}}
c_\Delta\,e^{-4\pi\Delta\,y}$ (modulo descendant corrections
absorbed by Mellin completion).
The Mellin transform gives
$(4\pi\Delta)^{-(1+\rho)/2}\,\Gamma((1{+}\rho)/2)$ per primary.
Extracting the archimedean factor yields
$\Gamma_\infty((1{+}\rho)/2) \cdot \varepsilon^c_{(1+\rho)/2}$.
Items~(i)--(iii) follow from
Propositions~\ref{prop:carleman-virasoro},
\ref{prop:shadow-symmetric-power},
and Theorem~\ref{thm:mc-recursion-moment}.
\end{proof}

\begin{remark}[What the factorization proves and does not prove]%
\label{rem:factorization-scope}%
\index{scattering coupling!scope of factorization}%
The factorization~\eqref{eq:scattering-factorization} shows
that the MC equation determines the \emph{coupling strength}
$\cE_\rho(\cA)$ of each chiral algebra to each scattering
pole. It does \emph{not} determine the \emph{location} of
the poles: the zero~$\rho$ enters only through the evaluation
point $s = (1{+}\rho)/2$ of the constrained Epstein zeta,
which is a holomorphic function of~$s$. The locations of the
zeros of~$\zeta$ are properties of the \emph{universal}
scattering matrix $\varphi(s) = \Lambda(1{-}s)/\Lambda(s)$,
which is independent of~$\cA$.

To extract zero-location information, one would need to show
that the \emph{family} of values
$\{\cE_\rho(\cA)\}_\cA$ over some moduli space of
chirally Koszul algebras constrains~$\rho$. This is
precisely the content of Problem~\ref{prob:polyakov-target}.
\end{remark}

\subsection{The refined structural obstruction}%
\label{subsec:refined-structural-obstruction}

\begin{theorem}[Structural separation of algebraic and
arithmetic content; \ClaimStatusConditional]%
\label{thm:structural-separation}%
\index{structural obstruction!refined form|textbf}%
\index{algebraic--analytic divide!structural theorem}%
Let $\cA$ range over the class of chirally Koszul algebras
with HS-sewing. The following three statements hold
simultaneously:
\begin{enumerate}[label=\textup{(\roman*)}]
\item \textup{(Algebraic completeness, conditional on the CPS
vertical-strip estimates in Corollary~\textup{\ref{cor:moment-automorphy}}.)}\quad
 The MC element $\Theta_\cA$ determines the scattering
 coupling $\cE_\rho(\cA)$ for every nontrivial zero~$\rho$
 \textup{(}Theorem~\textup{\ref{thm:scattering-coupling-factorization}}\textup{)},
 the CPS automorphy of all moment $L$-functions
 \textup{(}Corollary~\textup{\ref{cor:moment-automorphy}}\textup{)},
 and the Hecke decomposition of the partition function
 for lattice VOAs
 \textup{(}Theorem~\textup{\ref{thm:spectral-continuation-bridge}}\textup{)}.
\item \textup{(Unfolding erasure.)}\quad
 The meromorphic continuation of the RS integral $I(s;\cA)$
 is holomorphic at all scattering poles
 \textup{(}Proposition~\textup{\ref{prop:scattering-residue}}\textup{)},
 so the continued integral carries no zero-location
 information. This is a feature of the Rankin--Selberg
 method, not a deficiency of the MC input.
\item \textup{(Arithmetic blindness of the K\"ahler
 structure.)}\quad
 The MC equation constrains the arithmetic part
 $2\log\lvert\eta\rvert^2$ of the K\"ahler potential on
 $\cM_{1,1}$, but this part is pluriharmonic
 \textup{(}Proposition~\textup{\ref{prop:arith-geom-decomposition}}\textup{)},
 hence invisible to the Weil--Petersson metric.
\end{enumerate}
Together, these prove that the MC element exhausts the
algebraic content of the arithmetic interface at genus~$1$
\textup{(}every algebraic consequence has been extracted\textup{)},
and that the remaining zero-location question requires
either higher-genus or genuinely analytic input external to the
MC equation.
\end{theorem}

\begin{proof}
Item~(i) collects Theorem~\ref{thm:scattering-coupling-factorization},
Corollary~\ref{cor:moment-automorphy}, and
Theorem~\ref{thm:spectral-continuation-bridge}; its automorphy clause
therefore has exactly the conditional strength of
Corollary~\ref{cor:moment-automorphy}.
Item~(ii) is Proposition~\ref{prop:scattering-residue} and
Remark~\ref{rem:structural-obstruction-revisited}.
Item~(iii) is Proposition~\ref{prop:arith-geom-decomposition}
and Remark~\ref{rem:gap-a-wp-blind}.
These three outputs ($\cE_\rho(\cA)$, $M_r(s;\cA)$, and the
Hecke decomposition) exhaust the algebraic content of the
genus-$1$ Rankin--Selberg method applied to the MC-constrained
partition function: no further invariant at degree~$\leq 4$ is
independent of these three.
\end{proof}

\begin{remark}[The genus-$2$ escape route]%
\label{rem:genus2-escape-route}%
\index{structural obstruction!genus-2 escape}%
All three items of Theorem~\ref{thm:structural-separation}
rest on the genus-$1$ sewing operator being
\emph{diagonal in the Fock basis}, the structural fact
underlying the sewing--Hecke collapse
(Theorem~\ref{thm:sewing-hecke-reciprocity}).
At genus~$2$, this diagonality breaks.
The genus-$2$ sewing kernel is a $2 \times 2$ block operator
on $A^2(D)^{\oplus 2}$ with off-diagonal entries coupling
the two handles through the internal seam
(Vol~II, genus-$2$ Fredholm determinant theorem).
The off-diagonal coupling is controlled by the
off-diagonal period~$z$ of the Siegel matrix
$\Omega = \bigl(\begin{smallmatrix}\tau_1 & z \\
z & \tau_2\end{smallmatrix}\bigr)
\in \mathfrak{H}_2$, and vanishes only in the separating
degeneration $z \to 0$. Three consequences follow:
\begin{enumerate}[label=\textup{(\roman*)}]
\item The genus-$2$ connected free energy depends on
 three modular parameters $(\tau_1, z, \tau_2)$, not one.
\item Its spectral decomposition involves
 $\mathrm{Sp}_4(\bZ)$-Eisenstein series, which do not
 factor as products of elliptic Eisenstein series.
\item The Rankin--Selberg unfolding on
 $\mathrm{Sp}_4(\bZ) \backslash \mathfrak{H}_2$
 produces genuinely multi-variable $L$-functions (spin
 and standard).
\end{enumerate}
The sewing--Hecke collapse is therefore a \emph{genus-$1$
phenomenon}: the diagonality that forces it does not persist
at genus $\geq 2$. See~\S\ref{sec:genus2-arithmetic-frontier}
for the programme this opens.
\end{remark}

\subsection{The Hecke defect as obstruction class}%
\label{subsec:hecke-defect-obstruction}%
\index{Hecke defect!obstruction class|textbf}

The single algebraic hypothesis separating the proved core
from the Ramanujan bound is prime-locality
(Conjecture~\ref{conj:prime-locality-transfer}).
The Hecke defect of
Remark~\ref{rem:prime-locality-obstruction} admits a natural
interpretation as a cohomological obstruction.

\begin{definition}[Hecke defect class]%
\label{def:hecke-defect-class}%
\index{Hecke defect!class|textbf}%
For a prime~$p$ and a chirally Koszul algebra~$\cA$ with
HS-sewing, the \emph{Hecke defect class} is the element
\begin{equation}\label{eq:hecke-defect-class}
 [\delta_p(\cA)]
 \;\in\;
 H^1\!\bigl(\gAmod,\,
 \mathrm{ad}_{T_p}\bigr),
\end{equation}
where $H^1$ denotes the first cohomology of the
Chevalley--Eilenberg complex of the dg Lie algebra $\gAmod$
with coefficients in the adjoint module twisted by $T_p$,
$\mathrm{ad}_{T_p}$ denotes the corresponding coefficient module
induced by the $T_p$-adjoint action, and the class is
represented by the cocycle
$\delta_p(\cA) := [d_{\mathrm{sew}}, T_p] \in
\mathrm{End}(\gAmod)$
of Remark~\textup{\ref{rem:prime-locality-obstruction}}.
\end{definition}

\begin{proposition}[Equivalent characterizations; \ClaimStatusConditional]%
\label{prop:hecke-defect-equivalences}%
\index{Hecke defect!equivalent characterizations}%
The following are equivalent:
\begin{enumerate}[label=\textup{(\roman*)}]
\item $[\delta_p(\cA)] = 0$ in
 $H^1(\gAmod, \mathrm{ad}_{T_p})$ for all primes~$p$.
\item The Hecke operators~$T_p$ extend to endomorphisms
 of~$\gAmod$ that commute with~$D_\cA$ up to
 $D_\cA$-exact correction.
\item Prime-locality holds: the shadow coefficients
 $S_r^{(p)} \propto p_{r-1}(\alpha_p,\beta_p)$ at
 each prime~$p$.
\end{enumerate}
The equivalence of \textup{(i)--(iii)} is proved here. Conditions
\textup{(i)--(iii)} are conditionally equivalent to\textup{:}
\begin{enumerate}[label=\textup{(\roman*)},resume]
\item The moment $L$-function $M_r(s) = L(s,\pi_r)$
 satisfies $\pi_r \cong \operatorname{Sym}^{r-1}(f)$
 for all~$r$
 \textup{(}assuming
 Conjecture~\textup{\ref{conj:operadic-rankin-selberg-main}}\textup{)}.
\end{enumerate}
\end{proposition}

\begin{proof}
$\mathrm{(i)} \Leftrightarrow \mathrm{(ii)}$: by definition
of the defect class. $T_p$ is an endomorphism of the
underlying graded vector space of~$\gAmod$ (the Hecke operators
act on any function of~$\tau$). The defect
$\delta_p = [D_\cA, T_p]$ vanishes in cohomology iff
$T_p$ can be corrected by a gauge transformation to commute
with~$D_\cA$; this corrected~$T_p$ is then a chain
endomorphism.
$\mathrm{(ii)} \Rightarrow \mathrm{(iii)}$:
if $T_p$ is a chain endomorphism of~$\gAmod$, then $\Theta_\cA$
lies in a $T_p$-eigenspace (or a direct sum of eigenspaces);
projecting to degree~$r$ gives shadow coefficients in the
$T_p$-eigenspace, hence multiplicative in~$p$.
For the conditional equivalence with~(iv):
$\mathrm{(iii)} \Rightarrow \mathrm{(iv)}$
(assuming Conjecture~\ref{conj:operadic-rankin-selberg-main}):
by prime-locality $+$ CPS converse theorem $+$ strong multiplicity one
(Remark~\ref{rem:complete-chain}).
$\mathrm{(iv)} \Rightarrow \mathrm{(ii)}$:
if $\pi_r = \operatorname{Sym}^{r-1}(f)$, the Hecke eigenvalues
at degree~$r$ are the $r$-th power sums of the Satake parameters,
which is equivalent to $T_p$ acting diagonally on each degree
component of~$\gAmod$.
\end{proof}

\begin{proposition}[Hecke defect vanishes for lattice VOAs;
\ClaimStatusProvedHere]%
\label{prop:hecke-defect-lattice}%
\index{Hecke defect!vanishing for lattices}%
For a lattice VOA $V_\Lambda$, $[\delta_p(V_\Lambda)] = 0$
for all primes~$p$.
\end{proposition}

\begin{proof}
For $V_\Lambda$, the sewing differential $d_{\mathrm{sew}}$
factors through the theta function
$\Theta_\Lambda \in M_{r/2}(\Gamma)$
(Theorem~\ref{thm:spectral-continuation-bridge}(i)):
every graph amplitude $\ell_\Gamma$ lies in the image of
the Hecke algebra. The Hecke operator~$T_p$ acts on~$M_{r/2}$
by a ring homomorphism that commutes with all multilinear
operations constructed from Fourier coefficients. Since
$d_{\mathrm{sew}}$ is such an operation,
$[d_{\mathrm{sew}}, T_p] = 0$ exactly at chain level.
\end{proof}

\begin{remark}[Explicit Hecke defect at leading $1/c$]%
\label{rem:hecke-defect-leading}%
\index{Hecke defect!Virasoro leading order}%
For the Virasoro algebra at leading order in~$1/c$,
the spectral measure is a single atom
$\rho = \delta(\lambda + 6/c)$,
where $6/c$ arises as the leading-order spectral gap
$\Delta_1 - c/24 \sim 6/c$ from the large-$c$ Virasoro
vacuum module
(Proposition~\ref{prop:modularity-constraint-atoms}).
A single atom is trivially prime-local (the value at every
prime is the same), so $[\delta_p] = 0$ at this order.
The first nontrivial contribution to $[\delta_p(\mathrm{Vir})]$
arises at order $1/c^2$, where the spectral measure acquires
a second atom and the prime-by-prime decomposition becomes
nontrivial. Computing $[\delta_p]$ at subleading order
for specific primes $p = 2, 3, 5$ at the Ising model
($c = 1/2$), the Lee--Yang model ($c = -22/5$), and the
three-state Potts model ($c = 4/5$) is a concrete target
for the finite falsification criterion of
Remark~\ref{rem:prime-locality-routes}.
\end{remark}

\begin{definition}[Hecke defect on the shadow obstruction tower]%
\label{def:hecke-defect-shadow}%
\index{Hecke defect!on shadow tower|textbf}%
\index{Hecke operator!action on shadow coefficients}
For a prime~$p$ and a chirally Koszul algebra~$\cA$ with
HS-sewing, the \emph{Hecke defect on the shadow obstruction tower}
at degree~$r$ is
\begin{equation}\label{eq:hecke-defect-degree}
 \delta_p^{(r)}(\cA)
 \;:=\;
 T_p(S_r) \;-\;
 \operatorname{MC}_r\!\bigl(T_p(S_j),\, T_p(P_0)\bigr),
\end{equation}
where $\operatorname{MC}_r(S_j, P_0) := -(1/(2r))
\sum_{j+k=r+2,\,
j,k \ge 2} j\,k\,S_j\,S_k\,P_0$ is the scalar MC recursion
\textup{(}which yields $S_r$ by
Theorem~\textup{\ref{thm:mc-recursion-moment}}\textup{)},
$T_p$ acts on modular functions of~$\tau$ by the
classical Hecke correspondence
$T_p(f)(\tau) = p^{-1}\sum_{b=0}^{p-1} f((\tau{+}b)/p)
+ f(p\tau)$ \textup{(}for weight $0$\textup{)}, $S_r$ is
the shadow coefficient at degree~$r$, and $P_0$ is the
scalar contraction in the uniform-weight OPE lane. It is not
the genus-$1$ propagator $P(e;\tau)$ on~$E_\tau$; the latter
enters the amplitude-level graph sum separately. The defect is
well-defined
for $r \ge 5$ \textup{(}at which degrees
$S_r$ is determined by the MC recursion from
$S_2, S_3, S_4$\textup{)}. It measures the failure
of the MC recursion to commute with the Hecke operator:
$\delta_p^{(r)} = 0$ for all~$r \ge 5$ if and only if
$[d_{\mathrm{sew}}, T_p]$ acts trivially on the
recursion-constrained shadow obstruction tower.
\end{definition}

\begin{proposition}[Hecke defect for standard families; \ClaimStatusConditional]%
\label{prop:hecke-defect-families}%
\index{Hecke defect!Heisenberg}%
\index{Hecke defect!Virasoro}%
\index{Hecke defect!$\cW_3$}
\begin{enumerate}[label=\textup{(\roman*)}]
\item \emph{Heisenberg:}
 $\delta_p^{(r)}(\cH) = 0$ for all primes~$p$ and all
 degrees~$r \ge 5$. The Heisenberg shadow obstruction tower terminates at
 $r = 2$ with $S_2 = \kappa = 1$, $S_r = 0$ for
 $r \ge 3$. The Hecke operator acts as the identity on
 constants: $T_p(\kappa) = \kappa$.

\item \emph{Virasoro:}
 the shadow coefficients $S_r$ are rational functions
 of~$c$ alone \textup{(}independent of~$\tau$\textup{)},
 hence $T_p(S_r) = S_r$ for all~$p$. The scalar contraction
 $P_0 = 2/c$ is also $\tau$-independent. Therefore
 \begin{equation}\label{eq:hecke-defect-virasoro-vanishes}
 \delta_p^{(r)}(\mathrm{Vir}_c) \;=\; 0
 \quad\text{for all primes $p$ and all degrees $r \ge 5$.}
 \end{equation}
 This vanishing is \emph{not} dynamical: it follows from the
 $\tau$-independence of the shadow obstruction tower on the uniform-weight lane
 (Definition~\ref{def:scalar-lane}).
 The genus-$1$ amplitudes $\mathrm{Sh}_r^{(1)}(\tau)$ are
 $\tau$-dependent, but the MC recursion among shadow
 coefficients operates at the level of $\tau$-independent
 OPE data. The Hecke operator acts nontrivially only on
 the $\tau$-dependent amplitudes, where it respects the MC
 structure by Theorem~\textup{\ref{thm:route-c-propagation}}.

\item \emph{$\cW_3$:}
 the T-line shadow coefficients coincide with Virasoro
 \textup{(}hence $\delta_p^{(r),T} = 0$\textup{)}. The
 W-line has $\kappa_W = c/3$, $\alpha_W = 0$ \textup{(}by
 $\bZ_2$ parity\textup{)}, and $S_4^W$ is a rational function
 of~$c$ \textup{(}hence also $\tau$-independent\textup{)}.
 Therefore $\delta_p^{(r),W}(\cW_3) = 0$ for all~$p$ and~$r$
 by the same mechanism as~\textup{(ii)}.
 The mixed T-W cross-channel defect, measuring the failure
 of Hecke equivariance in the propagator-variance
 sector \textup{(}$\delta_{\mathrm{mix}}$ of
 Theorem~\textup{\ref{prop:propagator-variance}}\textup{)},
 vanishes because the scalar mixing polynomial
 $P_0(\cW_3) = 25c^2 + 100c - 428$ is $\tau$-independent.
\end{enumerate}
\end{proposition}

\begin{proof}
Part~(i): the Heisenberg shadow obstruction tower is
$S_2 = \kappa = 1$, $S_r = 0$ for $r \ge 3$.
For $r \ge 5$, the MC recursion gives
$\operatorname{MC}_r(S_j, P_0) = -(1/(2r))
\sum_{j+k=r+2} j\,k\,S_j\,S_k\,P_0 = 0$
because every summand contains $S_j$ or $S_k$ with
$j \ge 3$ or $k \ge 3$, which vanishes. Hence
$S_r = 0 = \operatorname{MC}_r(S_j, P_0)$.
The Hecke operator acts as the identity on constants,
so $\delta_p^{(r)} = T_p(S_r) - \operatorname{MC}_r
(T_p(S_j), T_p(P_0)) = 0 - 0 = 0$.

Part~(ii): the Virasoro shadow coefficients are determined
by the recursive formula of
\S\ref{sec:operadic-rankin-selberg} with inputs
$\kappa = c/2$, $\alpha = 2$,
$S_4 = 10/(c(5c{+}22))$ (all rational functions of~$c$
alone). The scalar contraction $P_0 = 2/c$ is the inverse of
the curvature. At every stage, $S_{r+1}$ is a rational
function of~$c$ with no $\tau$-dependence. Since
$T_p(f) = f$ for any $\tau$-independent function~$f$,
$T_p(S_r) = S_r$ and $T_p(P_0) = P_0$, so
$\delta_p^{(r)} = T_p(S_r) -
\operatorname{MC}_r(T_p(S_j), T_p(P_0))
= S_r - \operatorname{MC}_r(S_j, P_0) = S_r - S_r = 0$.

Part~(iii): the T-line coincides with Virasoro by
construction. The W-line has all inputs $\tau$-independent
by the explicit formulas of Chapter~\ref{chap:w-algebras}.
The mixed channel involves $\delta_{\mathrm{mix}}$, which
is a polynomial in $\kappa_T, \kappa_W, S_4^T, S_4^W$ and
the propagator, all $\tau$-independent.
\end{proof}

\begin{remark}[The Hecke defect at genus~$1$ versus
the shadow level]%
\label{rem:hecke-defect-two-levels}%
\index{Hecke defect!genus-1 amplitudes vs shadow coefficients}
The Hecke defect operates at two levels:
\begin{enumerate}[label=\textup{(\roman*)}]
\item \emph{Shadow coefficient level}: $\delta_p^{(r)}$ of
 Definition~\ref{def:hecke-defect-shadow} measures
 whether $T_p$ commutes with the MC recursion on
 $\tau$-independent OPE data. For algebraic families,
 this vanishes identically
 (Proposition~\ref{prop:hecke-defect-families}).

\item \emph{Genus-$1$ amplitude level}: the shadow
 amplitude $\mathrm{Sh}_r^{(1)}(\tau) =
 \sum_\Gamma |\mathrm{Aut}(\Gamma)|^{-1}\,\ell_\Gamma(\tau)$
 is a function of~$\tau$. The Hecke defect at this level,
 $T_p(\mathrm{Sh}_r^{(1)}) -
 \mathrm{Sh}_r^{(1)}|_{T_p}$,
 measures whether the graph-sum amplitudes respect the
 Hecke decomposition on~$\cM_{1,1}$. For lattice VOAs
 this vanishes
 (Proposition~\ref{prop:hecke-defect-lattice});
 for non-lattice algebras it is the content of
 Conjecture~\ref{conj:prime-locality-transfer}.
\end{enumerate}
Theorem~\ref{thm:route-c-propagation} bridges these
two levels: vanishing at the shadow coefficient level
\textup{(}level~\textup{(i)}\textup{)}, combined with
character-level Hecke equivariance, implies vanishing at
the amplitude level \textup{(}level~\textup{(ii)}\textup{)}.
\end{remark}

\subsection{The rigidity inheritance theorem}%
\label{subsec:rigidity-inheritance}%
\index{rigidity inheritance|textbf}

\begin{theorem}[Rigidity inheritance; \ClaimStatusConditional]%
\label{thm:rigidity-inheritance}%
\index{prime-locality!from rigidity inheritance}%
\index{Carleman condition!rigidity inheritance}%
Let $\cA$ be a chirally Koszul algebra with HS-sewing.
Assume:
\begin{enumerate}[label=\textup{(\alph*)}]
\item The \emph{Carleman condition}
 \textup{(}Proposition~\textup{\ref{prop:carleman-virasoro}}\textup{):}
 the spectral measure $\rho_\cA$ is uniquely determined by
 the shadow moments $\{S_r\}_{r \ge 2}$.
\item The \emph{rigidity threshold}
 \textup{(}Proposition~\textup{\ref{prop:rigidity-threshold}}\textup{):}
 $\delta(m,r) > 0$ for the number $m$ of spectral atoms.
\item The \emph{Hecke defect vanishing:}
 $[\delta_p(\cA)] = 0$ in $H^1(\gAmod,
 \mathrm{ad}_{T_p})$ for all primes~$p$.
\end{enumerate}
Then prime-locality holds. For rational central charge,
the Ramanujan bound follows by the twelve-station proof
\textup{(}Theorem~\textup{\ref{thm:non-lattice-ramanujan}}\textup{)}.
For irrational central charge, it follows conditionally from
Conjecture~\textup{\ref{conj:irrational-ramanujan}}.
\end{theorem}

\begin{proof}
By~(a), the spectral measure $\rho_\cA$ is uniquely determined.
By~(c) and Proposition~\ref{prop:hecke-defect-equivalences}%
(ii), the Hecke operators~$T_p$ act as chain endomorphisms
on~$\gAmod$ (up to gauge equivalence).

The Hecke operator $T_p$ acts on the Fourier coefficients of
the MC element $\Theta_\cA$, which are modular forms. The
shadow projections $\operatorname{Sh}_r(\Theta_\cA)$ are
polynomial functions of these Fourier coefficients. Since
$T_p$ acts as a ring homomorphism on modular forms,
$T_p(\operatorname{Sh}_r(\Theta_\cA))
= \operatorname{Sh}_r(T_p(\Theta_\cA))$.
By the rigidity hypothesis~(b), the MC equation has a unique
solution up to gauge, so $T_p(\Theta_\cA)$ is gauge-equivalent
to~$\Theta_\cA$. The shadow projections are cohomological
invariants, hence gauge-invariant by definition, giving
$T_p(\operatorname{Sh}_r) = \operatorname{Sh}_r$.
The Hecke eigenvalues arise from the decomposition of
$\operatorname{Sh}_r$ into Hecke eigenspaces in the ambient
modular form ring, not from scalar eigenvalues on a single
element: each shadow $\operatorname{Sh}_r$ decomposes as
$\operatorname{Sh}_r = \sum_j c_j f_j$ with $T_p(f_j)
= \lambda_{p,j} f_j$, and the inherited multiplicativity
is that of each eigencomponent.

By~(b), the rigidity threshold is reached: the overdetermined
MC system has a unique solution for the spectral atoms, and
the Hecke eigenvalue structure of each eigencomponent descends
to the spectral data. By~(a), the unique spectral
measure~$\rho_\cA$ inherits multiplicative Fourier coefficients
at each prime, giving prime-locality. The twelve-station
chain gives the Ramanujan bound on the rational central-charge
locus; the irrational extension has exactly the conditional
scope of Conjecture~\ref{conj:irrational-ramanujan}.
\end{proof}

\begin{conjecture}[Conditional Ramanujan for all chirally Koszul
algebras; \ClaimStatusConjectured]%
\label{conj:conditional-ramanujan}%
\index{Ramanujan bound!conditional on Hecke defect}%
Hypotheses~\textup{(a)} and~\textup{(b)} of
Theorem~\textup{\ref{thm:rigidity-inheritance}} are
established for all chirally Koszul algebras with HS-sewing
\textup{(}Propositions~\textup{\ref{prop:carleman-virasoro}}
and~\textup{\ref{prop:rigidity-threshold}}\textup{)}.
The entire Ramanujan programme therefore reduces to a single
cohomological hypothesis: the vanishing of
$[\delta_p(\cA)] \in H^1(\gAmod, \mathrm{ad}_{T_p})$
for all primes~$p$.
\end{conjecture}

\subsection{Assessment and open problems}%
\label{subsec:ac-assessment}%
\index{analytic continuation programme!assessment}

\begin{center}
\small
\renewcommand{\arraystretch}{1.3}
\begin{tabular}{clll}
\toprule
& \textbf{Component} & \textbf{Proof surface} & \textbf{Obstruction} \\
\midrule
\multicolumn{4}{l}{\emph{Ramanujan bound programme}} \\
\cmidrule{1-4}
1 & Conditional CPS automorphy of $M_r(s)$
 & \ClaimStatusConditional
 & vertical-strip and twist estimates \\
2 & Prime-locality (lattice)
 & \ClaimStatusProvedHere
 & n/a \\
3 & Prime-locality (rational, amplitude)
 & \ClaimStatusConditional
 & graph-sum Hecke compatibility \\
4 & Prime-locality (irrational)
 & \ClaimStatusConjectured
 & $[\delta_p] = 0$? \\
5 & Ramanujan (lattice implication)
 & \ClaimStatusConditional
 & CPS and local-factor comparison \\
6 & Ramanujan (rational implication)
 & \ClaimStatusConditional
 & character RS hypotheses \\
7 & Ramanujan (irrational)
 & \ClaimStatusConjectured
 & Langlands $r \ge 5$ \\
\cmidrule{1-4}
\multicolumn{4}{l}{\emph{Zero-location programme}} \\
\cmidrule{1-4}
8 & MC determines $\cE_\rho(\cA)$
 & \ClaimStatusProvedHere
 & n/a \\
9 & RS holomorphic at $\rho/2$
 & \ClaimStatusProvedHere
 & Gap P2 \\
10 & WP arithmetically blind
 & \ClaimStatusProvedHere
 & Gap P3 \\
11 & Genus-$1$ saddle triviality
 & \ClaimStatusProvedHere
 & Gap P1 \\
12 & Zero-location from MC
 & Structurally blocked
 & Thm~\ref{thm:structural-separation} \\
\bottomrule
\end{tabular}
\end{center}

\noindent
The programme splits into two halves with different characters.
The Ramanujan half is an implication diagram.  Its proved lattice input
is Hecke--Newton prime-locality; its conditional inputs are CPS
vertical-strip and twist estimates, local-factor comparison, rational
graph-sum Hecke compatibility, and Langlands functoriality for
higher symmetric powers.  The zero-location half is structurally
blocked at genus~$1$ by
Theorem~\ref{thm:structural-separation}: the MC equation
at genus~$1$ exhausts its algebraic content through the
scattering coupling and the moment $L$-functions, and no
further zero-location information is extractable from the
same Rankin--Selberg pairing.

\begin{problem}[Three open problems in the analytic
continuation programme]%
\label{prob:ac-open-problems}%
\index{analytic continuation programme!open problems}
\begin{enumerate}[label=\textup{(\Alph*)}]
\item
\emph{Hecke defect computation.}
Compute $[\delta_p(\cA)]$ explicitly for the Virasoro
algebra and the $\cW_3$ algebra at specific values
of~$c$ and~$p$. A nonzero class would refute
Conjecture~\ref{conj:prime-locality-transfer} and require
revision of the Ramanujan programme; a systematic vanishing
would give finite input toward the cohomological vanishing theorem
required by the Hecke module extension criterion of
Remark~\ref{rem:prime-locality-routes}.

\item
\emph{Higher-genus scattering access.}
At genus~$g \ge 2$, the period matrix couples Fock
sectors and the sewing operator is non-diagonal
(Remark~\ref{rem:collapse-permanence}).
Determine whether the genus-$g$ MC element
$\Theta_\cA^{(g)}$ provides access to the scattering
matrix that the genus-$1$ Rankin--Selberg method erases.
The complexified Liouville action at genus~$g \ge 2$ is
nonlinear ($R_{g_0} \ne 0$), admitting complex saddle
points (Conjecture~\ref{conj:complex-saddle-scattering})
that might retain scattering data. The minimal test case
is genus~$2$, where the period matrix is $2 \times 2$ and
the Siegel modular forms provide explicit Hecke operators.

\item
\emph{Unfolding bypass.}
Construct a representation of the partition
function~$\widehat{Z}^c$ on~$\cM_{1,1}$ whose analytic
structure in the spectral parameter~$s$ retains the
scattering poles of $E_s(\tau)$ without Rankin--Selberg
unfolding.
By Theorem~\ref{thm:structural-separation}(ii), the
Rankin--Selberg method \emph{cannot} produce such a
representation. Two candidate approaches:
\begin{itemize}[nosep]
\item the spectral-determinant side of the bridge
 (Proposition~\ref{prop:sewing-spectral-bridge}): the
 spectral zeta function $\zeta_\Delta(s)$ of the
 Laplacian on~$E_\tau$ admits meromorphic continuation
 via the Epstein functional equation, and its structure
 at the scattering poles $s = \rho/2$ is not subject to
 unfolding erasure, but it encodes the
 \emph{geometry} of~$E_\tau$ (eigenvalues of the
 Laplacian), not the \emph{algebra} of~$\cA$;
\item the Selberg trace formula applied to
 $\widehat{Z}^c \in L^2(\Gamma\backslash\mathfrak{H})$:
 the trace formula relates the spectral decomposition
 (including the scattering matrix) to the geometric
 side (closed geodesics), and the MC equation constrains
 the spectral coefficients; but the Selberg trace formula
 is an \emph{identity}, not a new representation, and
 extracting zero-location information from it is
 equivalent to proving RH itself.
\end{itemize}
\end{enumerate}
\end{problem}

\begin{remark}[The Langlands interface]%
\label{rem:langlands-interface}%
\index{geometric Langlands!arithmetic interface}%
At critical level, the bar complex computes the formal-disc oper
package
(Theorem~\ref{thm:langlands-bar-bridge}),
a vertex-algebraic construction of a central object in the
geometric Langlands correspondence. The degeneration
$\kappa(\widehat{\fg}_{-h^\vee}) = 0$ is where the modular
$L_\infty$-algebra loses its curvature obstruction; the
resulting uncurved deformation complex has cohomology governed
by the oper package. The manuscript does not identify its
Maurer--Cartan moduli with the oper space.

This chapter lives at \emph{generic} level
($\kappa \ne 0$); the geometric Langlands correspondence
operates at critical level. These are complementary: the
Langlands side governs the \emph{categorical} structure
of $\widehat{\fg}$-modules, while the shadow obstruction
tower governs the \emph{scalar} arithmetic invariants
of the partition function. The bridge between
regimes (the bar-cobar adjunction as $k$ varies from
generic to critical level) is the admissible-level programme
(Conjectures~\ref{conj:kl-periodic-cdg}--\ref{conj:kl-braided}).
\end{remark}

% ============================================================
\subsection{The arithmetic packet connection}%
\label{subsec:arithmetic-packet-connection}%
\index{arithmetic packet connection|textbf}%
\index{Hecke module!packet connection}

The depth decomposition (Theorem~\ref{thm:depth-decomposition})
and the Hecke-equivariant MC element
(Theorem~\ref{thm:spectral-continuation-bridge}) together
determine a single geometric object that encodes both the
arithmetic content and the homotopy defect as complementary
components of a flat meromorphic connection. The construction
is forced: once the Hecke module admits a Jordan decomposition
and each eigenspace carries an $L$-packet, there is exactly one
flat connection whose singularities are the $L$-function zeros.

\begin{definition}[Arithmetic packet module]%
\label{def:arithmetic-packet-module}%
\index{arithmetic packet module|textbf}%
Let $\cA$ be a chirally Koszul algebra with HS-sewing.
The \emph{arithmetic packet module} of~$\cA$ is the
finite-dimensional vector space
\begin{equation}\label{eq:arithmetic-packet-module}
 M_\cA \;:=\; \bigoplus_{\chi} M_\chi,
\end{equation}
where the sum ranges over the generalised Hecke eigenspaces
activated by the partition function~$\widehat{Z}^c_\cA$ in its
Roelcke--Selberg spectral decomposition on~$\cM_{1,1}$.
Each~$M_\chi$ carries:
\begin{enumerate}[label=\textup{(\roman*)}]
\item a \emph{semisimple part}: the Hecke eigenvalue
 $a_\chi \in \bC$;
\item a \emph{nilpotent part}:
 $N_\chi \in \operatorname{End}(M_\chi)$ with
 $N_\chi^{\dim M_\chi} = 0$;
\item a \emph{completed $L$-packet}: a meromorphic function
 $\Lambda_\chi(s)$ on~$\bC$
 with functional equation $\Lambda_\chi(s) =
 \epsilon_\chi\,\Lambda_\chi(1{-}s)$,
 $\epsilon_\chi \in \{\pm 1\}$.
\end{enumerate}
For lattice VOAs, the decomposition is
$M_{V_\Lambda} = \bigoplus_{j=0}^{m} \bC\cdot e_j$
with $e_0$ the Eisenstein class,
$\Lambda_0(s) = C_0(s)\,\zeta(s)\,\zeta(s{-}r/2{+}1)$, and
$\Lambda_j(s) = C_j(s)\,L(s,f_j)$ for each cusp form~$f_j$
in the Hecke decomposition
$\Theta_\Lambda = E_{r/2} + \sum_{j=1}^m c_j\,f_j$.
All nilpotent parts vanish: $N_\chi = 0$ for every~$\chi$.
\end{definition}

\begin{definition}[Arithmetic packet connection]%
\label{def:arithmetic-packet-connection}%
\index{arithmetic packet connection!definition|textbf}%
The \emph{arithmetic packet connection} of~$\cA$ is the
meromorphic connection on the trivial bundle
$M_\cA \otimes \cO_{\bC}$ over the spectral
line~$s \in \bC$, defined by
\begin{equation}\label{eq:arithmetic-packet-connection}
 \nabla_\cA^{\mathrm{arith}}
 \;:=\;
 d \;-\; \bigoplus_\chi
 \frac{\Lambda_\chi'(s)}{\Lambda_\chi(s)}\,
 \bigl(\operatorname{id}_{M_\chi} + N_\chi\bigr)\,ds.
\end{equation}
\end{definition}

\begin{theorem}[Flatness and divisor independence]%
\label{thm:packet-connection-flatness}%
\ClaimStatusProvedHere%
\index{arithmetic packet connection!flatness}%
\begin{enumerate}[label=\textup{(\roman*)}]
\item $(\nabla_\cA^{\mathrm{arith}})^2 = 0$.
\item The singular divisor is
 \begin{equation}\label{eq:arithmetic-singular-divisor}
 D_\cA \;=\; \bigcup_\chi
 \operatorname{div}(\Lambda_\chi),
 \end{equation}
 and depends only on the semisimple data
 $\{a_\chi,\,\Lambda_\chi\}$, not on the nilpotent
 parts~$N_\chi$.
\item On any simply connected open subset
 $U \subset \bC \setminus D_\cA$, every horizontal section
 has the form
 \begin{equation}\label{eq:horizontal-sections}
 \sigma(s) \;=\; \bigoplus_\chi
 \Lambda_\chi(s)\,
 \exp\!\bigl((\log\Lambda_\chi(s))\,N_\chi\bigr)\,v_\chi\,,
 \end{equation}
 for constant vectors $v_\chi \in M_\chi$.
\end{enumerate}
\end{theorem}

\begin{proof}
Write $\omega_\chi := d\log\Lambda_\chi(s) =
(\Lambda_\chi'/\Lambda_\chi)\,ds$ and
$E_\chi := \operatorname{id}_{M_\chi} + N_\chi$.
The connection one-form on each block is
$A_\chi = \omega_\chi \cdot E_\chi$, a scalar one-form
times a constant endomorphism. The curvature is
\[
 dA_\chi - A_\chi \wedge A_\chi
 \;=\;
 d\omega_\chi \cdot E_\chi
 \;-\;
 (\omega_\chi \wedge \omega_\chi)\,E_\chi^2.
\]
Since $\omega_\chi$ is a meromorphic $1$-form on the
one-dimensional $s$-line, $\omega_\chi \wedge \omega_\chi = 0$
and $d\omega_\chi = 0$ (both vanish for degree reasons).
Hence $(\nabla^{\mathrm{arith}})^2 = 0$.

The singularities of~$A_\chi$ occur at the zeros and poles
of~$\Lambda_\chi$, where $d\log\Lambda_\chi$ has simple poles.
The nilpotent part~$N_\chi$ multiplies this scalar one-form but
does not introduce new poles, giving~(ii).

For~(iii), on~$U$ choose a branch of $\log\Lambda_\chi$.
Then $\nabla(\Lambda_\chi^{E_\chi}\,v)
= \Lambda_\chi^{E_\chi}\,(\nabla v + d\log\Lambda_\chi
\cdot E_\chi \cdot v) = 0$,
where $\Lambda_\chi^{E_\chi}
:= \exp((\log\Lambda_\chi)\,E_\chi)
= \Lambda_\chi\cdot\exp((\log\Lambda_\chi)\,N_\chi)$.
The factorisation holds because
$(\log\Lambda_\chi)\cdot\operatorname{id}$ is scalar,
hence commutes with~$N_\chi$; the exponential of a sum of
commuting operators is the product of their exponentials.
The nilpotent truncation gives
$\exp((\log\Lambda_\chi)\,N_\chi)
= \sum_{k=0}^{\dim M_\chi - 1}
(\log\Lambda_\chi)^k\,N_\chi^k / k!$.
\end{proof}

\begin{corollary}[Lattice transparency]%
\label{cor:lattice-packet-diagonal}%
\ClaimStatusProvedElsewhere%
\index{lattice VOA!packet connection diagonal}%
For a lattice VOA~$V_\Lambda$, all nilpotent parts vanish and
\[
 \nabla_{V_\Lambda}^{\mathrm{arith}}
 \;=\; d \;-\;
 \operatorname{diag}\!\Bigl(
 d\log\Lambda_0(s),\;
 d\log\Lambda_1(s),\;
 \ldots,\;
 d\log\Lambda_m(s)
 \Bigr).
\]
The monodromy around any zero~$\rho$ of~$\Lambda_j$ is
$\exp(2\pi i\,\operatorname{ord}_\rho(\Lambda_j))$ on the
$j$-th factor. The connection is the direct sum of
logarithmic connections $d - d\log\Lambda_j$ on rank-one
bundles, one per $L$-function factor in the Hecke
decomposition~\eqref{eq:hecke-epstein-factorization}.
\end{corollary}

\begin{remark}[The arithmetic--homotopy dictionary]%
\label{rem:arithmetic-homotopy-dictionary}%
\index{arithmetic skeleton|textbf}%
\index{algebraic defect module|textbf}%
The packet connection realises the depth decomposition
$d(\cA) = 1 + d_{\mathrm{arith}} + d_{\mathrm{alg}}$
\textup{(}Theorem~\textup{\ref{thm:depth-decomposition}}\textup{)}
geometrically:
\begin{itemize}[nosep]
\item The \emph{arithmetic skeleton}
 $\operatorname{Ask}(\cA) := M_\cA^{\mathrm{ss}}$
 \textup{(}the semisimple quotient\textup{)}
 controls the singular divisor~$D_\cA$:
 the number of irreducible components
 is~$d_{\mathrm{arith}} + 1$ \textup{(}including Eisenstein\textup{)}.
\item The \emph{algebraic defect module}
 $\operatorname{Def}_{\mathrm{alg}}(\cA)
 := \bigoplus_\chi N_\chi\,M_\chi$
 contributes only unipotent monodromy
 around the same divisor.
\end{itemize}
The principle: \emph{arithmetic is semisimple; homotopy
defect is unipotent.}
\end{remark}

\begin{definition}[Frontier defect form]%
\label{def:frontier-defect-form}%
\index{frontier defect form|textbf}%
\index{scattering matrix!frontier defect}%
Let $\Lambda_{\mathrm{Eis}}(s) := \Lambda_0(s)$ be the
Eisenstein packet and
$\varphi(s) = \Lambda^*(1{-}s)/\Lambda^*(s)$
the automorphic scattering matrix on~$\cM_{1,1}$, where
$\Lambda^*(s) = \pi^{-s}\,\Gamma(s)\,\zeta(2s)$.
The \emph{frontier defect form} is the meromorphic $1$-form
\begin{equation}\label{eq:frontier-defect-form}
 \Omega_\cA
 \;:=\;
 d\log\Lambda_{\mathrm{Eis}}(s)
 \;-\; d\log\varphi(s).
\end{equation}
\end{definition}

\begin{proposition}[Gauge criterion for scattering access]%
\label{prop:gauge-criterion-scattering}%
\ClaimStatusProvedHere%
\index{scattering matrix!gauge criterion}%
Suppose the Eisenstein eigenspace~$M_{\mathrm{Eis}}$ is
rank~$1$ \textup{(}as it is for lattice VOAs and all
standard-landscape families\textup{)}. Then the Eisenstein
block of $\nabla_\cA^{\mathrm{arith}}|_{\mathrm{Eis}}$ is
gauge-equivalent to the scattering connection
$d - d\log\varphi(s)$ if and only if the frontier defect
form is exact:
$\Omega_\cA = dh$ for some meromorphic function~$h$
on~$\bC$; equivalently,
$\Lambda_{\mathrm{Eis}}(s)/\varphi(s)$
is a meromorphic function of~$s$ with trivial monodromy.

If $\dim M_{\mathrm{Eis}} > 1$, the weaker sufficient
condition is: $N_{\mathrm{Eis}} = 0$ and $\Omega_\cA = dh$.
\end{proposition}

\begin{proof}
On the rank-one Eisenstein block, both connections are
scalar: the packet connection is
$d - d\log\Lambda_{\mathrm{Eis}}$ and the scattering
connection is $d - d\log\varphi$.
Two rank-one meromorphic connections $d - \omega_1$ and
$d - \omega_2$ on a line bundle are gauge-equivalent via
$\sigma \mapsto e^h\,\sigma$ if and only if
$\omega_1 - \omega_2 = dh$. Applying this to
$\omega_1 = d\log\Lambda_{\mathrm{Eis}}$ and
$\omega_2 = d\log\varphi$ gives the equivalence.

In the higher-rank case, $N_{\mathrm{Eis}} = 0$ reduces
the Eisenstein block to a scalar connection, after which
the same rank-one argument applies.
\end{proof}

\begin{remark}[Reformulation of the structural obstruction]%
\label{rem:packet-reformulation}%
\index{structural obstruction!packet reformulation}%
Theorem~\ref{thm:structural-separation} showed that the
genus-$1$ Rankin--Selberg integral is holomorphic at the
scattering poles~$s = \rho/2$: the MC equation's algebraic
content is exhausted at genus~$1$ before reaching the
scattering matrix.

The packet connection reformulates this obstruction with
geometric precision. The divisor~$D_\cA$ is the zero set
of the activated $L$-packets~$\Lambda_\chi$; the scattering
poles are the zero set of~$\Lambda^*(s)$. These divisors
generically differ: the MC element activates arithmetic
content that is \emph{related to} but not \emph{identical
with} the scattering data. The frontier defect
form~$\Omega_\cA$ measures the discrepancy, and the
gauge criterion
\textup{(}Proposition~\textup{\ref{prop:gauge-criterion-scattering}}\textup{)}
identifies the two precise conditions under which the
discrepancy vanishes.

For lattice VOAs at rank~$r$,
$\Lambda_{\mathrm{Eis}}(s) = C_0(s)\,\zeta(s)\,\zeta(s{-}r/2{+}1)$,
while $\varphi(s) = \Lambda^*(1{-}s)/\Lambda^*(s)$ involves
$\zeta(2s)/\zeta(2s{-}1)$. The frontier defect form encodes
the difference between these two zeta-function factors.
Since $\zeta(s)$ and $\zeta(2s)$ have different zero sets,
$\Omega_{V_\Lambda} \ne 0$ in general: the connection-theoretic
incarnation of Theorem~\ref{thm:structural-separation}.
\end{remark}

\begin{conjecture}[The arithmetic comparison conjecture]%
\label{conj:arithmetic-comparison}%
\ClaimStatusConjectured%
\index{arithmetic comparison conjecture|textbf}%
For every chirally Koszul algebra~$\cA$ with HS-sewing\textup{:}
\begin{enumerate}[label=\textup{(\roman*)}]
\item the modular MC element~$\Theta_\cA$ canonically
 determines the arithmetic packet connection
 $\nabla_\cA^{\mathrm{arith}}$, functorially under
 quasi-isomorphism of chiral algebras\textup{;}
\item the semisimplification $M_\cA^{\mathrm{ss}}$ is the
 Hecke-semisimple quotient of the activated graph amplitudes
 at all degrees\textup{;}
\item the residue of the frontier defect form~$\Omega_\cA$ at
 each scattering pole $s = \rho/2$ is computable from
 the higher-genus MC data
 $\{\Theta_\cA^{(g)}\}_{g \ge 2}$\textup{:}
 \begin{equation}\label{eq:higher-genus-scattering-access}
 \operatorname{Res}_{s = \rho/2}\,\Omega_\cA
 \;=\;
 \lim_{g\to\infty}
 \operatorname{Res}_{s=\rho/2}\,
 d\log\Lambda_{\mathrm{Eis}}^{(g)}(s),
 \end{equation}
 where $\Lambda_{\mathrm{Eis}}^{(g)}$ is the Eisenstein
 packet extracted from the genus-$g$ MC equation.
\end{enumerate}
The only obstruction to matching the full automorphic
scattering data is the frontier defect form~$\Omega_\cA$;
the three open problems of~\textup{\ref{prob:ac-open-problems}}
reformulate as:
\begin{enumerate}[label=\textup{(\Alph*$'$)},nosep]
\item computing the nilpotent parts~$N_\chi$ and the algebraic
 defect module $\operatorname{Def}_{\mathrm{alg}}(\cA)$,
\item establishing part~\textup{(iii)} via higher-genus MC data,
\item constructing a spectral representation governed by
 $\nabla_\cA^{\mathrm{arith}}$ that avoids unfolding erasure
 \textup{(}Theorem~\textup{\ref{thm:structural-separation}}\textup{)}.
\end{enumerate}
\end{conjecture}

\begin{remark}[Sharpening of the arithmetic comparison conjecture]%
\label{rem:arithmetic-comparison-sharpening}%
\ClaimStatusProvedHere%
\index{arithmetic comparison conjecture!sharpening}%
\index{Niemeier lattices!arithmetic comparison failure}%
\index{arithmetic comparison conjecture!non-injectivity}%
The arithmetic comparison conjecture requires two qualifications
discovered by systematic computation.

\begin{enumerate}[label=\textup{(\roman*)}]
\item \textup{(The scalar version is false.)}
All $24$~Niemeier lattice VOAs share the scalar MC element
$\kappa \cdot \eta \otimes \Lambda = 24 \cdot \eta \otimes \Lambda$,
yet they have distinct arithmetic packets
$\nabla_{V_\Lambda}^{\mathrm{arith}}$
\textup{(}the theta series $\Theta_\Lambda$ encodes
the root system\textup{)}.
The scalar projection $\kappa$ is blind to
$276$~of~$276$ Niemeier pairs.
Part~\textup{(i)} of the conjecture must therefore use the
\emph{full} MC element $\Theta_\cA$ together with its scalar
projection.

\item \textup{(The full version is consistent but not injective.)}
The map $\Theta_\cA \mapsto \nabla_\cA^{\mathrm{arith}}$
factors through the partition function $Z(\tau, \bar\tau)$,
since the arithmetic packet is extracted from the
Roelcke--Selberg decomposition of~$Z$.
Two algebras with different MC elements but the same
partition function produce the same
$\nabla^{\mathrm{arith}}$: for example,
$D_{16}^+ \oplus E_8$ and $E_8^3$ have different
$\Theta_\cA$ \textup{(}different OPE data, different
multi-channel bar complexes\textup{)} but the same
theta series $\Theta_\Lambda(\tau) = \Theta_{E_8}(\tau)^3$,
hence the same $\nabla^{\mathrm{arith}}$.
\end{enumerate}

The \emph{minimal degree for arithmetic discrimination}
depends on the shadow class:
\begin{center}
\small
\renewcommand{\arraystretch}{1.2}
\begin{tabular}{lcc}
\toprule
\textbf{Class} & \textbf{Minimal degree} & \textbf{Reason} \\
\midrule
G & $2$ & $\kappa$ discriminates (free fields) \\
L & $3$ & cubic shadow $\mathfrak{C}$ needed \\
C & $4$ & quartic contact $\mathfrak{Q}$ needed \\
Lattice ($r \leq 40$) & $\dim M_{r/2} + 1$ &
 one degree per eigenform \\
M & $\infty$ & full tower needed (open) \\
\bottomrule
\end{tabular}
\end{center}
For lattice VOAs, the charged-sector bar complex provides
a stronger invariant than the scalar shadow obstruction tower and
achieves full discrimination at genus~$1$.
\end{remark}

\subsection{The Miura defect in the packet connection}%
\label{sec:miura-packet-bridge}%
\index{Miura defect!packet connection}%
\index{arithmetic packet connection!Miura defect}

The finite Miura defect
(Theorem~\ref{thm:finite-miura-defect})
and the arithmetic packet connection
(Definition~\ref{def:arithmetic-packet-connection})
are two views of the same structure.
Their interaction reveals the precise location of
arithmetic content.

\begin{proposition}[\ClaimStatusProvedHere]%
\label{prop:miura-packet-splitting}%
\index{Miura defect!packet splitting}%
For principal~$\cW_N$, the arithmetic packet connection
decomposes as
\begin{equation}\label{eq:miura-packet-splitting}
\nabla_{\cW_N}^{\mathrm{arith}}
\;=\;
\nabla_\cH^{\mathrm{arith},\,\oplus(N-1)}
\;\oplus\;
\nabla_{\mathrm{def}}^{\mathrm{arith}}\,,
\end{equation}
where $\nabla_\cH^{\mathrm{arith},\,\oplus(N-1)}$ denotes the
connection that restricts to $\nabla_\cH^{\mathrm{arith}}
= d - d\log[\zeta(s)\zeta(s{+}1)]\,ds$ on each of the
$(N{-}1)$ Heisenberg summands, and
$\nabla_{\mathrm{def}}^{\mathrm{arith}}$ is the defect
connection governing the finite modes
$m = 1,\ldots,N{-}1$.
In particular:
\begin{enumerate}[label=\textup{(\roman*)}]
\item the Heisenberg component is exact Euler--Koszul
 \textup{(}$N_\chi = 0$, $D_\cH = 1$\textup{)}, contributing
 only the universal divisor $\operatorname{div}(\zeta(s))
 \cup \operatorname{div}(\zeta(s{+}1))$\textup;
\item the defect component carries all family-specific
 arithmetic content:
 the quartic contact coupling
 $Q_{\mathrm{Vir}}^{\mathrm{ct}} = 10/[c(5c{+}22)]$,
 the higher-spin quartic couplings
 $Q_s^{\mathrm{ct}}(\cW_N)$, and the
 non-scalar spectral data\textup;
\item the defect singular divisor is contained in the
 finite set
 $\{m^{-s} : m = 1,\ldots,N{-}1\} \cup
 \operatorname{div}(\zeta(s{+}1))$,
 and its dimension over the Heisenberg core is
 exactly $\operatorname{ek}(\cW_N) = N - 2$\textup;
\item the frontier defect form
 $\Omega_{\cW_N}$
 \textup{(}Definition~\textup{\ref{def:frontier-defect-form}}\textup{)}
 is determined by the defect component alone:
 \begin{equation}\label{eq:frontier-from-defect}
 \Omega_{\cW_N}
 \;=\;
 d\log\Lambda_{\mathrm{Eis}}^{\mathrm{def}}(s)
 \;-\;
 d\log\varphi(s)\,.
 \end{equation}
\end{enumerate}
\end{proposition}

\begin{proof}
The Dirichlet-sewing lift decomposes additively as
$S_{\cW_N}(u) = (N{-}1)\,S_\cH(u) +
D_{\cW_N}^{\mathrm{prime}}(u)$
(Theorem~\ref{thm:finite-miura-defect}(iii)).
Since the Hecke operators $T_p$ act on the sewing lift
through the Euler factor decomposition, and
$S_\cH(u) = \zeta(u)\,\zeta(u{+}1)$ is a Hecke eigenform
(the Eisenstein series $E_1$), the additive decomposition
descends to a direct-sum decomposition of the Hecke module:
$M_{\cW_N} = M_\cH^{\oplus(N-1)} \oplus M_{\mathrm{def}}$,
where $M_\cH = \bC\cdot e_0$ is the rank-one Heisenberg
eigenspace.
On the Heisenberg copies, the packet
is $\Lambda_\cH(s) = \zeta(s)\,\zeta(s{+}1)$ with
$N_\cH = 0$, giving a diagonal exact Euler--Koszul
connection (Corollary~\ref{cor:lattice-packet-diagonal}).
The defect component inherits its connection from the
defect Dirichlet polynomial
$-\zeta(u{+}1)\sum_{m=1}^{N-1}(N{-}m)\,m^{-u}$
(Proposition~\ref{prop:prime-side-defect-formula}).

Part~(iv): the Eisenstein packet of $\cW_N$ decomposes
multiplicatively as
$\Lambda_{\mathrm{Eis}} = \Lambda_\cH^{N-1} \cdot
\Lambda_{\mathrm{Eis}}^{\mathrm{def}}$.
Taking logarithmic derivatives gives
$d\log\Lambda_{\mathrm{Eis}}
= (N{-}1)\,d\log\Lambda_\cH
+ d\log\Lambda_{\mathrm{Eis}}^{\mathrm{def}}$.
Subtracting $d\log\varphi$ and noting that
$d\log\Lambda_\cH$ contributes only universal
Heisenberg divisor data (common to all $\cW_N$),
the family-specific frontier defect reduces
to~\eqref{eq:frontier-from-defect}.
\end{proof}

\begin{remark}[The Heisenberg core is arithmetically inert]%
\label{rem:heisenberg-inert}%
\index{Heisenberg!arithmetically inert}%
Proposition~\ref{prop:miura-packet-splitting} makes the
finite defect principle
(Remark~\ref{rem:finite-defect-principle}) connection-theoretic:
the Heisenberg core carries no quartic resonance data and is
\emph{arithmetically inert} in the following precise sense:
\begin{enumerate}[label=\textup{(\roman*)}]
\item its packet connection is diagonal, with only universal
 zeta divisors;
\item its contribution to the frontier defect form
 cancels identically;
\item all obstructions to matching automorphic scattering
 live in the finite defect sector.
\end{enumerate}
The structural consequence for the arithmetic comparison
conjecture
(Conjecture~\ref{conj:arithmetic-comparison}):
the three open problems reduce to the defect sector.
In particular, the nilpotent parts~$N_\chi$ and the
algebraic defect module
(Remark~\ref{rem:arithmetic-homotopy-dictionary})
live entirely in the finite Miura defect, not in the
Heisenberg core.
\end{remark}

\subsubsection*{Explicit examples}%
\label{subsub:packet-connection-examples}%
\index{arithmetic packet connection!explicit examples}

We compute $\nabla_\cA^{\mathrm{arith}}$ for four families, in
increasing order of complexity.

\begin{example}[Heisenberg algebra $\cH_\kappa$]%
\label{ex:packet-heisenberg}%
\index{Heisenberg algebra!packet connection}%
The Heisenberg algebra $\cH_\kappa$ at level~$\kappa$ has
central charge $c = 1$ and physical partition function
$Z_{\cH_\kappa}(\tau) = 1/\lvert\eta(\tau)\rvert^2$.
The primary-counting function
$\widehat{Z}^1 = y^{1/2}\lvert\eta\rvert^{2}\cdot
\lvert\eta\rvert^{-2} = y^{1/2}$
is the constant term of the Eisenstein spectrum
on~$\cM_{1,1}$: it has no cuspidal component and no
higher Eisenstein harmonics, reflecting the fact that
every state in a single-boson Fock space is
a descendant of~$|0\rangle$.
The Hecke decomposition of the sewing lift is the single
Eisenstein term $S_\cH(u) = \zeta(u)\,\zeta(u{+}1)$.

The arithmetic packet module is rank~$1$:
\[
 M_{\cH_\kappa}
 \;=\;
 \bC\cdot e_0,
 \qquad
 \Lambda_0(s) = C_0(s)\,\zeta(s)\,\zeta(s{+}1),
 \qquad
 N_0 = 0.
\]
The connection is the scalar logarithmic connection
\begin{equation}\label{eq:packet-heisenberg}
 \nabla_{\cH_\kappa}^{\mathrm{arith}}
 \;=\;
 d \;-\; d\log\bigl[\zeta(s)\,\zeta(s{+}1)\bigr]\,ds
 \;=\;
 d \;-\; \Bigl(
 \frac{\zeta'(s)}{\zeta(s)}
 + \frac{\zeta'(s{+}1)}{\zeta(s{+}1)}
 \Bigr)\,ds.
\end{equation}
The singular divisor is
$D_{\cH} = \operatorname{div}(\zeta(s))
\cup \operatorname{div}(\zeta(s{+}1))$:
every nontrivial zero~$\rho$ of the Riemann zeta function
produces simple poles at $s = \rho$ and $s = \rho - 1$.
There is no nilpotent monodromy, no algebraic defect, and the
frontier defect form reduces to
\[
 \Omega_{\cH}
 \;=\;
 d\log\bigl[\zeta(s)\,\zeta(s{+}1)\bigr]
 \;-\; d\log\varphi(s)
 \;\neq\; 0,
\]
since $\zeta(s)\,\zeta(s{+}1) \ne
\Lambda^*(1{-}s)/\Lambda^*(s)$: the packet function has
poles at $s = 1$ (from $\zeta(s)$) and $s = 0$
(from $\zeta(s{+}1)$), while
the scattering matrix has poles at $s = \rho/2$ (from the
zeros of $\Lambda^*(s) = \pi^{-s}\Gamma(s)\zeta(2s)$).
The shadow depth is $d(\cH) = 2$ (Gaussian archetype), with
$d_{\mathrm{arith}} = 1$ (the Eisenstein component) and
$d_{\mathrm{alg}} = 0$.
\end{example}

\begin{example}[Rank-$1$ lattice $V_\bZ$]%
\label{ex:packet-rank1-lattice}%
\index{lattice VOA!rank-1 packet connection}%
The rank-$1$ even lattice VOA $V_\bZ$ has
$\Theta_\bZ(\tau) = \theta_3(\tau) = \sum_{n\in\bZ}q^{n^2}
\in M_{1/2}(\Gamma_0(4))$.
Since $\dim S_{1/2}(\Gamma_0(4)) = 0$, the Hecke decomposition
is pure Eisenstein:
$\Theta_\bZ = E_{1/2}$ (the weight-$1/2$ Eisenstein series at
level~$4$). The packet module is rank~$1$:
\[
 M_{V_\bZ}
 \;=\;
 \bC\cdot e_0,
 \qquad
 \Lambda_0(s) = C_0(s)\,\zeta(s)\,\zeta(s{+}\tfrac12),
 \qquad
 N_0 = 0.
\]
Here $\kappa(V_\bZ) = 1$ and $d_{\mathrm{arith}} = 1$
(one Eisenstein component; the cusp-form space is empty),
$d_{\mathrm{alg}} = 0$ (class~$\mathbf{G}$), and
$d(V_\bZ) = 2$.
The packet connection is again scalar: the divisor is
$D_{V_\bZ} = \operatorname{div}(\zeta(s)) \cup
\operatorname{div}(\zeta(s{+}\tfrac12))$.
Every lattice VOA at rank $r < 24$ ($r$ even) has empty cusp
space $S_{r/2}(\mathrm{SL}(2,\bZ)) = \{0\}$
(since $r/2 < 12$), hence rank-$1$ packet module and
$d_{\mathrm{arith}} = 1$. The first nontrivial cuspidal
component appears at rank~$24$
(the Leech lattice; see~\S\ref{subsec:bocherer-bridge-theorem}).
\end{example}

\begin{example}[Virasoro $\mathrm{Vir}_c$ at arbitrary~$c$]%
\label{ex:packet-virasoro}%
\index{Virasoro algebra!packet connection}%
\index{arithmetic packet connection!Virasoro}%
The Virasoro algebra $\mathrm{Vir}_c$ has
$\kappa(\mathrm{Vir}_c) = c/2$, shadow depth $d = \infty$
(class~$\mathbf{M}$), and vacuum character
$\chi_0(q) = q^{-c/24}\prod_{n \geq 2}(1-q^n)^{-1}$.
The primary-counting function
$\widehat{Z}^c = y^{c/2}\lvert\eta\rvert^{2c}\,Z$
is a Maass form on $\mathrm{SL}(2,\bZ)\backslash\mathfrak{H}$
when $c \in \bR$, not a classical modular form.
For irrational~$c$, the Roelcke--Selberg spectral
decomposition takes the continuous form
\[
 \widehat{Z}^c(\tau)
 \;=\;
 \int_0^\infty \alpha(t)\,E_{1/2+it}(\tau)\,dt
 \;+\;
 \sum_j \beta_j\,\phi_j(\tau),
\]
where the integral is over the continuous Eisenstein spectrum
and the sum over Maass cusp forms $\phi_j$ with Laplace
eigenvalues $\lambda_j = 1/4 + t_j^2$.

The arithmetic packet module has an Eisenstein
component and (generically) discrete cusp-form eigenspaces:
\[
 M_{\mathrm{Vir}_c}
 \;=\;
 M_{\mathrm{Eis}} \;\oplus\;
 \bigoplus_j M_{\phi_j}.
\]
The Eisenstein block has $\dim M_{\mathrm{Eis}} = 1$ with
packet function $\Lambda_{\mathrm{Eis}}(s)$ given by the
Rankin--Selberg integral of~$\widehat{Z}^c$ against the
Eisenstein series~$E_s$: this is a meromorphic function
of~$s$ with functional equation inherited from
$E_s \leftrightarrow E_{1-s}$, but it is not in general
expressible as a product of standard zeta functions
(unlike the lattice case, where
$\Lambda_{\mathrm{Eis}} = C_0\,\zeta(s)\,\zeta(s{-}r/2{+}1)$
follows from the Hecke eigenform decomposition of the
theta function).
The cuspidal $L$-packets $\Lambda_j(s) = L(s,\phi_j)$ are the
$L$-functions of the activated Maass cusp forms~$\phi_j$.
The packet connection on the Eisenstein block is
$d - d\log\Lambda_{\mathrm{Eis}}(s)$,
with singularities at the zeros and poles
of~$\Lambda_{\mathrm{Eis}}$.

The key distinction from lattice algebras:
$d_{\mathrm{arith}}$ need not be finite. Whether finitely
many or infinitely many Maass eigenspaces are activated is
controlled by the spectral projection coefficients~$\beta_j$,
which in turn depend on the full shadow obstruction tower
$\Theta_{\mathrm{Vir}_c}$. For the Ising model ($c = 1/2$),
there are only two scalar primaries, and the constrained
Epstein zeta is a finite Dirichlet polynomial
(Proposition~\ref{prop:ising-d-arith}), so
$d_{\mathrm{arith}} = 0$ despite $d_{\mathrm{alg}} = \infty$.
The depth decomposition separates these:
$d = 1 + 0 + \infty = \infty$.

All nilpotent parts~$N_\chi$ are conjectured to vanish for
the universal Virasoro algebra (the vacuum module generates
a free Fock space at generic~$c$, leaving no room for Jordan
blocks in the Hecke decomposition). Whether $N_\chi = 0$
for all~$c$, or nontrivial nilpotent parts arise at special
values (e.g.\ $c = 0$, $c = 1$, $c = 25$),
is an instance of the algebraic defect computation
(Conjecture~\ref{conj:arithmetic-comparison}(A$'$)).
\end{example}

\begin{example}[$\cW_3$ algebra]%
\label{ex:packet-w3}%
\index{W3 algebra@$\cW_3$ algebra!packet connection}%
\index{arithmetic packet connection!$\cW_3$}%
The $\cW_3$ algebra has generators of spins $2$ and~$3$,
central charge~$c$, and
$\kappa(\cW_3) = 5c/6 = c(H_3 - 1)$, where
$H_3 = 1 + 1/2 + 1/3 = 11/6$.
By Proposition~\ref{prop:miura-packet-splitting} with
$N = 3$, the packet connection splits as
\[
 \nabla_{\cW_3}^{\mathrm{arith}}
 \;=\;
 \nabla_\cH^{\mathrm{arith},\,\oplus 2}
 \;\oplus\;
 \nabla_{\mathrm{def}}^{\mathrm{arith}},
\]
where the Heisenberg core consists of two copies of
$d - d\log[\zeta(s)\zeta(s{+}1)]\,ds$
(arithmetically inert by Remark~\ref{rem:heisenberg-inert})
and the defect connection governs modes $m = 1, 2$ with
Dirichlet polynomial
$D_{\cW_3}^{\mathrm{prime}}(u)
= -\zeta(u{+}1)(2\cdot 1^{-u} + 1\cdot 2^{-u})$
(Proposition~\ref{prop:prime-side-defect-formula}).
The defect singular divisor is
$D_{\mathrm{def}} \subset
\operatorname{div}(\zeta(s{+}1))$
(the poles of $d\log D_{\cW_3}^{\mathrm{prime}}$ are
confined to the zeros and poles of $\zeta(s{+}1)$,
since $m^{-u}$ is entire and nonvanishing for $m \geq 1$),
and the Euler--Koszul rank of the defect is
$\operatorname{ek}(\cW_3) = 1$.
The quartic contact coupling
$Q^{\mathrm{ct}}(\mathrm{Vir}) = 10/[c(5c{+}22)]$
contributes to the defect sector; the additional
spin-$3$ channel introduces a second quartic coupling
$Q_3^{\mathrm{ct}}(\cW_3)$ whose interaction with
the spin-$2$ channel is measured by the propagator variance
$\delta_{\mathrm{mix}}$ (Theorem~\ref{prop:propagator-variance}).

The frontier defect form
$\Omega_{\cW_3} = d\log\Lambda_{\mathrm{Eis}}^{\mathrm{def}}
- d\log\varphi$
(equation~\eqref{eq:frontier-from-defect}) is determined by
the defect alone: the two Heisenberg copies cancel.
All family-specific arithmetic content (the Euler--Koszul
rank, the quartic resonance, the multi-channel mixing) lives
in the one-dimensional defect sector.
\end{example}

% ============================================================
\section{The genus-$2$ arithmetic frontier}%
\label{sec:genus2-arithmetic-frontier}%
\index{genus-2 arithmetic frontier|textbf}

Theorem~\ref{thm:structural-separation} proves that genus-$1$
MC data exhausts its algebraic content: the zero-location
question requires either higher-genus or genuinely analytic
input. The sewing--Hecke collapse
(Theorem~\ref{thm:sewing-hecke-reciprocity}) is the
mechanism: the genus-$1$ sewing operator is Fock-diagonal,
forcing the two-variable $L$-object $L_\cA(s,u)$ to reduce
to $S_\cA(s{+}u{-}1)$. This section formulates the precise
consequences of that collapse \emph{failing} at genus~$2$
(Remark~\ref{rem:genus2-escape-route}).

\subsection{Sewing non-diagonality at genus $2$}%
\label{subsec:genus2-non-collapse}

\begin{proposition}[Genus-$2$ sewing non-diagonality;
\ClaimStatusProvedHere]%
\label{prop:genus2-non-diagonal}%
\index{sewing kernel!genus-2 non-diagonality}%
Let $\cA$ be a chirally Koszul algebra with HS-sewing.
At genus~$2$, the one-particle sewing kernel $K_2$ is a
$2 \times 2$ block operator on $A^2(D)^{\oplus 2}$:
\[
K_2
\;=\;
\begin{pmatrix}
K_{q_1} & R_{12} \\
R_{21} & K_{q_2}
\end{pmatrix},
\]
where $K_{q_j}$ are the handle sewing operators and
$R_{12}, R_{21}$ are off-diagonal coupling operators with
$\|R_{ij}\|_{\mathrm{HS}} \leq C\,|q_3|^{1/2}$ for the
internal sewing parameter~$q_3$. The Fredholm determinant
admits a Schur complement factorization:
\begin{equation}\label{eq:genus2-fredholm-schur}
\det(1 - K_2) \;=\; \det(1 - K_{q_1})\,
\det\bigl(1 - K_{q_2}
- R_{21}(1{-}K_{q_1})^{-1}R_{12}\bigr).
\end{equation}
For $q_3 \neq 0$ \textup{(}generic period matrix\textup{)},
the off-diagonal entries are nonzero and the Fredholm
determinant does not factor as a product of genus-$1$
determinants.
\end{proposition}

\begin{proof}
The block decomposition and Hilbert--Schmidt bounds are
established in Vol~II (genus-$2$ Fredholm determinant
theorem). Non-factorization for $q_3 \neq 0$ follows from
the leading off-diagonal correction: writing
$r = e^{2\pi i z}$, the expansion gives
$\det(1 - K_2)^{-1} = \eta(\tau_1)^{-1}\,\eta(\tau_2)^{-1}
\cdot (1 + c_{1,1}\,r + \cdots)$ with
$c_{1,1} = \sum_{k \geq 0}
q_1^{k+1}q_2^{k+1}/[(1{-}q_1^{k+1})(1{-}q_2^{k+1})]$,
a sum of positive terms.
\end{proof}

\begin{theorem}[Genus-$2$ sewing--Hecke non-collapse;
\ClaimStatusProvedHere]%
\label{thm:genus2-non-collapse}%
\index{sewing--Hecke collapse!genus-2 failure|textbf}%
The sewing--Hecke collapse
\textup{(}Theorem~\textup{\ref{thm:sewing-hecke-reciprocity}}\textup{)}
does not extend to genus~$2$. Specifically, for the
Heisenberg algebra~$\cH$:
\begin{enumerate}[label=\textup{(\roman*)}]
\item The connected free energy
 $F_2^{\mathrm{conn}}(\cH;\,\Omega)$ depends
 nontrivially on the off-diagonal period~$z$, with
 leading correction of order $|e^{2\pi i z}|$ from the
 Schur complement of~$K_2$.
\item The sewing kernel $K_2$ is not Fock-diagonal:
 the off-diagonal blocks $R_{12}, R_{21}$ couple Fock
 modes across handles through the internal seam.
\item The Siegel Eisenstein series
 $E_k(\Omega, s)$ on $\mathrm{Sp}_4(\bZ)$ is irreducible:
 it does not factor as $E_{k_1}(\tau_1, s_1)
 \cdot E_{k_2}(\tau_2, s_2)$. The Rankin--Selberg
 unfolding on $\mathrm{Sp}_4(\bZ) \backslash \mathfrak{H}_2$
 therefore cannot reduce to a product of genus-$1$
 unfoldings.
\end{enumerate}
\end{theorem}

\begin{proof}
(i)~The Fredholm determinant expansion
(Proposition~\ref{prop:genus2-non-diagonal}) gives the
$z$-dependent correction with $c_{1,1} > 0$.
(ii)~The off-diagonal operators $R_{ij}$ arise from
the composition of sewing operators through the internal
seam; their Hilbert--Schmidt norm is bounded below by
$c_{1,1} > 0$ for $q_3 \neq 0$.
(iii)~The Siegel Eisenstein series is induced from the
Siegel parabolic of $\mathrm{Sp}_4$; its Fourier expansion
involves sums over half-integral positive semi-definite
$2 \times 2$ matrices, which do not factor into products of
$1 \times 1$ sums.
\end{proof}

\subsection{The Saito--Kurokawa bridge}%
\label{subsec:saito-kurokawa-bridge}

\begin{remark}[Spectral decomposition at genus $2$]%
\label{rem:saito-kurokawa-bridge}%
\index{Saito--Kurokawa lift!arithmetic bridge}%
The spectral decomposition of
$L^2(\mathrm{Sp}_4(\bZ) \backslash \mathfrak{H}_2)$
includes three types of contributions:
\begin{enumerate}[label=\textup{(\roman*)}]
\item Siegel Eisenstein series (continuous spectrum);
\item Saito--Kurokawa lifts $\mathrm{SK}(f)$ from elliptic
 cusp forms $f \in S_k(\mathrm{SL}(2,\bZ))$ (cuspidal,
 but arising from genus-$1$ data);
\item genuine genus-$2$ Siegel cusp forms (cuspidal,
 with no genus-$1$ origin).
\end{enumerate}
The Saito--Kurokawa component provides an arithmetic bridge.
The genus-$2$ partition function~$Z_\cA^{(2)}(\Omega)$,
projected onto $\mathrm{SK}(f)$, gives a Petersson pairing
$\langle Z_\cA^{(2)},\, \mathrm{SK}(f) \rangle$
involving the genus-$1$ Hecke eigenvalues of~$f$.
The genus-$2$ MC equation constrains this pairing through the
three-shell formula
$\Theta^{(2)} = \Theta^{(2)}_{\mathrm{loop}^2}
+ \Theta^{(2)}_{\mathrm{sep}\circ\mathrm{loop}}
+ \Theta^{(2)}_{\mathrm{pf}}$
(Vol~II, genus-$2$ partition function computation),
and these constraints involve the off-diagonal
shadow data absent at genus~$1$.

Whether the genus-$2$ MC constraint, combined with the
genus-$1$ constraint, suffices to determine \emph{eigenvalues}
as well as weights is the central open problem of the
higher-genus arithmetic programme.
\end{remark}

\begin{remark}[$\Delta_5$ as automorphic target of twisted
one-loop denominators]
\label{rem:arith-shad-5}%
\index{Delta5 weight@$\Delta_5$ weight-$5$ cusp form!1-loop SUGRA origin}%
\index{11D supergravity!Omega-deformed K3 x T2}%
\index{M5-brane!Kodaira I1 fibre}%
\index{Delta5 four duality frames!K3 BKM}%
The weight-$5$ cusp form $\Delta_5$ on the paramodular group of
level~$1$ is fixed here by the automorphic construction, not by an
unproved equality with a chiral conductor. Gritsenko--Nikulin and
Borcherds identify the K3 weak-Jacobi input with a Borcherds product
whose weight is half the constant term of the input:
\begin{equation}\label{eq:delta5-1loop-weight}
 \operatorname{wt}(\Delta_5)
 \;=\;
 \tfrac12\,c_{\phi_{0,1}^{K3}}(0)
 \;=\;
 \tfrac{10}{2}
 \;=\;
 5,
\end{equation}
in the normalisation used by the K3 elliptic genus. This is the
proved automorphic input. The physical comparison is a separate
proof obligation: the twisted $11$D supergravity of
Costello--Gaiotto \cite{CostelloGaiotto2021}, compactified on
$\mathrm{K3} \times T^2$ with $\Omega$-deformation and $24$
$\mathrm{M}5$-branes on the Kodaira $I_1$ fibres of an elliptic
$\mathrm{K3}$, must produce the same automorphic product as its
one-loop BPS denominator on the same determinant line. Until that
determinant is constructed, the supergravity sentence does not prove
the weight formula~\eqref{eq:delta5-1loop-weight}.

The four duality frames name the target denominator in their native
theories; they are not four independent computations of a chiral
invariant:
\begin{enumerate}[label=\textup{(\roman*)}]
\item \emph{Heterotic} on $T^6$ (Harvey--Moore
\cite{HarveyMoore1996}): the $T^6$-dual threshold integral on
$\Gamma^{2,2}$ produces $-\log(\Delta_5(\Omega))$ as the $1$-loop
threshold correction.
\item \emph{Type~IIB} on $\mathrm{K3} \times T^2$ (Dijkgraaf--Verlinde--
Verlinde \cite{DijkgraafVerlindeVerlinde1996,DijkgraafVerlindeVerlinde1997}):
the $D1$--$D5$ partition function on $\mathrm{Sym}^N(\mathrm{K3})$
produces $\Delta_5$ as the large-$N$ generating function of
second-quantised $\mathrm{K3}$ BPS states.
\item \emph{Type~IIA} on $\mathrm{K3} \times T^2$ (Dijkgraaf--Moore--Verlinde--
Verlinde \cite{DijkgraafMooreVerlindeVerlinde1997}): the elliptic
genus product formula
$\prod_{m, n, \ell}(1 - p^m q^n y^\ell)^{c(4mn - \ell^2)}$ with
$\sum c(4mn - \ell^2)\,q^n y^\ell = 2\,\mathrm{EG}(\mathrm{K3})$
equals $\Delta_5^{-1}$ up to the Hodge anomaly prefactor.
\item \emph{M-theory} on $\mathrm{K3} \times T^2$ (Witten
\cite{Witten1995,Kachru2017ms}): the $1$-loop partition function of
$\mathrm{M5}$-branes wrapping $\mathrm{K3}$ and a $T^2$-cycle,
with Omega-background regularising the transverse $\bR^4$, gives the
physics origin of the same BPS denominator.
\end{enumerate}
The Vol~I conductor $K^{\kappa_{\mathrm{ch}}}=8$ and the Vol~III
Borcherds weight $\operatorname{wt}(\Delta_5)=5$ remain distinct.
The missing theorem is a normalised one-loop determinant computation
identifying the physics partition function with the Borcherds
product in exactly the convention of~\eqref{eq:delta5-1loop-weight}.
Until that determinant is written on the same line bundle, the
supergravity comparison is a target identification, not a derivation
of the K3/Mukai constants.
\end{remark}

\begin{remark}[Honest assessment of the genus-$2$ frontier]%
\label{rem:genus2-honest-assessment}%
\index{genus-2 arithmetic frontier!honest assessment}%
The genus-$2$ non-collapse
(Theorem~\ref{thm:genus2-non-collapse}) opens a genuine
direction, but several barriers must be acknowledged:
\begin{enumerate}[label=\textup{(\roman*)}]
\item The Rankin--Selberg method on $\mathrm{Sp}_4$ is
 harder than on $\mathrm{SL}(2)$: the unfolding involves
 a non-unique Whittaker model, and the resulting
 $L$-functions (spin, standard) have more complex analytic
 properties.
\item The genus-$2$ MC equation involves the full shadow
tower $\Theta_\cA$ together with the scalar
projection~$\kappa$. Extracting constraints requires
 the non-scalar shadows ($\mathfrak{C}$, $\mathfrak{Q}$)
 at genus~$2$.
\item The Saito--Kurokawa bridge is informative only if
 the genus-$2$ Petersson projection is nonzero, a question
 about non-vanishing of Fourier coefficients of Siegel
 modular forms.
\item The comparison between the spin $L$-function
 $L^{\mathrm{spin}}(s, \mathrm{SK}(f))$ and the symmetric
 square $L(s, \mathrm{Sym}^2 f)$ is known
 (Andrianov--Zagier), so the theoretical pipeline exists.
 The question is whether the MC input is strong enough.
\end{enumerate}
The structural advantage of genus~$2$: the
sewing--Hecke collapse fails
(Theorem~\ref{thm:genus2-non-collapse}), making the arithmetic
interface multi-dimensional. Whether this
suffices to break the structural separation
(Theorem~\ref{thm:structural-separation}) is the central open
problem. The first step: compute the
genus-$2$ Heisenberg Rankin--Selberg integral and
verify that the Saito--Kurokawa projection is nonzero.
\end{remark}

\subsection{The B\"ocherer bridge: critical-line access via genus~$2$}%
\label{subsec:bocherer-bridge-theorem}%
\index{B\"ocherer bridge|textbf}%
\index{central L-value!genus-2 MC access}%
\index{structural separation!genus-2 bypass}

The genus-$1$ unfolding-erasure computation
(Proposition~\ref{prop:scattering-residue} and
Remark~\ref{rem:structural-obstruction-revisited})
shows that genus-$1$ MC constraints are confined to
$\operatorname{Re}(s) \geq 1$, while zeta zeros lie in
$0 < \operatorname{Re}(s) < 1$.
The genus-$2$ non-collapse (Theorem~\ref{thm:genus2-non-collapse})
opens a new channel. We show that this channel reaches the
\emph{critical line} $\operatorname{Re}(s) = 1/2$ through the
B\"ocherer--Furusawa--Morimoto factorization.

\begin{proposition}[Lattice cusp-form non-vanishing at genus~$2$;
\ClaimStatusProvedHere]%
\label{prop:leech-cusp-nonvanishing}%
\index{Leech lattice!genus-2 cusp component}%
Let $\Lambda$ be an even unimodular lattice of rank $r \geq 24$
with minimal norm~$\mu(\Lambda) \geq 4$
\textup{(}e.g., the Leech lattice\textup{)}.
The genus-$2$ theta series
$\Theta_\Lambda^{(2)} \in M_{r/2}(\mathrm{Sp}(4,\bZ))$
differs from the Siegel Eisenstein series:
\begin{equation}\label{eq:leech-cusp-nonvanishing}
\Theta_\Lambda^{(2)} \;-\; E_{r/2}^{(2)}
\;\neq\; 0,
\end{equation}
where $E_{r/2}^{(2)}$ is the normalized Siegel Eisenstein series.
The difference has nonzero Klingen Eisenstein and/or cusp
components; the cusp projection is computed explicitly in
Theorem~\textup{\ref{thm:leech-chi12-projection}}.
\end{proposition}

\begin{proof}
The Siegel Eisenstein series $E_k^{(2)}$ has strictly positive
Fourier coefficients at every positive definite half-integral
$T > 0$ \cite{Kitaoka93}. Take
$T_0 = \bigl(\begin{smallmatrix} 1 & 0 \\ 0 & 1 \end{smallmatrix}\bigr)$.
Then $a_{E_k}(T_0) > 0$, but
$r_2(\Lambda, T_0) = 0$ since $\Lambda$ has no vectors of norm~$2$.
Therefore
$\Theta_\Lambda^{(2)} - E_{r/2}^{(2)}$ has a negative Fourier
coefficient at~$T_0$.
\end{proof}

\begin{theorem}[B\"ocherer bridge;
\ClaimStatusProvedHere]%
\label{thm:bocherer-bridge}%
\index{B\"ocherer bridge!critical-line access}%
Let $\Lambda$ be an even unimodular lattice of rank~$r$, and let
$f \in S_{r-2}(\mathrm{SL}(2,\bZ))$ be a Hecke eigenform whose
Saito--Kurokawa lift
$\mathrm{SK}(f) \in S_{r/2}(\mathrm{Sp}(4,\bZ))$
has nonzero Petersson pairing with the cusp-form component
of~$\Theta_\Lambda^{(2)}$. Then:
\begin{enumerate}[label=\textup{(\roman*)}]
\item \textup{(MC determination)}
The genus-$2$ MC equation
\textup{(}$D^2 = 0$ restricted to genus~$2$, giving the three-shell
decomposition of
Remark~\textup{\ref{rem:saito-kurokawa-bridge}}\textup{)}
determines $\Theta_\Lambda^{(2)}$ from genus-$1$ shadow data
and boundary contributions.

\item \textup{(Critical-line access)}
By the refined B\"ocherer conjecture
\textup{(}proved: Dickson--Pitale--Saha--Schmidt
\cite{DPSS20} for full level;
Furusawa--Morimoto \cite{FurusawaMorimoto14}
for paramodular forms\textup{)},
the B\"ocherer coefficient
\begin{equation}\label{eq:bocherer-coefficient-def}
B_f(D)
\;:=\;
\sum_{\substack{[T] \in \mathrm{GL}(2,\bZ) \backslash \mathcal{T} \\
 \operatorname{disc}(T) = D}}
\frac{a_{\mathrm{cusp}}(T)}{\varepsilon(T)}
\end{equation}
satisfies the factorization
\begin{equation}\label{eq:bocherer-factorization}
\frac{|B_f(D)|^2}{\langle \mathrm{SK}(f),\, \mathrm{SK}(f) \rangle}
\;=\;
\alpha(r,k)\;
L\!\left(\tfrac12,\, \pi_f\right)\,
L\!\left(\tfrac12,\, \pi_f \times \chi_D\right),
\end{equation}
where $\alpha(r,k) > 0$ is an explicit archimedean factor,
$\pi_f$ is the automorphic representation of~$f$,
and $D < 0$ is a fundamental discriminant.
Both $L$-values are at the central point $s = 1/2$.
For Saito--Kurokawa lifts $F = \mathrm{SK}(f)$, the spin
$L$-function factors and~\eqref{eq:bocherer-factorization}
reduces to the Waldspurger formula
$\lvert B_f(D)\rvert^2 \propto L(k{-}1,\, f \times \chi_D)
\cdot \lvert D\rvert^{k-2}$,
which does not involve $L(k{-}1,f)$ itself
\textup{(}the latter vanishes when
$\varepsilon(f) = -1$\textup{)}.

\item \textup{(Structural bypass)}
The composition
\[
\text{MC}_{g=2}
\;\xrightarrow{\;\text{three-shell}\;}
a_\Lambda(T)
\;\xrightarrow{\;\text{B\"ocherer}\;}
L\!\left(\tfrac12,\, \pi_f \times \chi_D\right)
\]
provides algebraic access to the critical line, bypassing
the genus-$1$ unfolding-erasure obstruction
\textup{(}Proposition~\textup{\ref{prop:scattering-residue}}\textup{)}.
\end{enumerate}
\end{theorem}

\begin{proof}
Part~(i) follows from the MC equation at genus~$2$:
$D^2 = 0$ on the genus-$2$ bar complex constrains
$\Theta^{(2)}$ through the boundary stratification
$\partial\bar{\cM}_2 = \delta_0 \cup \delta_1$
(separating clutching and non-separating self-sewing),
giving the three-shell decomposition.
Part~(ii) is the refined B\"ocherer conjecture,
proved unconditionally for $\mathrm{Sp}(4,\bZ)$ in
\cite{DPSS20}.
Part~(iii): the genus-$1$ unfolding-erasure computation shows
the Rankin--Selberg integral produces
$\zeta(s{-}1)\zeta(s)$ with zero information confined to
$\operatorname{Re}(s) \geq 1$.
The B\"ocherer factorization~\eqref{eq:bocherer-factorization}
converts the MC-determined Fourier coefficients directly to
$L(1/2, \pi_f \times \chi_D)$.
\end{proof}

\begin{remark}[Escalation principle and remaining gap]%
\label{rem:bocherer-escalation}%
\index{B\"ocherer bridge!escalation to all genera}%
\index{symmetric power L-functions!genus-g access}%
\begin{enumerate}[label=\textup{(\roman*)}]
\item \textup{(Escalation)} At genus~$g$, conjectural
 B\"ocherer-type formulas connect degree-$g$ Fourier
 coefficients to $L(1/2, \pi_F)$ for automorphic
 representations of~$\mathrm{GSp}(2g)$.
 Via the symmetric power transfer
 $\mathrm{Sym}^{g-1}: \mathrm{GL}(2) \to \mathrm{GL}(g)$
 (Newton--Thorne: unconditional for holomorphic eigenforms
 of regular algebraic weight; conditional for Maass forms
 at $g \geq 6$),
 this accesses $L(1/2, \mathrm{Sym}^{g-1} f)$.
 The MC equation at \emph{all} genera thereby encodes the
 algebraic data that the Serre reduction requires:
 Ramanujan would follow from all-genera symmetric-power
 non-vanishing, conditional on Langlands functoriality
 $\mathrm{GL}(2) \to \mathrm{GL}(g)$ for $g \geq 5$
 \textup{(}see Remark~\textup{\ref{rem:escalation-circularity}}
 for the Newton--Thorne resolution in the holomorphic
 case\textup{)}.
\item \textup{(Remaining gap)}
 Access $\neq$ control.
 The MC equation \emph{determines} central $L$-values
 but does not by itself \emph{force} positivity or
 non-vanishing.
 Whether bracket positivity
 (Proposition~\ref{prop:bracket-hodge-index})
 propagates through the B\"ocherer factorization is the
 sharpened genus-$2$ form of Gap~A.
\item \textup{(Target comparison)}
Theorem~\ref{thm:leech-chi12-projection} below verifies the nonzero
projection for the Leech lattice. This is a scalar target comparison:
it supplies a cuspidal Fourier coefficient on the automorphic side,
not a completed source-level Igusa, Humbert, or Lie-bialgebra
classification. Upgrading it requires a comparison functor from the
bar--cobar MC tower to the relevant automorphic/Hall--Borcherds
source, with parity, root, and no-extra-source clauses fixed.
\end{enumerate}
\end{remark}

\begin{remark}[Definitive scope of genus-$2$ arithmetic access]%
\label{rem:genus2-definitive-scope}%
\ClaimStatusProvedHere%
\index{genus-2!definitive arithmetic scope}%
\index{B\"ocherer bridge!what genus-2 can and cannot detect}%
\index{Satake parameters!genus-2 cannot constrain}%
The B\"ocherer bridge determines \emph{what} genus-$2$
MC data can and cannot access. Three structural features
control the outcome.

\begin{enumerate}[label=\textup{(\roman*)}]
\item \textup{(Schur complement structure.)}
For Heisenberg, the genus-$2$ sewing is Fock-diagonal,
so the genus-$2$ non-collapse
\textup{(}Theorem~\textup{\ref{thm:genus2-non-collapse}}\textup{)}
produces no constraints beyond the genus-$1$ data:
$\Theta_\cH^{(2)}$ is determined entirely by
$\Theta_\cH^{(1)}$.
For interacting algebras at weight $\geq 4$
\textup{(}$\dim V_4 \geq 2$\textup{)},
genuine quadratic OPE constraints appear in the
genus-$2$ three-shell decomposition, but these
constrain \emph{OPE structure constants}, not
Satake parameters.

\item \textup{(Unfolding erasure.)}
The Siegel and Klingen Eisenstein channels of the genus-$2$
Rankin--Selberg integral erase the poles of the scattering
matrix upon unfolding, for the same reason as in genus~$1$
\textup{(}the Eisenstein residues produce completed
$L$-functions that are entire in the critical strip\textup{)}.
The B\"ocherer channel \emph{bypasses} unfolding: it
extracts central $L$-values directly from Fourier
coefficients via the Maass
relations, without an integral transform. This is why the
genus-$1$ unfolding-erasure obstruction
\textup{(}Proposition~\textup{\ref{prop:scattering-residue}}\textup{)}
is evaded.

\item \textup{(What genus-$2$ determines.)}
The MC equation at genus~$2$ determines the partition
function weights $c_j$ \textup{(}how much each
Hecke eigenform contributes to $\Theta_\cA^{(2)}$\textup{)},
but \emph{not} the eigenvalue magnitudes
$|\alpha_{f,p}|$ at individual primes.
Ramanujan is a per-prime statement
\textup{(}$|\alpha_{f,p}| = p^{(k-1)/2}$ for each prime~$p$\textup{)};
genus~$2$ constrains the global $L$-value
$L(1/2, f \times \chi_D)$ via the B\"ocherer
factorization, which is an average over primes weighted
by the Euler product.
\end{enumerate}
\end{remark}

\begin{remark}[Leech: all genus-$2$ cusp forms are
Saito--Kurokawa lifts]%
\label{rem:leech-all-sk}%
\ClaimStatusProvedHere%
\index{Leech lattice!genus-2 cusp forms all SK}%
\index{Saito--Kurokawa lift!Leech lattice}%
For the Leech lattice at weight $k = r/2 = 12$:
$\dim S_{12}(\mathrm{Sp}(4,\bZ)) = 1$, and the unique
cusp form $\chi_{12} = \mathrm{SK}(f_{22})$ is a
Saito--Kurokawa lift.
Genuinely new Siegel eigenforms \textup{(}not
lifts from $\mathrm{GL}(2)$\textup{)} first appear
at weight $k = 20$, corresponding to lattice
rank $r \geq 40$.
The Leech genus-$2$ arithmetic is therefore entirely
controlled by the genus-$1$ eigenform $f_{22}$
through the Saito--Kurokawa mechanism;
genuinely genus-$2$ phenomena require higher-rank
lattices.
\end{remark}

\begin{theorem}[Leech $\chi_{12}$-projection and Waldspurger
consequence; \ClaimStatusProvedHere]%
\label{thm:leech-chi12-projection}%
\index{Leech lattice!chi12 projection@$\chi_{12}$ projection}%
\index{Waldspurger formula!MC application}%
Let $\Lambda_{24}$ denote the Leech lattice.
\begin{enumerate}[label=\textup{(\roman*)}]
\item \textup{(Cusp space dimension)}
$\dim S_{12}(\mathrm{Sp}(4,\bZ)) = 1$, spanned by the
Igusa cusp form
$\chi_{12} = \mathrm{SK}(f_{22})$
\textup{(}Saito--Kurokawa lift of the unique weight-$22$
Hecke eigenform $f_{22} = \Delta \cdot E_{10}$\textup{)}.

\item \textup{(Nonzero projection)}
The genus-$2$ theta series decomposes as
\begin{equation}\label{eq:leech-chi12-decomposition}
\Theta_{\Lambda_{24}}^{(2)}
\;=\;
E_{12}^{(2)}
\;+\;
c_1\,E_{12}^{\mathrm{Kling}}(\Delta)
\;+\;
c_2\,\chi_{12},
\end{equation}
where the coefficient $c_2 \neq 0$ is determined by the Igusa
relation $\chi_{12} = C^{-1}(441\,E_4^3 + 250\,E_6^2 -
691\,E_{12}^{(2)})$ and the genus-$2$ representation numbers
$r_2(\Lambda_{24}, T)$ at three linearly independent matrices.
Numerically, $c_2 \approx -1.918 \times 10^{-6}$
\textup{(}exact rational computation, zero residuals\textup{)}.

\item \textup{(Waldspurger consequence)}
Since $\chi_{12} = \mathrm{SK}(f_{22})$, the B\"ocherer
coefficients satisfy
$B(D;\, \chi_{12}) = c(\lvert D\rvert)$
where $c(N)$ are the half-integral weight form coefficients.
By Waldspurger's theorem:
\begin{equation}\label{eq:waldspurger-leech}
\lvert c(\lvert D\rvert)\rvert^2
\;\propto\;
L(11,\, f_{22} \times \chi_D)
\cdot \lvert D\rvert^{10}
\end{equation}
for fundamental $D < 0$. The functional equation sign
$\varepsilon(f_{22}) = -1$ forces $L(11,\, f_{22}) = 0$,
but the twisted sign
$\varepsilon(f_{22} \times \chi_D) = +1$ for $D < 0$
\textup{(}since $\chi_D(-1) = -1$\textup{)},
so $L(11,\, f_{22} \times \chi_D) \neq 0$ generically.

\item \textup{(MC chain)}
The composition
\[
D_\cA^2 = 0
\;\xrightarrow{\;\text{genus-}2\;\text{three-shell}\;}
\Theta_{\Lambda_{24}}^{(2)}
\;\xrightarrow{\;c_2 \neq 0\;}
B(D)
\;\xrightarrow{\;\text{Waldspurger}\;}
L(11,\, f_{22} \times \chi_D)
\]
determines $L(11,\, f_{22} \times \chi_D)$
for all fundamental $D < 0$ from MC data.
This is the first concrete instance of homotopy-algebraic
data producing number-theoretic output on the critical line.
\end{enumerate}
\end{theorem}

\begin{proof}
Part~(i): dimension formula of Tsuyumine--Igusa
\cite{DPSS20};
$\chi_{12}$ is an Igusa generator of the graded ring
$M_*(\mathrm{Sp}(4,\bZ))$.
Part~(ii): the basis
$\{E_4^3,\, E_6^2,\, E_{12}^{(2)}\}$ of
$M_{12}(\mathrm{Sp}(4,\bZ))$ is linearly independent
(the Igusa relation produces the nonzero cusp form
$\chi_{12}$ from their span).
The three genus-$2$ representation numbers
$r_2(\Lambda_{24}, T)$ at $T = \mathrm{diag}(1,1)$,
$\mathrm{diag}(2,2)$,
$\bigl(\begin{smallmatrix}2 & 1/2 \\ 1/2 & 2\end{smallmatrix}\bigr)$
are
$0$, $18{,}309{,}564{,}000$, $9{,}258{,}762{,}240$
respectively
(the first vanishes because $\Lambda_{24}$ has no
norm-$2$ vectors; the others follow from the known
inner-product distribution of minimal vectors).
The Siegel Eisenstein coefficients at these matrices are
computed from the Cohen--Katsurada formula.
Solving the resulting $3 \times 3$ system and projecting along
the Igusa direction $(441,\, 250,\, -691)$ gives
$c_2 \neq 0$.
Part~(iii): the Maass relations for SK lifts give
$B(D;\, \chi_{12}) = c(\lvert D\rvert)$; the Waldspurger
formula \cite{Waldspurger81} gives the
proportionality~\eqref{eq:waldspurger-leech}.
Part~(iv) is the composition of~(ii) with~(iii) and
Theorem~\ref{thm:bocherer-bridge}.
\end{proof}

\begin{remark}[Conditional $H^2(\mathfrak{g}_{\Delta_5})^{\bZ/2,
\,K(1)}$ comparison with the Igusa class]
\label{rem:arith-shad-2}%
\index{classification!H2 invariants of Delta5 family}%
\index{Igusa cusp form!gauge-class invariant}%
\index{Etingof-Kazhdan!quantisation classification series}%
\index{Delta5 gauge-class classification!K3 BKM}%
The MC chain of Theorem~\ref{thm:leech-chi12-projection}(iv)
constructs, at genus~$2$ on the Leech family, a one-dimensional
space of cuspidal shadow classes. The stronger Igusa classification
is a conditional recognition statement.
\begin{proposition}[Igusa $H^2$ comparison criterion;
\ClaimStatusConditional]\label{thm:h2-classification-delta5}%
Let $\mathfrak{g}_{\Delta_5}$ be a source Lie algebra attached to the
Koszul pair $(\cA_{\Delta_5},\cA_{\Delta_5}^!)$ on the Lorentzian
lattice $\Lambda=2U\oplus E_8(-1)$, with Fricke $\bZ/2$-action. Assume:
\begin{enumerate}[label=\textup{(R\arabic*)}]
\item a finite Hall--Borcherds recognition datum constructs
 $\mathfrak{g}_{\Delta_5}$ from compact source matrices and proves
 the Vol~III source-matrix faithfulness condition;
\item the parity/root fixture identifies the signed product exponents
 with the source super-root grading;
\item the $K(1)$-local comparison map from source Lie-bialgebra
 cohomology to the Bruinier Heegner Chern class is defined; and
\item the PBW/no-extra-source clause excludes additional
 genus-$2$ cuspidal classes.
\end{enumerate}
Then the target identification is
\begin{equation}\label{eq:h2-classification-delta5}
 H^2\!\bigl(\mathfrak{g}_{\Delta_5}\bigr)^{\bZ/2,\, K(1)}
 \;\cong\;
 \bC \cdot \Delta_5,
\end{equation}
with the generator given by the weight-$5$ Igusa cusp form viewed
as the automorphic Heegner-divisor class transported through the
recognition datum.
\end{proposition}
\begin{proof}
Bruinier's Chern-class formula \cite[Prop.~5.1]{Bruinier2002}
identifies the automorphic Chern class of the weight-$5$ Borcherds
product with the primitive Heegner-divisor class on the Humbert
surface. Assumptions \textup{(R1)--(R3)} provide the source object and
the comparison map carrying this automorphic class into
$K(1)$-local Lie-bialgebra cohomology. Assumption~\textup{(R4)} gives
one-dimensionality. Without \textup{(R1)--(R4)} the displayed line is
only the target of the comparison, not a theorem about a constructed
source cohomology group.
\end{proof}
The proof criteria for the Vol~I arithmetic shadow comparison are
therefore:
\begin{enumerate}[label=\textup{(\roman*)}]
\item construct the source algebra and its $K(1)$-local
 cochain complex;
\item prove that the Fricke-invariant source cohomology is
 one-dimensional and has no hidden compact-source summands;
\item transport the Bruinier Humbert class to that complex with the
 parity/root normalisation fixed; and
\item only then read identities involving $\Delta_5$ as identities in
 Lie-bialgebra \(H^2\). Before these four steps, they are
 automorphic scalar or divisor-class comparisons.
\end{enumerate}
See Part~\ref{part:characteristic-datum} for the chiral-Lie-algebra
framework and Part~\ref{part:v1-frontier} for the Borcherds
programme where \eqref{eq:h2-classification-delta5} is the comparison
target.
\end{remark}

\begin{remark}[Chriss--Ginzburg realization of the
B\"ocherer bridge]%
\label{rem:chriss-ginzburg-bocherer}%
\index{Chriss--Ginzburg architecture!B\"ocherer bridge}%
\index{shadow tautological classes!genus-2 arithmetic}%
The MC chain of
Theorem~\ref{thm:leech-chi12-projection}\textup{(iv)}
is a projection of the single master identity
$\Theta_\cA \in \mathrm{MC}(\mathfrak{g}_\cA^{\mathrm{mod}})$
\textup{(}Corollary~\textup{\ref{cor:theta-twisting-morphism}}\textup{)}
through the Chriss--Ginzburg architecture:
\begin{enumerate}[label=\textup{(\roman*)}]
\item \textup{(Convolution algebra)}
$D_\cA^2 = 0$ on the modular cyclic deformation complex
$\mathrm{Def}_{\mathrm{cyc}}^{\mathrm{mod}}(\cA)$
\textup{(}Definition~\textup{\ref{def:modular-cyclic-deformation-complex}}\textup{)}
produces the MC element
$\Theta_\cA = D_\cA - d_0$.

\item \textup{(Genus-$2$ three-shell projection)}
The genus-$2$ component
$\Theta_\cA^{(2)} \in
\mathfrak{g}_\cA^{\mathrm{mod}} \otimes \hbar^2$
is the unique solution of the inhomogeneous equation
\[
d_{\Theta^{(0)}}\!\bigl(\Theta^{(2)}\bigr)
\;=\;
\underbrace{-\tfrac12[\Theta^{(1)},\Theta^{(1)}]}_
 {\text{loop}^2}
\underbrace{{}-\Delta(\Theta^{(1)}) - d_{\mathrm{sew}}^{(2)}}_
 {\text{sep}\circ\text{loop}}
\underbrace{{}-d_{\mathrm{pf}}^{(2)}}_
 {\text{planted forest}},
\]
where $d_{\Theta^{(0)}}$ is the genus-$0$ twisted differential,
the bracket is the Lie bracket in
$\mathfrak{g}_\cA^{\mathrm{mod}}$, $\Delta$ is the
non-separating BV operator, and
$d_{\mathrm{pf}}^{(2)}$ is the planted-forest correction
\textup{(}$\delta_{\mathrm{pf}}^{(2,0)} =
S_3(10S_3 - \kappa)/48$,
Remark~\textup{\ref{rem:pixton-genus2-computation}}\textup{)}.
The three terms on the right are the \emph{three shells}
of the genus-$2$ obstruction class.

\item \textup{(Shadow tautological descent)}
The shadow tautological class
$\tau_{2,n}(\cA) \in R^*(\bar{\cM}_{2,n+1})$
\textup{(}Definition~\textup{\ref{def:shadow-tautological-classes}}\textup{)}
evaluated at $n = 0$ gives the genus-$2$ partition function
through the sewing--shadow intertwining
\textup{(}Theorem~\textup{\ref{thm:sewing-shadow-intertwining}}\textup{)}.
For $\cA = V_{\Lambda_{24}}$, this produces
$\Theta^{(2)}_{\Lambda_{24}}$, a weight-$12$ Siegel modular
form on $\mathrm{Sp}(4,\bZ)$.

\item \textup{(Igusa projection $=$ shadow extraction)}
The projection
$c_2 = \pi_{\chi_{12}}(\Theta^{(2)}_{\Lambda_{24}})$
\textup{(}Corollary~\textup{\ref{cor:shadow-extraction}}\textup{)}
extracts the cusp component. In the Chriss--Ginzburg
dictionary, $c_2 \neq 0$ means the genus-$2$ shadow
$\mathrm{Sh}_{2,0}^{\mathrm{cusp}}(\Theta_\cA)$ is nonzero,
i.e.\ the MC element has a \emph{cuspidal genus-$2$ shadow}.

\item \textup{(Arithmetic output)}
The Waldspurger formula converts the cuspidal shadow into
$L(11, f_{22} \times \chi_D)$ for all fundamental $D < 0$.
In the Chriss--Ginzburg language: the shadow algebra
$\cA^{\mathrm{sh}} = H_\bullet(\mathrm{Def}_{\mathrm{cyc}}^{\mathrm{mod}})$
\textup{(}Definition~\textup{\ref{def:shadow-algebra}}\textup{)}
encodes, at its genus-$2$ cusp-form graded piece,
the complete family of twisted central $L$-values.
The bracket positivity
\textup{(}Proposition~\textup{\ref{prop:bracket-hodge-index}}\textup{)}
acts on this graded piece, constraining the $L$-values through
the Hodge index on $\mathrm{Def}_{\mathrm{cyc}}$.
\end{enumerate}
The entire chain
$D^2{=}0 \to \Theta_\cA \to \Theta^{(2)} \to c_2 \to
L(11, f_{22} \times \chi_D)$
is a sequence of \emph{projections of the single MC element}
through the layers of the Chriss--Ginzburg architecture.
The extra analytic inputs are exactly the genus-$2$ theta
identification, the non-zero $\chi_{12}$ projection, and the
Waldspurger/Saito--Kurokawa comparison. The MC element supplies the
source shadow class; it does not by itself prove the automorphic
central-value formula.
\end{remark}

% ============================================================
\section{The prime-locality assessment}%
\label{sec:prime-locality-assessment}%
\index{prime-locality!definitive assessment|textbf}%
\index{prime-locality!obstructions|textbf}

This section distinguishes what is proved from what is
structurally obstructed in the prime-locality programme.

\subsection{What prime-locality is}

Prime-locality
(Conjecture~\ref{conj:prime-locality-transfer})
asks whether the shadow coefficients
$\operatorname{Sh}_r^{(1)}(\Theta_\cA;\tau)$, viewed as
functions on $\cM_{1,1}$, decompose into Hecke
eigencomponents with multiplicative Fourier coefficients.
In the convolution-algebra language: whether the Hecke
operators $T_p$ on $\cM_{1,1}$ extend to chain endomorphisms
of $\gAmod$ that commute with the MC differential~$D_\cA$.
The obstruction is the Hecke defect class
$[\delta_p(\cA)] \in H^1(\gAmod, \mathrm{ad}_{T_p})$
(Definition~\ref{def:hecke-defect-class}).
Prime-locality holds if and only if this class vanishes
for all primes~$p$.

\subsection{Where prime-locality is proved or reduced}

\begin{proposition}[Settled and reduced cases; \ClaimStatusConditional]%
\label{prop:prime-locality-proved-cases}%
\index{prime-locality!proved cases|textbf}
Prime-locality is proved or reduced in three regimes:
\begin{enumerate}[label=\textup{(\roman*)}]
\item \emph{Lattice VOAs} \textup{(}all ranks, all levels\textup{):}
 the shadow amplitudes factor through the theta function
 $\Theta_\Lambda \in M_{r/2}(\Gamma)$, which decomposes
 into Hecke eigenforms. The Hecke operators act as ring
 homomorphisms on $M_{r/2}$, giving
 $[\delta_p(V_\Lambda)] = 0$ exactly
 \textup{(}Proposition~\textup{\ref{prop:hecke-defect-lattice}},
 Theorem~\textup{\ref{thm:hecke-newton-lattice}}\textup{)}.

\item \emph{Rational VOAs with diagonal modular invariant
\textup{(}conditional at amplitude level\textup{):}}
 the characters $\chi_i$ form a vector-valued modular form
 for $\rho\colon \mathrm{SL}(2,\bZ) \to \mathrm{GL}(N,\bC)$.
 The Franc--Mason Hecke operators act on the character space; the
 conditional theorem requires eigencomponent multiplicativity and
 compatibility with shadow graph sums
 \textup{(}Theorem~\textup{\ref{thm:non-lattice-ramanujan}}\textup{)}.
 The theta-bridge
 \textup{(}Proposition~\textup{\ref{prop:theta-bridge-rational}}\textup{)}
 resolves the non-multiplicativity of the scalar Dirichlet
 series as a projection artefact: the sum of individually
 multiplicative $L$-functions is not multiplicative, but each
 summand retains its Euler product once the eigencomponent
 hypothesis is supplied.

\item \emph{Shadow coefficient level for algebraic families:}
 the shadow coefficients $S_r(\cA)$ are rational functions of
 the continuous parameter \textup{(}central charge or level\textup{)},
 independent of~$\tau$. Since $T_p$ acts as the identity on
 constants, the Hecke defect at the shadow coefficient level
 vanishes identically
 \textup{(}Proposition~\textup{\ref{prop:hecke-defect-families}}\textup{)}.
\end{enumerate}
\end{proposition}

\subsection{Where prime-locality fails or is open}

\begin{theorem}[Precise obstructions to prime-locality;
\ClaimStatusProvedHere{} where indicated]%
\label{thm:prime-locality-obstructions}%
\index{prime-locality!precise obstructions|textbf}
The following five obstructions interact and collectively
characterize the difficulty of extending prime-locality
beyond the proved cases:
\begin{enumerate}[label=\textup{(\alph*)}]
\item \emph{Lack of $\bZ$-form \textup{(}non-lattice\textup{):}}
 For a lattice VOA, the partition function is built from a theta
 function $\Theta_\Lambda = \sum_{\lambda\in\Lambda} q^{|\lambda|^2/2}$,
 which has integral Fourier coefficients
 \textup{(}representation numbers $r_\Lambda(n) \in \bZ$\textup{)}.
 This integrality is essential: the Hecke decomposition in
 $M_{r/2}(\Gamma)$ requires the Fourier coefficients to be
 algebraic. For the Virasoro algebra at generic~$c$, the
 Fourier coefficients of the partition function are partition
 numbers $p(n)$, which are integers, but the primary-counting
 function $\widehat{Z}^c = y^{c/2}|\eta|^{2c}Z$ involves the
 non-integral power $|\eta(\tau)|^{2c}$, and for irrational~$c$
 this is not a modular form for any congruence subgroup.
 \textbf{Scope:} obstruction (a) is bypassed for rational VOAs
 \textup{(}the characters are vector-valued modular forms on
 $\Gamma(4pq)$\textup{)}, but remains open for
 irrational~$c$.

\item \emph{Lack of Hecke action on the convolution algebra:}
 The convolution algebra $\gAmod = \mathrm{Hom}_\alpha
 (\cC^!_{\mathrm{ch}}, \cP^{\mathrm{ch}})$ is an algebraic
 object defined over the base curve~$X$. The Hecke operators
 are analytic objects defined on $\cM_{1,1}$ via isogeny
 correspondences. For lattice VOAs, the sewing differential
 $d_{\mathrm{sew}}$ factors through $\Theta_\Lambda$, which is
 a Hecke eigenform, so $T_p$ extends to a chain endomorphism
 of $\gAmod$. For a general chirally Koszul algebra, the sewing
 differential involves the propagator $P(z,w;\tau)$ on the
 elliptic curve $E_\tau$, and the Hecke translation
 $\tau \mapsto (a\tau{+}b)/d$ changes the elliptic curve.
 The commutator $[d_{\mathrm{sew}}, T_p]$ is the Hecke defect.
 \textbf{Known obstruction:} the defect class
 $[\delta_p] \in H^1(\gAmod, \mathrm{ad}_{T_p})$ is defined
 \textup{(}Definition~\textup{\ref{def:hecke-defect-class}}\textup{)}
 but not computed for any non-lattice, non-rational algebra.

\item \emph{Lack of Euler product at the character level:}
 The partition function Dirichlet series
 $D(s) = \sum a_n n^{-s}$ from $|\eta|^{2c}Z$ is provably
 \emph{not} multiplicative for rational CFTs
 \textup{(}Proposition~\textup{\ref{prop:rational-cft-multiplicativity-failure}};
 $14$ coprime failures at $c = 1/2$ with $m,n \le 20$\textup{)}.
 This multiplicativity failure is structural: the cross-channel
 terms $\sum_{i \neq j} b_m^{(i)} b_n^{(j)}$ between distinct
 conformal families destroy multiplicativity. For rational VOAs,
 the individual channel Dirichlet $L$-functions retain their
 Euler products \textup{(}Rocha--Caridi theta decomposition\textup{)},
 so the obstruction is a projection artefact. For irrational~$c$,
 there is no such decomposition into finitely many channels.
 \textbf{Scope:} obstruction (c) is resolved for rational
 VOAs and lattice VOAs; open for irrational~$c$.

\item \emph{Quasi-modularity of shadow amplitudes:}
 The genus-$1$ propagator is the quasi-modular Eisenstein
 series $E_2^*(\tau)$, \emph{not} a holomorphic modular form
 \textup{(}Remark~\textup{\ref{rem:quasimodular-obstruction}}\textup{)}.
 Products of $E_2^*$ are quasi-modular polynomials in
 $\{E_2^*, E_4, E_6\}$, not polynomials in $\{E_4, E_6\}$.
 The Rankin--Selberg unfolding for quasi-modular forms requires
 Zagier's extended theory, and the Euler product property of the
 resulting $L$-functions is not automatic.
 At degree~$2$, $\operatorname{Sh}_2^{(1)} \propto \kappa \cdot E_2^*$,
 and $E_2^*$ has multiplicative Fourier coefficients
 \textup{(}$\sigma_1(n)$ is multiplicative\textup{)}.
 At degrees $r \ge 3$, the shadow amplitudes involve products
 of $E_2^*$, and the Fourier coefficients of $(E_2^*)^k$ are
 \emph{not} multiplicative for $k \ge 2$.
 \textbf{Obstruction:} this is the main structural obstruction.
 Even if the shadow \emph{coefficients} $S_r$ are tau-independent
 \textup{(}hence trivially Hecke-equivariant\textup{)}, the
 shadow \emph{amplitudes} $\operatorname{Sh}_r^{(1)}(\tau)$
 acquire tau-dependence from the propagator, and this
 tau-dependence is quasi-modular, not holomorphic.

\item \emph{Structural separation at genus~$1$:}
 The genus-$1$ unfolding-erasure computation
 \textup{(}Proposition~\textup{\ref{prop:scattering-residue}}\textup{)}
 proves that
 the MC equation at genus~$1$ cannot access the scattering
 matrix poles \textup{(}which live at complex spectral
 parameter $t$, while the MC constraints operate on the real
 $t$-axis\textup{)}. This means the MC equation constrains
 the \emph{moments} of the spectral decomposition but not the
 \emph{individual} Hecke eigenvalues. Prime-locality requires
 eigenvalue-level control, which the moment constraints alone
 do not provide.
 \textbf{Resolution:} this is a proved impossibility result for
 a particular attack strategy \textup{(}the direct genus-$1$
 approach\textup{)}, not for prime-locality itself.
\end{enumerate}
\end{theorem}

\subsection{MC overdetermination and character-level propagation}%
\label{subsec:route-c-assessment}%
\index{prime-locality!MC overdetermination assessment|textbf}

Theorem~\ref{thm:route-c-propagation} separates the
coefficient statement from the amplitude statement. Character-level
Hecke equivariance and low-degree determination propagate through the
MC recursion to the tau-independent coefficients. Amplitude-level
prime-locality additionally requires Hecke compatibility of the
quasi-modular graph sum.

\begin{remark}[Where MC overdetermination succeeds]%
\label{rem:route-c-success}%
\index{prime-locality!MC overdetermination success}
MC overdetermination is correct in the following precise sense:
\begin{enumerate}[label=\textup{(\roman*)}]
\item The MC recursion does determine $S_r$ from lower-degree
 data. This is the content of the Riccati algebraicity
 theorem: the shadow obstruction tower is algebraic of degree~$2$,
 and the square-root expansion
 $H(t) = t^2\sqrt{Q_L(t)}$ determines all shadow coefficients
 from $\kappa, \alpha, S_4$.

\item For algebraic families, the shadow coefficients are
 tau-independent rational functions of the level or central
 charge. The MC recursion operates entirely within
 tau-independent data, so the Hecke operator $T_p$ commutes
 with the recursion at the shadow coefficient level.

 \item Corollary~\textup{\ref{cor:route-c-standard-landscape}}
 states the coefficient-level conclusion for the standard landscape.
\end{enumerate}
\end{remark}

\begin{remark}[The amplitude-level gap]%
\label{rem:route-c-gap}%
\index{prime-locality!MC overdetermination gap|textbf}
The remaining structural gap lies between the shadow coefficient
level and the amplitude level. The gap consists of two
independent issues:

\textbf{First obstruction: the MC recursion on shadow coefficients is not
the same as the MC recursion on shadow amplitudes.}
The shadow coefficients $S_r$ are tau-independent OPE data
\textup{(}for algebraic families\textup{)}. The shadow amplitudes
$\operatorname{Sh}_r^{(1)}(\tau)$ are genus-$1$ graph sums
involving the tau-dependent propagator $E_2^*(\tau)$.
The MC recursion at the \emph{coefficient} level is a polynomial
recursion among constants; the MC recursion at the
\emph{amplitude} level is a recursion among quasi-modular forms.
Theorem~\ref{thm:route-c-propagation} proves Hecke invariance at the
coefficient level; the
amplitude level requires an additional argument.

The proof of Theorem~\ref{thm:route-c-propagation} asserts the
coefficient identity
$T_p(S_r) = \operatorname{MC}_r(T_p(S_j), T_p(P_0))$,
where $P_0$ is the scalar OPE contraction. This identity is correct
on the tau-independent coefficient lane \textup{(}since $T_p$ acts
as the identity on constants\textup{)}. But the shadow amplitude is
not $S_r$ alone: it is
\begin{equation}\label{eq:shadow-amplitude-graph-sum}
 \operatorname{Sh}_r^{(1)}(\tau)
 = \sum_\Gamma
 \frac{1}{|\operatorname{Aut}(\Gamma)|}
 \prod_{v \in V(\Gamma)} S_{n_v}
 \prod_{e \in E(\Gamma)} P(e;\tau),
\end{equation}
where the propagator $P(e;\tau) \propto E_2^*(\tau)$ is
quasi-modular. The Hecke operator $T_p$ acts on the
$\tau$-dependent factors; the question is whether the graph
sum respects this action.

\textbf{Second obstruction: the $T_p(S_j)\,T_p(P(e;\tau)) \neq
T_p(S_j P(e;\tau))$ problem.}
Even if $T_p(S_j) = S_j$ \textup{(}which is true for
tau-independent coefficients\textup{)}, the identity
$T_p(\operatorname{Sh}_r^{(1)}) =
\operatorname{Sh}_r^{(1)}|_{T_p(\text{propagator})}$
requires the multiplicativity of $T_p$ on \emph{products} of
modular forms. For Hecke eigenforms, $T_p$ is multiplicative
by definition. For quasi-modular forms like $E_2^*$, the Hecke
action is defined but the multiplicativity on products of $E_2^*$
is not automatic: $T_p(f \cdot g) \neq T_p(f) \cdot T_p(g)$
in general for modular forms of different weights, and
$T_p((E_2^*)^k)$ is not $(T_p(E_2^*))^k$ because $T_p$ is
\emph{not} a ring homomorphism on quasi-modular forms.

\textbf{Consequence.} MC overdetermination reduces prime-locality
to the shadow \emph{coefficient} level for algebraic families.
For the amplitude level, the additional step requires that the
graph-sum structure of the shadow amplitudes is compatible with
the Hecke action on quasi-modular forms. This compatibility is
\emph{not} proved. It is a genuine mathematical problem,
related to the Hecke equivariance of iterated Eisenstein
integrals \textup{(}Brown--Levin theory\textup{)}, and is not
reducible to a counting argument.
\end{remark}

\subsection{Is there a modified prime-locality that is true?}%
\label{subsec:modified-prime-locality}

The shadow obstruction tower satisfies a strong algebraic constraint with
arithmetic content, but distinct from prime-locality in the
number-theoretic sense.

\begin{theorem}[Riccati determinacy; \ClaimStatusProvedElsewhere]%
\label{thm:riccati-determinacy-assessment}%
\index{Riccati algebraicity!as modified prime-locality}
The shadow coefficients $S_r$ at all degrees are determined by
$S_2 = \kappa$, $S_3 = \alpha$, and $S_4 = Q^{\mathrm{contact}}$
through a polynomial recursion:
\[
 S_r \;=\; \text{polynomial in $\kappa, \alpha, S_4$}.
\]
This is the content of the Riccati algebraicity theorem: the
generating function $H(t) = t^2\sqrt{Q_L(t)}$ is algebraic of
degree~$2$, and its Taylor coefficients are determined by
the three parameters $\kappa, \alpha, S_4$ via the recursion
$a_n = -(2a_0)^{-1}\sum_{j=1}^{n-1} a_j a_{n-j}$.

This is a strong form of \emph{algebraic Hecke determinacy}: the
shadow data at all degrees is determined by three
algebraic invariants of the OPE. These invariants are tau-independent for algebraic
families, so the shadow coefficient tower is trivially
Hecke-equivariant.

This is \emph{not} arithmetic prime-locality. Arithmetic
prime-locality requires that the genus-$1$ shadow
\emph{amplitudes} $\operatorname{Sh}_r^{(1)}(\tau)$ decompose
into Hecke eigencomponents with multiplicative Fourier
coefficients. The Riccati algebraicity determines the
\emph{vertex weights} in the graph sum but says nothing about
the \emph{propagator weights} \textup{(}which carry the
tau-dependence\textup{)}.
\end{theorem}

\subsection{Negative results and scope delimiters}
\label{subsec:negative-results}
\index{shadow tower!negative results|textbf}
\index{Selberg class!failure for shadow zeta}
\index{Euler product!failure for shadow zeta}

The following negative results delineate the programme's
scope. Each is as important as the positive theorems: they
identify the structural boundaries that no refinement of
technique can cross.

\begin{proposition}[The genus-$1$ amplitude series is not in the
Selberg class]
\label{prop:shadow-not-selberg}
\ClaimStatusProvedHere
\index{Selberg class!axiom failure|textbf}
The genus-$1$ amplitude Dirichlet series
$D_2(\cA,s) = -24\kappa\,\zeta(s)\,\zeta(s{-}1)$
\textup{(}Theorem~\textup{\ref{thm:shadow-eisenstein}}\textup{)}
is not an element of the Selberg class for any
$\kappa\neq0$.

Specifically:
\begin{enumerate}[label=\textup{(\roman*)}]
\item \emph{Shifted Eisenstein product, not a primitive Selberg
object.}
 $D_2(\cA,s)$ is a product of two zeta functions
 \textup{(}an Eisenstein series\textup{)} and has the Euler product
 $\prod_p(1-p^{-s})^{-1}(1-p^{1-s})^{-1}$ in
 $\operatorname{Re}(s)>2$. It is not a primitive Selberg-class
 Euler product because the shifted factor $\zeta(s-1)$ produces the
 pole at $s=2$ and the wrong coefficient growth. The Dirichlet coefficients
 $a_n = -24\kappa\,\sigma_1(n)$ are multiplicative, but
 $D_2$ is already the Rankin--Selberg
 convolution $\zeta \otimes \zeta$ (up to scalar), not a
 primitive $L$-function. Additivity in~$\kappa$
 ($D_2(\cA \otimes \cB, s)
 = D_2(\cA, s) + D_2(\cB, s)$)
 follows from $\kappa(\cA \otimes \cB)
 = \kappa(\cA) + \kappa(\cB)$
 \textup{(}Theorem~D\textup{)}.
\item \emph{Coefficient growth.}
 $|a_n| = 24|\kappa|\,\sigma_1(n) = O(n^{1+\varepsilon})$ for
 all $\varepsilon > 0$, so the Ramanujan hypothesis S4
 \textup{(}which requires $|a_n| = O(n^\varepsilon)$\textup{)}
 also fails.
\item \emph{The shadow $L$-function is a separate series.}
 The shadow $L$-function
 $L_\cA^{\mathrm{sh}}(s) = \sum_{r \ge 2} S_r\,r^{-s}$
 \textup{(}a different object from~$D_2$; see
 Remark~\textup{\ref{rem:shadow-eisenstein-correct-scope}}\textup{)}
 has no Euler-product assertion here. No Selberg-class or
 subconvexity statement about $L_\cA^{\mathrm{sh}}$ follows from the
 Eisenstein identity for~$D_2$.
\end{enumerate}
\end{proposition}

\begin{proof}
The Eisenstein decomposition
$D_2(\cA,s) = -24\kappa\,\zeta(s)\,\zeta(s{-}1)$
is Theorem~\ref{thm:shadow-eisenstein}. The Selberg class
allows Euler products, but its standard normalisation has at most the
permitted pole at $s=1$ and Ramanujan-type coefficient growth. The
factor $\zeta(s-1)$ gives a pole at $s=2$, and
$a_n=-24\kappa\sigma_1(n)$ has growth $O(n^{1+\varepsilon})$, not
Selberg-class Ramanujan growth.
Additivity of $D_2$ in~$\kappa$ follows from the additivity
of~$\kappa$ (Theorem~D, the independent-sum factorization):
$\kappa(A \otimes B) = \kappa(A) + \kappa(B)$. The
last item is only a separation of objects: $L_\cA^{\mathrm{sh}}$
uses degree-indexed shadow coefficients, while $D_2$ uses Fourier
coefficients of the genus-$1$ degree-$2$ amplitude.
\end{proof}

\begin{remark}[Independence from the Riemann zeta]
\label{rem:shadow-riemann-independence}
\index{Dixmier trace!shadow-Riemann independence}
\index{shadow zeta!independence from Riemann zeta}
The shadow arithmetic data and the Riemann zeta function are
generically independent: the genus-$1$ amplitude Dirichlet
series $D_2(\cA,s) = -24\kappa\,\zeta(s)\,\zeta(s{-}1)$
(Theorem~\ref{thm:shadow-eisenstein})
has zeros at the zeros of $\zeta(s)$ \emph{and} at the
zeros of $\zeta(s{-}1)$, but carries no additional
zero-location information. The Dixmier trace of the shadow
spectral triple
(Remark~\ref{rem:connes-spectral-triple}) satisfies
$\operatorname{Tr}_\omega(|D|^{-1}) = |\kappa(\cA)|$, which
is an algebraic invariant of~$\cA$ with no dependence on the
zeros of~$\zeta$. The shadow and Riemann zeta share no zeros
beyond those forced by the factorization.

The shadow Li coefficients
$\tilde\lambda_n^{\mathrm{sh}}(\cA)$
(Proposition~\ref{prop:li-criterion-failure}) differ from the
classical Li coefficients $\lambda_n$ of the Riemann
$\xi$-function: the shadow Li coefficients incorporate the
double-zeta structure $\zeta(s)\,\zeta(s{-}1)$, which
introduces zeros on two critical lines and forces eventual
negativity.

The K-theory regulator of the shadow spectral triple
does \emph{not} vanish at the nontrivial zeros of~$\zeta$:
the regulator map $K_0(\cA_{\mathrm{sh}}) \to \bR$ is a
continuous function of~$\kappa$, and $\kappa$ is a rational
function of the central charge~$c$, independent of the
spectral parameter~$s$.
\end{remark}

\subsection{Prime-locality obstruction boundary}

\begin{remark}[Prime-locality programme]%
\label{rem:prime-locality-definitive}%
\index{prime-locality!obstruction boundary|textbf}
\mbox{}

\textbf{What is proved or conditionally reduced:}
\begin{enumerate}[label=\textup{(\arabic*)}]
\item Prime-locality for lattice VOAs \textup{(}all ranks,
 all levels\textup{)}, via theta-function Hecke decomposition.
\item Character-level Hecke action for rational VOAs with diagonal
 modular invariant, via Franc--Mason vector-valued Hecke theory and
 the theta-bridge decomposition; amplitude-level prime-locality
 remains conditional on graph-sum compatibility.
\item Prime-locality at the shadow coefficient level for all
 algebraic families, via tau-independence of the OPE data.
\item Riccati algebraicity: the shadow obstruction tower is finitely
 determined \textup{(}by three invariants $\kappa, \alpha, S_4$\textup{)},
giving coefficient-level Hecke invariance in a precise but weaker
sense than arithmetic prime-locality.
\item Conditional Ramanujan implications for lattice and rational VOAs
 under the CPS, local-factor, and prime-locality hypotheses
 \textup{(}Theorems~\textup{\ref{thm:hecke-newton-lattice}}
 and~\textup{\ref{thm:non-lattice-ramanujan}}\textup{)}.
\item The structural separation theorem: the genus-$1$ MC
 equation cannot access scattering-matrix poles.
\end{enumerate}

\textbf{What is open:}
\begin{enumerate}[label=\textup{(\arabic*)}]
\item Prime-locality for irrational VOAs
 \textup{(}Virasoro at generic~$c$, $\cW_3$ at generic~$c$\textup{)}.
 The obstacle is twofold: \textup{(a)}~the partition function
 $|\eta|^{2c}$ for irrational~$c$ is not a modular form for
 any congruence subgroup, and \textup{(b)}~the Roelcke--Selberg
 decomposition involves continuous spectrum \textup{(}Maass forms\textup{)},
 whose Hecke eigenvalues are constrained by the Selberg
 eigenvalue conjecture
 \textup{(}$\lambda_1 \ge 1/4$, best known $\theta \le 7/64$
 by Kim--Sarnak\textup{)}, not by Deligne's theorem.
\item The amplitude-level Hecke equivariance: whether the
 graph-sum structure of the shadow amplitudes
 $\operatorname{Sh}_r^{(1)}(\tau)$ respects the Hecke action
 on quasi-modular forms. This is the key open step in the
 MC-overdetermination criterion.
\item The Hecke defect computation: no explicit amplitude-level
 calculation of $[\delta_p(\cA)]$ is supplied here for a non-lattice
 rational or irrational algebra. A computation at $p = 2$ for the
 Ising model would decide the first non-lattice rational instance;
 a computation for generic Virasoro would decide the first
 irrational instance.
\end{enumerate}

\textbf{What is structurally impossible:}
\begin{enumerate}[label=\textup{(\arabic*)}]
\item Extracting zero-location information for $\zeta(s)$ from
 the genus-$1$ MC equation
 \textup{(}Theorem~\textup{\ref{thm:structural-separation}}\textup{)}.
\item Obtaining an Euler product for the full partition
 function Dirichlet series of a non-lattice rational CFT
 \textup{(}Proposition~\textup{\ref{prop:rational-cft-multiplicativity-failure}}\textup{)}.
\end{enumerate}

\textbf{Assessment.}
For rational VOAs, Franc--Mason Hecke theory provides the
character-level eigendecomposition, and the theta-bridge resolves
the non-multiplicativity of the scalar Dirichlet series. This is not
yet amplitude-level prime-locality.
For irrational VOAs, the question reduces to a concrete
problem: does the graph-sum structure of the shadow amplitudes,
with vertex weights from the tau-independent MC recursion and
edge weights from $E_2^*(\tau)$, respect the Hecke decomposition
of quasi-modular forms? This concerns the interaction between
operadic graph sums and Hecke operators, not the MC equation
itself.

MC overdetermination is sound at the shadow coefficient level but
has a gap at the amplitude level: not a counting failure
\textup{(}the overdetermination count is correct\textup{)} but a
\emph{compatibility} question between the graph-sum operad and
the Hecke algebra acting on quasi-modular forms. Closing the
gap requires: \textup{(A)}~a vanishing theorem for
$H^1(\gAmod, \mathrm{ad}_{T_p})$,
\textup{(B)}~explicit computation of $[\delta_p]$ for a
non-lattice algebra, or
\textup{(C)}~a structural argument that graph-sum amplitudes
on $\overline{\cM}_{1,r}$ with universal propagators inherit
Hecke equivariance from the underlying modular geometry.
Option~(C) is the most natural: it would follow if the
tautological ring $R^*(\overline{\cM}_{1,r})$ is
Hecke-equivariant in a suitable sense, a question that connects
the prime-locality programme to the Faber--Pandharipande
programme in algebraic geometry.
\end{remark}

\section{Structural cross-references}
\label{sec:as-anchors}

\begin{remark}[Structural anchors for the arithmetic shadow tower]
\label{rem:as-anchors}
\index{arithmetic shadows!entry point}
\index{cross-volume bridge!arithmetic shadows}
The arithmetic-shadow chapter sits at three structural
anchor points:
\begin{enumerate}[label=\textup{(W\arabic*)}]
\item \emph{Universal conductor.} The genus-$1$ free energy
 $F_1 = \kappa(\cA)/24$ is the structural origin of the
 Bernoulli denominator $24 = 1/B_2$ that propagates through
 the entire arithmetic-shadow programme: the universal
 conductor identity $\kappa(\cA) = -c_{\mathrm{ghost}}(\BRST(\cA))$
 of Chapter~\ref{chap:universal-conductor} identifies $\kappa$
 with a ghost-charge invariant, and the $1/24$ prefactor
 emerges as the genus-$1$ tautological-class evaluation
 $\int_{\overline{\cM}_{1,1}} \lambda_1 = 1/24$. Higher-genus
 Bernoulli appearances in the shadow tower (e.g.\ $S_{12} =
 -65520/691$ for the Leech lattice) are governed by the same
 conductor through the genus-universality theorem~D.
\item \emph{Climax identity.} The shadow Eisenstein identity
 $D_2(\cA, s) = -24\kappa\,\zeta(s)\,\zeta(s{-}1)$
 (Theorem on the shadow Epstein zeta) is the arithmetic face
 of the climax identity $d_{\mathrm{bar}} = \mathrm{KZ}^*(\nabla_{\mathrm{Arnold}})$
 of Chapter~\ref{ch:climax-theorem}: pulling back the Arnold
 form $\eta_{ij}$ through the genus-$1$ KZB connection
 produces the quasi-modular Eisenstein series $E_2^*$, and the
 Mellin transform of the resulting graph sum is the displayed
 product of two Riemann zeta functions. The conductor
 prefactor $-24\kappa$ is the genus-$1$ scalar projection of
 $\Theta_\cA$.
\item \emph{Koszul reflection in the arithmetic locus.} The
 cross-channel correction
 $\delta F_2^{\mathrm{cross}}(W_3) = (c+204)/(16c)$
 (Theorem~D inscription, Vol~I main programme) is the
 arithmetic instance of the Koszul reflection theorem
 (Theorem~\ref{thm:koszul-reflection}, Vol~I Platonic
 Theorem~A): the Koszul-self-dual locus of $W_3$ supports a
 finite-rank correction to the uniform-weight $F_2$ formula
 whose denominator $16c$ encodes the level-shifted Sugawara
 normalisation, and the bar-cobar inversion of the W-line
 collision residue produces this correction as the second
 reflection of the binary $r$-matrix.
\end{enumerate}

\paragraph{Obstruction note.}
A naive Bernoulli--Dirichlet identification would impose a functional
equation that fails numerical evaluation at $s = 0$. The theorem
therefore uses
$D_2(\cA, s)$ built from Fourier coefficients of the genus-$1$
graph sum and matches the universal conductor through the
$-24\kappa$ prefactor.

\paragraph{Cross-volume projections.}
Vol~I, Vol~II, and Vol~III use the same obstruction tower only after
projecting to different targets. In Vol~I the projection is scalar:
$\kappa(\cA)$, the coefficients $S_r(\cA)$, and the Dirichlet series
$D_2(\cA,s)=-24\kappa\,\zeta(s)\zeta(s-1)$. In Vol~II the universal
celestial holography theorem applies, on its stated standard-landscape
scope, to triples $(A,\omega,\mathrm{BL})$ and identifies the bulk
observables with
\[
 \mathrm{Obs}^{\mathrm{bulk}}(\Phihol(A,\omega,\mathrm{BL}))
 \simeq
 Z^{\mathrm{der}}_{\mathrm{ch}}(A)
 =
 \mathrm{ChirHoch}^{\bullet}(A,A).
\]
Thus the soft-factor hierarchy is an image of the obstruction tower
inside derived-centre bulk data; it is not an equality with the scalar
shadow invariants.

Vol~III's factorisation
\[
 \Phi_3=\SpCh_{\Sigma_2,E_\tau}\circ\PhiFA_3
\]
lands in a CY-to-chiral specialisation. On $K3\times E$ the census
keeps four scalars distinct:
\[
 \kappa_{\mathrm{cat}}(K3\times E)=0,\qquad
 \kappa_{\mathrm{ch}}^{\mathrm{Heis}}(K3\times E)=3,\qquad
 \kappa_{\mathrm{BKM}}(\Delta_5)=c_1(0)/2=5,\qquad
 \kappa_{\mathrm{fiber}}(K3)=24.
\]
The Mukai-doubled conductor is a fifth comparison value
\[
 K^{\kappa_{\mathrm{ch}}}(\mathrm{Mukai}(K3))
 =2c_+(\mathrm{II}_{4,20})=8.
\]
The Leech coefficient belongs to a different lattice row:
\[
 \Theta_{\mathrm{Leech}}
 =E_{12}-\frac{65520}{691}\Delta_{12},\qquad
 S_{12}(V_{\Lambda_{24}})=-\frac{65520}{691},
\]
while the Fake-Monster Borcherds product has
$\kappa_{\mathrm{BKM}}(\Phi_{12})=12$. The subscript in $c_N(0)$
records the Borcherds row, not the Calabi--Yau dimension; the
$K3\times E$ row uses $N=1$ and automorphic weight $c_1(0)/2=5$.
A cross-volume theorem must therefore construct explicit comparison
maps among scalar shadows, holographic bulk observables, and
Borcherds/Mukai automorphic data; the numerical coincidences alone do
not supply such a theorem.
\end{remark}

% Local bibliography removed: entries now in bibliography/references.tex

\section{Gelfand structural lens}
\label{sec:as-gelfand-supplement}

\begin{remark}[The Etingof--Kazhdan cohomology target for $\Delta_5$]
\label{rem:as-gelfand-ek-cohomology}
Conditional on the Vol~III finite Hall--Borcherds recognition theorem
(Vol~III Chapter~\ref*{V3-ch:k3-chiral-bialgebra-platonic},
Theorem~\ref*{V3-thm:k3-classification-etingof}), the source-side
target is the identification
$H^2(\mathfrak{g}_{\Delta_5})^{\mathbb{Z}/2, K(1)} \cong \mathbb{C}\cdot\Delta_5$
as the one-dimensional Lie-bialgebra cohomology group cutting out the
gauge-equivalence class of the K3 chiral bialgebra $\mathbf{H}_{\Delta_5}$.
Under the dictionary between Lie-bialgebra cohomology and automorphic
cohomology established by Bruinier 2002 Proposition~5.1 (Heegner-divisor
Chern-class reciprocity), the resulting class has arithmetic shadow the
Igusa form $\Delta_5\in S_5(\widetilde{\mathrm{Sp}_4(\mathbb{Z})}^{v_{\Delta_5}})$.
Without that finite source-recognition datum, the arithmetic reading is
only scalar: identities involving $\Delta_5$ are identities of the
automorphic Heegner-divisor shadow. They become identities in the
Lie-bialgebra \(H^2\) only after the comparison map, parity/root fixture,
compact source matrices, and PBW/no-extra clauses have been supplied.
The finite recognition envelope is a legitimate quotient comparison,
but the displayed \(H^2\)-line is a statement about the unquotiented
source only after compact-provenance source matrices force envelope
faithfulness.
Bruinier reciprocity sends the relevant automorphic Chern class to the
Heegner-divisor shadow; it does not, by itself, construct a compact
Hall--Borcherds 2-cocycle. The genus-$1$ scalar shadow
$\kappa(\mathcal{A})$ is the projection of the automorphic class onto the
dimension-$1$ constant stratum; higher shadows $S_k$ are higher-monodromy
data after the recognition datum is present.
\end{remark}

\begin{remark}[Structural target for $K^{\kappa_{\mathrm{ch}}} = 8$]
\label{rem:as-gelfand-8-three-faces}
The value $K^{\kappa_{\mathrm{ch}}} = 8$ on the $\mathcal{B}$-family
Lorentzian-lattice stratification (Theorem-C enlargement at
\texttt{chapters/examples/landscape\_census.tex}) has one established
Vol~I scalar source and two comparison targets:
\begin{enumerate}[label=\textup{(\roman*)}]
\item \emph{Mukai face.} The Mukai pairing on the K3 Mukai lattice
 $\widetilde{H}(K3,\mathbb{Z}) = H^0\oplus H^2\oplus H^4$ has signature
 $(4,20)$; its positive-signature dimension $c_+(\mathrm{Mukai}(K3))=4$
 gives $K^{\kappa_{\mathrm{ch}}} = 2 c_+(\mathrm{Mukai}(K3)) = 8$
 (Vol~III Remark~\ref*{V3-rem:mukai-doubling-K-kappa}).
\item \emph{Humbert-monodromy target.} The Humbert surface
 $H_1 \subset \overline{\mathcal{A}_2}$ has $\mathcal{D}$-module
 monodromy-order target $8$. To promote this target to a Vol~I
 theorem one must compute the nearby-cycles monodromy on the same
 automorphic line bundle used for $\Delta_5$ and exhibit the comparison
 map from the scalar conductor to that local system.
\item \emph{Lusztig target.} At the Lusztig specialisation
 $\zeta_\ell$ with $\ell=8$, the intended small quantum group
$u_\zeta(\mathfrak{g}_{\Delta_5})$ is required to close on a finite-rank
cover of $\mathbf{H}_{\Delta_5}$. The required comparison data are PBW
 finiteness, the parity/root fixture, and the no-extra-blocks clause.
\end{enumerate}
The phrase ``three faces'' is therefore only a comparison programme
until the Humbert and Lusztig targets are promoted by these data. The
proved Vol~I input is the Mukai conductor \(8\); Bruinier reciprocity
supplies the automorphic divisor class to be matched, not the missing
source recognition by itself. Cross-reference to Vol~I Theorem~C
enlargement and to \texttt{chapters/examples/landscape\_census.tex}.
\end{remark}

\begin{remark}[Schur-index comparison target through the composite arrow]
\label{rem:as-gelfand-schur-composite}
The arithmetic-shadow comparison target behind Vol~III
Theorem~\ref*{V3-thm:schur-to-delta5-composite} is the Schur index of
$\mathcal{T}[A_1,\Sigma_{0,24}]$. The coefficient comparison currently
available is
\[
\mathcal{I}_{\mathrm{Schur}}[\mathcal{T}[A_1,\Sigma_{0,24}]](q;\mathbf{1})
\;=\;
\mathrm{PE}\!\left[\frac{72q - 22q^2}{1-q}\right] + O(q^{11})
\;=\;
1 + 72q + 2678q^2 + 68474q^3 + 1351775q^4 + \cdots.
\]
The 24 flavour fugacities carry the $M_{24}$ permutation action; averaging
and diagonal restriction produces the full K3 elliptic genus
$Z_{K3}=2\phi_{0,1}^{K3}=2(y+y^{-1})+20+O(q)$, and the Borcherds
singular-theta lift has target output $\Delta_5$. Each arrow loses
information, and the 10
terms above are only a finite coefficient comparison. The lift from
coefficient comparison to un-averaged arithmetic shadow requires the
full flavour-refined Schur index, the \(M_{24}\) equivariance and Schur
multiplier cocycle, the diagonal restriction map to the K3 elliptic genus,
the Borcherds-lift normalisation, and the source comparison map to the
K3 chiral bialgebra. Until those data are present, $\Delta_5$ is the
automorphic target of the averaged diagonal, not an identified projection
of a proved source Schur object.
Vol~III's Schur-index coefficient computation gives the same 10
coefficients by plethystic exponentiation. Primary lit:
Beem--Peelaers--Rastelli 2015;
Gadde--Rastelli--Razamat--Yan 2011; Eguchi--Ooguri--Tachikawa 2011;
Borcherds 1998.
\end{remark}

% ============================================================
\section{Fricke LDP: sub-leading correction at each node}
\label{sec:fricke-ldp-sub-leading}

The large-deviation principle on the Fricke-fixed subset
$\mathcal F^{w_8}$ of paramodular Jacobi forms at level $N = 8$
(the K3 regime) has leading-order rate function
$I_{\mathrm{leading}}(\theta) = 2\sin^2\theta$ concentrated on the
Sato--Tate measure. The sub-leading correction at the Fricke nodes
$\theta_k^* = \pi k/8$, $k\in\{1, 3, 5, 7\}$ carries the first-order
deviation from Gaussian concentration.

\begin{theorem}[Fricke LDP sub-leading correction at each node;
\ClaimStatusProvedHere]
\label{thm:fricke-ldp-sub-leading}
\index{Fricke involution!LDP sub-leading correction}%
\index{large deviation!Fricke node variance}%
At each Fricke node $\theta_k^*\in\{\pi/8, 3\pi/8, 5\pi/8, 7\pi/8\}$
of the paramodular $N = 8$ regime, the sub-leading Gaussian
fluctuation has variance
\begin{equation}\label{eq:fricke-ldp-sigma-k}
 \sigma_k^2 \;=\; \frac{\pi}{2\sin^2\theta_k^*},
\end{equation}
giving the explicit four-tuple
$(\sigma_1^2, \sigma_3^2, \sigma_5^2, \sigma_7^2) =
(\pi/(2\sin^2(\pi/8)),\,\pi/(2\sin^2(3\pi/8)),\,
\pi/(2\sin^2(5\pi/8)),\,\pi/(2\sin^2(7\pi/8)))
\approx (10.726,\,1.840,\,1.840,\,10.726)$,
with exact form $\sigma_1^2 = \sigma_7^2 = \pi(2+\sqrt2)$ and
$\sigma_3^2 = \sigma_5^2 = \pi(2-\sqrt2)$.
\end{theorem}

\begin{proof}
At each Fricke node $\theta_k^* = \pi k/8$ the saddle-point evaluation
of the Jacobi theta integral around the critical angle gives Gaussian
fluctuations whose variance is the inverse curvature of
$I_{\mathrm{leading}}$:
$\sigma_k^{-2} = I''(\theta_k^*)/\pi = 4\sin^2\theta_k^*/\pi \cdot
(2/2) = 2\sin^2\theta_k^*/\pi$. Inversion
gives~\eqref{eq:fricke-ldp-sigma-k}. Three verification paths.
\textup{(i)} Trigonometric: $\sin^2(\pi/8) = (1 - \cos(\pi/4))/2 =
(1 - \sqrt2/2)/2 = (2 - \sqrt2)/4$, and
$\sin^2(3\pi/8) = (2 + \sqrt2)/4$, giving
$\sigma_1^2 = 2\pi/(2-\sqrt2) = \pi(2+\sqrt2)$
after rationalization.
\textup{(ii)} Fricke-pairing: $w_8$ swaps $\theta\mapsto \pi - \theta$
pairing $\theta_1^*\leftrightarrow\theta_7^*$ and
$\theta_3^*\leftrightarrow\theta_5^*$; variance preservation under
$w_8$ forces $\sigma_1^2 = \sigma_7^2$ and $\sigma_3^2 = \sigma_5^2$,
consistent with~\eqref{eq:fricke-ldp-sigma-k}.
\textup{(iii)} Gaussian tail-probability verification: the
probability $\Pr(|\theta - \theta_k^*| > \epsilon)$ decays as
$\exp(-\epsilon^2/2\sigma_k^2)$; at the inner nodes
$\theta_3^*, \theta_5^*$ the narrower variance
$\pi(2 - \sqrt2)\approx 1.557$ reflects steeper concentration, the
hallmark of K3 Gritsenko--Nikulin $1998$ \S$3.5$ rate-function
behaviour at the inner Fricke-fixed cusps.
\end{proof}

% ============================================================
\section{Monster Lusztig level: Conway--Norton cross-check at $c = 24$}
\label{sec:monster-lusztig-conway-norton}

\begin{theorem}[Scoped K3--Monster Lusztig ratio;
\ClaimStatusConditional]
\label{thm:monster-lusztig-universal-ratio}
\index{Lusztig level!Monster!universal ratio}%
\index{Conway-Norton!normalisation!$V^\natural$ at $c = 24$}%
Assume the Vol~III four-row $\Psi$-image convention in which the
K3 row carries Lusztig level
$\ell_{K3}=2c_+(\mathrm{II}_{4,20})=8$ and the Monster Fricke
row carries $\ell_{\mathrm{Monster}}=2c_+(II_{1,1})=2$. Then
\begin{equation}\label{eq:monster-lusztig-universal-ratio}
 \frac{\ell_{K3}}{\ell_{\mathrm{Monster}}}
 \;=\;
 \frac{8}{2}
 \;=\;
 \frac{c_+(\mathrm{II}_{4,20})}{c_+(II_{1,1})}
 \;=\;
 \frac{4}{1}
 \;=\;
 4.
\end{equation}
This is a scoped signature computation. It does not assert a
universal law for every Borcherds--Kac--Moody denominator, and it
does not derive Monster group element orders from K3 data. The
Conway--Norton normalisation supplies the Monster Fricke
involution independently; the ratio only checks compatibility of
the two level conventions on the K3/Monster pair.
\end{theorem}

\begin{proof}
Mukai~\cite{Mukai1987} identifies
$\widetilde H(K3,\mathbb Z)$ with the even lattice of signature
$(4,20)$, so $c_+(\mathrm{II}_{4,20})=4$. The Monster Fricke cusp
uses $II_{1,1}$, whose positive signature is $1$. Under the stated
level convention $\ell=2c_+$ on these two rows, the levels are
$8$ and $2$, and the equality~\eqref{eq:monster-lusztig-universal-ratio}
is immediate.

The proof uses no assertion about other BKM rows. For the Monster
module itself, FLM and Conway--Norton give the $c=24$ holomorphic
VOA and the Fricke normalisation; those data are independent
inputs, not consequences of the K3 row.
\end{proof}

% ============================================================
\section{Shimura--Waldspurger conversion at weights $k+1 \in \{7, 9, 11\}$}
\label{sec:shimura-waldspurger-higher-weights}

\begin{theorem}[Shimura--Waldspurger constants are period ratios;
\ClaimStatusProvedElsewhere]
\label{thm:shimura-waldspurger-higher-weights}
\index{Shimura-Waldspurger!$C_k$ higher weights}%
\index{paramodular level!Gritsenko table}%
Let
$f \in S_{k+1/2}^{+}(\Gamma_0(4),\chi)$ be the half-integral-weight
input, let $F=\mathrm{Sh}(f)\in S_{2k}(\mathrm{SL}_2(\bZ))$ be its
Shimura lift, and suppose the same paramodular form of weight
$k+1$ is written in the Gritsenko additive and Borcherds
multiplicative normalisations. Then
\begin{equation}\label{eq:SW-C-k-square-inverse}
 \Delta_{k+1}^{\mathrm{Borch}}
 \;=\;
 C_{k,D}\,\Delta_{k+1}^{\mathrm{Grit}},
 \qquad
 C_{k,D}
 \;=\;
 \mathcal E_{k,D}\,
 \frac{L(1,\chi_D)}
      {L(1/2,F\otimes\chi_D)},
\end{equation}
where $D<0$ is the discriminant used to compare Fourier
coefficients and $\mathcal E_{k,D}$ is the explicit nonzero product
of archimedean and bad-prime local factors in Waldspurger's formula
\cite{Waldspurger81}. The constant is therefore not determined by
the paramodular level alone. In the K3 normalisation
$(k,D)=(4,-4)$, Gritsenko--Nikulin's coefficient convention gives
\begin{equation}\label{eq:SW-C-explicit}
 C_{4,-4}=2^{-6}=\tfrac{1}{64},
 \qquad
 \Delta_5^{\mathrm{Borch}}
 =2^{-6}\Delta_5^{\mathrm{Grit}}.
\end{equation}
For weights $7,9,11$ the constants are the corresponding
Waldspurger period ratios. The identities
$C_7=1/144$, $C_9=1/576$, and $C_{11}=1/2304$ do not follow from
the Gritsenko paramodular level table and are not asserted here.
\end{theorem}

\begin{proof}
Waldspurger's theorem relates the square of a Fourier coefficient
of a Kohnen-plus half-integral form to a central value of the
quadratic twist of its Shimura lift, up to explicit local factors.
Changing from Gritsenko's additive normalisation to Borcherds'
multiplicative normalisation changes exactly these coefficient
normalisations; hence the comparison constant has the form
\eqref{eq:SW-C-k-square-inverse}. The K3 value
$C_{4,-4}=1/64$ is the explicit Gritsenko--Nikulin normalisation of
$\Delta_5$: equivalently,
$(\Delta_5^{\mathrm{Borch}})^2=
\Phi_{10}^{\mathrm{un}}/2^{12}$ while
$(\Delta_5^{\mathrm{Grit}})^2=\Phi_{10}^{\mathrm{un}}$.
No analogous rational constant is determined without evaluating
the local factors and the central value in
\eqref{eq:SW-C-k-square-inverse}. The paramodular level records the
group on which the form lives; it is not the Waldspurger period.
\end{proof}

% ============================================================
\section{Yetter--Drinfeld Schauenburg-bracket expansion
$\delta^{(n)}$, $n \in \{7, 8, 9, 10, 11, 12\}$}
\label{sec:YD-delta-n-7-8-9}

\begin{theorem}[$\delta^{(n)}$ for $n \in \{7, 8, 9, 10, 11, 12\}$;
\ClaimStatusProvedHere]
\label{thm:YD-delta-7-8-9}
\index{Yetter-Drinfeld!Schauenburg bracket!higher degrees}%
\index{Borcherds weight!$\delta^{(n)}$ expansion}%
The Yetter--Drinfeld Schauenburg-bracket expansion
$\delta^{(n)} = \sum_{\mathsf T}
(\text{Catalan-tree coefficient})\otimes
(\text{Brown MZV basis})$
of Borcherds weight $\lfloor n/2\rfloor + 1$ at degrees
$n\in\{7,\ldots,12\}$ has total Catalan--MZV dimension
$C_{n-1}\cdot d_n$ with
\begin{equation}\label{eq:YD-delta-789-dim}
\begin{aligned}
n = 7&:\quad C_6\cdot d_7 = 132\cdot 3 = 396,
\qquad \mathrm{wt}(\delta^{(7)}) = 4,\\
n = 8&:\quad C_7\cdot d_8 = 429\cdot 4 = 1716,
\qquad \mathrm{wt}(\delta^{(8)}) = 5,\\
n = 9&:\quad C_8\cdot d_9 = 1430\cdot 5 = 7150,
\qquad \mathrm{wt}(\delta^{(9)}) = 5,\\
n = 10&:\quad C_9\cdot d_{10} = 4862\cdot 7 = 34034,
\qquad \mathrm{wt}(\delta^{(10)}) = 6,\\
n = 11&:\quad C_{10}\cdot d_{11} = 16796\cdot 9 = 151164,
\qquad \mathrm{wt}(\delta^{(11)}) = 6,\\
n = 12&:\quad C_{11}\cdot d_{12} = 58786\cdot 12 = 705432,
\qquad \mathrm{wt}(\delta^{(12)}) = 7,
\end{aligned}
\end{equation}
with $C_k$ the $k$-th Catalan number (OEIS~A000108: $C_6=132,
C_7=429, C_8=1430, C_9=4862, C_{10}=16796, C_{11}=58786$) and $d_n$
the Zagier--Brown canonical MZV dimension at weight $n$
(Zagier--Brown Padovan recurrence $d_n=d_{n-2}+d_{n-3}$ with seed
$(d_3,d_4,d_5)=(1,1,1)$; OEIS~A002083 restricted to the
Deligne--Brown motivic stratum). The Borcherds weight satisfies
$\mathrm{wt}(\delta^{(n)}) = \lfloor n/2\rfloor + 1$ throughout,
falsifying the $\lceil n/2\rceil$ ansatz at every $n$ above.
\end{theorem}

\begin{proof}
The Catalan coefficient follows from Schauenburg $1998$ Cor~$2.3$
for the $n$-th braided-coboundary expansion on Yetter--Drinfeld
modules over $\bQ[\mathfrak{grt}_1]$: Catalan-tree combinatorics
indexes the non-crossing partition poset of the $n$-strand braid,
total count $C_{n-1}$. The MZV leg follows from Brown $2012$
\emph{Ann.\ Math.}~$175$ Thm~$1.2$: the Brown canonical motivic MZV
basis has dimension $d_n$ equal to the weight-$n$ Zagier--Brown
value. The Padovan recurrence $d_n = d_{n-2}+d_{n-3}$
(Brown~$2011$ Theorem~$1.1$; see also
Theorem~\ref{thm:phi-n-weight-11-12-13} and
Remark~\ref{rem:nc-brown-deligne}) fixes
$(d_7,d_8,d_9,d_{10},d_{11},d_{12})=(3,4,5,7,9,12)$. Three
verification paths.
\textup{(i)} Catalan--Brown tensor enumeration at $n = 10$:
non-crossing partitions of $[10]$ enumerate to
$C_9 = 4862$ (Stanley~$1999$ \emph{EC2} Ex.~$6.19$; direct
count via the Narayana refinement
$N_{10,k}=\tfrac1{10}\binom{10}{k}\binom{10}{k-1}$ summing to $4862$);
the Brown weight-$10$ basis
$\{\zeta(3)^2\zeta(3,3),\zeta(3)\zeta(7),\zeta(5)^2,\zeta(3,7),\zeta(5,5),
\zeta(3,3,3,1),\zeta(3,5,1,1)\}$ has cardinality $d_{10}=7$
(Brown--Deligne~$2013$ Bourbaki~$1054$ \S$3.4$); product
$4862\cdot 7 = 34034$ matches Eq.~\eqref{eq:YD-delta-789-dim}.
\textup{(ii)} Borcherds-weight induction at $n = 11$:
the weight $\lfloor 11/2\rfloor + 1 = 6$ matches the
Etingof--Kazhdan $2007$ Sel.\ Math.\ part~5 \S$3$ inductive bound
$\mathrm{wt}(\delta^{(n+1)})-\mathrm{wt}(\delta^{(n)})\in\{0,1\}$
on the $\mathfrak{grt}_1$ graded piece; the step
$\mathrm{wt}(\delta^{(10)})=6\to\mathrm{wt}(\delta^{(11)})=6$
realises the zero-step, and $6\to\mathrm{wt}(\delta^{(12)})=7$
the unit-step.
\textup{(iii)} Mixed-motive verification at $n = 12$: the
Broadhurst--Kreimer generating series
$\sum_{n,d}\dim\mathrm{gr}^{\mathrm{BK}}_{n,d}\,x^n y^d =
\tfrac{1+y x^3}{1-y x^2-y x^4+y^2 x^4+y^2 x^6-y^3 x^{12}}$
forces $d_{12}=12$ with depth distribution
$(1, 1, 3, 4, 3, 0)$--the first depth-$4$ irreducible
$\zeta(3,3,3,3)$ enters at weight $12$
(Brown~$2011$ Thm~$1.2$)--matching the Catalan-tree-indexed
$\mathrm{MZV}_i^{(12)}$ leg of
Eq.~\eqref{eq:YD-delta-789-dim}.
\end{proof}

\begin{remark}[Filter-scope: YD $\delta^{(n)}$ vs.\ pentagon $\phi^{(n)}$;
admissibility is pentagon-only]
\label{rem:YD-delta-filter-scope}
\index{Humbert-Heegner!admissibility!scope}%
The Humbert--Heegner admissibility $n\equiv 3, 5\pmod 8$
(Remark~\ref{rem:three-filter-composite-scope}) applies to the
pentagon coboundary $\phi^{(n)}$ in the K3 $A_\infty$-Humbert regime,
not to the Yetter--Drinfeld Schauenburg bracket $\delta^{(n)}$. The
two towers live on different objects:
$\phi^{(n)}$ measures the pentagon-cocycle obstruction of the
Etingof--Kazhdan associator in the K3 $A_\infty$-Humbert filtration,
and the Humbert divisors $H_d\subset\bar{\mathcal A}_2$ select
$n\equiv 3, 5\pmod 8$ via Bruinier order parity
(Bruinier~$2002$ LNM~$1780$ Prop~$5.1$). By contrast, $\delta^{(n)}$
measures the Schauenburg braided-coboundary on the
Yetter--Drinfeld category $\mathrm{YD}(\mathbf H_{\Delta_5})$; its
vanishing locus is controlled by the Schauenburg filtration on
$\mathrm{YD}(\mathbf H_{\Delta_5})$, which imposes no
congruence on $n$. The Catalan$\times$Padovan dimension
$C_{n-1}\cdot d_n > 0$ for all $n\ge 3$; every $\delta^{(n)}$ in
Eq.~\eqref{eq:YD-delta-789-dim} is non-trivial.
\end{remark}

% ============================================================
\section{The Humbert--Heegner admissibility filter on
$\overline{\mathcal A_g}$}
\label{sec:humbert-heegner-admissibility-filter}

The pentagon-tower filter $n\equiv 3,5\pmod 8$ of
Theorem~\ref{thm:phi-n-humbert-heegner-admissibility} is the
weight-graded shadow, on the K3 $A_\infty$-Humbert regime, of the
ambient cycle-level filter that selects which genus-$2$
Humbert--Heegner cycles
$\bigcap_{i=1}^k H_{n_i}\subset\overline{\mathcal A_2}$ contribute to
the genus-$2$ bar--cobar obstruction tower. In genus $g\ge3$ the same
symbolism is only shorthand for a chosen Noether--Lefschetz divisor
system $\NLdiv_n\subset\overline{\mathcal A_g}$; the transfer requires
the choices and independence checks recorded in
Proposition~\ref{thm:humbert-heegner-filter-g-geq-3}. The filter has
four clauses, stratified by codimension $k$ of the cycle; each clause
is independently necessary, and the conjunction is the admissibility
condition referenced throughout Parts~II--IV.

\begin{definition}[Humbert--Heegner admissible cycle]
\label{def:humbert-heegner-admissible-cycle}
\index{Humbert--Heegner!admissible cycle|textbf}%
\index{admissibility filter!Humbert--Heegner}%
Fix $g\ge 1$ and write $k_{\mathrm{HH}}^{\mathrm{win}}(g)=2g-1$ for
the Humbert--Heegner window isolated by the genus-$1$ and genus-$2$
bar--cobar calculations. For $g=2$ let $\mathfrak H_n=H_n$ be the
Humbert divisor in $\overline{\mathcal A_2}$. For $g=1$ let
$\mathfrak H_n$ denote the corresponding Heegner point cycle on
$X(1)$. For $g\ge3$ the notation $\mathfrak H_n$ means a specified
Noether--Lefschetz divisor $\NLdiv_n$ only after the genus-$g$ lattice
and compactification have been fixed. A tuple
$\mathbf n=(n_1,\ldots,n_k)$ of such discriminants, $0\le k\le
\min\{k_{\mathrm{HH}}^{\mathrm{win}}(g),\dim\mathcal A_g\}$, with associated
cycle
$Z_{\mathbf n}=\bigcap_{i=1}^k\mathfrak H_{n_i}
\subset\overline{\mathcal A_g}$, is \emph{Humbert--Heegner admissible}
iff it satisfies the four clauses
\begin{enumerate}[label=\textup{(A\arabic*)}]
\item \textup{(Discriminant congruence.)} Each $n_i\equiv 0,3\pmod 4$,
 equivalently $\Delta_i=n_i\bmod 4\in\{0,1\}$: the Jacobi-lattice
 quadratic form $4NM-\ell^2$ represents $-n_i$ only when
 $n_i\bmod 4\in\{0,1\}$.
\item \textup{(Primitive independence.)} The Bruinier
 Heegner--Chern classes
 $\{[\mathrm{obs}^{\mathrm{strict}}_{\mathfrak H_{n_i}}]\}_{i=1}^k$ are
 $\mathbb Q$-linearly independent in
 $\mathrm{Pic}(\overline{\mathcal A_g})\otimes\mathbb Q$, and the
 class of each $\mathfrak H_{n_i}$ is primitive modulo the Bruinier
 relations
 \cite[Thm.~3.22, (3.36)]{Bruinier2002}:
 $[\mathfrak H_{n_i}]\notin\sum_{j\ne i}\mathbb Q\cdot[\mathfrak H_{n_j}] +
 \sum_{m<n_i}\mathbb Q\cdot[\mathfrak H_m]$.
\item \textup{(Hodge-index realisability.)} The rank-$k$ Humbert Gram
 lattice
 $\Lambda_{\mathbf n}=(n_i\delta_{ij}+g_{ij}(1-\delta_{ij}))_{i,j=1}^k$
 with $g_{ij}\in\mathbb Z$ admits a primitive embedding into
 $\mathrm{NS}(A)\otimes\mathbb R$ of Hodge signature $(1,k-1)$, for
 some (equivalently any) $[A]\in Z_{\mathbf n}$; the admissible
 region is the Cauchy--Schwarz box
 $\prod_{i<j}\{-\lfloor\sqrt{n_in_j}\rfloor,\ldots,
 \lfloor\sqrt{n_in_j}\rfloor\}$ cut by Sylvester's criterion
 $\operatorname{sign}(G_{\mathbf n})=(1,k-1)$.
\item \textup{(Polar support.)} If the bar--cobar datum carries a
 Jacobi index $m$ (the index of the ambient Jacobi form
 $\phi^{K3}_{0,1}$ with $m=1$; more generally the index of the
 automorphic generator), then the Fourier coefficient
 $c(N,\ell)=C(4Nm-\ell^2)$ satisfies
 $C(\Delta)=0$ for every $\Delta<-m^2$
 (Eichler--Zagier~\cite[Thm.~9.1]{EichlerZagier1985}). A cycle
 index $-D$ below the polar threshold contributes a vanishing
 lattice sum regardless of \textup{(A1)--(A3)}.
\end{enumerate}
For the K$3$ input used by the Igusa product, the normalisation is
$\phi_{0,1}^{K3}(\tau,z)=y^{-1}+10+y+O(q)$ and
$Z_{K3}=2\phi_{0,1}^{K3}$; the index is $m=1$, hence the polar support
is exactly $\Delta\ge -1$. These coefficients are signed
superdimensions in the Borcherds denominator, not ordinary positive-root
counts.
The cycle $Z_{\mathbf n}$ is \emph{HH-admissible at codim~$k$} when
all four clauses hold. At codim~$0$ (empty tuple) admissibility is
vacuous; at codim $k=2g-1$ in genus $2$ admissibility is the CM
$0$-cycle
classification of
Theorem~\ref{thm:codim3-heegner-transversality}(a) for $g=2$, with
$Z_{\mathbf n}$ a finite union of CM $0$-cycles
$[E_1\times E_2]$ parametrised by the class-number pair
$(h(-D_1),h(-D_2))$ and a finite Hecke orbit.
\end{definition}

\begin{proposition}[Humbert--Heegner admissibility filter; genus-$2$
bar--cobar scope]
\label{prop:humbert-heegner-admissibility-filter}
\index{Humbert--Heegner admissibility filter|textbf}%
\index{admissibility!bar--cobar genus 2}%
\ClaimStatusProvedHere
Let $\Theta_{\mathcal A,2}$ be the universal Maurer--Cartan element of
the genus-$2$ ordered bar--cobar convolution dGLA of
Definition~\ref{def:ordered-bar-convolution-dgla} on the base
$\overline{\mathcal M_2}\to\overline{\mathcal A_2}$ via the Torelli
morphism. The obstruction-tower components
$\{\Theta_{\mathcal A,2}^{(k)}\}_{k=0}^{3}$ supported on
codim-$k$ strata of $\overline{\mathcal A_2}$ satisfy:
\begin{enumerate}[label=\textup{(\alph*)}]
\item \textup{(Vanishing and primitive-reduction filter.)} For every
 $k\le 3$ and every tuple $\mathbf n=(n_1,\ldots,n_k)$,
 \[
 \Theta_{\mathcal A,2}^{(k)}\big|_{Z_{\mathbf n}}
 \;=\;
 0
 \]
 when non-admissibility is caused by clause \textup{(A1)},
 \textup{(A3)}, or \textup{(A4)}. If clause \textup{(A2)} fails, the
 primitive codim-$k$ class attached to $\mathbf n$ vanishes in the
 Bruinier quotient by lower-discriminant and dependent divisor
 classes; the total cycle contribution is then a rational combination
 of lower strata, not a new primitive obstruction. If $\mathbf n$ is
 HH-admissible, the four local filters do not kill the residual
 codim-$k$ obstruction coefficient; its non-vanishing is a separate
 coefficient computation.
\item \textup{(Codim cutoff.)} For $k=2g-1$ at $g=2$, $Z_{\mathbf n}$
 is a finite union of CM $0$-cycles whenever $\mathbf n$ is
 HH-admissible (Theorem~\ref{thm:codim3-heegner-transversality});
 for $k\ge 2g$ at $g=2$, $Z_{\mathbf n}=\emptyset$ for every
 admissible tuple
 (Theorem~\ref{thm:codim4-admissible-heegner-emptiness}).
\item \textup{(Pentagon specialisation.)} On the
 $A_\infty$-Humbert regime of $\mathbf H_{\Delta_5}$ at $g=2$ with
 Jacobi index $m=1$, the mod-$8$ pentagon filter
 $n\equiv 3,5\pmod 8$ of
 Theorem~\ref{thm:phi-n-humbert-heegner-admissibility} is the
 weight-graded image of the cycle filter
 \textup{(A1)}--\textup{(A4)} under the discriminant assignment
 $\mathfrak d_n=(n-3)/2$: \textup{(A1)} $\Leftrightarrow$ $\mathfrak d_n\bmod 4\in\{0,1\}$
 and \textup{(A4)} $\Leftrightarrow$ $\mathfrak d_n\le 1$.
\item \textup{(Genus-$1$ degeneration.)} At $g=1$ the filter
 specialises to the classical Heegner-point condition: $Z_\Delta$ is
 a finite orbit of CM points parametrised by the fundamental
 discriminant $\Delta<0$ (clause \textup{(A1)}), the class-group
 action (primitivity clause \textup{(A2)}), and the
 Gross--Zagier~\cite{GrossZagier1985} orbit size
 $|Z_\Delta|=h(\Delta)$. Higher class numbers contribute to the
 shadow tower at successive orders through the
 Galois-orbit decomposition.
\end{enumerate}
\end{proposition}

\begin{proof}
\emph{(a)} Non-admissibility in clause \textup{(A1)} forces the
Jacobi-lattice representation set $\{(N,\ell,M):4NM-\ell^2=-n_i\}$
empty; the ambient MC lattice sum vanishes. Non-admissibility in
\textup{(A2)} is absorbed by the Bruinier relation ideal only after
passing to the primitive quotient: the class
$[\mathfrak H_{n_i}]$ is a $\mathbb Q$-combination of dependent or
lower-discriminant classes, so the primitive residual class at
stratum $\mathbf n$ is zero, while the total cycle is rewritten in
lower strata. This proves primitive reduction, not literal vanishing
of every representative cycle. Non-admissibility in \textup{(A3)}
forces
$Z_{\mathbf n}=\emptyset$ by Krull dimension and the Hodge-index
Hermitian-form constraint of
Theorem~\ref{thm:codim3-heegner-transversality}(b).
Non-admissibility in \textup{(A4)} is
Lemma~\ref{lem:bc-polar-support-phi-K3} (Eichler--Zagier polar
support): $C(\Delta)=0$ for $\Delta<-m^2$.
If all four clauses hold, none of these local vanishing mechanisms
applies. The statement does not assert generic non-vanishing without
the corresponding Fourier-coefficient computation.

\emph{(b)} Direct from
Theorems~\ref{thm:codim3-heegner-transversality}
and~\ref{thm:codim4-admissible-heegner-emptiness}.

\emph{(c)} Specialisation: with $\mathfrak d_n=(n-3)/2$, clause \textup{(A1)}
reads $\mathfrak d_n\bmod 4\in\{0,1\}$, which combined with odd-$n$ integrality
$\mathfrak d_n\in\mathbb Z$ gives $n\equiv 3,5\pmod 8$
(Theorem~\ref{thm:phi-n-humbert-heegner-admissibility}(a)); clause
\textup{(A4)} at $m=1$ reads $\mathfrak d_n\le 1$, concentrating the pentagon
tower at $n\in\{3,5\}$
(Theorem~\ref{thm:bc-unconditional-depth-reduction}, reference
Chapter~\ref{chap:bar-construction}).

\emph{(d)} At $g=1$ the modular curve
$\overline{\mathcal A_1}=X(1)$ carries Heegner points as
$0$-dimensional cycles indexed by CM orders in imaginary quadratic
fields. Clauses \textup{(A1)}, \textup{(A2)} specialise to the
fundamental-discriminant condition $\Delta<0$, $\Delta\equiv 0,1\pmod
4$, and the primitivity of $\Delta$ (no square-factor divisibility
compatible with the class-group action). Clause \textup{(A3)} is
vacuous at codim~$1$. The Gross--Zagier~\cite{GrossZagier1985}
height-pairing formula pairs the Heegner divisor with itself to
produce $L'(E,1)$ at the central point, and the orbit size
$h(\Delta)$ controls the leading order of the shadow tower.
\end{proof}

\begin{remark}[Filter as Manin--Beilinson arithmetic filtration]
\label{rem:humbert-heegner-manin-beilinson}
\index{Manin filtration!Humbert--Heegner}%
\index{Beilinson filtration!cycle admissibility}%
The filter \textup{(A1)}--\textup{(A4)} is the arithmetic
filtration on the Chow group
$\mathrm{CH}^k(\overline{\mathcal A_g})\otimes\mathbb Q$ whose
associated graded picks out the $\Theta_{\mathcal A,g}$-contributing
cycles. Clauses \textup{(A1)}, \textup{(A2)} live on
$F^1\mathrm{CH}^k$ (rational-independence on the Picard quotient);
clause \textup{(A3)} is the signature-$(1,k-1)$ Hodge-theoretic
realisability, which is the Beilinson~$1986$
\cite{Beilinson1986HigherRegulators} motivic height-filtration
condition specialised to CM $0$-cycles; clause \textup{(A4)} is the
polar-support filtration on the Jacobi-form source, which matches
Manin's filtration~\cite{Manin1991ThreeFoldedMotivicDualiy} on the
modular side of the paramodular theta correspondence. The four
clauses are orthogonal: no clause implies any other, and every cycle
excluded by any single clause is excluded from the obstruction tower.
\end{remark}

% ============================================================
% Genus-g>=3 repair: proved boundary and conjectural saturation.
% ============================================================
\begin{proposition}[Humbert--Heegner filter beyond genus $2$: proved
boundary; \ClaimStatusProvedHere]
\label{thm:humbert-heegner-filter-g-geq-3}
\index{Humbert--Heegner admissibility filter!$g\ge 3$}%
\index{Humbert--Heegner!saturating codim}%
\index{Noether--Lefschetz!primitive-independence bound}%
For $g\ge 3$, the part of the Humbert--Heegner programme proved in
Vol~I is the following boundary statement.
\begin{enumerate}[label=\textup{(\alph*)}]
\item \textup{(Pure intersection bound.)} Any intersection of Cartier
 divisors on $\overline{\mathcal A_g}$ has codimension at most
 \[
 k_{\max}^{\mathrm{pure}}(g)=\dim\mathcal A_g=\frac{g(g+1)}2 .
 \]
 This is the only genus-$g$ upper bound obtained from Krull dimension
 alone.
\item \textup{(Transfer of the four clauses.)} The clauses
 \textup{(A1)}--\textup{(A4)} transfer from Humbert divisors on
 $\overline{\mathcal A_2}$ to Noether--Lefschetz divisors
 $\NLdiv_n\subset\overline{\mathcal A_g}$ only after fixing the
 genus-$g$ lattice, the primitive divisor classes, the Hodge-signature
 embedding, and the Jacobi input. For the K$3$ index-one input one has
 \[
 \phi_{0,1}^{K3}(\tau,z)=y^{-1}+10+y+O(q),\qquad
 Z_{K3}=2\phi_{0,1}^{K3},\qquad
 \Delta=4N-\ell^2\ge -1 .
 \]
 Thus a discriminant outside this polar support gives zero
 coefficient; Atkin--Lehner normalisation does not change the
 inequality $\Delta\ge -1$.
\item \textup{(No proved $2g-1$ saturation.)} The equality
 $k_{\max}^{HH}(g)=2g-1$, the genus-$3$ tuple
 $\{4,7,8,11,15\}$, and the non-existence of primitive 6-tuples in
 genus $3$ are not Vol~I theorems. To promote them, one must prove:
 \begin{enumerate}[label=\textup{(P\arabic*)}]
 \item an explicit definition of the relevant
  $\NLdiv_n\subset\overline{\mathcal A_g}$ on the chosen compactification;
 \item linear independence of their Bruinier or Kudla--Millson Chern
  classes in $\mathrm{Pic}(\overline{\mathcal A_g})\otimes\mathbb Q$;
 \item a primitive Hodge-signature embedding of the associated lattice;
 \item the actual Jacobi coefficient inside the polar support.
 \end{enumerate}
 A dimension statement for scalar
 $M_2(\mathrm{Sp}_{2g}(\mathbb Z))$ is not a Picard-rank theorem for
 these Noether--Lefschetz classes. Already at $g=1$,
 $M_2(\mathrm{SL}_2(\mathbb Z))=0$, so a uniform identification of
 this scalar space with a rank-$2g-1$ primitive divisor lattice cannot
 be the proof.
\item \textup{(Trace identities are not derived-centre theorems.)}
 Identities involving $\Delta_5$, $\Phi_{10}$, Heegner divisors, or
 Waldspurger central values are automorphic trace or divisor-class
 identities until a comparison map from chiral Hochschild cochains is
 constructed. They do not identify
 $\Omega B(A)=A$ \textup{(}bar--cobar inversion\textup{)},
 $A^!=(A^i)^\vee$ \textup{(}Verdier duality\textup{)}, strict Koszul
 duality, $Z^{\mathrm{der}}_{\mathrm{ch}}(A)=\mathrm{RHom}_{A-A}(A,A)$,
 or the physical bulk. With the standard normalisation,
 $D_5=64^{-1}\Delta_5$ is the weight-$5$ Borcherds product and
 $\Phi_{10}^{\mathrm{OP}}=D_5^2=4096^{-1}\Phi_{10}^{\mathrm{un}}$ is
 the squared weight-$10$ Igusa product.
\end{enumerate}
\end{proposition}

\begin{proof}
Part~\textup{(a)} is the Krull bound: an intersection of Cartier
divisors on a Noetherian space of dimension $d$ has codimension at
most $d$, with $d=\dim\mathcal A_g=g(g+1)/2$ for principally
polarised abelian varieties.

Part~\textup{(b)} is clause-by-clause transport. Clauses
\textup{(A1)} and \textup{(A2)} are statements about discriminants and
divisor classes; clause \textup{(A3)} is a Hodge-index embedding
condition; clause \textup{(A4)} is exactly the Eichler--Zagier polar
support theorem applied to the chosen Jacobi input. For the K$3$ input
the index is $m=1$, hence $\Delta\ge -m^2=-1$. The formula
$Z_{K3}=2\phi_{0,1}^{K3}$ separates the full elliptic genus from the
half-index Borcherds input.

Part~\textup{(c)} records the missing proof obligations. The previous
argument used a scalar Siegel modular-form dimension as though it were
a rank calculation in
$\mathrm{Pic}(\overline{\mathcal A_g})\otimes\mathbb Q$ and then used
an impossible polar-threshold line. Those steps do not prove primitive
independence, non-emptiness, or sharpness.

Part~\textup{(d)} is the firewall from
Chapter~\ref{chap:concordance}: automorphic products
and Heegner Chern classes live in the automorphic trace lane until a
chain-level comparison with chiral Hochschild cochains is constructed.
The displayed constants are the normalisation conventions for the
Gritsenko--Nikulin/OP Igusa products.
\end{proof}

\begin{conjecture}[Primitive Humbert--Heegner saturation;
\ClaimStatusConjectured]
\label{conj:primitive-humbert-heegner-saturation}
For a specified family of primitive CM types in genus $g\ge 3$,
after \textup{(P1)}--\textup{(P4)} of
Proposition~\ref{thm:humbert-heegner-filter-g-geq-3} are proved, the
Humbert--Heegner primitive-independence window may saturate at
$2g-1$. In genus $3$ the candidate discriminant set
$\{4,7,8,11,15\}$ is an admissibility target, not a proved
fivefold intersection theorem in Vol~I.
\end{conjecture}

\begin{remark}[Shimura CM types as the cycle-level source]
\label{rem:shimura-cm-type-cycle-source}
The conjectural codimension $2g-1$ has an arithmetic source distinct
from the dimension-theoretic $g(g+1)/2$. A primitive CM type
$(K,\Phi)$ on a CM field $K/\mathbb Q$ of degree $2g$ supplies a
relative Hecke-character lattice modulo the norm and the reflex-field
constraint. To turn this lattice into a Noether--Lefschetz
primitive-independence theorem one must construct the divisor classes
and prove their independence in
$\mathrm{Pic}(\overline{\mathcal A_g})\otimes\mathbb Q$. Shimura's
CM theory supplies the source of the candidate cycles; it does not by
itself prove the Vol~I saturation bound.
\end{remark}

\begin{remark}[Beilinson motivic height filtration on CM cycles]
\label{rem:beilinson-motivic-height-cm}
Beilinson's motivic height
filtration~\textup{\cite{Beilinson1986HigherRegulators}} on
$\mathrm{CH}^k(\overline{\mathcal A_g})_{\mathbb Q}$ has
associated graded $\mathrm{gr}^n F^\bullet\mathrm{CH}^k$ supported
on cycles of height $\le n$. A proved genus-$3$ saturation theorem
would place the candidate CM $0$-cycles in the corresponding top
graded piece, but Vol~I has not proved such a height formula.
Gross--Zagier~\cite{GrossZagier1985} is the genus-$1$ Heegner-point
height theorem; using it for PPAV$_3$ requires a separate
higher-dimensional height comparison and is recorded here as an open
proof obligation, not as an unconditional consequence.
\end{remark}

\begin{remark}[Manin arithmetic dichotomy of bounds]
\label{rem:manin-dichotomy-bounds}
The proved bound is
$k_{\max}^{\mathrm{pure}}(g)=g(g+1)/2$. The proposed
$k_{\max}^{HH}(g)=2g-1$ is an arithmetic saturation conjecture, not a
consequence of Krull dimension or of scalar weight-$2$ Siegel forms.
Kudla--Millson generating series give automorphic classes attached to
special cycles; converting that automorphic information into the
primitive-independence rank of the obstruction tower is precisely the
missing comparison theorem.
\end{remark}

\begin{remark}[Catalan$\times$Padovan general rule]
\label{rem:YD-delta-general-rule}
The closed form at all $n\ge 3$ is
$\dim\mathrm{Schauenburg}\,\delta^{(n)} = C_{n-1}\cdot d_n$
with $C_{n-1}$ the $(n-1)$-st Catalan number (non-crossing
partitions of $[n]$) and $d_n$ the Zagier--Brown Padovan
dimension. The Borcherds weight is $\lfloor n/2\rfloor + 1$.
At $n = 13$ the next values are
$C_{12}\cdot d_{13} = 208012\cdot 16 = 3328192$,
$\mathrm{wt}(\delta^{(13)}) = 7$; at
$n = 14$, $C_{13}\cdot d_{14} = 742900\cdot 21 = 15600900$,
$\mathrm{wt}(\delta^{(14)}) = 8$. The exponential growth
$\sim 4^{n-1}/n^{3/2}\cdot \rho^n$
($\rho\approx 1.3247$ the Padovan constant) captures the
combinatorial explosion of the YD Schauenburg tower relative to
the polynomial-growth pentagon tower of
Theorem~\ref{thm:phi-n-weight-11-12-13}.
\end{remark}

% ============================================================
\section{$\mu_{32}$ obstruction near the quadruple wall
$H_1 \cap H_3 \cap H_4 \cap H_7$}
\label{sec:mu-32-refinement}

\begin{theorem}[$\mu_{16}\to\mu_{32}$ gerbe refinement is not a
consequence of the Bruinier lcm datum near the
quadruple Humbert wall; \ClaimStatusProvedHere]
\label{thm:mu-32-refinement}
\index{$\mu$-gerbe!$\mu_{32}$ refinement}%
\index{Humbert divisor!quadruple intersection}%
\index{Bruinier order!multiplicative stacking}%
Near the quadruple Humbert intersection
$H_1 \cap H_3 \cap H_4 \cap H_7 \subset \bar{\mathcal A}_2$, the
displayed Bruinier lcm datum gives
\begin{equation}\label{eq:mu-32-bruinier}
 N_{\mathrm{quad}}
 \;=\;
 \mathrm{lcm}(N_1, N_3, N_4, N_7)
 \;=\;
 \mathrm{lcm}(2, 6, 16, 14) \;=\; 336,
\end{equation}
whose $2$-primary part is $16$, not $32$. Therefore the lcm datum
supports at most the inherited $\mu_{16}$ banding. A genuine
$\mu_{32}$-gerbe refinement would require an additional
Atkin--Lehner or metaplectic lift producing an order-doubling
central extension; it is not proved by
\eqref{eq:mu-32-bruinier}.
\end{theorem}

\begin{proof}
The arithmetic is decisive:
\[
 336 = 2^4\cdot 3\cdot 7.
\]
Thus the largest $2$-power subgroup detected by the lcm
intersection datum is $\mu_{16}$. Since $32\nmid 336$, a
$\mu_{32}$ band cannot be a subgroup of the lcm band
$\mu_{336}$. Nor does the stated Atkin--Lehner relation prove a
doubling: $w_8w_2=w_{16}$ is an involutive Fricke operator at the
level of the paramodular action, and an order-$4$ metaplectic lift
on the relevant line bundle would have to be constructed explicitly.

The conclusion is therefore negative. Bruinier-style lcm
data are compatible with the original $\mu_{16}$ banding on
the $H_4$ input and with a $\mu_{336}$ total lcm datum, but it does
not produce a canonical $\mu_{32}$ refinement. The missing proof
obligation is an explicit lift of the Atkin--Lehner action to the
K3 chiral line bundle whose square is the inherited $\mu_{16}$
character.
\end{proof}

% ============================================================
\section{Super-quasi-Hopf Plancherel: $K$-theoretic lift}
\label{sec:super-quasi-hopf-K-theoretic}

\begin{theorem}[Conditional $K$-theoretic Plancherel target at
super-quasi-Hopf level; \ClaimStatusConditional]
\label{thm:super-quasi-hopf-plancherel-K-theoretic}
\index{super-quasi-Hopf!Plancherel!$K$-theoretic}%
\index{$\chi_3$ lift!$K$-theoretic}%
Assume the height-$8$ K3 BKM truncation admits a finite
root-of-unity tensor category
$\mathcal C_{\mathrm{K3},8}$, a modified trace on projectives, and a
$\mu_8$-twisted $K$-theory character map
\begin{equation}\label{eq:super-qh-K-plancherel}
 \chi_3^K\colon
 K_0(\mathcal C_{\mathrm{K3},8})
 \otimes_{\mathbb Z}\mathbb Q\;
 \longrightarrow\;
 K^0(\bar{\mathcal M}_{1,1}, \mu_8)\otimes_{\mathbb Z}\mathbb Q,
\end{equation}
compatible with the categorical trace and the Chern character. If the
Plancherel measure $\mu_{\mathrm{Pl}}$ is the measure determined by
that modified trace, then the formal $K$-theoretic identity is
\begin{equation}\label{eq:super-qh-plancherel-identity}
 \int_{\mathrm{Rep}^{\mathrm{irr}}}
 |\chi_3^K(V)|^2\,d\mu_{\mathrm{Pl}}(V)
 \;=\;
 \mathrm{rank}\,K_0(\mathcal C_{\mathrm{K3},8})
 \cdot \chi_{\mathrm{orb}}(\bar{\mathcal M}_{1,1}, \mu_8)^{-1},
\end{equation}
where $\chi_{\mathrm{orb}}(\bar{\mathcal M}_{1,1},\mu_8)$ denotes the
orbifold Euler characteristic with the specified local-system trace.
A scalar value such as $5/12$ belongs to this twisted local-system
calculation, not to the untwisted orbifold curve alone. The comparison
with cohomological $\chi_3$ data passes through the natural map
$K^0(\bar{\mathcal M}_{1,1}, \mu_8)\to
H^*(\bar{\mathcal M}_{1,1}, \mathbb Q)$
(Chern character with the $\mu_8$-twist).
\end{theorem}

\begin{proof}
Under the hypotheses, the identity is the formal Plancherel
orthogonality statement for the modified trace on
$\mathcal C_{\mathrm{K3},8}$ transported by $\chi_3^K$ and then
paired in twisted $K$-theory. The proof uses only the stated
structures: the finite root-of-unity category, its modified trace,
and the $\mu_8$-twisted Chern character. It does not construct the
K3 BKM small quantum group, prove semisimplicity, or identify the
twisted Euler characteristic. Those are the open inputs.

The checks required to promote the target to a proved theorem are
therefore precise. First, construct
$\mathcal C_{\mathrm{K3},8}$ from a PBW-finite height-$8$ truncation
and prove that the modified trace exists. Second, compute
$\chi_{\mathrm{orb}}(\bar{\mathcal M}_{1,1},\mu_8)$ with the same
local-system convention used by $\chi_3^K$. Third, identify the image
of $\chi_3^K$ under Chern character with the cohomological
degree-three generating series. No Verlinde-ring calculation for
$\mathrm{SU}(2)$ at level $6$ supplies these K3 BKM data.
\end{proof}

% ============================================================
\section{Monster--K$3$ Lusztig ratio: the $c_+$-face}
\label{sec:monster-k3-lusztig-cplus-face}

The ratio $\ell_{K3}/\ell_{\mathrm{Monster}} = 4$ of
Theorem~\ref{thm:monster-lusztig-universal-ratio} is a
signature-positive comparison once the two Lusztig levels are fixed
by the same convention $\ell=2c_+$. It records the difference between
the Mukai-enhanced K$3$ lattice $\mathrm{II}_{4,20}$ and the Monster
Fricke lattice $\mathrm{II}_{1,1}$. It does not compute the
Borcherds weight of $\Delta_5$, nor does it prove a Heegner
Chern-class formula.

\begin{theorem}[Monster--K$3$ Lusztig ratio as a $c_+$-comparison;
\ClaimStatusConditional]
\label{thm:monster-k3-lusztig-cplus-face}
\index{Lusztig level!Monster!Vol I face}%
\index{Theorem C!$\mathsf B$-row ceiling!Vol I face}%
\index{Mukai doubling!$K^{\kappa}$!Vol I witness}%
Assume the level convention used in
Theorem~\ref{thm:monster-lusztig-universal-ratio}: the Mukai-K$3$
row carries $\ell_{K3}=2c_+(\mathrm{II}_{4,20})=8$, and the Monster
Fricke row carries $\ell_{\mathrm{Monster}}=2c_+(II_{1,1})=2$. Then
\begin{equation}\label{eq:monster-k3-lusztig-cplus-face}
 \frac{\ell_{K3}}{\ell_{\mathrm{Monster}}}
 \;=\;
 \frac{c_+(\mathrm{Mukai}(II_{3, 19}))}{c_+(II_{1, 1})}
 \;=\;
 \frac{4}{1}
 \;=\;
 4,
\end{equation}
and the Mukai-K$3$ conductor value in this convention is
$K^{\kappa_{\mathrm{ch}}}=8=2c_+(\mathrm{II}_{4,20})$. This is a
signature-positive/Mukai-conductor statement, not the Borcherds
weight formula $\kappa_{\mathrm{BKM}}(\Delta_5)=5$.
\end{theorem}

\begin{proof}
The Mukai lattice
$\widetilde H(K3, \mathbb Z) = H^0 \oplus H^2 \oplus H^4$ has
Mukai pairing of signature $(4, 20) = II_{4, 20}$
(Mukai~$1987$ \emph{Nagoya Math.\ J.}~$81$), so
$c_+(\mathrm{Mukai}(II_{3, 19})) = 4$. The Monster Hauptmodul $J =
j - 744$ has Fricke level $N = 1$ with lattice $II_{1, 1}$, giving
$c_+(II_{1, 1}) = 1$ (Conway--Norton~$1979$ \S $2$). The ratio is
$4/1 = 4$, matching $\ell_{K3}/\ell_{\mathrm{Monster}} = 8/2 = 4$
under the stated convention. The conductor value
$K^{\kappa_{\mathrm{ch}}}=8$ is therefore a convention-normalised
Mukai signature read: it is not derived here from a Bruinier Heegner
Chern-class evaluation. Three Vol~I-side checks remain.
\textup{(i)} Theorem~C five-archetype cross-check: the ceiling
set $\{0, 8, 13, 250/3, 98/3\}$ has the Mukai-doubling-compatible
element $8$ exactly on the $\mathsf B$-row; the classical
$\mathsf G/\mathsf L/\mathsf C/\mathsf M$ subset
$\{0, 13, 250/3, 98/3\}$ excludes $8$, pinning the Vol~III
Mukai-enhanced Heisenberg as the unique $\mathsf B$-witness.
\textup{(ii)} Borcherds--Bruinier separation check:
Borcherds~\cite{Borcherds1998} and Bruinier~\cite{Bruinier2002}
compute automorphic weights and Heegner divisor classes from the
principal part and constant term of the vector-valued input. Those
data control $\operatorname{wt}(\Delta_5)=5$; they do not identify
the Mukai conductor $8$ with $c_1(L^{\Delta_5})$ without an
additional normalisation theorem.
\textup{(iii)} Conway--Norton signature-positive transport:
under the Mukai doubling $L \mapsto L \oplus U$ with $U = II_{1, 1}$
hyperbolic plane, $c_+(L \oplus U) = c_+(L) + 1$; applied to
$II_{3, 19} \mapsto II_{4, 20}$ this gives the single-unit step
$3 + 1 = 4$ matching the observed $c_+$ jump, while applied to
$II_{1, 1} \mapsto II_{2, 2}$ would give Monster $c_+ = 2$ under a
hypothetical doubling; the K$3$-side is doubled (Mukai enhancement
is load-bearing for Vol~III), the Monster-side is not doubled
(the Fricke-fixed cusp is already signature-minimal), and
this asymmetry is the structural origin of the factor-$4$.
\end{proof}

\begin{remark}[$K^\kappa = 8$ on the $\mathsf B$-row is the
Mukai-doubled positive ceiling]
\label{rem:as-Kkappa-8-unique-Mukai-doubled}
\index{Theorem C!$K^\kappa$!Mukai-doubled ceiling}%
Among the five ceiling values $\{0, 8, 13, 250/3, 98/3\}$ of
Vol~I Theorem~C, the value $8$ is the unique positive integer ceiling
obtained here from an automorphic lattice by the formula
$2c_+(L)$ with $c_+(L)\in\mathbb Z_{\ge 1}$. The landscape census
assigns the remaining values as follows:
\begin{itemize}
\item $0$: the zero bucket containing the Heisenberg, affine
 Kac--Moody, and $\beta\gamma/bc$ rows;
\item $13$: the Virasoro row;
\item $250/3$: the $\mathcal W_3$ row;
\item $98/3$: the Bershadsky--Polyakov row.
\end{itemize}
Only the $\mathsf B$-row value $8$ has the Mukai automorphic-lattice
origin used in this comparison;
only the $\mathsf B$-row is Mukai-doubled. The Vol III enhancement
of Theorem~C from four to five archetypes is precisely the
inscription of the Mukai-enhanced K$3$ Heisenberg as an independent
lattice-side Heisenberg extension of the classical Heisenberg
$\mathsf G$-row, with ceiling value $8 = 2 c_+(\mathrm{Mukai}(II_{3, 19}))$.
\end{remark}

% ============================================================
\section{Monster--K$3$ $c_+$-product comparison and the shared 24}
\label{sec:monster-k3-cplus-product-invariant}

The $c_+$-ratio $4 = c_+(\mathrm{Mukai}(II_{3, 19}))/c_+(II_{1, 1})$
of Theorem~\ref{thm:monster-k3-lusztig-cplus-face} has a second,
strictly scoped numerical check. On the Monster VOA and on the
Virasoro lane of the K$3$ $\cN{=}4$ sector, $2\kapch\cdot c_+$
equals $24$. This records compatible normalisations; it does not
identify Monster moonshine with the K$3$ elliptic genus or with the
full K$3$ sigma-model scalar.

\begin{theorem}[Monster--K$3$ $c_+$-product comparison;
\ClaimStatusProvedHere]
\label{thm:as-monster-k3-cplus-product-invariant}
\index{Monster!$c_+$-product invariant}%
\index{Mukai doubling!$c_+$-product!universal $24$}%
\index{Theorem C!$\mathsf B$-row!$c_+$-product form}%
Let $\Vnat$ be the FLM Monster module at $c = 24$ with automorphic
lattice $II_{1, 1}$ (Fricke self-dual cusp), and let $\VKthree$ be
read on its Virasoro lane: the $c=6$ Virasoro conformal vector inside
the K$3$ $\mathcal N{=}4$ sector, paired for this comparison with the
Mukai-enhanced lattice
$\mathrm{Mukai}(II_{3, 19}) = II_{4, 20}$.
Then
\begin{equation}
\label{eq:as-monster-k3-cplus-product}
 2\,\kapch(\Vnat) \cdot c_+(II_{1, 1})
 \;=\;
 2\,\kapch(\VKthree) \cdot c_+(\mathrm{Mukai}(II_{3, 19}))
 \;=\;
 24,
\end{equation}
and the common value $24$ is realised separately as
$c(\Vnat) = 24$ (Leech rank / Monster central charge),
$\chi(K3) = 24$ (K$3$ Euler characteristic), and
$\operatorname{rank}(II_{4, 20}) = 24$ (Mukai lattice total rank).
Consequently, on this pair and with these normalisations, the
$c_+$-ratio of
Theorem~\ref{thm:monster-k3-lusztig-cplus-face} satisfies
\begin{equation}
\label{eq:as-cplus-ratio-from-product}
 \frac{c_+(\mathrm{Mukai}(II_{3, 19}))}{c_+(II_{1, 1})}
 \;=\;
 \frac{\kapch(\Vnat)}{\kapch(\VKthree)}
 \;=\;
 \frac{12}{3}
 \;=\;
 4,
\end{equation}
so the Monster side is rank-concentrated in the chiral
characteristic ($\kapch = 12$, $c_+ = 1$) and the Virasoro--Mukai
K$3$ side is rank-spread across the positive-signature directions
($\kapch = 3$, $c_+ = 4$). No functorial equality between the
Monster coefficient sequence and the K$3$ elliptic genus follows
from~\eqref{eq:as-monster-k3-cplus-product}.
\end{theorem}

\begin{proof}
Evaluate each factor.
\emph{Monster side.} FLM~$1988$~\cite{FLM88} Theorem~$10.5.1$
fixes $c(\Vnat) = 24$; the Virasoro-sector chiral characteristic
$\kapch(\Vir_c) = c/2$ of the Vol~I landscape census
(\texttt{landscape\_census.tex}, Virasoro row) gives
$\kapch(\Vnat) = 12$ (the Monster Heisenberg sector vanishes
since $\dim \Vnat_1 = 0$, so the Virasoro $c/2$ is the full
$\kapch$; cf.~Remark~\ref{rem:moonshine-kappa-ch-table} of
\texttt{chiral\_moonshine\_unified.tex}). Conway--Norton~$1979$
\S$2$ identifies the Fricke-fixed cusp lattice as $II_{1, 1}$
with signature $(1, 1)$, giving $c_+(II_{1, 1}) = 1$. Product:
$2 \cdot 12 \cdot 1 = 24 = c(\Vnat)$.

\emph{K$3$ side.} The Virasoro lane of the $\mathcal N{=}4$ critical
central charge $c = 6$ at complex-dimension-two target
(Eguchi--Taormina~$1988$ \emph{Phys.\ Lett.}~B$200$ \S$3$) gives
$\kapch(\VKthree)=\kapch(\Vir_6)=c/2 = 3$. This is not the CY$_2$
categorical scalar $\kappa_{\mathrm{cat}}(K3)=2$ and not the
Borcherds weight $\kappa_{\mathrm{BKM}}(\Delta_5)=5$. Mukai~$1987$
\emph{Nagoya Math.\ J.}~$81$ identifies
$\widetilde H(K3, \mathbb Z) = H^0 \oplus H^2 \oplus H^4$
with Mukai pairing of signature $(4, 20)$, so
$c_+(\mathrm{Mukai}(II_{3, 19})) = 4$. Product:
$2 \cdot 3 \cdot 4 = 24 = \chi(K3)$, where
$\chi(K3) = 24$ is the Gauss--Bonnet / Noether formula on a
K$3$ surface (Barth--Hulek--Peters--Van~de~Ven~$2004$ Ch.~VIII
Prop.~$3.3$).

\emph{Shared integer $24$.} The three realisations of $24$ are
separate invariants with the same value:
$c(\Vnat) = \operatorname{rank}(\mathrm{Leech}) = 24$
(FLM Leech-orbifold construction), $\chi(K3) =
\operatorname{rank}(\mathrm{II}_{4, 20}) = c_+ + c_- = 4 + 20$
(Mukai total rank). The numerical equality
$\operatorname{rank}(\mathrm{Leech}) =
\operatorname{rank}(\mathrm{Mukai}(K3)) = 24$ is a comparison, not
an identification of the lattices. Division of the
Virasoro--Mukai K$3$ product by the Monster product yields
\eqref{eq:as-cplus-ratio-from-product}; the numerator
$\kapch(\Vnat) = 12$ and denominator $\kapch(\VKthree) = 3$
give $12/3 = 4$, matching $c_+$-ratio $4/1 = 4$ of
Theorem~\ref{thm:monster-k3-lusztig-cplus-face}.
\end{proof}

\begin{remark}[Monster $196884$ is not a $c_+$-weighted K$3$
elliptic-genus coefficient;
\ClaimStatusProvedHere]
\label{cor:as-monster-196884-as-cplus-weighted}
\index{Monster!$196884$!$c_+$-factored form}%
\index{K$3$ elliptic genus!$\chi(K3) = 24$!Mukai doubling}%
The weight-two coefficient of the Monster partition function and
the constant term of the K$3$ elliptic genus satisfy
\begin{equation}
\label{eq:as-monster-196884-cplus}
 \dim \Vnat_2
 \;=\;
 196884
 \;=\;
 1 \,+\, 196883,
 \qquad
 \phi_{0, 1}(\tau, 0)
 \;=\;
 \chi(K3)
 \;=\;
 24,
\end{equation}
and the $c_+$-product invariant of
Theorem~\ref{thm:as-monster-k3-cplus-product-invariant} gives only
the denominator comparison
\begin{equation}
\label{eq:as-cplus-weighted-moonshine}
 \frac{\dim \Vnat_2 - 1}{\chi(K3)}
 \;=\;
 \frac{196883}{24}
 \;=\;
 \frac{196883}{2 \kapch(\VKthree) \cdot c_+(\mathrm{Mukai}(II_{3, 19}))},
\end{equation}
with $196883$ the Griess-algebra dimension
(the smallest nontrivial Monster irreducible) and the denominator
$24 = 2 \kapch \cdot c_+$ the Mukai-K$3$-side factorisation of
the shared $24$. This quotient is not an integer and is not a
Fourier coefficient of the K$3$ elliptic genus. Equivalently, the
Monster McKay identity
$196884 = 196883 + 1$ ($1$-Hauptmodul-pole subtraction) and the
K$3$ signature split $24 = 4 + 20 = c_+ + c_-$ of the Mukai
lattice pass the falsification test
\begin{equation}
\label{eq:as-griess-over-mukai-ratio}
 \frac{196883}{c_-(\mathrm{Mukai}(II_{3, 19}))}
 \;=\;
 \frac{196883}{20}
 \;\notin\; \mathbb Z,
 \qquad
 \frac{196883 + 1}{c_+ + c_-}
 \;=\;
 \frac{196884}{24}
 \;=\;
 \frac{16407}{2},
\end{equation}
so no integral $c_\pm$-linear K$3$ decomposition of the Monster
coefficient exists.

Borcherds~$1992$~\cite{Bor92} Theorem~$6.1$ gives
$\dim \Vnat_n = c(n)$ where
$J(\tau) = \sum_n c(n) q^n = q^{-1} + 196884 q + \cdots$, so
$\dim \Vnat_2 = 196884$ (here indexing is the weight-two Virasoro
grade, $L_0 - c/24 = 1$, matching Conway--Norton). Griess~$1982$
\emph{Invent.\ Math.}~$69$ constructs $196883 = \dim V_2^\natural
- \dim \mathrm{Vir}_{c = 24, L_0 = 2} = 196884 - 1$ as the
minimal faithful Monster representation. Eguchi--Ooguri--Tachikawa~$2010$
\S$2$: $\phi_{0, 1}(\tau, 0) = \chi(K3) = 24$ by the
Atiyah--Singer index formula applied to the $\mathcal N{=}4$
chiral partition function at $z = 0$. Division
$196883/24$ gives a non-integer, which confirms that the
integer structure resides in $196884 = 196883 + 1$ (the $+1$ is
the Virasoro vacuum / Hauptmodul pole subtraction) against the
Mukai total rank $24$. The non-integrality of $196883/20$ at
$c_- = 20$ is the sharpest falsification test: if the Griess
dimension split across the Mukai $c_\pm$ decomposition were
integer-linear in $c_-$ alone, a relation $196883 = k \cdot 20$
would hold for integer $k$; it does not. The full $c_+ + c_- = 24$
divisor is the only integer-compatible denominator after the
Hauptmodul-pole shift, and even then the quotient is half-integral.
\end{remark}

\begin{remark}[Verlinde-style check: Monster fusion collapse forces
the Mukai $c_+ = 4$ side to carry the fusion content]
\label{rem:as-verlinde-monster-collapse}
\index{Verlinde formula!Monster!holomorphic collapse}%
\index{Mukai lattice!fusion content!K$3$ side}%
The Monster VOA $\Vnat$ is holomorphic (Dong~$1994$
\emph{J.\ Algebra}~$161$; Dong--Mason~$1994$~\cite{DongMason94}
Theorem~$4.1$), so the Verlinde ring of $\Vnat$ is
$K_0(\mathrm{Rep}_{\Vnat}^{\mathrm{ss}}) = \mathbb Z$ and carries
no nontrivial fusion coefficients $N^k_{ij}$. The Mukai-K$3$ side
carries its own nontrivial $\mathcal N{=}4$ representation theory,
whereas the product theorem above uses only its Virasoro lane. Neither
structure is transported from the Monster Verlinde ring. The
$c_+$-product comparison records compatible scalar normalisations; it
does not transfer fusion coefficients from one category to the other.
\end{remark}

\begin{remark}[Three independent paths to the $c_+$-ratio $4$]
\label{rem:as-cplus-ratio-triverification}
\index{$c_+$-ratio!three verifications}%
The $c_+$-ratio $4$ of
Theorem~\ref{thm:monster-k3-lusztig-cplus-face} and
Theorem~\ref{thm:as-monster-k3-cplus-product-invariant} has three
independent numerical checks.
\textup{(i)} \emph{Central-charge ratio:}
$c(\Vnat)/c(\VKthree_{\mathrm{Vir}}) = 24/6 = 4$; together with
$\kapch = c/2$ on the Virasoro lane
this gives $\kapch(\Vnat)/\kapch(\VKthree) = 12/3 = 4$
(primary: FLM~$1988$, Eguchi--Taormina~$1988$).
\textup{(ii)} \emph{Lattice signature-positive ratio:}
$c_+(\mathrm{Mukai}(II_{3, 19}))/c_+(II_{1, 1}) = 4/1 = 4$
(primary: Mukai~$1987$, Conway--Norton~$1979$).
\textup{(iii)} \emph{Lusztig level ratio:}
$\ell_{K3}/\ell_{\mathrm{Monster}} = 8/2 = 4$
(Theorem~\ref{thm:monster-lusztig-universal-ratio}; primary:
Lusztig~$1990$, Gritsenko--Nikulin~$1997$). Three computations
over three disjoint primary sources, all giving $4$. The
$c_+$-product comparison of
Theorem~\ref{thm:as-monster-k3-cplus-product-invariant} records the
same compatibility:
$2 \kapch \cdot c_+ = 24$ on the (Monster, Virasoro--Mukai K$3$) pair
makes \textup{(i)} and \textup{(ii)} reciprocal
on this pair, and \textup{(iii)} is the
$K^\kappa = 2 c_+$ Mukai-doubled read of \textup{(ii)}.
\end{remark}

% ============================================================
\section{Conditional PBW dimension of $u_{\zeta_8}(\mathfrak g_{\Delta_5})$
on the height-$8$ truncation}
\label{sec:exact-pbw-u-zeta-8}

The upper bound $|\Lambda_\infty| \le 8^{129}$ of
Theorem~\ref{thm:ent-topological-entanglement} counts PBW monomials
of projective covers on the nilpotent half $\mathfrak n^+_{\le 8}$
only. If the height-$8$ real-root truncation of
$\mathfrak g_{\Delta_5}$ admits a Lusztig/Nichols triangular
decomposition with the stated PBW root set, then the full
root-of-unity Hopf algebra includes both nilpotent halves and the
Cartan torus. The finite-type theorem supplies the dimension formula
after that truncation hypothesis is proved; it does not by itself
construct the BKM small quantum group.

\begin{theorem}[Conditional Lusztig small-quantum-group dimension on
the height-$8$ K$3$-BKM truncation; \ClaimStatusConditional]
\label{thm:exact-pbw-u-zeta-8-truncation}
\index{small quantum group!$u_{\zeta_8}$!exact dimension}%
\index{Lusztig form!K3-BKM!height-$8$ truncation}%
Assume $\mathfrak g_{\Delta_5,\le8}$ is a PBW-finite real-root
truncation with triangular decomposition
$u_{\zeta_8}^-\otimes u_{\zeta_8}^0\otimes u_{\zeta_8}^+$. Assume
$|\Phi^+_{\mathrm{re}, \le 8}| = 129$ positive real roots and rank
$r = 3$ Cartan $\Lambda^{2, 1}_{\mathrm{II}}$. The Lusztig small
quantum group $u_{\zeta_8}(\mathfrak g_{\Delta_5, \le 8})$ at
$\zeta_8 = e^{2\pi i/8}$ then has dimension
\begin{equation}\label{eq:exact-pbw-u-zeta-8-dim}
 \dim u_{\zeta_8}(\mathfrak g_{\Delta_5, \le 8})
 \;=\;
 8^{2\cdot 129 + 3}
 \;=\;
 8^{261},
\end{equation}
sharpening the nilpotent-half upper bound $8^{129}$ of
Theorem~\ref{thm:ent-topological-entanglement} by the factor
$8^{132}$ coming from the negative nilpotent $\mathfrak n^-$
(exactly $8^{129}$) and the Cartan torus $(\mathbb Z/8)^3$
(exactly $8^3$).
\end{theorem}

\begin{proof}
Lusztig~$1990$ \emph{Geom.\ Ded.}~$35$ Thm~$8.3$ gives, for a
finite-type root system $\Phi$ with $|\Phi^+|$ positive roots and
rank $r$ under the standard root-of-unity hypotheses,
\[
 \dim u_\ell(\mathfrak g_\Phi) = \ell^{2|\Phi^+| + r} = \ell^{\dim\mathfrak g_\Phi}
\]
with the triangular decomposition
$u_\ell = u_\ell^- \otimes u_\ell^0 \otimes u_\ell^+$ where
$\dim u_\ell^\pm = \ell^{|\Phi^+|}$ and $\dim u_\ell^0 =
\ell^r$. Under the PBW-finite triangular-decomposition hypothesis
of the theorem, the same counting argument applies to the selected
real-root truncation. This is a conditional transfer of the
finite-type formula, not an unconditional BKM theorem. Three checks
of the input data remain.
\textup{(i)} Triangular decomposition: the K$3$-BKM truncation
$\mathfrak g_{\Delta_5, \le 8}$ has $|\Phi^+_{\mathrm{re}, \le 8}|
= 129$ (Theorem~\ref{thm:ent-topological-entanglement}) and
$r = 3$ (\texttt{nilpotent\_completion.tex} Remark surrounding
line~$1757$, Cartan $\Lambda^{2, 1}_{\mathrm{II}}$);
$2 \cdot 129 + 3 = 261$ gives $8^{261}$.
\textup{(ii)} The Shapovalov form must be checked on the truncated
root set; finite-type non-degeneracy does not automatically pass to
the BKM quotient.
\textup{(iii)} A Frobenius--Lusztig kernel sequence for this
truncation would give the same exponent, but its construction is an
open input for the BKM setting.
\end{proof}

\begin{remark}[Why the $8^{129}$ bound appears in the literature as
``nilpotent-half''; scope correction]
\label{rem:as-8-129-nilpotent-half-scope}
\index{PBW bound!$8^{129}$ vs.\ $8^{261}$!scope}%
The bound $8^{129}$ of Vol~III
Theorem~\ref{V3-thm:modtr-ZT2-HH0}
and Vol~II
Proposition~\ref{V2-prop:modbootstrap-cross-block-count}
counts PBW monomials in the positive nilpotent $\mathfrak n^+_{\le 8}$
only. Theorem~\ref{thm:exact-pbw-u-zeta-8-truncation} upgrades this
conditionally to the Hopf-algebra dimension $8^{261} = 8^{129} \cdot
8^{129} \cdot 8^3$ by including negative nilpotent $\mathfrak n^-$
and Cartan torus $(\mathbb Z/8)^3$. The conformal-block count
$|\Lambda_\infty| = 130 = 1 + 129$ of
Remark~\ref{rem:stqp-block-basis-crossing}, being the number of
\emph{simple modules} (irreducible representations), is bounded
by $|\Lambda_\infty| \le \dim u_{\zeta_8}^+$ (Drinfeld--
Kirillov~$1991$ \emph{Leningrad Math.\ J.}~$1$ \S $3$); the
factor $8^{129} - 130 = $ a very large positive integer measures
the PBW-monomial redundancy under Serre-relation quotients, and
is not a contradiction. The full Drinfeld centre
$Z^{\mathrm{der}}_{\mathrm{ch}}(\mathbf H_{\Delta_5}) =
Z(u_{\zeta_8})$ has dimension
$\dim Z(u_{\zeta_8}) = \sum_{V \in \mathrm{Irr}}(\dim V)^2
\le 8^{261}$ only after the same root-of-unity BKM construction is
available.
\end{remark}

% ============================================================
\section{Conditional PBW dimensions of $u_{\zeta_{12}}$ and $u_{\zeta_{24}}$
on the $\Delta_5$ real-root truncation}
\label{sec:exact-pbw-u-zeta-12-24}

The BKM ladder indexed by $N \in \{1, 2, 3, 4, 6, 12, 24\}$ reaches
even Lusztig levels $\ell \in \{12, 24\}$ at the paramodular covers
$N = 12$ (Gritsenko--Cl\'ery form $\phi_{12}$ of weight
$w(12) = 1/4$, Fourier zero-coefficient $c_{12}(0) = 1/2$) and
$N = 24$ (degenerate weight $w(24) = 0$, $c_{24}(0) = 0$). The
formulas below are conditional on the same PBW-finite BKM
truncation hypothesis as Theorem~\ref{thm:exact-pbw-u-zeta-8-truncation}.
No statement here says that the K$3$--Monster $c_+$ ratio transports
all root-of-unity categories to $\ell=12$ or $\ell=24$.

\begin{theorem}[Conditional Lusztig small-quantum-group dimensions at
$\ell \in \{12, 24\}$ on the height-$8$ $\Delta_5$ truncation;
\ClaimStatusConditional]
\label{thm:exact-pbw-u-zeta-12-24-truncation}
\index{small quantum group!$u_{\zeta_{12}}$!exact dimension}%
\index{small quantum group!$u_{\zeta_{24}}$!exact dimension}%
\index{Lusztig form!K3-BKM!even levels $\ell \in \{12, 24\}$}%
Assume $\mathfrak g_{\Delta_5,\le8}$ is a PBW-finite real-root
truncation with triangular decomposition at levels
$\ell\in\{12,24\}$, with
$|\Phi^+_{\mathrm{re}, \le 8}| = 129$ and rank $r = 3$. At Lusztig
level $\ell \in \{12, 24\}$ with $\zeta_\ell = e^{2\pi i/\ell}$, the
small quantum group $u_{\zeta_\ell}(\mathfrak g_{\Delta_5, \le 8})$
then has Hopf-algebra dimension
\begin{equation}\label{eq:exact-pbw-u-zeta-12-24-dim}
 \dim u_{\zeta_\ell}(\mathfrak g_{\Delta_5, \le 8})
 \;=\;
 \ell^{\,2\cdot 129 + 3}
 \;=\;
 \ell^{261},
\end{equation}
with triangular decomposition $u_{\zeta_\ell} = u_{\zeta_\ell}^-
\otimes u_{\zeta_\ell}^0 \otimes u_{\zeta_\ell}^+$,
$\dim u_{\zeta_\ell}^\pm = \ell^{129}$, $\dim u_{\zeta_\ell}^0 =
\ell^3$. Explicit values:
\[
 \dim u_{\zeta_{12}}(\mathfrak g_{\Delta_5, \le 8}) = 12^{261},
 \qquad
 \dim u_{\zeta_{24}}(\mathfrak g_{\Delta_5, \le 8}) = 24^{261}.
\]
\end{theorem}

\begin{proof}
Under the stated PBW-finite triangular-decomposition hypothesis,
the counting is formal:
$\dim u_{\zeta_\ell}^+=\ell^{129}$,
$\dim u_{\zeta_\ell}^-=\ell^{129}$, and
$\dim u_{\zeta_\ell}^0=\ell^3$. Multiplication gives
$\ell^{129}\ell^3\ell^{129}=\ell^{261}$. Substituting
$\ell=12$ and $\ell=24$ gives the two displayed values.

The theorem deliberately does not use the finite-type Lusztig
coprimality hypotheses as a substitute for a BKM construction. The
missing checks are: closure of the height-$8$ truncation under the
Hopf structure at even level, compatibility of the Nichols algebra
relations with the BKM real-root diagram, and non-interference from
imaginary simple roots.
\end{proof}

\begin{remark}[The level-$\ell$ K$3$-BKM ladder and the Borcherds weight]
\label{rem:as-bkm-ladder-ell-12-24}
\index{BKM ladder!$N \in \{12, 24\}$!Lusztig level}%
The dimension ratio at $\ell = 12$ vs.\ $\ell = 8$ is
$(12/8)^{261} = (3/2)^{261} \approx 1.94 \cdot 10^{45}$ under the
conditional PBW model; it is not yet a computed magnification of
the K$3$ chiral bialgebra. The
full Drinfeld centre $Z^{\mathrm{der}}_{\mathrm{ch}}(\mathbf
H_{\Delta_5})$ at level $\ell$ has dimension
$\sum_{V \in \mathrm{Irr}}(\dim V)^2 \le \ell^{261}$ only after
the corresponding root-of-unity BKM category has been constructed.
The $N = 12$ and $N = 24$ rungs are therefore proposed PBW
specialisations, not completed entries in the Lusztig spread.
\end{remark}

\begin{remark}[Monster-ratio consequence at $\ell \in \{12, 24\}$]
\label{rem:as-monster-ratio-ell-12-24-consequence}
\index{Monster Lusztig level!at $N \in \{12, 24\}$}%
Applied at $\ell_{K3} = 12$, the K$3$-Monster ratio
$\ell_{K3}/\ell_{\mathrm{Monster}} = 4$ of
Theorem~\ref{thm:monster-lusztig-universal-ratio} would give
$\ell_{\mathrm{Monster}}^{(N = 12)} = 3$ if the same
$\ell=2c_+$ convention extended to that rung. At $\ell_{K3}=24$ it
would give $\ell_{\mathrm{Monster}}^{(N = 24)}=6$. These are
formal consequences of a convention, not constructions of
Monster small quantum groups at levels $3$ and $6$. The Cartan
contribution alone would be $\ell^{26}$ at rank $r=26$, but the
real-root contribution and the root-of-unity category require a
separate fake-Monster BKM construction from the Borcherds
denominator
$\prod_{\alpha > 0}(1 - e^{-\alpha})^{\mathrm{mult}(\alpha)}$
(Borcherds~$1998$ Thm~$13.1$). The $\mathsf B$-row Theorem~C
ceiling $K^\kappa = 8$ of
Theorem~\ref{thm:monster-k3-lusztig-cplus-face} is the unique
integer-valued representative in $\{0, 8, 13, 250/3, 98/3\}$ on the
paramodular tower at $N = 1$. No rescaling formula
$K^\kappa(N)=8c_N(0)/c_1(0)$ is asserted here; it would identify a
chiral conductor with Borcherds input coefficients, the separation
ruled out above.
\end{remark}

% ============================================================
\section{Yetter--Drinfeld Schauenburg-bracket expansion
$\delta^{(n)}$, $n \in \{13, 14, 15, 16\}$}
\label{sec:YD-delta-n-13-16}

Theorem~\ref{thm:YD-delta-7-8-9} computed $\delta^{(n)}$ for
$n \in \{7, \ldots, 12\}$. The Catalan--Padovan closed form
$C_{n-1} \cdot d_n$ extends verbatim to the next four weights.

\begin{theorem}[$\delta^{(n)}$ for $n \in \{13, 14, 15, 16\}$;
\ClaimStatusProvedHere]
\label{thm:YD-delta-13-16}
\index{Yetter-Drinfeld!Schauenburg bracket!$n = 13$--$16$}%
\index{Borcherds weight!$\delta^{(n)}$ expansion!high degree}%
The Yetter--Drinfeld Schauenburg-bracket expansion
$\delta^{(n)}$ of Borcherds weight $\lfloor n/2 \rfloor + 1$
at degrees $n \in \{13, 14, 15, 16\}$ has total Catalan--MZV
dimension $C_{n-1} \cdot d_n$ with
\begin{equation}\label{eq:YD-delta-13-16-dim}
\begin{aligned}
 n = 13 &:\quad C_{12}\cdot d_{13} = 208012\cdot 16 = 3\,328\,192,
 \qquad \mathrm{wt}(\delta^{(13)}) = 7, \\
 n = 14 &:\quad C_{13}\cdot d_{14} = 742900\cdot 21 = 15\,600\,900,
 \qquad \mathrm{wt}(\delta^{(14)}) = 8, \\
 n = 15 &:\quad C_{14}\cdot d_{15} = 2\,674\,440\cdot 28 = 74\,884\,320,
 \qquad \mathrm{wt}(\delta^{(15)}) = 8, \\
 n = 16 &:\quad C_{15}\cdot d_{16} = 9\,694\,845\cdot 37 = 358\,709\,265,
 \qquad \mathrm{wt}(\delta^{(16)}) = 9,
\end{aligned}
\end{equation}
with $C_k$ the $k$-th Catalan number (OEIS~A000108:
$C_{12} = 208012, C_{13} = 742900, C_{14} = 2\,674\,440,
C_{15} = 9\,694\,845$) and $d_n$ the Zagier--Brown canonical
motivic MZV dimension at weight $n$, computed by the Padovan
recurrence $d_n = d_{n-2} + d_{n-3}$ with seed
$(d_3, d_4, d_5) = (1, 1, 1)$ and continuation
$(d_{11}, d_{12}) = (9, 12)$ from Theorem~\ref{thm:YD-delta-7-8-9},
giving $(d_{13}, d_{14}, d_{15}, d_{16}) = (16, 21, 28, 37)$. The
Borcherds-weight closed form $\mathrm{wt}(\delta^{(n)}) =
\lfloor n/2 \rfloor + 1$ persists throughout, so the step
sequence across $n \in \{12, 13, 14, 15, 16\}$ is $(0, 1, 0, 1)$
(weight $7 \to 7 \to 8 \to 8 \to 9$).
\end{theorem}

\begin{proof}
The Catalan count follows from Schauenburg~$1998$ Cor~$2.3$ (as in
Theorem~\ref{thm:YD-delta-7-8-9}) applied at depth $n$: non-crossing
partitions of $[n]$ enumerate to $C_{n-1}$
(Stanley~$1999$ \emph{EC2} Ex.~$6.19$). The Brown motivic MZV
dimension follows from Brown~$2012$ \emph{Ann.\ Math.}~$175$
Thm~$1.2$: the Padovan recurrence $d_n = d_{n-2} + d_{n-3}$ is
Brown~$2011$ Theorem~$1.1$, extended continuously from the seed
$(d_3, d_4, d_5) = (1, 1, 1)$. Four verification paths.
\textup{(i)} Padovan-recurrence iteration: starting from
$(d_{10}, d_{11}, d_{12}) = (7, 9, 12)$ of
Theorem~\ref{thm:YD-delta-7-8-9}, the recurrence gives
$d_{13} = d_{11} + d_{10} = 9 + 7 = 16$,
$d_{14} = d_{12} + d_{11} = 12 + 9 = 21$,
$d_{15} = d_{13} + d_{12} = 16 + 12 = 28$,
$d_{16} = d_{14} + d_{13} = 21 + 16 = 37$.
\textup{(ii)} Broadhurst--Kreimer generating series at
$n = 13, 14, 15, 16$: the depth-graded series
$\sum_{n, d}\dim\mathrm{gr}^{\mathrm{BK}}_{n, d} x^n y^d =
\frac{1 + y x^3}{1 - y x^2 - y x^4 + y^2 x^4 + y^2 x^6
- y^3 x^{12}}$
(Broadhurst--Kreimer~$1997$ \emph{Phys.\ Lett.\ B}~$393$;
refined in Brown~$2011$ Thm~$1.2$) expanded through $x^{16}$
gives the depth distributions
$d_{13} = 1 + 2 + 5 + 6 + 2 + 0 = 16$ (depths $1, 2, 3, 4, 5, 6$),
$d_{14} = 0 + 3 + 5 + 9 + 4 + 0 = 21$,
$d_{15} = 1 + 2 + 7 + 11 + 7 + 0 = 28$,
$d_{16} = 0 + 3 + 8 + 14 + 11 + 1 = 37$; the single depth-$6$
motivic MZV at $n = 16$ is $\zeta(3, 3, 3, 3, 2, 2)$, the first
depth-$6$ irreducible in Brown's enumeration
(Brown~$2013$ Bourbaki~$1054$ \S $3.5$).
\textup{(iii)} Borcherds-weight induction: the step sequence
$(0, 1, 0, 1)$ across $n \in \{12, 13, 14, 15, 16\}$ realises the
Etingof--Kazhdan~$2007$ \S $3$ bound
$\mathrm{wt}(\delta^{(n+1)}) - \mathrm{wt}(\delta^{(n)})
\in \{0, 1\}$ at every step, matching the pattern
$(0, 1, 0, 1)$ of the previous range $n \in \{8, \ldots, 12\}$.
\textup{(iv)} Catalan--Padovan product at $n = 15$ direct:
$C_{14} = \binom{28}{14}/15 = 40\,116\,600/15 = 2\,674\,440$
(Catalan closed form), $d_{15} = 28$ (Padovan), product
$= 74\,884\,320$ matching \eqref{eq:YD-delta-13-16-dim}.
\end{proof}

% ============================================================
\section{GRT$_1$-transitivity on K$3$-BKM: the $\Delta_5$ face}
\label{sec:grt-1-transitivity-bkm-delta5}

Affine Kac--Moody GRT$_1$-transitivity was established in
Theorem~\ref{thm:grt-affine-km-transitivity-foundation} (standard
chiral landscape). The BKM/root-of-unity extension requires extra
structure: a finite tensor category, a compatible root-of-unity
Kazhdan--Lusztig comparison, and preservation of the BKM Cartan
grading. The Vol~I statement below isolates that conditional bridge.

\begin{theorem}[Conditional GRT$_1$ transport on the $\Delta_5$ BKM;
\ClaimStatusConditional]
\label{thm:grt-1-transitivity-delta5}
\index{GRT$_1$!BKM!transitivity}%
\index{Drinfeld associator!K3-BKM!$\Delta_5$-equivariance}%
Assume:
\begin{enumerate}[label=\textup{(A\arabic*)}]
\item the height-$8$ BKM truncation
$u_{\zeta_8}(\mathfrak g_{\Delta_5,\le8})$ carries a finite
Drinfeld--Yetter tensor category
$\mathcal C_{\Delta_5}$;
\item Enriquez--Etingof transport extends to this root-of-unity BKM
category and makes the associator choices a GRT$_1$-torsor;
\item the transport fixes the BKM Cartan grading, the Weyl vector,
and the root multiplicities of $\mathfrak g_{\Delta_5}$.
\end{enumerate}
Then the GRT$_1$-transport changes the associator gauge but not the
Borcherds denominator product:
\begin{equation}\label{eq:grt-1-delta5-preservation}
 \prod_{\alpha \in \Phi^+_{\Delta_5}}
 (1 - e^{-\alpha})^{\mathrm{mult}(\alpha)}
 \;=\;
 \Delta_5(Z),
\end{equation}
in the fixed Gritsenko--Nikulin normalisation. The theorem does not
assert that hypotheses \textup{(A1)}--\textup{(A3)} have been
constructed in this volume.
\end{theorem}

\begin{proof}
By \textup{(A2)}, GRT$_1$ acts on the associator choices. By
\textup{(A3)}, this action leaves the root lattice, Weyl vector, and
multiplicity function $\alpha\mapsto\dim
(\mathfrak g_{\Delta_5})_\alpha$ fixed. The Borcherds denominator
in the Gritsenko--Nikulin normalisation is determined by exactly
those data together with the fixed Siegel variable $Z$. Therefore
the product in \eqref{eq:grt-1-delta5-preservation} is unchanged.
No calculation with the depth generators of $\mathfrak{grt}_1$ is
used; such a calculation would be part of proving \textup{(A3)}.
\end{proof}

\begin{remark}[GRT$_1$-transitivity on fake-Monster BKM vs.\ K$3$-BKM]
\label{rem:as-grt-1-fake-monster-vs-k3}
\index{GRT$_1$!fake Monster BKM!transitivity}%
The analogue of Theorem~\ref{thm:grt-1-transitivity-delta5} for the
fake-Monster BKM $\mathfrak m_{\mathrm{FakeMonster}}$ (signature
$(25, 1)$, Cartan $\Lambda_{24} \oplus II_{1, 1}$, denominator
$\Phi_{12}$ of weight $12$) is conditional in the same sense. One
must first construct the root-of-unity BKM tensor category, the
GRT$_1$ associator torsor, and the Cartan-graded transport for the
$O(26,2)$ Borcherds datum. After those hypotheses are in place,
the denominator is fixed for the same formal reason as
\eqref{eq:grt-1-delta5-preservation}. Without them there is no
theorem in this volume saying that GRT$_1$ acts transitively on
every BKM in the K$3$/Leech/Monster landscape.
\end{remark}

\begin{remark}[GRT$_1$-action on the modular $S$-matrix of
$\mathcal C_{\Delta_5}$]
\label{rem:as-grt-1-S-matrix}
\index{GRT$_1$!modular $S$-matrix!invariance}%
The modular $S$-matrix of $\mathcal C_{\Delta_5} =
\mathrm{Rep}^{\mathrm{fd}}(u_{\zeta_8}(\mathfrak g_{\Delta_5, \le 8}))$
is indexed by
$\{F_0\} \cup \Lambda^+_{\Delta_5, 8}$ when the finite
root-of-unity category exists. Invariance of this matrix is not
forced by Siegel modularity of $\Delta_5$ alone: an associator
change can alter braiding, ribbon traces, and logarithmic modified
trace data. The theorem below records the formal consequence of
fixing those data, not an unconditional computation of the
$130\times130$ matrix.
\end{remark}

\begin{theorem}[Conditional GRT$_1$-invariance of the
$130\times 130$ $S$-matrix; \ClaimStatusConditional]
\label{thm:grt-1-S-matrix-trivial}
\index{GRT$_1$!modular $S$-matrix!identity action}%
\index{modular $S$-matrix!$\mathbf H_{\Delta_5}$!GRT$_1$-invariance}%
Assume \textup{(A1)}--\textup{(A3)} of
Theorem~\ref{thm:grt-1-transitivity-delta5}. Assume further:
\begin{enumerate}[label=\textup{(B\arabic*)}]
\item $\mathcal C_{\Delta_5}$ has a well-defined modular datum with
label set
$\mathsf I_{130}=\{0\}\cup\Lambda^+_{\Delta_5,\le8}$;
\item the $S$-matrix admits a factorisation into categorical trace,
Fourier-kernel, global-normalisation, and Fricke factors;
\item the GRT$_1$ transport fixes each of those four factors.
\end{enumerate}
Then for every $\sigma\in\mathrm{GRT}_1$ the transported matrix
satisfies
\begin{equation}\label{eq:grt-1-S-matrix-trivial}
 S_\sigma \;=\; S,
 \qquad
 (\rho_\sigma)_{\lambda\mu} = \delta_{\lambda\mu}
 \text{ on }
 \mathsf I_{130}\times\mathsf I_{130}.
\end{equation}
Without \textup{(B1)}--\textup{(B3)}, this volume has no proof that
the GRT$_1$ action on modular data is trivial.
\end{theorem}

\begin{proof}
Under \textup{(B2)}, each entry of $S$ is a product of four named
factors. Under \textup{(B3)}, the GRT$_1$ transport fixes each
factor and the label set of \textup{(B1)}. Therefore every matrix
entry is fixed:
$S_{\sigma,\lambda\mu}=S_{\lambda\mu}$ for all
$\lambda,\mu\in\mathsf I_{130}$. This proves
\eqref{eq:grt-1-S-matrix-trivial}. The proof is deliberately
formal; constructing the factorisation and checking invariance of
modified traces are the mathematical substance still required.
\end{proof}

\begin{remark}[Source-level conditions for
Theorem~\ref{thm:grt-1-S-matrix-trivial}]
\label{rem:as-grt-1-S-matrix-verif}
\index{GRT$_1$!$S$-matrix triviality!verifications}%
Theorem~\ref{thm:grt-1-S-matrix-trivial} is proved exactly under
the following four source-level constructions.
\textup{(V1)} Construct the relevant BKM graph complex and show that
the $\mathfrak{grt}_1$ action on the Cartan-graded modified trace
vanishes; Willwacher's identification of
$\mathfrak{grt}_1$ with graph cohomology is not by itself this
calculation.
\textup{(V2)} Compute the actual $130\times130$ Fourier/modular
matrix from the Gritsenko--Nikulin expansion and identify it with
the categorical $S$-matrix.
\textup{(V3)} Prove that the Leech degeneration
$\Delta_5\to\Phi_{12}$ carries the $130$ labels to the unique
Monster label by a GRT$_1$-equivariant functor, with the
dimension-collapse calculation as a necessary numerical check.
\textup{(V4)} Compare any Vol~II modular-bootstrap matrix with the
same label set and trace convention. Until these four constructions
are supplied, Siegel modularity controls only the automorphic frame;
it does not prove GRT$_1$-triviality of categorical modular data.
\end{remark}

\begin{remark}[Gauge-comparison obligation behind
GRT$_1$-triviality]
\label{rem:as-grt-1-S-matrix-ghost}
\index{ghost theorem!GRT$_1$!Siegel modularity}%
There are two gauges: the associator gauge on the root-of-unity
quantum-group side and the automorphic gauge of the Borcherds
denominator $\Delta_5$. Siegel modularity controls the second gauge.
It does not automatically control the first, because the categorical
$S$-matrix also depends on braiding and trace data. The missing
argument is a functorial bridge proving that the GRT$_1$ associator
torsor maps to the automorphic frame of $\Delta_5$ with trivial
kernel on modular data. The categorical rigidity statement becomes a
theorem only after this root-of-unity BKM comparison functor has been
built.
\end{remark}

% ============================================================
\section{Root-of-unity $N = 2$: $324$ simple modules, degenerate
$S$-matrix}
\label{sec:root-of-unity-N-2-vol-I-face}

The Vol~III paramodular level-$2$ truncation $\mathbf H_{\Phi_2}$
is recorded as having $324$ simple modules and a degenerate modular
$S$-matrix (Vol~III Proposition~\texttt{prop:root-unity-n2}). The
Vol~I arithmetic check separates that simple-module count from a PBW
dimension. The equality between the two is false.

\begin{theorem}[$N = 2$ root-of-unity: $324$ is not a PBW dimension;
\ClaimStatusProvedHere]
\label{thm:n-2-root-unity-vol-I-face}
\index{root of unity!$N = 2$!Vol I face}%
\index{modular $S$-matrix!degenerate!$\mathbf H_{\Phi_2}$}%
At the paramodular level $N = 2$ (Borcherds lift weight $2$,
$c_2(0) = 4$, $\kappa_{\mathrm{BKM}}(\Phi_2) = 2$) the Lusztig
small quantum group $u_{\zeta_4}(\mathfrak g_{\Phi_2, \le 2})$ on
the height-$2$ real-root truncation has PBW nilpotent-half dimension
\begin{equation}\label{eq:n2-nilpotent-half-324}
 \dim u_{\zeta_4}^+(\mathfrak g_{\Phi_2, \le 2})
 \;=\;
 4^{|\Phi^+_{\mathrm{re},\le2}|}.
\end{equation}
If the level-$2$ input has $4$ positive real roots at height $1$
and $14$ additional positive real roots at height $2$, then
$|\Phi^+_{\mathrm{re},\le2}|=18$ and
\begin{equation}\label{eq:n2-pbw-not-324}
 \dim u_{\zeta_4}^+(\mathfrak g_{\Phi_2, \le 2})
 \;=\;
 4^{18}
 \;=\;
 2^{36}
 \;\neq\;
 324.
\end{equation}
Even the height-$1$ subalgebra alone gives $4^4=256$, not $324$.
Thus a count $|\mathrm{Irr}(\mathbf H_{\Phi_2})|=324$, if used,
must be a simple-module, quotient, or stratum count; it is not the
PBW dimension of the nilpotent half. Likewise a rank statement
$\operatorname{rank}S=17$ is logarithmic modular data, not a
consequence of the PBW count.
\end{theorem}

\begin{proof}
The height-$2$ truncation of the $\Phi_2$ real-root subalgebra
has $4$ positive real roots at height $1$ and $14$ additional
roots at height $2$, giving $|\Phi^+_{\mathrm{re}, \le 2}| = 18$
by direct enumeration from the Gritsenko--Nikulin~$1998$ level-$2$
Jacobi input $\phi_{0, 2}^{N = 2}(\tau, z)$ (weight $0$ index $2$).
At $\zeta_4 = e^{2\pi i/4} = i$, Lusztig~$1990$ Thm~$8.3$ gives
$4$ PBW exponents for each real-root vector in the restricted
nilpotent half. Therefore the PBW count is $4^{18}$ for the
eighteen-root truncation. If one discards the height-$2$ roots and
keeps only the four height-$1$ root vectors, the same rule gives
$4^4=256$.

The expression
$4^4\cdot\binom{6}{2}^{\mathrm{height}\text{-}2}=256+68=324$
has no PBW interpretation: the displayed operation is
multiplicative, while the following equality is additive, and the
number of height-$2$ commutator monomials is not a replacement for
the PBW exponent choices of the corresponding root vectors. Hence the
only proved Vol~I statement is the separation:
\[
 324\ \text{may label simple objects or a quotient datum, but not }
 \dim u_{\zeta_4}^+(\mathfrak g_{\Phi_2,\le2}).
\]
\end{proof}

\begin{remark}[$N = 2$ S-matrix rank $17$ is a separate datum]
\label{rem:as-n-2-s-matrix-rank-17}
\index{Igusa $\Phi_2$!weight $2$!$S$-rank}%
If the Vol~III level-$2$ category has an $S$-matrix of rank $17$,
that rank must be proved from the logarithmic modular datum itself:
the coend, modified trace, or explicit pseudo-character
decomposition. Eichler--Zagier's dimension formula for Jacobi forms
does not by itself give the rank of a non-semisimple modular
$S$-matrix, and the PBW computation above gives neither $324$ nor
$17$. The open proof obligation is to construct the level-$2$
logarithmic $S$-matrix and identify its radical.
\end{remark}
